\documentclass[UTF8,12pt,a4paper,fontset=fandol]{ctexbook}


% 支持从项目根目录或当前笔记目录编译。
\makeatletter
\def\input@path{{../}{../../}{../../../}{../../../../}}
\makeatother
\usepackage{researchnotes}


\title{高等代数习题（谢启鸿）笔记}
\author{}
\date{\today}


\begin{document}


\maketitle
\tableofcontents
\clearpage


\chapter{行列式}
\label{chap:determinants}

% 研读来源：高等代数/textbooks/高等代数习题（谢启鸿）/高代习题-1.pdf，
% PDF 第 1--53 页，即原书第 1 章 §1.1--§1.12。
行列式（determinant）把一个方阵对应到一个数，既能判断方程组是否具有唯一解，又能把消元、排列符号、多项式因子与递推关系联系起来。本章按教材的十二节整理基本概念和计算方法，选择能够说明方法及其条件的例题展开推导，并在基础训练中处理易混的符号和操作。计算时应先观察零元素、重复部分、行列和及参数结构，再决定如何化简，而不是立即展开全部乘积。

除另有说明外，元素取自实数域或复数域，$n$ 为正整数，$I_n$ 为 $n$ 阶单位矩阵，$\mathbf1$ 为适当维数的全一列向量。方阵记为 $A=(a_{ij})$，行列式记为 $\det A$ 或 $|A|$；绝对值另写为 $|\det A|$。约定空乘积和零阶行列式均为 $1$，以便降阶公式在最低阶处保持一致。

\section{基本概念}
\label{sec:xie-ex-det-basics}

\subsection{排列定义与余子式展开}

\begin{definition}
\label{def:xie-ex-det-leibniz}
设 $S_n$ 为 $1,\ldots,n$ 的所有排列组成的集合。排列 $\sigma=(\sigma(1),\ldots,\sigma(n))$ 的\emph{逆序数}（inversion number）为
\[
N(\sigma)=\#\{(i,j):1\leq i<j\leq n,\ \sigma(i)>\sigma(j)\},
\]
其符号为 $\operatorname{sgn}(\sigma)=(-1)^{N(\sigma)}$。$n$ 阶方阵的行列式定义为
\begin{equation}
\label{eq:xie-ex-det-leibniz}
\det A=\sum_{\sigma\in S_n}\operatorname{sgn}(\sigma)
\prod_{i=1}^{n}a_{i,\sigma(i)}.
\end{equation}
每个乘积从每行、每列各取一个元素，共有 $n!$ 个带符号的项。特别地，$\det(a_{11})=a_{11}$。
\end{definition}

\begin{definition}
\label{def:xie-ex-det-cofactor}
删除 $A$ 的第 $i$ 行和第 $j$ 列，余下元素保持原有次序组成的行列式，称为元素 $a_{ij}$ 的\emph{余子式}（minor），记为 $M_{ij}$；$C_{ij}=(-1)^{i+j}M_{ij}$ 称为该元素的\emph{代数余子式}（cofactor）。这里用 $C_{ij}$ 区分矩阵 $A$ 与教材常用的代数余子式记号 $A_{ij}$。当 $n=1$ 时，$M_{11}=C_{11}=1$。
\end{definition}

\begin{theorem}[按一行或一列展开]
\label{thm:xie-ex-det-expansion}
对任意固定的行指标 $i$ 和列指标 $j$，
\begin{equation}
\label{eq:xie-ex-det-expansion}
\det A=\sum_{k=1}^{n}a_{ik}C_{ik}
=\sum_{k=1}^{n}a_{kj}C_{kj}.
\end{equation}
因而可以从一阶开始，按任一固定行或列递归定义行列式，所得结果与排列定义相同。
\end{theorem}

\begin{proof}
在式~\eqref{eq:xie-ex-det-leibniz} 中，把选中 $a_{ij}$ 的项归为一组。将第 $i$ 行和第 $j$ 列分别移到首位，需要 $i-1$ 次和 $j-1$ 次相邻对换，故这组的符号因子为 $(-1)^{i+j}$，剩余排列的求和为 $M_{ij}$。按固定行中的不同列分组，就得到行展开；列展开同理。递归定义与排列定义的等价性再由阶数归纳得到。
\end{proof}

\subsection{性质与行列变换}

\begin{proposition}
\label{prop:xie-ex-det-properties}
行列式对每行、每列分别是线性的；对换两行或两列，行列式变号；转置不改变行列式。某行或某列全零、两行或两列相同或成比例时，行列式为零。上三角或下三角矩阵的行列式等于主对角元素之积。
\end{proposition}

\begin{proof}
排列展开的每项恰含指定行中的一个元素，因此对该行具有可加性和齐次性，列的情形相同。对换两行相当于把排列复合一个对换，故符号改变。两行相同时，对换后矩阵不变而行列式变号，在实数域或复数域中便有行列式为零；成比例的情形再用齐次性。转置的项与原矩阵的逆排列项对应，逆排列具有相同符号。对于上三角矩阵，非零项要求 $\sigma(i)\geq i$；又因 $\sum_i\sigma(i)=\sum_i i$，只能处处相等，故只剩对角乘积。下三角情形由转置得到。
\end{proof}

用 $R_i,K_j$ 表示第 $i$ 行和第 $j$ 列。$R_i\leftarrow R_i+cR_k$（$i\ne k$）保持行列式不变，因为拆开后增加的行列式有两行成比例；列变换同理。把某行乘以 $c$ 会使行列式乘以 $c$，交换两行则变号，消元时必须记录这些因子。\emph[caution]{多重线性是对一行或一列分别线性，并不意味着 $\det(A+B)=\det A+\det B$。}把全部元素乘以 $c$，所得行列式为 $c^n\det A$，而不是 $c\det A$。

\begin{proposition}[分块三角性质]
\label{prop:xie-ex-det-basic-blocks}
若 $P,Q$ 为方阵，$X$ 为相应大小的矩阵，则 $\det\begin{pmatrix}P&X\\0&Q\end{pmatrix}=\det P\det Q$；分块下三角矩阵也有相同结论，不要求两个对角块可逆。
\end{proposition}

\begin{proof}
设 $P$ 为 $p$ 阶、$Q$ 为 $q$ 阶。非零排列项中，最后 $q$ 行只能选择最后 $q$ 列，因而必须恰好占满这些列；最前 $p$ 行便只能选择最前 $p$ 列。两组内部排列独立，符号为各自符号之积，分别求和即得两个对角块行列式的乘积。下三角情形由转置得到。
\end{proof}

\begin{proposition}
\label{prop:xie-ex-det-cross-cofactor}
同一行元素与另一行的对应代数余子式的乘积之和为零。完整地写成
\begin{equation}
\label{eq:xie-ex-det-cross-cofactor}
\sum_{j=1}^{n}a_{ij}C_{kj}
=\begin{cases}\det A,&i=k,\\0,&i\ne k.\end{cases}
\end{equation}
列也有相应公式 $\sum_i a_{ij}C_{i\ell}=\delta_{j\ell}\det A$，其中 $\delta_{j\ell}$ 在 $j=\ell$ 时为 $1$，否则为 $0$。
\end{proposition}

\begin{proof}
$i=k$ 时就是展开公式。$i\ne k$ 时，以第 $i$ 行替换第 $k$ 行，新矩阵有两行相同，行列式为零；删除第 $k$ 行后余子式不受替换影响，按这一行展开便得到左端。列公式同理。
\end{proof}

\subsection{Cramer 法则}

\begin{theorem}[Cramer 法则]
\label{thm:xie-ex-det-cramer}
设 $Ax=b$ 是 $n$ 个方程、$n$ 个未知量的线性方程组。若系数行列式 $D=\det A\ne0$，则方程组有唯一解
\begin{equation}
\label{eq:xie-ex-det-cramer}
x_j=\frac{D_j}{D},\qquad 1\leq j\leq n,
\end{equation}
其中 $D_j$ 是以常数列 $b$ 替换 $A$ 的第 $j$ 列后所得的行列式。
\end{theorem}

\begin{proof}
若 $x$ 为解，把 $b=\sum_k x_kA_{\cdot k}$ 代入 $D_j$ 并按替换列拆开，除 $k=j$ 外的项都有重复列，故 $D_j=x_jD$。这证明任何解都必须有上述形式。为证明存在性，按替换列展开得 $D_j=\sum_k b_kC_{kj}$，由式~\eqref{eq:xie-ex-det-cross-cofactor}，
\[
\sum_j a_{ij}\frac{D_j}{D}
=\frac1D\sum_k b_k\sum_j a_{ij}C_{kj}=b_i.
\]
因此给出的向量满足所有方程，且解唯一。
\end{proof}

$D\ne0$ 是比值公式的前提。若 $D=0$，不能仅凭这一点判定无解或只有零解：$0x=0$ 有无穷多解，$0x=1$ 无解。对齐次方程组，$D\ne0$ 保证只有零解；系数行列式为零时应进一步消元，或使用后续的秩理论。

\section{降阶法}
\label{sec:xie-ex-det-reduction}

降阶法先用行列变换制造零元素，再沿零最多的行或列展开。连续作差时操作次序也是推导的一部分：要让每个新行等于两个原行之差，通常应从下向上操作，避免前一行已经被改写。

% 原书例 1.1；引入平移参数准确描述降阶后的同类结构。
\begin{example}
\label{ex:xie-ex-det-binomial}
设 $s$ 为非负整数，求
\[
P_n(s)=\det\left(\binom{s+i+j-2}{j-1}\right)_{i,j=1}^{n}.
\]
特别地，$P_n(0)$ 的第一行和第一列均为 $1$。
\end{example}

\begin{solve}
从 $i=n$ 到 $i=2$ 作 $R_i\leftarrow R_i-R_{i-1}$。第一列除首项外全零，且对 $i,j\geq2$，
\[
\binom{s+i+j-2}{j-1}-\binom{s+i+j-3}{j-1}
=\binom{s+i+j-3}{j-2}.
\]
按第一列展开并重新编号，得到 $P_n(s)=P_{n-1}(s+1)$。由 $P_1(s)=1$ 得 $P_n(s)=1$。平移参数 $s$ 用于说明降阶后的结构，不能未经核对便把新行列式视为原式的简单缩小。
\end{solve}

% 原书例 1.4。
\begin{example}
\label{ex:xie-ex-det-arrow}
计算\emph{爪型行列式}（arrowhead determinant）
\[
H=\begin{vmatrix}
a_1&b_2&b_3&\cdots&b_n\\
c_2&a_2&0&\cdots&0\\
c_3&0&a_3&\cdots&0\\
\vdots&\vdots&\vdots&\ddots&\vdots\\
c_n&0&0&\cdots&a_n
\end{vmatrix},\qquad n\geq2.
\]
\end{example}

\begin{solve}
若 $a_2,\ldots,a_n$ 全非零，把第 $j$ 列的 $-c_j/a_j$ 倍加到第一列，$2\leq j\leq n$，即得上三角矩阵，所以
\begin{equation}
\label{eq:xie-ex-det-arrow-nonzero}
H=\left(a_1-\sum_{j=2}^{n}\frac{b_jc_j}{a_j}\right)\prod_{j=2}^{n}a_j.
\end{equation}
不作非零假设时，仍有
\begin{equation}
\label{eq:xie-ex-det-arrow-polynomial}
H=a_1\prod_{j=2}^{n}a_j-
\sum_{i=2}^{n}b_ic_i\prod_{\substack{2\leq j\leq n\\j\ne i}}a_j.
\end{equation}
由排列定义，非零项或者全部取对角元素，或者只交换指标 $1$ 与某个 $i$；后一种排列的符号为负。其余排列必取到某个非对角的零元素。这也直接证明了不含分母的公式，包括部分 $a_i=0$ 的情形。
\end{solve}

% 原书例 1.5，使用对角修正量 d_i。
\begin{example}
\label{ex:xie-ex-det-common-rows}
设 $B$ 的非对角元素为 $b_{ij}=x_j$，主对角元素为 $b_{ii}=x_i+d_i$。求 $\det B$。
\end{example}

\begin{solve}
写成 $B=\operatorname{diag}(d_1,\ldots,d_n)+\mathbf1(x_1,\ldots,x_n)$。按列拆分，至少两列都选公共部分的项具有成比例列，因而为零；恰在第 $i$ 列选择公共部分时，其余列的对角元素迫使该列只取第 $i$ 项。因此
\begin{equation}
\label{eq:xie-ex-det-common-rows}
\det B=\prod_j d_j+\sum_i x_i\prod_{j\ne i}d_j.
\end{equation}
各 $d_i\ne0$ 时可写成 $\prod_i d_i(1+\sum_i x_i/d_i)$；有零参数时应使用前一个多项式公式。教材主对角元素为 $x_i-a_i$ 的模板对应于 $d_i=-a_i$。
\end{solve}

% 原书例 1.7。
\begin{proposition}[添加一行一列]
\label{prop:xie-ex-det-border}
设 $u,v$ 为 $n$ 维列向量，$z$ 为一个数，则
\begin{equation}
\label{eq:xie-ex-det-border}
\det\begin{pmatrix}A&u\\v^{\mathsf T}&z\end{pmatrix}
=z\det A-\sum_{i,j=1}^{n}C_{ij}u_iv_j.
\end{equation}
此式不要求 $A$ 可逆。
\end{proposition}

\begin{proof}
先按最后一列展开，最后一个元素贡献 $z\det A$。对 $u_i$ 的余子式再按其中的最后一行展开，$v_j$ 对应的余子式为 $M_{ij}$，两次展开的符号之积为 $(-1)^{i+n+1}(-1)^{n+j}=(-1)^{i+j+1}$，故其贡献为 $-u_iv_jC_{ij}$。
\end{proof}

\section{求和法}
\label{sec:xie-ex-det-summing}

各行元素和相同时，把所有列加到一列，就产生一个常数列；各列元素和相同时则合并行。这种操作利用整体的和，不要求原矩阵已有大量零元素。

\begin{proposition}
\label{prop:xie-ex-det-row-sum}
设 $n\geq2$，$A$ 的各行元素和均为 $\lambda$，则
\begin{equation}
\label{eq:xie-ex-det-row-sum}
\det A=\lambda\det(a_{ij}-a_{1j})_{i,j=2}^{n}.
\end{equation}
特别地，$\lambda=0$ 时 $\det A=0$。
\end{proposition}

\begin{proof}
把第 $2$ 至第 $n$ 列加到第一列，第一列变成 $\lambda\mathbf1$。提出 $\lambda$，再把第一行从其余行中减去，按第一列展开即得。提出因子只用了齐次性，没有除以 $\lambda$，所以也包括 $\lambda=0$。
\end{proof}

% 原书例 1.8、1.10、1.11。
\begin{example}
\label{ex:xie-ex-det-equal-sums}
求主对角元素为零、其余元素为 $1$ 的行列式，并计算 $\det(bI_n+\mathbf1a^{\mathsf T})$，其中 $a=(a_1,\ldots,a_n)^{\mathsf T}$。
\end{example}

\begin{solve}
第一个矩阵各行和为 $n-1$，合并列、提出该因子后，用第一行消去其余行的第一列，得到 $\det(\mathbf1\mathbf1^{\mathsf T}-I_n)=(-1)^{n-1}(n-1)$。第二个矩阵各行和为 $b+\sum_i a_i$，相同操作使降阶矩阵成为 $bI_{n-1}$，故
\begin{equation}
\label{eq:xie-ex-det-scalar-rank-one}
\det(bI_n+\mathbf1a^{\mathsf T})=b^{n-1}\left(b+\sum_i a_i\right).
\end{equation}
这些多项式公式包括 $b=0$ 或行和为零的情形，$n=1$ 可直接核对。另如三次方程 $t^3+pt+q=0$ 的三个根按循环次序组成三阶矩阵，Vieta 关系给出根之和为零，故行列式也为零。
\end{solve}

% 原书例 1.12。
\begin{example}
\label{ex:xie-ex-det-cyclic-rows}
设 $n\geq2$，第一行为 $(1,2,\ldots,n)$，以后每行由上一行循环右移一位得到，求所得行列式。
\end{example}

\begin{solve}
从下向上作 $R_i\leftarrow R_i-R_{i-1}$，$i=n,\ldots,2$，第 $i$ 行在第 $i-1$ 列为 $n-1$，其余位置为 $-1$。再把前 $n-1$ 列加到最后一列，该列只有首项 $n(n+1)/2$ 非零。按该列展开，剩下的矩阵为 $nI_{n-1}-\mathbf1\mathbf1^{\mathsf T}$，其行列式由式~\eqref{eq:xie-ex-det-scalar-rank-one} 得到 $n^{n-2}$。故结果为
\[
(-1)^{n+1}\frac{n(n+1)}2\,n^{n-2}
=(-1)^{n-1}\frac{n+1}{2}\,n^{n-1}.\qedhere
\]
\end{solve}

\section{递推法与数学归纳法}
\label{sec:xie-ex-det-recurrence}

递推法把 $n$ 阶行列式与低阶的同类行列式联系起来，再由初值确定结果。数学归纳法则在知道或猜到通式之后，验证它满足递推式和初值；猜测本身不是证明。写递推关系时必须同时说明适用阶数和所需初值。

% 原书例 1.13。
\begin{proposition}[三对角行列式]
\label{prop:xie-ex-det-tridiagonal}
设 $D_k$ 为下列矩阵左上角的 $k$ 阶行列式：
\[
\begin{pmatrix}
a_1&c_1&0&\cdots&0\\
b_1&a_2&c_2&\ddots&\vdots\\
0&b_2&a_3&\ddots&0\\
\vdots&\ddots&\ddots&\ddots&c_{n-1}\\
0&\cdots&0&b_{n-1}&a_n
\end{pmatrix}.
\]
则
\begin{equation}
\label{eq:xie-ex-det-tridiagonal}
D_0=1,\quad D_1=a_1,\quad
D_n=a_nD_{n-1}-b_{n-1}c_{n-1}D_{n-2}\quad(n\geq2).
\end{equation}
\end{proposition}

\begin{proof}
按最后一列展开，$a_n$ 贡献 $a_nD_{n-1}$；$c_{n-1}$ 的代数余子式带负号，其余子式再按最后一行展开，只留下 $b_{n-1}D_{n-2}$。当 $n=2$ 时，公式为 $a_1a_2-b_1c_1$，与 $D_0=1$ 相容。
\end{proof}

% 原书例 1.14。
\begin{example}
\label{ex:xie-ex-det-constant-tridiagonal}
若主对角元素全为 $a$，上、下副对角元素分别全为 $c,b$，求 $D_n$。
\end{example}

\begin{solve}
递推式为 $D_n=aD_{n-1}-bcD_{n-2}$。设 $\alpha,\beta$ 为 $r^2-ar+bc=0$ 的两个复根。当 $\alpha\ne\beta$ 时，
\begin{equation}
\label{eq:xie-ex-det-tridiagonal-distinct}
D_n=\frac{\alpha^{n+1}-\beta^{n+1}}{\alpha-\beta}.
\end{equation}
把右端代入递推式，并核对 $D_0=1,D_1=\alpha+\beta=a$ 即可证明。当 $\alpha=\beta=a/2$，即 $a^2=4bc$ 时，应另用
\begin{equation}
\label{eq:xie-ex-det-tridiagonal-repeated}
D_n=(n+1)(a/2)^n\quad(n\geq1),\qquad D_0=1.
\end{equation}
这一式也可直接核对递推式和初值，包括 $a=0,bc=0$ 的退化情形。\emph[caution]{重根时不能继续使用以 $\alpha-\beta$ 为分母的式子。}
\end{solve}

% 原书例 1.15、1.19。
\begin{example}
\label{ex:xie-ex-det-induction}
求两类三对角行列式：第一类的主对角元素为 $\cos\theta,2\cos\theta,\ldots,2\cos\theta$，两条副对角线均为 $1$；第二类的主对角元素为 $a_0+a_1,a_1+a_2,\ldots,a_{n-1}+a_n$，第 $i$ 个上、下副对角元素均为 $a_i$。
\end{example}

\begin{solve}
第一类记为 $T_n$，则 $T_0=1,T_1=\cos\theta$，且 $T_n=2\cos\theta\,T_{n-1}-T_{n-2}$。余弦加法公式给出相同递推关系，故 $T_n=\cos(n\theta)$。证明没有除以 $\sin\theta$，因此也包括 $\theta\in\pi\mathbb Z$。

第二类记为 $E_n$，则不论各 $a_i$ 是否为零，都有
\begin{equation}
\label{eq:xie-ex-det-weighted-tridiagonal}
E_n=\sum_{i=0}^{n}\prod_{\substack{0\leq j\leq n\\j\ne i}}a_j.
\end{equation}
设右端为 $P_n$，则 $P_n=a_nP_{n-1}+\prod_{j=0}^{n-1}a_j$，且 $P_{n-1}-a_{n-1}P_{n-2}=\prod_{j=0}^{n-2}a_j$。合并得
\[
P_n=(a_{n-1}+a_n)P_{n-1}-a_{n-1}^2P_{n-2},
\]
与 $E_n$ 的递推相同，再核对 $P_0=1,P_1=a_0+a_1$ 即得。若各 $a_i\ne0$，才可改写成 $a_0\cdots a_n\sum_{i=0}^{n}1/a_i$。
\end{solve}

\section{拆分法}
\label{sec:xie-ex-det-splitting}

有效的拆分不是把全部元素逐一展开，而是把列写成对角部分与公共方向之和，使含有重复列或成比例列的项自动消失。

% 原书例 1.22 及其推论，统一为两个任意向量。
\begin{theorem}[公共方向的增量]
\label{thm:xie-ex-det-rank-one-update}
设 $u,v$ 为 $n$ 维列向量，$t$ 为参数，$C_{ij}$ 为原矩阵 $A$ 的代数余子式，则
\begin{equation}
\label{eq:xie-ex-det-rank-one-update}
\det(A+tuv^{\mathsf T})=\det A+t\sum_{i,j=1}^{n}C_{ij}u_iv_j.
\end{equation}
特别地，全体元素同时增加 $t$ 时，
\begin{equation}
\label{eq:xie-ex-det-all-increment}
\det(A+t\mathbf1\mathbf1^{\mathsf T})=\det A+t\sum_{i,j=1}^{n}C_{ij}.
\end{equation}
这些公式不要求 $A$ 可逆。
\end{theorem}

\begin{proof}
第 $j$ 列为 $A_{\cdot j}+tv_ju$，按所有列拆分。至少两列都选增量的项具有成比例列，因而为零；不选增量的项为 $\det A$，只在第 $j$ 列选增量的项按该列展开为 $tv_j\sum_i u_iC_{ij}$，求和即得。
\end{proof}

参数次数至多为一的原因在于增量列共用一个方向 $u$，不能据此认定任何元素中含 $t$ 的行列式都对 $t$ 线性。类似地，例~\ref{ex:xie-ex-det-common-rows} 的拆分给出
\begin{equation}
\label{eq:xie-ex-det-diagonal-rank-one}
\det\bigl(\operatorname{diag}(d_1,\ldots,d_n)+uv^{\mathsf T}\bigr)
=\prod_i d_i+\sum_i u_iv_i\prod_{j\ne i}d_j.
\end{equation}
矩阵 $uv^{\mathsf T}$ 的各列都是同一向量的倍数，称为秩至多为一的矩阵；本处的推导只需这一具体结构。

% 原书例 1.25。
\begin{example}
\label{ex:xie-ex-det-identity-update}
设 $u=(u_1,\ldots,u_n)^{\mathsf T}$，求 $\det(I_n+uu^{\mathsf T})$。
\end{example}

\begin{solve}
在式~\eqref{eq:xie-ex-det-diagonal-rank-one} 中令各 $d_i=1,v=u$，得 $\det(I_n+uu^{\mathsf T})=1+\sum_i u_i^2$。实向量情形下结果严格为正；若 $u$ 为复向量，此处没有取共轭，不能将 $\sum_i u_i^2$ 当成平方模之和。
\end{solve}

% 原书例 1.23。
\begin{example}
\label{ex:xie-ex-det-uniform-triangle}
设 $U_n$ 主对角元素均为 $a$，上方元素均为 $b$，下方元素均为 $c$，求 $\det U_n$。
\end{example}

\begin{solve}
令 $F(t)=\det(U_n+t\mathbf1\mathbf1^{\mathsf T})$。$F$ 的次数至多为一，且 $F(-b)=(a-b)^n$、$F(-c)=(a-c)^n$，因为这时分别得到下三角和上三角矩阵。$b\ne c$ 时，两点确定 $F$，所以
\begin{equation}
\label{eq:xie-ex-det-uniform-triangle}
\det U_n=\frac{b(a-c)^n-c(a-b)^n}{b-c}.
\end{equation}
$b=c$ 时不能由这两个重合的取值求斜率，应直接使用式~\eqref{eq:xie-ex-det-scalar-rank-one}，得到 $(a-b)^{n-1}(a+(n-1)b)$。
\end{solve}

% 原书例 1.16。
\begin{example}
\label{ex:xie-ex-det-variable-diagonal}
将上一例的主对角元素改为 $x_1,\ldots,x_n$，求所得行列式 $D_n$。
\end{example}

\begin{solve}
同样设置 $F(t)$，则 $F(-b)=\prod_i(x_i-b),F(-c)=\prod_i(x_i-c)$，故
\begin{equation}
\label{eq:xie-ex-det-variable-diagonal}
D_n=\begin{cases}
\displaystyle\frac{b\prod_i(x_i-c)-c\prod_i(x_i-b)}{b-c},&b\ne c,\\
\displaystyle\prod_i(x_i-b)+b\sum_i\prod_{j\ne i}(x_j-b),&b=c.
\end{cases}
\end{equation}
第二种情形直接由式~\eqref{eq:xie-ex-det-diagonal-rank-one} 得到，不要求 $x_i\ne b$。拆分最后一列还得到递推式
\[
D_n=(x_n-b)D_{n-1}+b\prod_{i=1}^{n-1}(x_i-c).
\]
其中选择常数列的部分，可用该列消去其余列的公共部分，只留下对角元素 $x_i-c$。这也说明拆分法与递推法可以配合使用。
\end{solve}

% 原书例 1.24。
\begin{proposition}
\label{prop:xie-ex-det-low-degree}
设 $n\geq2$，多项式 $f_1,\ldots,f_n$ 的次数均不超过 $n-2$，则对任意 $x_1,\ldots,x_n$，$\det(f_j(x_i))_{i,j=1}^{n}=0$。
\end{proposition}

\begin{proof}
每个 $f_j$ 都是 $1,t,\ldots,t^{n-2}$ 的线性组合。按列拆分后，每个行列式的 $n$ 列只能从 $n-1$ 种列向量 $(x_1^k,\ldots,x_n^k)^{\mathsf T}$ 中选取，至少两列成比例，故各项都为零。次数上界不能随意改为 $n-1$：取 $f_j(t)=t^{j-1}$ 就会得到下一节中通常非零的行列式。
\end{proof}

\section{Vandermonde 行列式}
\label{sec:xie-ex-det-vandermonde}

\begin{theorem}[Vandermonde 行列式]
\label{thm:xie-ex-det-vandermonde}
对任意 $x_1,\ldots,x_n$，\emph{范德蒙行列式}（Vandermonde determinant）满足
\begin{equation}
\label{eq:xie-ex-det-vandermonde}
V_n(x_1,\ldots,x_n)
=\begin{vmatrix}
1&x_1&x_1^2&\cdots&x_1^{n-1}\\
1&x_2&x_2^2&\cdots&x_2^{n-1}\\
\vdots&\vdots&\vdots&\ddots&\vdots\\
1&x_n&x_n^2&\cdots&x_n^{n-1}
\end{vmatrix}
=\prod_{1\leq i<j\leq n}(x_j-x_i).
\end{equation}
因此它非零当且仅当各节点两两不同。
\end{theorem}

\begin{proof}
从最后一列向前作 $K_j\leftarrow K_j-x_1K_{j-1}$，$j=n,\ldots,2$。第一行变成 $(1,0,\ldots,0)$，其余行在后 $n-1$ 列的元素为 $(x_i-x_1)x_i^{j-2}$。按第一行展开并逐行提出 $x_i-x_1$，得
\[
V_n(x_1,\ldots,x_n)=\prod_{i=2}^{n}(x_i-x_1)V_{n-1}(x_2,\ldots,x_n).
\]
由 $V_1=1$ 递推即得。这里只提取因子，没有除以节点之差，所以节点重合时公式也成立。
\end{proof}

% 原书例 1.27。
\begin{proposition}[多项式列的化简]
\label{prop:xie-ex-det-polynomial-columns}
设 $f_k(t)$ 是次数为 $k$、最高次项系数为 $c_k\ne0$ 的多项式，$0\leq k\leq n-1$，则
\begin{equation}
\label{eq:xie-ex-det-polynomial-columns}
\det(f_{j-1}(x_i))_{i,j=1}^{n}
=\left(\prod_{k=0}^{n-1}c_k\right)V_n(x_1,\ldots,x_n).
\end{equation}
\end{proposition}

\begin{proof}
从常数列开始，利用前面的列依次消去每列的低次项，再提出该列的最高次项系数，剩下的就是标准 Vandermonde 行列式。前面各列的最高次项系数非零，保证消元系数可以选取。
\end{proof}

% 原书例 1.26。
\begin{example}
\label{ex:xie-ex-det-homogeneous-vandermonde}
求 $\det(a_i^{n-j}b_i^{j-1})_{i,j=1}^{n}$。
\end{example}

\begin{solve}
先设各 $a_i\ne0$，从第 $i$ 行提出 $a_i^{n-1}$，得到以 $b_i/a_i$ 为节点的 Vandermonde 行列式，因而
\begin{equation}
\label{eq:xie-ex-det-homogeneous-vandermonde}
\det(a_i^{n-j}b_i^{j-1})_{i,j=1}^{n}
=\prod_{i<j}(a_ib_j-a_jb_i).
\end{equation}
两端都是所有参数的多项式，故连续。让任意为零的 $a_i$ 由非零数趋近于零，就证明该公式也适用于零参数；最终公式本身没有任何分母。
\end{solve}

% 原书例 1.28、1.29。
\begin{example}
\label{ex:xie-ex-det-trigonometric-vandermonde}
设 $n\geq2$，计算 $\det(\cos((j-1)\theta_i))_{i,j=1}^{n}$ 和 $\det(\sin(j\theta_i))_{i,j=1}^{n}$。
\end{example}

\begin{solve}
令 $t=\cos\theta$。由余弦递推关系，$\cos(k\theta)$ 是 $t$ 的 $k$ 次多项式；$k=0$ 的系数为 $1$，$k\geq1$ 时最高次项系数为 $2^{k-1}$。由式~\eqref{eq:xie-ex-det-polynomial-columns}，
\begin{equation}
\label{eq:xie-ex-det-cosine}
\det(\cos((j-1)\theta_i))
=2^{(n-1)(n-2)/2}\prod_{i<j}(\cos\theta_j-\cos\theta_i).
\end{equation}
类似地，由正弦递推可构造多项式 $q_{j-1}$，使 $\sin(j\theta)=\sin\theta\,q_{j-1}(\cos\theta)$，其中 $q_{j-1}$ 的次数为 $j-1$、最高次项系数为 $2^{j-1}$。逐行提出 $\sin\theta_i$ 后得到
\begin{equation}
\label{eq:xie-ex-det-sine}
\det(\sin(j\theta_i))
=2^{n(n-1)/2}\prod_i\sin\theta_i
\prod_{i<j}(\cos\theta_j-\cos\theta_i).
\end{equation}
这里用的是乘法恒等式，没有除以 $\sin\theta_i$；因此某个正弦为零时也成立。余弦节点重复时，乘积和行列式同时为零。
\end{solve}

% 原书例 1.30。
\begin{corollary}
\label{cor:xie-ex-det-polynomial-uniqueness}
次数不超过 $m$ 的多项式若在 $m+1$ 个互异节点处取零值，就一定是零多项式。因此同次数范围内的多项式由这 $m+1$ 个节点值唯一确定。
\end{corollary}

\begin{proof}
把多项式的 $m+1$ 个系数视为未知量，节点处的零值给出一个系数矩阵为 Vandermonde 矩阵的齐次方程组。节点互异使其行列式非零，由 Cramer 法则只有零解。对两个具有相同节点值的多项式之差应用此结论，就得到唯一性。
\end{proof}

\section{升阶法}
\label{sec:xie-ex-det-bordering}

升阶法添加适当的行或列，把原式嵌入一个更容易识别的行列式。关键是同时证明新旧行列式的关系；若添加一行一列后按其展开恰好恢复原式，就不会引入未知的比例因子。缺少某个幂次列、含有公共常数项或具有少数几个公共列方向时，都适合考虑升阶。

% 原书例 1.31。
\begin{example}
\label{ex:xie-ex-det-one-plus-powers}
求 $D=\det(1+x_i^j)_{i,j=1}^{n}$。
\end{example}

\begin{solve}
构造
\[
B=\begin{pmatrix}
1&-1&-1&\cdots&-1\\
1&x_1&x_1^2&\cdots&x_1^n\\
\vdots&\vdots&\vdots&\ddots&\vdots\\
1&x_n&x_n^2&\cdots&x_n^n
\end{pmatrix}.
\]
把第一行从其余行中减去，再按第一列展开，得到 $\det B=D$。将第一行拆成 $(2,0,\ldots,0)-(1,1,\ldots,1)$，第一部分贡献 $2\det(x_i^j)=2\prod_i x_iV_n(x_1,\ldots,x_n)$；第二部分为节点 $1,x_1,\ldots,x_n$ 的 Vandermonde 行列式。因此
\begin{equation}
\label{eq:xie-ex-det-one-plus-powers}
D=\left(2\prod_i x_i-\prod_i(x_i-1)\right)V_n(x_1,\ldots,x_n).
\end{equation}
公式不要求节点互异，也不排除节点为 $0$ 或 $1$。
\end{solve}

% 原书例 1.32。
\begin{example}
\label{ex:xie-ex-det-missing-power}
在幂次 $0,1,\ldots,n$ 中删去第 $k$ 次幂，按剩余幂次递增组成 $n$ 列，求节点 $x_1,\ldots,x_n$ 对应的行列式 $W_k$，其中 $0\leq k\leq n$。
\end{example}

\begin{solve}
定义初等对称多项式
\[
e_r(x_1,\ldots,x_n)=\sum_{1\leq i_1<\cdots<i_r\leq n}x_{i_1}\cdots x_{i_r},
\qquad e_0=1.
\]
在完整的 $n+1$ 列幂次下添加最后一行 $(1,y,\ldots,y^n)$，得到
\[
V_{n+1}(x_1,\ldots,x_n,y)
=V_n(x_1,\ldots,x_n)\prod_{i=1}^{n}(y-x_i).
\]
左端按最后一行展开时，$y^k$ 的系数为 $(-1)^{n+k}W_k$；右端对应系数为 $(-1)^{n-k}e_{n-k}V_n$。比较系数，两个符号相同，故
\begin{equation}
\label{eq:xie-ex-det-missing-power}
W_k=e_{n-k}(x_1,\ldots,x_n)V_n(x_1,\ldots,x_n).
\end{equation}
边界 $k=n$ 恢复标准 Vandermonde 公式，$k=0$ 对应每行提出 $x_i$。
\end{solve}

% 原书例 1.33；直接描述两次升阶后的分块消元。
\begin{example}
\label{ex:xie-ex-det-offdiagonal-sums}
设 $n\geq2$，各 $a_i\ne0$，矩阵 $G$ 满足 $g_{ii}=0$、$g_{ij}=a_i+a_j$（$i\ne j$），求 $\det G$。
\end{example}

\begin{solve}
令 $a=(a_1,\ldots,a_n)^{\mathsf T}$，$D=-2\operatorname{diag}(a_1,\ldots,a_n)$，$U=(\mathbf1,a)$，$V=(a,\mathbf1)$，则 $G=D+UV^{\mathsf T}$。考虑 $n+2$ 阶矩阵
\[
B=\begin{pmatrix}I_2&-V^{\mathsf T}\\U&D\end{pmatrix}.
\]
以最上面两行为基准消去左下块，得到右下块 $D+UV^{\mathsf T}=G$，所以 $\det B=\det G$。另一方面，各 $a_i\ne0$ 保证 $D$ 的对角元非零，可以用最后 $n$ 列消去左下块，得到左上块 $I_2+V^{\mathsf T}D^{-1}U$。这里 $D^{-1}=\operatorname{diag}(-1/(2a_i))$，只是逐个对角元素取倒数。令 $S=\sum_i a_i,T=\sum_i1/a_i$，则
\[
I_2+V^{\mathsf T}D^{-1}U
=\begin{pmatrix}1-n/2&-S/2\\-T/2&1-n/2\end{pmatrix}.
\]
对两种消元后的分块三角矩阵，应用命题~\ref{prop:xie-ex-det-basic-blocks}，行列式都是对角块行列式之积，因而
\begin{equation}
\label{eq:xie-ex-det-offdiagonal-sums}
\det G=(-2)^{n-2}\left(\prod_i a_i\right)
\left((n-2)^2-\left(\sum_i a_i\right)\left(\sum_i\frac1{a_i}\right)\right).
\end{equation}
本式保留了题设的非零条件；若需代入零参数，必须先消去分母，不能直接使用 $1/a_i$。
\end{solve}

\section{求根法}
\label{sec:xie-ex-det-polynomial-roots}

行列式的元素为多项式时，其排列展开仍是多项式。求根法先估计次数，再找出使行列式出现相同行、相关列或零行的参数值，从而找到因子，最后确定尚未知道的常数。次数为 $d$ 的非零一元多项式，只有找齐按重数计算的 $d$ 个根并确定最高次项系数后，才能恢复整个多项式；只找出若干根不能排除额外因子。若次数估计为“至多 $d$”，还要防止最高次项相消。

因式定理的依据是多项式除法：$p(r)=0$ 当且仅当 $t-r$ 整除 $p(t)$。多个不同根可以依次提出因子；重根则须另行核对，例如检查导数或继续作整除。多参数问题可先把参数看作独立未定元，证明多项式恒等式，再代入特殊值；不能在参数已经相等时仍把相同根算作多个不同根。

% 原书例 1.34。
\begin{example}
\label{ex:xie-ex-det-vandermonde-roots}
用求根法再证 Vandermonde 公式。
\end{example}

\begin{solve}
把 $x_n$ 看成主未定元，$V_n$ 的次数至多为 $n-1$，$x_n=x_i$（$i<n$）时有两行相同，故含因子 $\prod_{i<n}(x_n-x_i)$。先把其余节点视为独立参数；按最后一行展开，$x_n^{n-1}$ 的系数为 $V_{n-1}$，所以
\[
V_n=V_{n-1}\prod_{i<n}(x_n-x_i).
\]
再对 $n$ 归纳得到式~\eqref{eq:xie-ex-det-vandermonde}。所得是多项式恒等式，因此任意节点重合的代入也成立。
\end{solve}

% 原书例 1.18 的两种解法；此处采用清分母后的因子比较。
\begin{theorem}[Cauchy 行列式]
\label{thm:xie-ex-det-cauchy}
设所有 $a_i+b_j\ne0$，则\emph{柯西行列式}（Cauchy determinant）为
\begin{equation}
\label{eq:xie-ex-det-cauchy}
\det\left(\frac1{a_i+b_j}\right)_{i,j=1}^{n}
=\frac{\prod_{i<j}(a_j-a_i)(b_j-b_i)}{\prod_{i,j=1}^{n}(a_i+b_j)}.
\end{equation}
在分母均非零的前提下，它非零当且仅当 $a_i$ 两两不同且 $b_j$ 两两不同。
\end{theorem}

\begin{proof}
逐行清分母，令 $B_{ij}=\prod_{k\ne j}(a_i+b_k)$，则左端等于 $\det B/\prod_{i,j}(a_i+b_j)$。多项式 $\det B$ 的总次数至多为 $n(n-1)$；交换两个 $a_i$ 相当于交换两行，交换两个 $b_j$ 相当于交换两列。因此 $a_i=a_j$ 或 $b_i=b_j$ 时它为零，所有一次因子 $a_j-a_i,b_j-b_i$ 均整除它。这些不同的一次因子在实数或复数系数多项式环中互不相伴，且均为素因子，故它们的乘积也整除 $\det B$。乘积总次数已为 $n(n-1)$，所以
\[
\det B=K\prod_{i<j}(a_j-a_i)(b_j-b_i)
\]
且 $K$ 为常数。为确定 $K$，在\emph{已经清分母的多项式}中令 $a_i=-b_i$，并取互异的 $b_i$。此时 $B$ 为对角矩阵，其对角元为 $\prod_{k\ne i}(b_k-b_i)$；两两配对后，对角元乘积恰等于右侧的两个差积，因此 $K=1$。此代入不能用于原来的分式矩阵，但用于多项式 $B$ 完全合法。
\end{proof}

% 原书例 1.36，按同一四元对称模式排列。
\begin{example}
\label{ex:xie-ex-det-four-factors}
求
\[
Q=\begin{vmatrix}
x&y&z&w\\y&x&w&z\\z&w&x&y\\w&z&y&x
\end{vmatrix}.
\]
\end{example}

\begin{solve}
行向量分别作全加、前两行减后两行、第一与第三行减第二与第四行、第一与第四行减第二与第三行，可以看出下列任一因子为零时，四行就有一个非平凡线性关系：
\[
x+y+z+w,\quad x+y-z-w,\quad x-y+z-w,\quad x-y-z+w.
\]
故这四个不同的一次多项式因子都整除 $Q$。总次数为四，且 $x^4$ 的系数为 $1$，所以
\[
Q=(x+y+z+w)(x+y-z-w)(x-y+z-w)(x-y-z+w).
\]
这是多参数的恒等式，某些参数取特殊值导致因子重合时仍然有效。
\end{solve}

\section{组合定义}
\label{sec:xie-ex-det-combinatorics}

组合定义适合追踪哪些排列项可能非零、某个参数的最高次数来自哪里，以及整数整除或共轭运算如何传递到行列式。它通常用于证明结构性结论，而不是直接计算全部 $n!$ 项。

% 原书例 1.39。
\begin{example}
\label{ex:xie-ex-det-sparse}
证明：若 $A$ 的零元素超过 $n^2-n$ 个，则 $\det A=0$。另求只在副对角线上可能非零的行列式。
\end{example}

\begin{solve}
第一个条件说明非零元素少于 $n$ 个。任一排列项必须从不同位置选择 $n$ 个元素，其中至少一个为零，因此所有项均为零；也可由抽屉原理得到至少有一整行为零。这个充分条件不能反过来使用，例如全一矩阵有 $n^2$ 个非零元素，$n\geq2$ 时行列式仍为零。

只在副对角线上可能非零时，唯一可能保留的排列为 $(n,n-1,\ldots,1)$，其逆序数为 $n(n-1)/2$，所以
\begin{equation}
\label{eq:xie-ex-det-antidiagonal}
\det A=(-1)^{n(n-1)/2}\prod_{i=1}^{n}a_{i,n+1-i}.
\end{equation}
其中某个副对角元素为零时，两端也同时为零。
\end{solve}

% 原书例 1.40。
\begin{proposition}
\label{prop:xie-ex-det-leading-characteristic}
对任意 $n$ 阶矩阵 $A$，$p_A(t)=\det(tI_n-A)$ 是最高次项系数为 $1$ 的 $n$ 次多项式，且 $t^{n-1}$ 的系数为 $-\sum_i a_{ii}$。
\end{proposition}

\begin{proof}
恒等排列贡献 $\prod_i(t-a_{ii})$，其最高两项为 $t^n-(\sum_i a_{ii})t^{n-1}$。非恒等排列至少移动两个指标，而参数 $t$ 只出现在主对角上，所以其余排列项的次数至多为 $n-2$，不影响最高两项。
\end{proof}

% 原书例 1.41、1.43。
\begin{proposition}
\label{prop:xie-ex-det-conjugate-skew}
对复矩阵，$\det\overline A=\overline{\det A}$，其中横线表示逐元素复共轭。若 $A^{\mathsf T}=-A$，称 $A$ 为\emph{反对称矩阵}（skew-symmetric matrix）；实数域或复数域上奇数阶反对称矩阵的行列式为零。
\end{proposition}

\begin{proof}
共轭保持有限和与乘积，排列系数 $\pm1$ 也不变，对式~\eqref{eq:xie-ex-det-leibniz} 逐项取共轭即得第一式。反对称情形由
\[
\det A=\det A^{\mathsf T}=\det(-A)=(-1)^n\det A
\]
得出 $n$ 为奇数时 $2\det A=0$，因而 $\det A=0$。这里的反对称使用转置，不能换成共轭转置后仍直接照用结论。
\end{proof}

% 原书例 1.42，以及例 1.20 和基础训练中的奇偶性。
\begin{example}
\label{ex:xie-ex-det-integrality}
设 $n\geq2$，整数矩阵 $A$ 的行列式 $D\ne0$，且 $D$ 整除每个元素，证明 $D=\pm1$。再说明元素均为 $1$ 或 $-1$ 的行列式具有怎样的整除性。
\end{example}

\begin{solve}
每个排列项是 $n$ 个元素之积，均被 $D^n$ 整除，所以 $D^n$ 整除它们的带符号之和 $D$。因 $D\ne0$，有 $D^{n-1}$ 整除 $1$，故整数 $D$ 只能为 $\pm1$。$n\geq2$ 不能删去：一阶矩阵 $(2)$ 就不满足结论。

若全体元素为 $\pm1$，将第一行从其余每行中减去，后 $n-1$ 行的元素都为 $0$ 或 $\pm2$，所以
\begin{equation}
\label{eq:xie-ex-det-sign-divisibility}
2^{n-1}\mid\det A.
\end{equation}
这比仅说行列式为偶数更精确。教材还给出界 $|\det A|\leq\frac23n!$（$n\geq3$）：$n=3$ 时先对各行、列乘 $-1$，可使第一列全为 $1$、第一行其余元素全为 $-1$，绝对值不变；将第一列加到其余两列后，按第一行展开，剩余二阶矩阵的元素均为 $0$ 或 $2$，故绝对值至多为 $4=\frac23\,3!$。若结论对 $n-1$ 成立，按第一行展开并用三角不等式，则绝对值至多为 $n\cdot\frac23(n-1)!=\frac23n!$，完成归纳。
\end{solve}

% 原书例 1.21 的组合定义证明。
\begin{theorem}[行列式求导]
\label{thm:xie-ex-det-differentiation}
设每个 $f_{ij}(t)$ 在区间 $J$ 上可微，$F(t)=\det(f_{ij}(t))$。令 $F_j(t)$ 表示只将第 $j$ 列逐元素求导、其余列保持不变所得的行列式，则
\begin{equation}
\label{eq:xie-ex-det-differentiation}
F'(t)=\sum_{j=1}^{n}F_j(t)
=\sum_{i,j=1}^{n}C_{ij}(t)f'_{ij}(t).
\end{equation}
也可逐行求导后求和，但不能把所有元素同时求导所得的一个行列式当成 $F'(t)$。
\end{theorem}

\begin{proof}
排列定义是有限和，每项是 $n$ 个可微函数之积。逐项使用乘积求导法则，再按被求导元素所在的列分组，就得到 $\sum_jF_j$。每个 $F_j$ 按第 $j$ 列展开后，其余子式都来自原矩阵，得到第二个表达式。整个过程只涉及有限和，不需要额外的极限交换条件。
\end{proof}

\section{Laplace 定理}
\label{sec:xie-ex-det-laplace}

\begin{definition}
\label{def:xie-ex-det-general-minor}
记 $[n]=\{1,\ldots,n\}$。设 $I=\{i_1<\cdots<i_k\}$、$J=\{j_1<\cdots<j_k\}$，$A[I,J]$ 表示保留这些行列、并按原次序排列的 $k$ 阶子矩阵；$\det A[I,J]$ 为相应的 $k$ 阶子式。记 $I^{\mathrm c}=[n]\setminus I$，$\Sigma I=\sum_{i\in I}i$。该子式的代数余子式定义为
\[
(-1)^{\Sigma I+\Sigma J}\det A[I^{\mathrm c},J^{\mathrm c}].
\]
当 $I=J$ 时，称 $\det A[I,I]$ 为\emph{主子式}（principal minor）。
\end{definition}

\begin{theorem}[Laplace 展开]
\label{thm:xie-ex-det-laplace}
固定 $k$ 行组成的指标集 $I$，则
\begin{equation}
\label{eq:xie-ex-det-laplace}
\det A=\sum_{\substack{J\subseteq[n]\\|J|=k}}
(-1)^{\Sigma I+\Sigma J}\det A[I,J]\det A[I^{\mathrm c},J^{\mathrm c}].
\end{equation}
固定列集、对行集求和也有同样公式。$k=1$ 时退化为一行或一列展开。
\end{theorem}

\begin{proof}
按所选 $k$ 行中使用的列集 $J$，把排列定义中的项分组。保持内部次序，将 $I,J$ 分别移到前 $k$ 个位置，所需相邻对换次数为
\[
\Sigma I-\frac{k(k+1)}2+\Sigma J-\frac{k(k+1)}2,
\]
与 $\Sigma I+\Sigma J$ 同奇偶。固定 $J$ 后，选中部分与补集部分的排列独立求和，分别给出 $\det A[I,J]$ 和 $\det A[I^{\mathrm c},J^{\mathrm c}]$，再乘这一符号并对 $J$ 求和即可。
\end{proof}

\begin{corollary}[分块三角与分块反对角]
\label{cor:xie-ex-det-blocks}
设 $P,Q$ 分别为 $p,q$ 阶方阵，$X,Y$ 为相应大小的任意矩阵，则
\begin{equation}
\label{eq:xie-ex-det-block-triangular}
\det\begin{pmatrix}P&X\\0&Q\end{pmatrix}
=\det\begin{pmatrix}P&0\\Y&Q\end{pmatrix}
=\det P\det Q,
\end{equation}
并且
\begin{equation}
\label{eq:xie-ex-det-block-antidiagonal}
\det\begin{pmatrix}0_{p\times q}&P\\Q&0_{q\times p}\end{pmatrix}
=(-1)^{pq}\det P\det Q.
\end{equation}
\end{corollary}

\begin{proof}
对上三角分块矩阵按最后 $q$ 行作 Laplace 展开，只有选中最后 $q$ 列的子式可能非零，得到对角块乘积；下三角情形同理。分块反对角式把最前 $q$ 列整体移到最后，需要 $pq$ 次相邻列对换，得到分块对角矩阵，因此增加符号 $(-1)^{pq}$。例如 $p=q=5$、$\det P=2,\det Q=3$ 时，分块反对角行列式为 $-6$，不能漏掉块交换的符号。
\end{proof}

% 原书例 1.44。
\begin{example}
\label{ex:xie-ex-det-plucker}
证明恒等式
\[
(ab'-a'b)(cd'-c'd)-(ac'-a'c)(bd'-b'd)
+(ad'-a'd)(bc'-b'c)=0.
\]
\end{example}

\begin{solve}
考虑四阶矩阵，其四行依次为 $(a,a',a,a')$、$(b,b',b,b')$、$(c,c',c,c')$、$(d,d',d,d')$。它有重复列，行列式为零。按前两列作 Laplace 展开，行集 $\{1,2\},\{1,3\},\{1,4\}$ 分别给出括号中的正、负、正三项；互补行集又各给出同样的一项。因此展开结果为所证左端的 $2$ 倍，除以 $2$ 即得。
\end{solve}

% 原书例 1.46。
\begin{proposition}[两个矩阵之和的行列式]
\label{prop:xie-ex-det-sum-minors}
若 $A,B$ 均为 $n$ 阶矩阵，则
\begin{equation}
\label{eq:xie-ex-det-sum-minors}
\det(A+B)=\sum_{k=0}^{n}
\sum_{\substack{I,J\subseteq[n]\\|I|=|J|=k}}
(-1)^{\Sigma I+\Sigma J}\det A[I,J]\det B[I^{\mathrm c},J^{\mathrm c}].
\end{equation}
$k=0,n$ 分别贡献 $\det B,\det A$，其余项是不能省略的混合项。
\end{proposition}

\begin{proof}
将每一行拆为来自 $A$ 与来自 $B$ 的两部分。固定从 $A$ 取行的指标集 $I$，对这一混合矩阵按 $I$ 作 Laplace 展开，得到对应的内层和。对所有 $I$ 求和，就把行的全部 $2^n$ 种选择计入一次。
\end{proof}

% 原书例 1.47。
\begin{corollary}[特征多项式的系数]
\label{cor:xie-ex-det-characteristic-coefficients}
设 $p_A(t)=\det(tI_n-A)=t^n+c_1t^{n-1}+\cdots+c_n$，则
\begin{equation}
\label{eq:xie-ex-det-characteristic-coefficients}
c_k=(-1)^k\sum_{\substack{I\subseteq[n]\\|I|=k}}\det A[I,I],
\qquad 1\leq k\leq n.
\end{equation}
特别地，$c_1=-\sum_i a_{ii}$，$c_n=(-1)^n\det A$。
\end{corollary}

\begin{proof}
在式~\eqref{eq:xie-ex-det-sum-minors} 中对 $-A+tI_n$ 展开。单位矩阵的补子式只有行、列指标相同才可能非零，这时等于 $t^{n-k}$，外部符号为 $(-1)^{2\Sigma I}=1$；$-A$ 的 $k$ 阶子式则贡献 $(-1)^k$，比较系数即得。
\end{proof}

\section{综合运用}
\label{sec:xie-ex-det-synthesis}

选择方法时先识别结构：零多就直接展开；相邻行列差简单就先作差；行列和相同就合并；主对角加公共方向就拆分；幂次或递增次数的多项式列就化为 Vandermonde；缺少幂次列就升阶；低次数参数式就找因子并比较系数；三对角等带状结构则优先找递推。方法可以组合，但每一步都应能说明行列式是否改变、乘了什么因子，以及是否引入非零条件。

% 原书例 1.9。
\begin{example}
\label{ex:xie-ex-det-row-complements}
令 $s_i=\sum_j a_{ij}$，$b_{ij}=s_i-a_{ij}$，证明 $\det B=(-1)^{n-1}(n-1)\det A$。
\end{example}

\begin{solve}
设 $A$ 的列为 $\alpha_1,\ldots,\alpha_n$，$s=\sum_j\alpha_j$，则 $B$ 的列为 $s-\alpha_j$。按列拆分，选取至少两个 $s$ 的项为零。不选 $s$ 的项为 $(-1)^n\det A$；只在第 $j$ 列选 $s$ 的项为 $(-1)^{n-1}\det A$，因为将 $s=\sum_k\alpha_k$ 代入该列后仅有 $k=j$ 不出现重复列。这样的项共有 $n$ 个，求和即得。$n=1$ 时两端也同为零。
\end{solve}

% 原书例 1.48。
\begin{example}
\label{ex:xie-ex-det-squared-diagonal}
设 $A(x)$ 的第 $j$ 列除第 $j$ 行外均为 $a_j^2$，而第 $(j,j)$ 个元素为 $(x-a_j)^2$，求 $\det A(x)$。
\end{example}

\begin{solve}
由 $(x-a_j)^2=a_j^2+(x^2-2a_jx)$，令 $d_j=x^2-2a_jx$，可写成 $A(x)=\operatorname{diag}(d_1,\ldots,d_n)+\mathbf1(a_1^2,\ldots,a_n^2)$。式~\eqref{eq:xie-ex-det-diagonal-rank-one} 给出
\[
\det A(x)=\prod_j(x^2-2a_jx)
+\sum_i a_i^2\prod_{j\ne i}(x^2-2a_jx).
\]
这个表达式在 $x=0$、$x=2a_i$ 或 $a_i=0$ 时均有效；若先改写为含 $1/d_i$ 的形式，就必须另外处理这些取值。
\end{solve}

% 原书例 1.49。
\begin{example}
\label{ex:xie-ex-det-cofactor-sum}
设 $A=(\alpha_1,\ldots,\alpha_n)$，证明
\[
\det(\alpha_1-\alpha_2,\alpha_2-\alpha_3,\ldots,
\alpha_{n-1}-\alpha_n,\mathbf1)=\sum_{i,j=1}^{n}C_{ij}.
\]
\end{example}

\begin{solve}
令 $F(t)=\det(A+t\mathbf1\mathbf1^{\mathsf T})$。从 $j=1$ 到 $j=n-1$ 依次作 $K_j\leftarrow K_j-K_{j+1}$，此时右侧的列尚未改写，所以
\[
F(t)=\det(\alpha_1-\alpha_2,\ldots,\alpha_{n-1}-\alpha_n,
\alpha_n+t\mathbf1).
\]
由最后一列的线性，$F(1)-F(0)$ 就是题目左端；由式~\eqref{eq:xie-ex-det-all-increment}，这一差又等于全部代数余子式之和。
\end{solve}

% 原书例 1.50。
\begin{theorem}[行列式的刻画]
\label{thm:xie-ex-det-characterization}
设实数域或复数域上的函数 $f$ 把每个 $n$ 阶方阵对应到一个数，并且对每一列分别线性、交换任意两列时变号，且 $f(I_n)=1$，则 $f(A)=\det A$。
\end{theorem}

\begin{proof}
设 $e_1,\ldots,e_n$ 为标准单位列向量，将第 $j$ 列写成 $\alpha_j=\sum_i a_{ij}e_i$，按各列的线性展开得
\[
f(A)=\sum_{k_1,\ldots,k_n=1}^{n}
a_{k_1,1}\cdots a_{k_n,n}f(e_{k_1},\ldots,e_{k_n}).
\]
若有两个 $k_j$ 相同，交换对应两列使函数值变号而输入不变，故函数值为零。其余项的 $(k_1,\ldots,k_n)$ 是排列，通过相邻交换恢复 $(1,\ldots,n)$，得到函数值 $(-1)^{N(k_1,\ldots,k_n)}f(I_n)$。于是剩下的求和正是按列选元素的排列定义，等于 $\det A$。
\end{proof}

\section{基础训练与方法核查}
\label{sec:xie-ex-det-exercises}

本节从教材的选择、填空与解答题中提取有代表性的训练，重点说明理由，不逐项抄录选择项。同类小题应回到行列式性质核对，而不是仅记住某个数值答案。例如，第一列移到末列要作 $n-1$ 次相邻交换，行列式乘以 $(-1)^{n-1}$；同时颠倒行序和列序的总符号为 $(-1)^{n(n-1)}=1$；将 $a_{ij}$ 换为 $(-1)^{i+j}a_{ij}$，等价于逐行、逐列提出符号，总符号也为 $1$。若换为 $r^{i-j}a_{ij}$ 且 $r\ne0$，行因子与列因子相消，行列式不变；若换为 $b_ib_ja_{ij}$，行列式则乘以 $\prod_i b_i^2$。

\subsection{余子式、操作次序与循环结构}

% 原书解答题 7，以及代数余子式求和的填空题。
\begin{example}
\label{ex:xie-ex-det-triangular-cofactors}
设 $A$ 为上三角矩阵，证明 $i<j$ 时 $M_{ij}=C_{ij}=0$，并说明如何计算某一行的代数余子式之和。
\end{example}

\begin{solve}
删除第 $i$ 行、第 $j$ 列且 $i<j$ 后，原来的前 $j-1$ 列在剩余矩阵中仍有 $j-1$ 列。上三角性使这些列的非零元素只能落在原前 $j-1$ 行；但其中第 $i$ 行已被删去，只剩 $j-2$ 行。每个排列项要从这 $j-1$ 列选取不同的行，必然选到零，因此 $M_{ij}=0$，从而 $C_{ij}=0$。不能将条件改成 $i>j$：例如二阶上三角矩阵的 $C_{21}=-a_{12}$ 未必为零。

若求第 $k$ 行的 $\sum_j C_{kj}$，把该行替换为全一行，再沿该行展开即可；代数余子式因删除了这一行而保持原值。这与将原来各行相加或求所有元素之和是不同操作。
\end{solve}

% 原书填空题 15，强调同时使用原列。
\begin{example}
\label{ex:xie-ex-det-cyclic-column-addition}
设 $n\geq2$，$A=(\alpha_1,\ldots,\alpha_n)$。同时令 $\beta_1=\alpha_1+\alpha_n$、$\beta_j=\alpha_j+\alpha_{j-1}$（$2\leq j\leq n$），求 $\det(\beta_1,\ldots,\beta_n)$。
\end{example}

\begin{solve}
按列拆分，每列选本列或前一原列。若某列选前一原列，而前一列仍选本列，就出现重复列；要避免重复，这种选择必须沿圆周传播，直到所有列都选前一原列。因此只有全部取本列或全部取前一原列的两项可能非零，得到
\[
\det(\beta_1,\ldots,\beta_n)
=\bigl(1+(-1)^{n-1}\bigr)\det A
=\begin{cases}2\det A,&n\text{ 为奇数},\\0,&n\text{ 为偶数}.\end{cases}
\]
题意是同时使用原列，不能逐列就地更新后再引用已改变的列；否则计算的会是另一矩阵。
\end{solve}

% 原书解答题 10、12。
\begin{example}
\label{ex:xie-ex-det-cyclic-bidiagonal}
设 $n\geq2$，$B$ 主对角元素均为 $a$，上副对角元素及左下角元素均为 $b$，其余元素为零。另设 $S$ 主对角元素为 $x$，上方元素为 $h$，下方元素为 $-h$。分别求其行列式。
\end{example}

\begin{solve}
$B$ 的非零排列项只有恒等排列和沿上副对角及左下角绕行一周的循环排列：只要一行选择上副对角元素，就会迫使下一行继续选择上副对角元素，直到绕回起点。因此
\[
\det B=a^n+(-1)^{n-1}b^n.
\]
对 $S$，在式~\eqref{eq:xie-ex-det-uniform-triangle} 中取 $a=x,b=h,c=-h$，$h\ne0$ 时得
\[
\det S=\frac{(x-h)^n+(x+h)^n}{2}.
\]
两端为多项式，或直接检查 $h=0$ 的对角情形，就知这一式也包括 $h=0$。
\end{solve}

\subsection{递推与结构矩阵}

% 原书解答题 8。
\begin{example}
\label{ex:xie-ex-det-continued-fraction}
令 $K(a_1,\ldots,a_n)$ 为主对角元素依次为 $a_1,\ldots,a_n$、上副对角均为 $1$、下副对角均为 $-1$ 的三对角行列式，约定 $K(\varnothing)=1$。将有限连分数表示为行列式之比。
\end{example}

\begin{solve}
按第一行展开，$n\geq2$ 时有
\[
K(a_1,\ldots,a_n)
=a_1K(a_2,\ldots,a_n)+K(a_3,\ldots,a_n),
\]
其中 $n=2$ 时最后一项为 $K(\varnothing)=1$；一阶值为 $K(a_1)=a_1$。若每个尾部 $K(a_k,\ldots,a_n)\ne0$（$2\leq k\leq n$），则各层分母都有定义，归纳得到
\begin{equation}
\label{eq:xie-ex-det-continued-fraction}
a_1+\cfrac1{a_2+\cfrac1{a_3+\cfrac1{\ddots+\cfrac1{a_n}}}}
=\frac{K(a_1,\ldots,a_n)}{K(a_2,\ldots,a_n)}.
\end{equation}
归纳步骤就是将递推式除以 $K(a_2,\ldots,a_n)$。这些尾部非零条件不能省去，否则原来的逐层除法可能没有意义。若所有 $a_i=1$，则 $K_n=K_{n-1}+K_{n-2}$，$K_0=K_1=1$，与 Fibonacci 数列相联系。
\end{solve}

% 原书解答题 13、14。
\begin{example}
\label{ex:xie-ex-det-max-distance}
分别求 $M=(\max\{i,j\})_{i,j=1}^{n}$ 和 $L=(|i-j|)_{i,j=1}^{n}$ 的行列式，第二个问题设 $n\geq2$。
\end{example}

\begin{solve}
对 $M$，从 $i=1$ 到 $i=n-1$ 依次作 $R_i\leftarrow R_i-R_{i+1}$。第 $i$ 行前 $i$ 个元素为 $-1$，其余元素为零，最后一行全为 $n$，所以成为下三角矩阵，得到
\begin{equation}
\label{eq:xie-ex-det-max}
\det M=(-1)^{n-1}n.
\end{equation}
对 $L$，先从下向上作 $R_i\leftarrow R_i-R_{i-1}$，$i=n,\ldots,2$；再从右向左作 $K_j\leftarrow K_j-K_{j-1}$，$j=n,\ldots,2$。直接由 $|i-j|$ 的差分可得新矩阵
\[
\begin{pmatrix}
0&1&1&\cdots&1\\
1&-2&0&\cdots&0\\
1&0&-2&\cdots&0\\
\vdots&\vdots&\vdots&\ddots&\vdots\\
1&0&0&\cdots&-2
\end{pmatrix}.
\]
由爪型公式~\eqref{eq:xie-ex-det-arrow-nonzero}，
\begin{equation}
\label{eq:xie-ex-det-distance}
\det L=(-2)^{n-1}\frac{n-1}{2}
=(-1)^{n-1}(n-1)2^{n-2}.
\end{equation}
这两题中的差分方向不同，均是为了确保被减去的行或列仍为原来的对象。
\end{solve}

\subsection{Vandermonde 方程组与升阶训练}

% 原书解答题 3、4；唯一解之外补出显式解。
\begin{example}
\label{ex:xie-ex-det-vandermonde-system}
设 $a_1,\ldots,a_n$ 两两不同，证明方程组
\[
\sum_{i=1}^{n}a_i^kx_i=b^k,\qquad k=0,1,\ldots,n-1,
\]
有唯一解，并求出解。再求以 $t,a_1,\ldots,a_{n-1}$ 为节点的 Vandermonde 行列式等于零时的全部 $t$。
\end{example}

\begin{solve}
系数矩阵是 Vandermonde 矩阵的转置，行列式非零，所以解唯一。定义
\[
\ell_i(t)=\prod_{j\ne i}\frac{t-a_j}{a_i-a_j}.
\]
它满足 $\ell_i(a_j)=\delta_{ij}$。对任意次数不超过 $n-1$ 的多项式 $q$，差 $q(t)-\sum_iq(a_i)\ell_i(t)$ 在 $n$ 个不同节点为零，由推论~\ref{cor:xie-ex-det-polynomial-uniqueness} 可知恒为零。取 $q(t)=t^k$ 并代入 $t=b$，就得到
\[
x_i=\ell_i(b)=\prod_{j\ne i}\frac{b-a_j}{a_i-a_j}.
\]
若 $b$ 恰等于某一节点，也可直接代入，得到该位置为 $1$、其余为零的解。第二问由 Vandermonde 差积得到全部根 $t=a_1,\ldots,a_{n-1}$，没有其他根。
\end{solve}

% 原书解答题 5。
\begin{example}
\label{ex:xie-ex-det-neighboring-powers}
设 $n\geq2$，矩阵第 $i$ 行依次为
\[
1,\quad x_i+1,\quad x_i^2+x_i,\quad\ldots,\quad x_i^{n-1}+x_i^{n-2}.
\]
求其行列式。
\end{example}

\begin{solve}
第 $j$ 列是 $x_i$ 的 $j-1$ 次首一多项式在各节点的值，直接由式~\eqref{eq:xie-ex-det-polynomial-columns} 得到 $V_n(x_1,\ldots,x_n)$。若逐列消去低次项，应从第二列向右进行：每列减去已经化成纯幂的前一列。这里恰好需要使用更新后的前列，和“所有列同时减去原前列”的操作不同。
\end{solve}

% 原书解答题 15。
\begin{example}
\label{ex:xie-ex-det-shifted-missing-power}
设 $n\geq2$，求第 $i$ 行为
\[
1,\quad x_i(x_i-a),\quad x_i^2(x_i-a),\quad\ldots,\quad x_i^{n-1}(x_i-a)
\]
的 $n$ 阶行列式 $D$。
\end{example}

\begin{solve}
添加缺少的一次多项式列，并添加最后一行对应于节点 $y$ 的值，得到列多项式为 $1,t-a,t(t-a),\ldots,t^{n-1}(t-a)$ 的 $n+1$ 阶行列式 $B(y)$。这些列多项式依次为 $0$ 至 $n$ 次的首一多项式，故
\[
B(y)=V_n(x_1,\ldots,x_n)\prod_i(y-x_i).
\]
按最后一行展开，分别代入 $y=a$ 和 $y=0$ 并相减：第一列的贡献相消，第二列贡献 $a(-1)^{n+3}D$，其余列在两个取值下均为零。因此
\[
aD=V_n(x_1,\ldots,x_n)\left(\prod_i x_i-\prod_i(x_i-a)\right).
\]
当 $a\ne0$ 时得到
\begin{equation}
\label{eq:xie-ex-det-shifted-missing-power}
D=\frac{\prod_i x_i-\prod_i(x_i-a)}{a}V_n(x_1,\ldots,x_n).
\end{equation}
当 $a=0$ 时，幂次恰好为 $0,2,\ldots,n$，由式~\eqref{eq:xie-ex-det-missing-power} 得到 $D=e_{n-1}(x_1,\ldots,x_n)V_n(x_1,\ldots,x_n)$。这也等于前式消去可约因子 $a$ 后的取值。
\end{solve}

本章的模板都应连同条件一起使用：展开式中的余子式来自哪一个原矩阵，行列变换引用的是原行列还是更新后的行列，提取因子是否被误当成除法，以及含参数公式是否覆盖零值和重根，都会直接影响结论。把结果代回一阶或二阶、检查对角或零参数情形，可以发现许多符号错误；这些核对帮助检验推导，但不能代替对一般阶数的证明。

\chapter{矩阵}
\label{chap:matrices}

% 教材来源：高等代数/textbooks/高等代数习题（谢启鸿）/高代习题-2.pdf。
% 覆盖原书第 2 章 §§2.1--2.12，印刷页 54--115；按主题整合，不逐题抄录。
矩阵把线性关系组织成可以运算的对象。本章围绕两类问题展开：怎样利用矩阵结构简化乘法、乘方和求逆，怎样把行列式或矩阵恒等式转化为更容易处理的问题。前半章建立基本运算、特殊矩阵与初等变换的工具，后半章通过伴随矩阵、迹、矩阵分解、Cauchy--Binet 公式和分块消元组织解题方法，最后说明如何将可逆情形的恒等式推广到一般矩阵。

除另有说明外，矩阵元素取自数域 $\mathbb F\subseteq\mathbb C$，$M_{m,n}(\mathbb F)$ 表示所有 $m\times n$ 矩阵，$M_n(\mathbb F)=M_{n,n}(\mathbb F)$，$I_n$ 表示 $n$ 阶单位矩阵，$0$ 的尺寸由上下文确定。记 $A^{\mathsf T}$ 为转置，$\overline A$ 为逐元素共轭，$A^{\mathsf H}=\overline A^{\mathsf T}$ 为共轭转置。本章沿用教材记号 $A^*$ 表示\emph{伴随矩阵}，不把它用于共轭转置；涉及平方和的非负性时，另行注明实数或复数条件。

\section{基本概念}
\label{sec:matrix-basic}

\subsection{矩阵及其运算}

\begin{definition}
\label{def:matrix-operations}
由 $mn$ 个数 $a_{ij}$ 按 $m$ 行、$n$ 列排成的阵列 $A=(a_{ij})$ 称为\emph{矩阵（matrix）}，其中 $1\le i\le m$、$1\le j\le n$。两个矩阵相等，当且仅当它们的尺寸相同且对应元素相等。同型矩阵的加法与数乘逐元素定义；若 $A\in M_{m,k}(\mathbb F)$、$B\in M_{k,n}(\mathbb F)$，则乘积 $AB\in M_{m,n}(\mathbb F)$ 定义为
\begin{equation}
 (AB)_{ij}=\sum_{\ell=1}^{k}a_{i\ell}b_{\ell j}.
 \label{eq:matrix-product-entry}
\end{equation}
\end{definition}

矩阵加法和数乘满足通常的交换、结合与分配律；矩阵乘法满足结合律和左右分配律，但一般不满足交换律。使用结合律前仍须检查各步尺寸，例如 $A\in M_{m,k}$、$B\in M_{k,\ell}$、$C\in M_{\ell,n}$ 时，$(AB)C=A(BC)\in M_{m,n}$。方阵的非负整数次幂规定为 $A^0=I_n$、$A^{r+1}=A^rA$，因此 $A^rA^s=A^{r+s}$。只有在 $AB=BA$ 时，才能直接使用通常的二项式展开 $(A+B)^k=\sum_{j=0}^k\binom{k}{j}A^{k-j}B^j$。

\begin{example}
\label{ex:matrix-noncommutative}
说明非零矩阵不能任意约去，且非零矩阵的乘积可能为零。
\end{example}
\begin{solve}
取二阶矩阵 $E_{ij}$，其第 $(i,j)$ 个元素为 $1$，其余为 $0$。有 $E_{12}E_{21}=E_{11}$ 而 $E_{21}E_{12}=E_{22}$，故乘法不交换；又有 $E_{11}E_{22}=0$，两因子却都非零。于是从 $AB=AC$ 只能得到 $A(B-C)=0$，若没有可逆性等附加条件，不能推出 $B=C$。即使两个加数均可逆，其和也可能不可逆，例如 $I_n+(-I_n)=0$。
\end{solve}

\begin{proposition}
\label{prop:matrix-transpose}
在各式运算有定义时，
\[
 (A+B)^{\mathsf T}=A^{\mathsf T}+B^{\mathsf T},\qquad
 (cA)^{\mathsf T}=cA^{\mathsf T},\qquad
 (AB)^{\mathsf T}=B^{\mathsf T}A^{\mathsf T},
\]
\[
 \overline{AB}=\overline A\,\overline B,\qquad
 \overline{cA}=\overline c\,\overline A,\qquad
 (AB)^{\mathsf H}=B^{\mathsf H}A^{\mathsf H}.
\]
\end{proposition}
\begin{proof}
加法和数乘的结论逐元素成立。对乘法，由式~\eqref{eq:matrix-product-entry}，$(AB)^{\mathsf T}$ 的第 $(i,j)$ 个元素为 $\sum_\ell a_{j\ell}b_{\ell i}$，恰是 $B^{\mathsf T}A^{\mathsf T}$ 的对应元素。共轭保持有限和与乘积，再与转置公式结合即得其余结论。
\end{proof}

\subsection{逆矩阵与伴随矩阵}

\begin{definition}
\label{def:matrix-inverse-adjugate}
若 $A\in M_n(\mathbb F)$ 存在同阶矩阵 $B$ 满足 $AB=BA=I_n$，则称 $A$ \emph{可逆（invertible）}或非奇异，$B$ 称为逆矩阵，记为 $A^{-1}$；否则称 $A$ 奇异。记 $C_{ij}=(-1)^{i+j}\det A[\widehat i,\widehat j]$ 为 $a_{ij}$ 的代数余子式，其中方括号表示删除第 $i$ 行与第 $j$ 列后所得矩阵，则
\[
 A^*=(C_{ji})_{1\le i,j\le n}
\]
称为 $A$ 的\emph{伴随矩阵（adjugate matrix）}。规定零阶行列式为 $1$，因而一阶矩阵 $(a)$ 的伴随为 $(1)$。
\end{definition}

\begin{theorem}
\label{thm:matrix-invertibility}
对任意 $n$ 阶矩阵 $A$，有
\begin{equation}
 AA^*=A^*A=(\det A)I_n.
 \label{eq:matrix-adjugate-identity}
\end{equation}
矩阵 $A$ 可逆当且仅当 $\det A\ne0$，此时 $A^{-1}=(\det A)^{-1}A^*$。对同阶方阵，若 $AB=I_n$ 或 $BA=I_n$ 中任一等式成立，则另一等式也成立。
\end{theorem}
\begin{proof}
$AA^*$ 的第 $(i,j)$ 个元素为 $\sum_k a_{ik}C_{jk}$：当 $i=j$ 时，这是 $\det A$ 按第 $i$ 行的展开；当 $i\ne j$ 时，这是将第 $j$ 行替换为第 $i$ 行所得行列式按第 $j$ 行的展开，因两行相同而为零。$A^*A$ 同理。若 $\det A\ne0$，式~\eqref{eq:matrix-adjugate-identity} 直接给出逆；若 $A$ 可逆，则由方阵乘积的行列式公式 $\det(AA^{-1})=\det A\det A^{-1}=1$ 得 $\det A\ne0$，该乘积公式亦由定理~\ref{thm:cauchy-binet} 推出。最后，$AB=I_n$ 蕴含 $\det A\det B=1$，故 $A$ 可逆且 $B=A^{-1}$；另一情形同理。
\end{proof}

逆矩阵唯一：若 $B,C$ 都是 $A$ 的逆，则 $B=B(AC)=(BA)C=C$。对可逆矩阵和非零标量 $c$，直接相乘验证可得
\begin{equation}
 (AB)^{-1}=B^{-1}A^{-1},\qquad
 (cA)^{-1}=c^{-1}A^{-1},\qquad
 (A^{\mathsf T})^{-1}=(A^{-1})^{\mathsf T}.
 \label{eq:matrix-inverse-rules}
\end{equation}
因此同阶可逆矩阵之积仍可逆；若一个方阵因子奇异，则乘积奇异。可逆性还能恢复消去律，例如 $A$ 可逆时，$AB=AC$ 蕴含 $B=C$。

\subsection{初等变换、初等矩阵与相抵}
\label{subsec:matrix-elementary-basic}

\begin{definition}
交换两行、以非零数乘某行、将某行的倍数加到另一行，分别称为三类\emph{初等行变换（elementary row operations）}；对应地定义初等列变换。对 $I_n$ 作一次初等变换得到\emph{初等矩阵（elementary matrix）}。同型矩阵 $A,B$ 若可经有限次初等行、列变换相互转化，则称它们\emph{相抵（equivalent）}，亦称等价。
\end{definition}

若 $i\ne j$，记 $P_{ij}$ 为交换 $I_n$ 第 $i,j$ 行所得矩阵，$D_i(c)=I_n+(c-1)E_{ii}$（$c\ne0$），$T_{ij}(c)=I_n+cE_{ij}$。左乘 $T_{ij}(c)$ 的含义是将第 $j$ 行的 $c$ 倍加到第 $i$ 行；右乘同一矩阵则将第 $i$ 列的 $c$ 倍加到第 $j$ 列。三类初等矩阵的行列式分别为 $-1,c,1$，其逆分别为 $P_{ij},D_i(c^{-1}),T_{ij}(-c)$。所以初等行变换是左乘适当初等矩阵，初等列变换是右乘适当初等矩阵，均保持矩阵的可逆性。

\begin{theorem}
\label{thm:matrix-equivalence-normal}
任意 $A\in M_{m,n}(\mathbb F)$ 均存在可逆矩阵 $P\in M_m(\mathbb F)$、$Q\in M_n(\mathbb F)$，使
\begin{equation}
 PAQ=\begin{pmatrix}I_r&0\\0&0\end{pmatrix}_{m\times n}.
 \label{eq:matrix-equivalence-normal}
\end{equation}
这里 $r$ 唯一，等于 $A$ 的最大非零子式阶数，称为 $A$ 的\emph{秩（rank）}；零矩阵的秩为 $0$。同型矩阵相抵当且仅当秩相同。
\end{theorem}
\begin{proof}
若矩阵非零，先将一个非零元素移到左上角并化为 $1$，再用行、列变换消去第一列与第一行其余元素，对剩余子块重复，最终得到式~\eqref{eq:matrix-equivalence-normal}。可逆行、列变换不改变所有 $s$ 阶子式是否全零：以行变换为例，新子式可按行列式的多线性展开为旧 $s$ 阶子式的线性组合，逆变换给出反向蕴含，列变换同理。标准型的最大非零子式阶数正好为 $r$，故 $r$ 唯一；两个秩相同的矩阵均可化为同一标准型，结论成立。
\end{proof}

对方阵，上述标准型给出可逆、秩为 $n$、相抵于 $I_n$、能表示为有限个初等矩阵之积之间的等价关系。把 $Ax=0$ 化为标准型对应的方程，还可知 $A$ 可逆当且仅当齐次方程只有零解，这也等价于其列向量线性无关；对转置应用同一结论得到行向量版本。若 $A$ 可逆，对增广矩阵 $(A\mid I_n)$ 只作行变换得到 $(I_n\mid A^{-1})$：所有行变换的乘积 $P$ 满足 $PA=I_n$，故右块恰为 $P=A^{-1}$。

\subsection{分块矩阵}

把相邻若干行、列分别组成一组，就把矩阵写成\emph{分块矩阵（block matrix）}。若 $A=(A_{ij})$、$B=(B_{jk})$ 的分块尺寸相容，则 $(AB)_{ik}=\sum_jA_{ij}B_{jk}$；这只是式~\eqref{eq:matrix-product-entry} 按索引分组。分块运算必须保留因子的先后次序，不能把矩阵块当作可交换的数。特别地，分块对角阵满足
\[
 \operatorname{diag}(A_1,\ldots,A_s)\operatorname{diag}(B_1,\ldots,B_s)
 =\operatorname{diag}(A_1B_1,\ldots,A_sB_s),
\]
其幂、逆逐块计算，逆存在当且仅当每一对角块可逆，行列式为各对角块行列式的乘积。分块行变换中的乘子在左侧，分块列变换中的乘子在右侧；详细的消元与降阶方法见第~\ref{sec:matrix-block-elimination} 节。

\subsection{Cauchy--Binet 公式}

\begin{theorem}
\label{thm:cauchy-binet}
设 $A\in M_{m,n}(\mathbb F)$、$B\in M_{n,m}(\mathbb F)$。若 $m>n$，则 $\det(AB)=0$；若 $m\le n$，则
\begin{equation}
 \det(AB)=\sum_{\substack{J\subseteq\{1,\ldots,n\}\\|J|=m}}
 \det A[:,J]\det B[J,:].
 \label{eq:cauchy-binet}
\end{equation}
这里 $J$ 中索引按递增次序排列，$A[:,J]$ 表示保留这些列，$B[J,:]$ 表示保留这些行。更一般地，$A\in M_{m,n}$、$B\in M_{n,p}$ 时，对大小均为 $r$ 的行、列索引集 $I,K$，
\begin{equation}
 \det\bigl((AB)[I,K]\bigr)
 =\sum_{\substack{J\subseteq\{1,\ldots,n\}\\|J|=r}}
 \det A[I,J]\det B[J,K],
 \label{eq:cauchy-binet-minor}
\end{equation}
$r>n$ 时右侧为空和，值为 $0$。这称为\emph{柯西--比内公式（Cauchy--Binet formula）}。
\end{theorem}
\begin{proof}
记 $A$ 的列为 $a_1,\ldots,a_n$，则 $AB$ 的第 $j$ 列为 $\sum_{k=1}^n b_{kj}a_k$。按列多线性展开，$\det(AB)$ 是所有有序索引 $(k_1,\ldots,k_m)$ 对应项 $\det(a_{k_1},\ldots,a_{k_m})\prod_{j=1}^m b_{k_jj}$ 之和。索引重复时行列式为零；若 $m>n$，每一项都如此。否则将互异索引按集合 $J=\{j_1<\cdots<j_m\}$ 分组，在组内将 $A$ 的列恢复递增顺序，置换符号与 $B$ 的系数乘积之和恰为 $\det B[J,:]$，得到式~\eqref{eq:cauchy-binet}。最后对 $A[I,:]B[:,K]$ 应用同一结论即得子式公式。
\end{proof}

当 $m=n$ 时求和只有一项，故 $\det(AB)=\det A\det B$；矩形矩阵则不能直接套用这个方阵公式。一般的子式公式还说明 $\operatorname{rank}(AB)\le\min\{\operatorname{rank}A,\operatorname{rank}B\}$，因为超过任一因子秩的阶数时，每个乘积项都为零。

\section{特殊矩阵}
\label{sec:matrix-special}

\subsection{标准单位向量与基础矩阵}

记 $e_i^{(n)}$ 为第 $i$ 个分量为 $1$、其余分量为 $0$ 的 $n$ 维标准单位列向量；尺寸明确时简写为 $e_i$。若 $A\in M_{m,n}$，则 $Ae_j^{(n)}$ 为第 $j$ 列，$(e_i^{(m)})^{\mathsf T}A$ 为第 $i$ 行，并有 $(e_i^{(m)})^{\mathsf T}Ae_j^{(n)}=a_{ij}$。所以证明两个同型矩阵相等，可以验证它们对全部标准单位向量的作用相同。

\begin{proposition}
\label{prop:matrix-units}
$n$ 阶\emph{基础矩阵（matrix units）} $E_{ij}=e_ie_j^{\mathsf T}$ 满足
\begin{equation}
 E_{ij}E_{k\ell}=\delta_{jk}E_{i\ell},\qquad
 A=\sum_{i,j=1}^n a_{ij}E_{ij},\qquad
 E_{ij}AE_{k\ell}=a_{jk}E_{i\ell},
 \label{eq:matrix-unit-rules}
\end{equation}
其中 $\delta_{jk}$ 在 $j=k$ 时为 $1$，否则为 $0$。
\end{proposition}
\begin{proof}
利用 $e_j^{\mathsf T}e_k=\delta_{jk}$ 及 $e_j^{\mathsf T}Ae_k=a_{jk}$ 即得第一、三式，第二式逐元素验证即可。特别地，$E_{ij}A$ 仅第 $i$ 行可能非零，该行为 $A$ 的第 $j$ 行；$AE_{ij}$ 仅第 $j$ 列可能非零，该列为 $A$ 的第 $i$ 列。
\end{proof}

\emph[caution]{基础矩阵与初等矩阵是不同的概念。}前者用于提取、组合元素，后者表示可逆的初等变换；例如 $n\ge2$ 时 $E_{11}$ 奇异，而初等矩阵总可逆。

\subsection{循环移位、幂零块与多项式友阵}

对 $n\ge2$，定义基础循环矩阵 $J$ 与幂零 Jordan 块 $N$ 为
\[
 J=\begin{pmatrix}0&1&&0\\ &0&\ddots&\\ &&\ddots&1\\1&&&0\end{pmatrix},\qquad
 N=\begin{pmatrix}0&1&&0\\ &0&\ddots&\\ &&\ddots&1\\0&&&0\end{pmatrix}.
\]
空白位置均为零；$n=1$ 时另定 $J=(1)$、$N=(0)$。$Je_1=e_n$、$Je_j=e_{j-1}$（$j\ge2$），故 $J^n=I_n$；$Ne_1=0$、$Ne_j=e_{j-1}$（$j\ge2$），故当 $1\le k<n$ 时，$N^k$ 恰在第 $k$ 条上副对角线上为 $1$、其余元素为零，且 $N^n=0$。这些结论从单位向量的作用即可得到，不必逐次进行完整矩阵乘法。

\begin{definition}
若存在正整数 $s$ 使 $A^s=0$，则称 $A$ 为\emph{幂零矩阵（nilpotent matrix）}。对首一多项式 $f(t)=t^n+a_1t^{n-1}+\cdots+a_n$，矩阵
\[
 C(f)=\begin{pmatrix}
 0&0&\cdots&0&-a_n\\
 1&0&\cdots&0&-a_{n-1}\\
 0&1&\cdots&0&-a_{n-2}\\
 \vdots&\vdots&\ddots&\vdots&\vdots\\
 0&0&\cdots&1&-a_1
 \end{pmatrix}
\]
称为它的\emph{友阵（companion matrix）}；其转置称为相应的 Frobenius 块。
\end{definition}

\begin{proposition}
\label{prop:matrix-companion}
友阵满足 $\det(tI_n-C(f))=f(t)$，因而 $C(f)$ 可逆当且仅当 $a_n\ne0$。
\end{proposition}
\begin{proof}
$tI_n-C(f)$ 的前 $n-1$ 列只有主对角元 $t$ 与其下方元素 $-1$ 可能非零。沿末列展开，第 $i<n$ 行的元素 $a_{n-i+1}$ 所对应代数余子式为 $t^{i-1}$：删除该行与末列后，前段贡献 $t^{i-1}$，后段的 $-1$ 因子与余子式符号抵消。末行贡献 $(t+a_1)t^{n-1}$，故展开结果为 $a_n+a_{n-1}t+\cdots+a_2t^{n-2}+(t+a_1)t^{n-1}=f(t)$。取 $t=0$ 得 $\det C(f)=(-1)^na_n$；$n=1$ 的结论直接成立。
\end{proof}

\subsection{对称、反对称与三角结构}

\begin{definition}
$A^{\mathsf T}=A$ 时称 $A$ 为\emph{对称矩阵（symmetric matrix）}，$A^{\mathsf T}=-A$ 时称其为\emph{反对称矩阵（skew-symmetric matrix）}。主对角线下方元素全为零的方阵称上三角矩阵；若其主对角元也全为零，则称严格上三角矩阵。下三角矩阵作对称定义。
\end{definition}

\begin{proposition}
\label{prop:matrix-symmetric-test}
任意方阵均唯一分解为
\begin{equation}
 A=\frac{A+A^{\mathsf T}}2+\frac{A-A^{\mathsf T}}2,
 \label{eq:matrix-symmetric-split}
\end{equation}
其中前项对称、后项反对称。又有 $x^{\mathsf T}Ax=0$ 对所有 $x\in\mathbb F^n$ 成立，当且仅当 $A$ 反对称；特别地，对称矩阵若满足这一条件，则必为零矩阵。
\end{proposition}
\begin{proof}
两部分的转置性质直接验证；若一个矩阵同时对称与反对称，则 $2A=0$，故为零，这给出分解的唯一性。反对称时，标量 $x^{\mathsf T}Ax$ 等于其转置 $x^{\mathsf T}A^{\mathsf T}x=-x^{\mathsf T}Ax$，因此为零。反之，取 $x=e_i$ 得 $a_{ii}=0$，再取 $x=e_i+e_j$ 得 $a_{ij}+a_{ji}=0$，于是 $A^{\mathsf T}=-A$。
\end{proof}

\begin{proposition}
\label{prop:matrix-triangular}
同阶上三角矩阵对加减、数乘、乘法及多项式运算封闭；其行列式为主对角元之积，可逆当且仅当各主对角元非零，此时逆仍为上三角矩阵，主对角元逐项取倒数。严格上三角 $n$ 阶矩阵 $U$ 满足 $U^n=0$。下三角情形相同。
\end{proposition}
\begin{proof}
若 $A,B$ 上三角，则 $(AB)_{ij}=\sum_k a_{ik}b_{kj}$ 中非零项必须满足 $i\le k\le j$，故 $i>j$ 时为零，且 $(AB)_{ii}=a_{ii}b_{ii}$。行列式的置换展开只有主对角项可能非零。可逆时逐列回代解 $Ax=e_j$，得到 $x_i=0$（$i>j$）和 $x_j=a_{jj}^{-1}$，故逆具有所述结构。严格上三角情形下，$(U^k)_{ij}$ 的非零项必须对应链 $i<i_1<\cdots<i_{k-1}<j$；当 $k=n$ 时不存在这样的索引链，故 $U^n=0$。
\end{proof}

对上三角矩阵，伴随矩阵也上三角，且第 $i$ 个主对角元为 $\prod_{j\ne i}a_{jj}$。可逆时这由 $A^*=(\det A)A^{-1}$ 得到；奇异时对 $A+tI_n$ 使用第~\ref{sec:matrix-perturbation} 节的摄动论证，再令 $t\to0$。相同结论可用于具有方形对角块的分块三角矩阵，但分块行列式的进一步化简仍须检查交换条件。

\section{矩阵运算中的结构方法}
\label{sec:matrix-calculation}

\begin{example}
\label{ex:matrix-positive-products}
设 $A\in M_{m,n}(\mathbb R)$、$B\in M_{n,p}(\mathbb R)$ 的元素均非负，且 $AB$ 的第 $i$ 行全为零。证明：或者 $A$ 的第 $i$ 行全为零，或者 $B$ 至少有一行全为零。
\end{example}
\begin{solve}
若 $A$ 的第 $i$ 行不全为零，选 $k$ 使 $a_{ik}>0$。对每个 $j$，$0=(AB)_{ij}=\sum_\ell a_{i\ell}b_{\ell j}$ 是非负数之和，故 $a_{ik}b_{kj}=0$，从而 $b_{kj}=0$。这对所有 $j$ 成立，因此 $B$ 的第 $k$ 行全为零。这里的非负条件排除了正负项相消，不能省略。
\end{solve}

\begin{proposition}
\label{prop:matrix-positive-zero}
实矩阵 $A$ 满足 $AA^{\mathsf T}=0$ 当且仅当 $A=0$；复矩阵 $A$ 满足 $AA^{\mathsf H}=0$ 当且仅当 $A=0$。
\end{proposition}
\begin{proof}
$AA^{\mathsf T}$ 的第 $i$ 个主对角元为 $\sum_j a_{ij}^2$，实数平方和为零迫使每项为零。复数情形用 $\sum_j|a_{ij}|^2$。若对复矩阵误用转置，结论可能失败，例如 $A=(1,\mathrm i)$ 非零却有 $AA^{\mathsf T}=0$。
\end{proof}

\begin{example}
\label{ex:matrix-center}
证明与所有 $n$ 阶矩阵乘法可交换的矩阵恰为纯量矩阵 $cI_n$。
\end{example}
\begin{solve}
设 $AE_{ii}=E_{ii}A$ 对每个 $i$ 成立，比较第 $i$ 行、第 $i$ 列可知所有非对角元都为零，故 $A=\operatorname{diag}(d_1,\ldots,d_n)$。再由 $AE_{ij}=E_{ij}A$ 得 $d_iE_{ij}=d_jE_{ij}$，因此所有 $d_i$ 相等。反之纯量矩阵与所有同阶矩阵可交换。
\end{solve}

\begin{example}
\label{ex:matrix-powers}
计算 $aI_n+N$ 与 $uv^{\mathsf T}$ 的正整数次幂，其中 $N^s=0$，$u,v\in\mathbb F^n$。
\end{example}
\begin{solve}
由于 $aI_n$ 与 $N$ 可交换，
\begin{equation}
 (aI_n+N)^k=\sum_{j=0}^{\min\{k,s-1\}}\binom{k}{j}a^{k-j}N^j.
 \label{eq:matrix-nilpotent-power}
\end{equation}
这里零次幂按多项式约定取 $1$，因此 $a=0$ 时也成立。对外积 $R=uv^{\mathsf T}$，中间的 $v^{\mathsf T}u$ 是标量，结合律给出 $R^2=(v^{\mathsf T}u)R$，进而
\begin{equation}
 R^k=(v^{\mathsf T}u)^{k-1}R\qquad(k\ge1).
 \label{eq:matrix-rank-one-power}
\end{equation}
例如 $u=(1,2,3)^{\mathsf T}$、$v=(1,2,4)^{\mathsf T}$ 时 $v^{\mathsf T}u=17$，故 $R^k=17^{k-1}R$。又若二阶矩阵 $A$ 满足某个正整数 $k$ 对应的 $A^k=0$，则 $\det A=0$；非零时 $A$ 的秩为 $1$，可写成 $uv^{\mathsf T}$，再由式~\eqref{eq:matrix-rank-one-power} 得 $v^{\mathsf T}u=0$，故总有 $A^2=0$。
\end{solve}

\begin{proposition}
\label{prop:matrix-circulant-polynomial}
同阶\emph{循环矩阵（circulant matrix）}
\[
 A=\begin{pmatrix}
 a_1&a_2&\cdots&a_n\\
 a_n&a_1&\cdots&a_{n-1}\\
 \vdots&\vdots&\ddots&\vdots\\
 a_2&a_3&\cdots&a_1
 \end{pmatrix}
\]
恰是基础循环矩阵 $J$ 的多项式 $A=a_1I_n+a_2J+\cdots+a_nJ^{n-1}$。所以同阶循环矩阵的积仍为循环矩阵，且彼此可交换。
\end{proposition}
\begin{proof}
由 $J$ 的移位作用逐行核对即可得到多项式表示。两个多项式相乘后，利用 $J^n=I_n$ 把幂次约化到 $0,\ldots,n-1$，因此乘积仍有同样表示；同一矩阵的两个多项式相互可交换。
\end{proof}

\begin{example}
\label{ex:matrix-row-sums}
设 $A\in M_n(\mathbb F)$ 每行元素之和均为 $c$，证明 $A^k$ 每行元素之和为 $c^k$；若 $A$ 可逆，则 $c\ne0$，且 $A^{-1}$ 每行元素之和为 $c^{-1}$。
\end{example}
\begin{solve}
令 $\mathbf1=(1,\ldots,1)^{\mathsf T}$。行和条件等价于 $A\mathbf1=c\mathbf1$，反复相乘得 $A^k\mathbf1=c^k\mathbf1$。若 $A$ 可逆而 $c=0$，则 $A\mathbf1=0$ 与 $\mathbf1\ne0$ 矛盾；$c\ne0$ 时原等式左乘 $c^{-1}A^{-1}$ 得 $A^{-1}\mathbf1=c^{-1}\mathbf1$。列和问题可对转置使用同一方法。
\end{solve}

\section{可逆矩阵的判定与求逆}
\label{sec:matrix-inverse-methods}

判定可逆性时，先观察已知条件能否给出非零行列式、非零核向量或一个直接的逆因子，再决定是否展开元素。Vandermonde 矩阵的行列式为各节点差之积，故节点两两不同时可逆；对满足多项式关系的矩阵，重组关系往往比计算行列式更直接。

\begin{example}
\label{ex:matrix-invertibility-tests}
分别证明以下结论：非零实矩阵若 $A^*=A^{\mathsf T}$，则 $A$ 可逆；实反对称矩阵 $K$ 的 $I_n-K$ 可逆；奇数阶实正交矩阵 $Q$ 若 $\det Q>0$，则 $I_n-Q$ 奇异。
\end{example}
\begin{solve}
第一问中有 $C_{ij}=a_{ij}$，选取非零行按该行展开，得到 $\det A=\sum_j a_{ij}^2>0$。第二问若 $(I_n-K)x=0$ 有非零实解，则 $Kx=x$，于是 $x^{\mathsf T}x=x^{\mathsf T}Kx=0$，与实平方和的正性矛盾，故齐次方程只有零解。第三问由 $QQ^{\mathsf T}=I_n$ 得 $\det Q=1$；又
\[
 \det(I_n-Q)=\det((I_n-Q)Q^{\mathsf T})
 =\det(Q^{\mathsf T}-I_n)=(-1)^n\det(I_n-Q),
\]
$n$ 为奇数，故该行列式为零。这里正交矩阵指满足 $QQ^{\mathsf T}=I_n$ 的实方阵。
\end{solve}

\begin{example}
\label{ex:matrix-polynomial-inverse}
设 $A^2-3A+2I_n=0$，求 $A^{-1}$ 与 $(A+I_n)^{-1}$，并证明 $A\ne I_n$ 时 $A-2I_n$ 奇异。
\end{example}
\begin{solve}
由 $A(3I_n-A)=2I_n$ 与 $(A+I_n)(4I_n-A)=6I_n$ 得
\[
 A^{-1}=\frac{3I_n-A}{2},\qquad
 (A+I_n)^{-1}=\frac{4I_n-A}{6}.
\]
又 $(A-I_n)(A-2I_n)=0$；若后一因子可逆，右乘其逆便得 $A=I_n$，与条件矛盾。一般地，若 $p(A)=0$ 且 $p(0)\ne0$，把常数项移到另一侧即可凑出 $A$ 的逆多项式。
\end{solve}

\begin{example}
若同阶方阵满足 $A+B=AB$，证明 $I_n-A$ 可逆且 $AB=BA$。
\end{example}
\begin{solve}
由 $(I_n-A)(I_n-B)=I_n$，定理~\ref{thm:matrix-invertibility} 保证 $(I_n-A)^{-1}=I_n-B$，因此反向乘积也为 $I_n$。展开反向乘积得 $BA=A+B=AB$。单侧逆推出双侧逆的结论依赖于这里两者都是同阶方阵。
\end{solve}

\begin{theorem}
\label{thm:matrix-inverse-transfer}
设 $U\in M_{m,n}(\mathbb F)$、$V\in M_{n,m}(\mathbb F)$。则 $I_m+UV$ 可逆当且仅当 $I_n+VU$ 可逆，此时
\begin{equation}
 (I_n+VU)^{-1}=I_n-V(I_m+UV)^{-1}U.
 \label{eq:matrix-inverse-transfer}
\end{equation}
\end{theorem}
\begin{proof}
设 $M=I_m+UV$ 可逆，由 $U(I_n+VU)=MU$ 得
\[
 \bigl(I_n-VM^{-1}U\bigr)(I_n+VU)
 =I_n+VU-VM^{-1}MU=I_n.
\]
于是右侧矩阵可逆且逆由式~\eqref{eq:matrix-inverse-transfer} 给出。交换 $U,V$ 的角色即得反向蕴含。
\end{proof}

\begin{theorem}
\label{thm:matrix-woodbury}
设 $A\in M_n(\mathbb F)$、$C\in M_m(\mathbb F)$ 可逆，$B\in M_{n,m}(\mathbb F)$、$D\in M_{m,n}(\mathbb F)$，且 $S=C^{-1}+DA^{-1}B$ 可逆，则
\begin{equation}
 (A+BCD)^{-1}
 =A^{-1}-A^{-1}BS^{-1}DA^{-1}.
 \label{eq:matrix-woodbury}
\end{equation}
这称为 Sherman--Morrison--Woodbury 公式。特别地，若 $u,v\in\mathbb F^n$ 且 $1+v^{\mathsf T}A^{-1}u\ne0$，则有 Sherman--Morrison 公式
\begin{equation}
 (A+uv^{\mathsf T})^{-1}
 =A^{-1}-\frac{A^{-1}uv^{\mathsf T}A^{-1}}{1+v^{\mathsf T}A^{-1}u}.
 \label{eq:matrix-sherman-morrison}
\end{equation}
\end{theorem}
\begin{proof}
令 $U=A^{-1}B$、$V=CD$，则 $I_m+VU=CS$ 可逆。由定理~\ref{thm:matrix-inverse-transfer}，
\[
 (I_n+UV)^{-1}
 =I_n-A^{-1}B(CS)^{-1}CD
 =I_n-A^{-1}BS^{-1}D.
\]
再由 $A+BCD=A(I_n+UV)$ 得第一式。取 $m=1$、$C=(1)$、$B=u$、$D=v^{\mathsf T}$ 得第二式。公式中的分母必须非零，矩阵因子的次序也必须保持。
\end{proof}

\begin{example}
\label{ex:matrix-diagonal-rank-one-inverse}
设 $D=\operatorname{diag}(d_1,\ldots,d_n)$、$d_i\ne0$，求 $D+\mathbf1\mathbf1^{\mathsf T}$ 的可逆条件和逆，并处理主对角元为零、其余元素全为 $1$ 的矩阵。
\end{example}
\begin{solve}
记 $s=1+\sum_{i=1}^n d_i^{-1}$。由定理~\ref{thm:matrix-inverse-transfer}，$D+\mathbf1\mathbf1^{\mathsf T}$ 可逆当且仅当 $s\ne0$；此时
\begin{equation}
 (D+\mathbf1\mathbf1^{\mathsf T})^{-1}
 =D^{-1}-\frac{D^{-1}\mathbf1\mathbf1^{\mathsf T}D^{-1}}s,
 \qquad (A^{-1})_{ij}=\frac{\delta_{ij}}{d_i}-\frac1{sd_id_j}.
 \label{eq:matrix-diagonal-rank-one-inverse}
\end{equation}
这里第二式记 $A=D+\mathbf1\mathbf1^{\mathsf T}$。若 $s=0$，则非零向量 $D^{-1}\mathbf1$ 满足 $A(D^{-1}\mathbf1)=0$，也可直接看出奇异。令 $D=-I_n$，当 $n\ge2$ 时得到
\[
 (-I_n+\mathbf1\mathbf1^{\mathsf T})^{-1}
 =-I_n+\frac1{n-1}\mathbf1\mathbf1^{\mathsf T}.
\]
$n=1$ 时原矩阵为零，不能使用这个分式。
\end{solve}

若 $A,B$ 均可逆，则
\[
 A+B=A(A^{-1}+B^{-1})B.
\]
所以 $A+B$ 与 $A^{-1}+B^{-1}$ 同时可逆，且在可逆时
\[
 (A^{-1}+B^{-1})^{-1}=B(A+B)^{-1}A=A(A+B)^{-1}B.
\]
后一种写法可由 $A(A+B)^{-1}B=A-A(A+B)^{-1}A$ 验证；这并不意味着各因子可以任意交换。

\section{初等变换及其应用}
\label{sec:matrix-elementary-applications}

\begin{proposition}
\label{prop:matrix-singular-witness}
对 $A\in M_n(\mathbb F)$，以下条件等价：$A$ 奇异；存在非零列向量 $x$ 使 $Ax=0$；存在非零同阶矩阵 $B$ 使 $AB=0$。
\end{proposition}
\begin{proof}
可逆时由左乘 $A^{-1}$ 可排除后两种情形。奇异时，由式~\eqref{eq:matrix-equivalence-normal} 取 $PAQ=\operatorname{diag}(I_r,0)$、$r<n$，令 $x=Qe_n\ne0$，便有 $PAx=0$，从而 $Ax=0$。再令 $B=xe_1^{\mathsf T}\ne0$ 即得 $AB=0$。若已知非零 $B$ 满足 $AB=0$，选取它的一个非零列也给出相应的 $x$。
\end{proof}

\begin{proposition}
\label{prop:matrix-transvection-reduction}
任意可逆 $n$ 阶矩阵仅用第三类初等行、列变换便可化为 $\operatorname{diag}(1,\ldots,1,\det A)$。任意 $n$ 阶矩阵均可表示为有限个形如 $I_n+cE_{ij}$ 的矩阵之积，此处允许 $i=j$，也允许因子奇异。
\end{proposition}
\begin{proof}
可逆矩阵若左上角为零，可把某个使该位置非零的行加到第一行；然后利用非零主元，仅通过行、列倍加消去第一列、第一行其余元素。剩余子块仍可逆，递归可将矩阵化为非零对角阵。对非零 $a,b$，变换
\[
 \begin{pmatrix}a&0\\0&b\end{pmatrix}
 \longrightarrow\begin{pmatrix}a&0\\1-a&b\end{pmatrix}
 \longrightarrow\begin{pmatrix}1&b\\1-a&b\end{pmatrix}
 \longrightarrow\begin{pmatrix}1&0\\1-a&ab\end{pmatrix}
 \longrightarrow\begin{pmatrix}1&0\\0&ab\end{pmatrix}
\]
依次对应 $R_2\gets R_2+(a^{-1}-1)R_1$、$R_1\gets R_1+R_2$、$C_2\gets C_2-bC_1$、$R_2\gets R_2-(1-a)R_1$，全是第三类变换。逐对合并对角元即得第一结论，$n=1$ 时无需变换。

一般矩阵的相抵标准型为 $D_r=\operatorname{diag}(I_r,0)$，且 $D_r=\prod_{i=r+1}^n(I_n-E_{ii})$。第三类初等矩阵已具有所需形式，第二类为 $I_n+(c-1)E_{ii}$；第一类初等矩阵由第一结论可表示为第三类矩阵与 $\operatorname{diag}(1,\ldots,1,-1)$ 的积。将 $A=P^{-1}D_rQ^{-1}$ 中各部分分解即可。奇异矩阵不能表示为可逆初等矩阵之积，故这里允许奇异因子是必要的。
\end{proof}

\begin{example}
\label{ex:matrix-geometric-triangular-inverse}
设 $N$ 为 $n$ 阶幂零 Jordan 块，求上三角矩阵
\[
 T=I_n+aN+a^2N^2+\cdots+a^{n-1}N^{n-1}
\]
的逆。
\end{example}
\begin{solve}
把有限等比和作为整体，相乘得 $(I_n-aN)T=T(I_n-aN)=I_n-a^nN^n=I_n$，故 $T^{-1}=I_n-aN$。从初等消元看，这相当于按适当次序用相邻行的倍数消去各条上副对角线；从结构看，只须识别有限几何级数。这一方法也说明任何幂零矩阵 $H$ 若 $H^s=0$，则 $(I_n-H)^{-1}=I_n+H+\cdots+H^{s-1}$。
\end{solve}

实际使用增广矩阵法时，每一步必须作用于整个增广矩阵，不能只消去左块而忘记同步更新右块。对有参数的矩阵，要记录主元非零的条件；若计算中出现新的分母，应回到原矩阵检查该条件是否确为可逆条件，例如例~\ref{ex:matrix-diagonal-rank-one-inverse} 中除 $d_i\ne0$ 外还必须有 $s\ne0$。

\section{伴随矩阵}
\label{sec:matrix-adjugate}

伴随矩阵由子式定义，对奇异矩阵仍有意义；公式 $A^*=(\det A)A^{-1}$ 则只适用于可逆矩阵。因此讨论一般恒等式时，应优先使用子式或相抵标准型，或者先证明可逆情形，再核实能否用摄动延拓。

\begin{theorem}
\label{thm:matrix-adjugate-rules}
对同阶方阵 $A,B$，有
\begin{equation}
 (AB)^*=B^*A^*,\qquad
 (A^{\mathsf T})^*=(A^*)^{\mathsf T}.
 \label{eq:matrix-adjugate-rules}
\end{equation}
当 $n\ge2$ 时，$(cA)^*=c^{n-1}A^*$ 对所有标量 $c$ 成立，且
\begin{equation}
 \det A^*=(\det A)^{n-1}.
 \label{eq:matrix-adjugate-determinant}
\end{equation}
当 $A$ 可逆时，$A^*$ 也可逆，$(A^*)^{-1}=(A^{-1})^*$。当 $n=2$ 时 $(A^*)^*=A$；当 $n\ge3$ 时
\begin{equation}
 (A^*)^*=(\det A)^{n-2}A.
 \label{eq:matrix-double-adjugate}
\end{equation}
\end{theorem}
\begin{proof}
转置与数乘公式直接由余子式得到。对 $n\ge2$，将式~\eqref{eq:cauchy-binet-minor} 应用于 $AB$ 的 $(n-1)$ 阶子式，所得求和索引集均为删除一个索引 $k$ 的集合，故
\[
 \bigl((AB)^*\bigr)_{ij}
 =(-1)^{i+j}\sum_{k=1}^n
 \det A[\widehat j,\widehat k]\det B[\widehat k,\widehat i]
 =\sum_{k=1}^n (B^*)_{ik}(A^*)_{kj}.
\]
$n=1$ 时两边均为 $(1)$。对可逆矩阵，将 $A^*=(\det A)A^{-1}$ 代入即可得到行列式公式及二次伴随公式；又由乘积伴随公式，$A^*(A^{-1})^*=(A^{-1}A)^*=I_n$，得到求逆公式。二阶二次伴随公式可直接逐元素验证。

为处理奇异情形，取 $PAQ=D_r$ 为相抵标准型。乘积伴随公式给出 $D_r^*=Q^*A^*P^*$，其中 $P^*,Q^*$ 可逆。若 $r\le n-2$，则 $D_r^*=0$，故 $A^*=0$；若 $r=n-1$，则 $D_r^*$ 仅末个主对角元为 $1$，故 $A^*$ 的秩为 $1$。因此 $n\ge2$ 时奇异 $A$ 的 $A^*$ 也奇异，得到式~\eqref{eq:matrix-adjugate-determinant}；$n\ge3$ 时上述两种情况下 $A^*$ 的秩均不超过 $n-2$，故 $(A^*)^*=0$，恰等于式~\eqref{eq:matrix-double-adjugate} 右侧。
\end{proof}

由证明同时得到 $n\ge2$ 时的重要判别：$\operatorname{rank}A=n$、$n-1$、不超过 $n-2$ 时，$\operatorname{rank}A^*$ 分别为 $n$、$1$、$0$。一阶情形需单独记住 $A^*=(1)$，不能照搬奇异性判别。若正整数 $m$ 满足 $A^m=I_n$，则 $(A^*)^m=(A^m)^*=I_n$，这也表明矩阵的有限阶关系可通过伴随乘积公式传递。

\begin{proposition}
\label{prop:matrix-block-adjugate}
设 $A\in M_m(\mathbb F)$、$B\in M_n(\mathbb F)$，则
\begin{equation}
 \begin{pmatrix}A&0\\0&B\end{pmatrix}^*
 =\begin{pmatrix}(\det B)A^*&0\\0&(\det A)B^*\end{pmatrix}.
 \label{eq:matrix-block-adjugate}
\end{equation}
\end{proposition}
\begin{proof}
左上块中的代数余子式等于 $A$ 中相应代数余子式乘 $\det B$，右下块同理。对非对角块的代数余子式，删除的行与列位于不同的对角块，剩下的行、列数失衡：例如保留了 $B$ 所在的 $n$ 行，却只有 $n-1$ 个相应列可供非零匹配，故该余子式为零。这给出全部块，而无须假设 $A,B$ 可逆。
\end{proof}

\begin{example}
\label{ex:matrix-cofactor-sum}
设 $A=(a_{ij})$ 的代数余子式为 $C_{ij}$，说明如何把 $\sum_{i,j}C_{ij}$ 转化为矩阵运算或行列式。
\end{example}
\begin{solve}
由于伴随是代数余子式矩阵的转置，$\sum_{i,j}C_{ij}=\mathbf1^{\mathsf T}A^*\mathbf1$；若 $A$ 可逆，还等于 $(\det A)\mathbf1^{\mathsf T}A^{-1}\mathbf1$。另一方法利用列多线性：$A+t\mathbf1\mathbf1^{\mathsf T}$ 的各列均为原列加 $t\mathbf1$，同时取两列以上的新增项时因列相同而为零，因此
\[
 \det(A+t\mathbf1\mathbf1^{\mathsf T})
 =\det A+t\sum_{i,j=1}^n C_{ij}.
\]
由此可以选择适当 $t$ 简化行列式，再求出余子式总和；该方法对奇异矩阵同样成立。
\end{solve}

\section{矩阵的迹}
\label{sec:matrix-trace}

\begin{definition}
方阵 $A=(a_{ij})\in M_n(\mathbb F)$ 的\emph{迹（trace）}定义为 $\operatorname{tr}A=\sum_{i=1}^n a_{ii}$。
\end{definition}

\begin{proposition}
\label{prop:matrix-trace-rules}
迹对矩阵加法与数乘线性，且 $\operatorname{tr}(A^{\mathsf T})=\operatorname{tr}A$。若 $A\in M_{m,n}$、$B\in M_{n,m}$，则
\begin{equation}
 \operatorname{tr}(AB)=\operatorname{tr}(BA).
 \label{eq:matrix-trace-cyclic}
\end{equation}
因此可循环移动乘积因子，且对可逆矩阵 $P$，$\operatorname{tr}(P^{-1}AP)=\operatorname{tr}A$。
\end{proposition}
\begin{proof}
线性与转置性质由定义给出，而 $\operatorname{tr}(AB)=\sum_{i=1}^m\sum_{j=1}^n a_{ij}b_{ji}=\sum_{j=1}^n\sum_{i=1}^m b_{ji}a_{ij}=\operatorname{tr}(BA)$。把若干相邻因子视为整体即可循环移动，取整体 $P^{-1}A$ 与 $P$ 便得到相似不变性。循环移动不允许任意交换两个相邻因子。
\end{proof}

\begin{example}
证明不存在 $A,B\in M_n(\mathbb F)$ 满足 $AB-BA=cI_n$，其中 $c\ne0$。
\end{example}
\begin{solve}
两边取迹得 $0=nc$，与 $\mathbb F$ 为数域、$n\ge1$、$c\ne0$ 矛盾。这里数域的特征为零这一条件是论证的一部分。
\end{solve}

\begin{proposition}
\label{prop:matrix-trace-positive}
复矩阵 $A\in M_{m,n}(\mathbb C)$ 满足
\begin{equation}
 \operatorname{tr}(AA^{\mathsf H})=\sum_{i,j}|a_{ij}|^2\ge0,
 \label{eq:matrix-trace-positive}
\end{equation}
等号当且仅当 $A=0$。实矩阵可将共轭转置换为转置。特别地，对实方阵有
\begin{equation}
 \operatorname{tr}(A^2)\le\operatorname{tr}(AA^{\mathsf T}),
 \label{eq:matrix-trace-square-bound}
\end{equation}
等号当且仅当 $A$ 对称。
\end{proposition}
\begin{proof}
第一式逐元素展开即可。对实方阵，利用线性、转置与循环性质得到
\[
 0\le\operatorname{tr}\bigl((A-A^{\mathsf T})(A-A^{\mathsf T})^{\mathsf T}\bigr)
 =2\operatorname{tr}(AA^{\mathsf T})-2\operatorname{tr}(A^2).
\]
由第一式，等号当且仅当 $A-A^{\mathsf T}=0$。
\end{proof}

例如，若实对称矩阵 $A_1,\ldots,A_s$ 满足 $\sum_jA_j^2=0$，取迹后每项均为 $\operatorname{tr}(A_jA_j^{\mathsf T})\ge0$，故每个 $A_j=0$。若实反对称矩阵 $K$ 满足 $KB=B$，则左乘 $B^{\mathsf T}$ 得 $B^{\mathsf T}KB=B^{\mathsf T}B$，左边反对称、右边对称，故均为零，再由式~\eqref{eq:matrix-trace-positive} 得 $B=0$；复数情形相应使用 $K^{\mathsf H}=-K$ 与共轭转置。

\begin{proposition}
\label{prop:matrix-trace-pairing}
若 $X\in M_n(\mathbb F)$ 满足 $\operatorname{tr}(XC)=0$ 对每个 $C\in M_n(\mathbb F)$ 成立，则 $X=0$。因此，若对所有 $C$ 有 $\operatorname{tr}(ABC)=\operatorname{tr}(CBA)$，则 $AB=BA$。
\end{proposition}
\begin{proof}
取 $C=E_{ji}$，则 $\operatorname{tr}(XE_{ji})=x_{ij}$，故所有元素均为零。第二问中由迹的循环性质把条件化为 $\operatorname{tr}((AB-BA)C)=0$，再应用第一结论。
\end{proof}

\begin{proposition}
\label{prop:matrix-trace-characterization}
若 $f:M_n(\mathbb F)\to\mathbb F$ 为 $\mathbb F$ 线性映射，且对任意 $A,B$ 满足 $f(AB)=f(BA)$、$f(I_n)=n$，则 $f(A)=\operatorname{tr}A$。
\end{proposition}
\begin{proof}
由 $E_{ij}E_{ji}=E_{ii}$ 与 $E_{ji}E_{ij}=E_{jj}$ 得 $f(E_{ii})=f(E_{jj})$，再由 $\sum_iE_{ii}=I_n$ 得每个 $f(E_{ii})=1$。若 $i\ne j$，则 $E_{ii}E_{ij}=E_{ij}$ 而 $E_{ij}E_{ii}=0$，所以 $f(E_{ij})=0$。最后按式~\eqref{eq:matrix-unit-rules} 展开 $A$ 并用线性性即得结论。
\end{proof}

\section{矩阵乘法与行列式计算}
\label{sec:matrix-factor-determinant}

计算行列式时可以寻找 $M=AB$，也可以构造容易求行列式的 $V,W$ 使 $MV=W$，再由 $\det M\det V=\det W$ 求解；后一方法只有在已知 $\det V\ne0$ 时才能除以它。若矩阵元素是两组变量的乘积、幂或和，往往可按变量将其分解；若分解通过较低维的中间空间，还可以直接判定高阶行列式为零。

\begin{example}
\label{ex:matrix-low-rank-factor}
设 $n\ge3$，证明 $M=(1+x_iy_j)_{i,j=1}^n$ 与 $K=(\cos(\alpha_i-\beta_j))_{i,j=1}^n$ 都奇异。
\end{example}
\begin{solve}
分别分解为
\[
 M=\begin{pmatrix}1&x_1\\\vdots&\vdots\\1&x_n\end{pmatrix}
 \begin{pmatrix}1&\cdots&1\\y_1&\cdots&y_n\end{pmatrix},\qquad
 K=\begin{pmatrix}\cos\alpha_1&\sin\alpha_1\\\vdots&\vdots\\\cos\alpha_n&\sin\alpha_n\end{pmatrix}
 \begin{pmatrix}\cos\beta_1&\cdots&\cos\beta_n\\\sin\beta_1&\cdots&\sin\beta_n\end{pmatrix}.
\]
两个乘积的中间维数都为 $2$，故秩均不超过 $2<n$，从而行列式为零。同理，若 $n$ 个多项式 $f_j$ 的次数均不超过 $n-2$，则评价矩阵 $(f_j(x_i))$ 可分解为 $n\times(n-1)$ 的幂矩阵乘 $(n-1)\times n$ 的系数矩阵，因此必奇异。
\end{solve}

\begin{example}
\label{ex:matrix-binomial-determinant}
对整数 $r\ge0$，求 $(r+1)$ 阶矩阵 $M=((x_i+y_j)^r)_{i,j=0}^r$ 的行列式。
\end{example}
\begin{solve}
二项式公式给出 $M=V_xDW_y$，其中 $(V_x)_{ik}=x_i^k$、$D=\operatorname{diag}(\binom r0,\ldots,\binom rr)$、$(W_y)_{kj}=y_j^{r-k}$，各索引从 $0$ 到 $r$。$V_x$ 为 Vandermonde 矩阵，$W_y$ 为幂次倒序的转置 Vandermonde 矩阵，故
\begin{equation}
 \det M=\left(\prod_{k=0}^r\binom rk\right)
 \prod_{0\le i<j\le r}(x_j-x_i)(y_i-y_j).
 \label{eq:matrix-binomial-determinant}
\end{equation}
即使节点重复，上式仍成立，此时对应的差为零。这一分解还解释了符号的来源：将 $r+1$ 个幂次逆序，共引入 $(-1)^{r(r+1)/2}$。
\end{solve}

\begin{example}
\label{ex:matrix-hankel-moments}
设 $s_k=\sum_{\ell=1}^n x_\ell^k$（$k\ge0$，其中 $s_0=n$），求 $H_q=(s_{i+j+q})_{i,j=0}^{n-1}$ 的行列式，其中 $q$ 为非负整数。
\end{example}
\begin{solve}
令 $V=(x_\ell^i)_{0\le i\le n-1,\,1\le\ell\le n}$，则
\[
 H_q=V\operatorname{diag}(x_1^q,\ldots,x_n^q)V^{\mathsf T},
\]
因而
\begin{equation}
 \det H_q=\left(\prod_{\ell=1}^n x_\ell^q\right)
 \prod_{1\le i<j\le n}(x_j-x_i)^2.
 \label{eq:matrix-hankel-determinant}
\end{equation}
$q=0$ 时对角块为 $I_n$，无须节点非零；节点为实数时 $H_0=VV^{\mathsf T}$，其行列式非负，且恰在节点两两不同时为正。$q=1$ 则给出整体平移一个下标的幂和行列式，额外出现因子 $\prod_\ell x_\ell$。
\end{solve}

\begin{example}
\label{ex:matrix-circulant-determinant}
设循环矩阵 $A=f(J)$，$f(t)=a_1+a_2t+\cdots+a_nt^{n-1}$，求 $\det A$。
\end{example}
\begin{solve}
在 $\mathbb C$ 中取全部 $n$ 次单位根 $\zeta_1,\ldots,\zeta_n$。对 $v_j=(1,\zeta_j,\ldots,\zeta_j^{n-1})^{\mathsf T}$，有 $Jv_j=\zeta_jv_j$，故 $Av_j=f(\zeta_j)v_j$。矩阵 $V=(v_1,\ldots,v_n)$ 为节点互异的 Vandermonde 矩阵，所以可逆，并且 $AV=V\operatorname{diag}(f(\zeta_1),\ldots,f(\zeta_n))$，从而
\begin{equation}
 \det A=\prod_{j=1}^n f(\zeta_j).
 \label{eq:matrix-circulant-determinant}
\end{equation}
即使 $A$ 的元素属于实数域，也可在复数域中完成计算，因为行列式值并不随扩充数域而改变。若某个 $f(\zeta_j)=0$，则 $v_j$ 同时给出一个非零核向量。
\end{solve}

\section{Cauchy--Binet 公式的应用}
\label{sec:matrix-cauchy-binet-applications}

\begin{proposition}
\label{prop:matrix-gram-minors}
设 $A\in M_{m,n}(\mathbb R)$，则 $AA^{\mathsf T}$ 的所有主子式均非负。更具体地，对任意大小为 $r$ 的行索引集 $I$，
\begin{equation}
 \det\bigl((AA^{\mathsf T})[I,I]\bigr)
 =\sum_{\substack{J\subseteq\{1,\ldots,n\}\\|J|=r}}\bigl(\det A[I,J]\bigr)^2.
 \label{eq:matrix-gram-minors}
\end{equation}
$r>n$ 时其值为 $0$；复矩阵 $AA^{\mathsf H}$ 的对应公式为绝对值平方之和。
\end{proposition}
\begin{proof}
由式~\eqref{eq:cauchy-binet-minor}，第二因子的子式为第一因子相应子式的转置行列式，故得到平方和；共轭转置时得到共轭乘积。若 $AA^{\mathsf T}=I_m$，则任意 $r\le m$ 个固定行所产生的 $r$ 阶子式平方和均为 $1$，这是同一公式的直接应用。
\end{proof}

\begin{proposition}
\label{prop:matrix-principal-minor-sums}
设 $A\in M_{m,n}(\mathbb F)$、$B\in M_{n,m}(\mathbb F)$，则 $AB$ 与 $BA$ 的 $r$ 阶主子式之和相等，其中 $1\le r\le\min\{m,n\}$。
\end{proposition}
\begin{proof}
对 $AB$ 的全部 $r$ 阶主子式使用式~\eqref{eq:cauchy-binet-minor}，得到
\[
 \sum_{|I|=r}\det\bigl((AB)[I,I]\bigr)
 =\sum_{|I|=r}\sum_{|J|=r}\det A[I,J]\det B[J,I].
\]
交换有限求和顺序并再次使用该公式，得到 $\sum_{|J|=r}\det\bigl((BA)[J,J]\bigr)$。$r=1$ 时正是迹的交换公式；更高阶主子式之和可视为它的推广。
\end{proof}

\begin{theorem}
\label{thm:matrix-lagrange-cauchy}
对实向量 $a,b\in\mathbb R^n$，有 Lagrange 恒等式
\begin{equation}
 \left(\sum_i a_i^2\right)\left(\sum_i b_i^2\right)
 -\left(\sum_i a_ib_i\right)^2
 =\sum_{i<j}(a_ib_j-a_jb_i)^2.
 \label{eq:matrix-lagrange}
\end{equation}
因此有\emph{柯西--施瓦茨不等式（Cauchy--Schwarz inequality）}
\begin{equation}
 \left(\sum_i a_ib_i\right)^2
 \le\left(\sum_i a_i^2\right)\left(\sum_i b_i^2\right),
 \label{eq:matrix-cauchy-schwarz}
\end{equation}
等号当且仅当 $a,b$ 线性相关，包括某个向量为零的情形。
\end{theorem}
\begin{proof}
对两行矩阵 $R=\begin{pmatrix}a^{\mathsf T}\\b^{\mathsf T}\end{pmatrix}$，$\det(RR^{\mathsf T})$ 展开即为左侧，Cauchy--Binet 公式则把它写成所有二阶子式的平方和，得到恒等式与不等式。等号当且仅当全部 $a_ib_j-a_jb_i=0$；若 $a\ne0$，取 $a_k\ne0$，即可得到 $b_i=(b_k/a_k)a_i$ 对所有 $i$ 成立，零向量情形另行直接成立。
\end{proof}

\begin{corollary}
\label{cor:matrix-determinant-cauchy}
若 $A,B\in M_{m,n}(\mathbb R)$，则
\begin{equation}
 \bigl(\det(AB^{\mathsf T})\bigr)^2
 \le\det(AA^{\mathsf T})\det(BB^{\mathsf T}).
 \label{eq:matrix-determinant-cauchy}
\end{equation}
当 $m\le n$ 时，等号当且仅当以所有 $m$ 阶子式为分量的两个向量 $(\det A[:,J])_J$、$(\det B[:,J])_J$ 线性相关；$m>n$ 时各行列式均为零。
\end{corollary}
\begin{proof}
$m\le n$ 时 Cauchy--Binet 公式把 $\det(AB^{\mathsf T})$ 写为上述两个子式向量的内积，把另外两个行列式写为各自平方和，再应用定理~\ref{thm:matrix-lagrange-cauchy}。$m>n$ 时由定理~\ref{thm:cauchy-binet} 得到零值。
\end{proof}

\section{分块初等变换与降阶公式}
\label{sec:matrix-block-elimination}

分块消元把若干元素的消去合并为一次矩阵运算。对
$M=\begin{pmatrix}A&B\\C&D\end{pmatrix}$，约定 $A\in M_m$、$D\in M_n$、$B\in M_{m,n}$、$C\in M_{n,m}$，于是每个乘法的尺寸都确定。第三类分块行变换对应左乘形如 $\begin{pmatrix}I_m&0\\X&I_n\end{pmatrix}$ 的矩阵，分块列变换对应右乘相应矩阵，它们的行列式均为 $1$。若交换两个含 $p,q$ 行的相邻分块行，行列式因子为 $(-1)^{pq}$，不能一概按普通两行交换的 $-1$ 处理。

\begin{theorem}
\label{thm:matrix-schur-determinant}
若 $A$ 可逆，则称 $S=D-CA^{-1}B$ 为关于块 $A$ 的\emph{舒尔补（Schur complement）}，并有
\begin{equation}
 \det\begin{pmatrix}A&B\\C&D\end{pmatrix}
 =\det A\det(D-CA^{-1}B).
 \label{eq:matrix-schur-a}
\end{equation}
若 $D$ 可逆，则相应地有
\begin{equation}
 \det\begin{pmatrix}A&B\\C&D\end{pmatrix}
 =\det D\det(A-BD^{-1}C).
 \label{eq:matrix-schur-d}
\end{equation}
\end{theorem}
\begin{proof}
$A$ 可逆时，左乘第三类分块初等矩阵得到
\[
 \begin{pmatrix}I_m&0\\-CA^{-1}&I_n\end{pmatrix}
 \begin{pmatrix}A&B\\C&D\end{pmatrix}
 =\begin{pmatrix}A&B\\0&D-CA^{-1}B\end{pmatrix}.
\]
左乘因子行列式为 $1$，右侧分块上三角，得到第一式。$D$ 可逆时用 $\begin{pmatrix}I_m&-BD^{-1}\\0&I_n\end{pmatrix}$ 左乘，消去右上块，得到第二式。
\end{proof}

\begin{corollary}
\label{cor:matrix-sylvester-determinant}
设 $U\in M_{m,n}(\mathbb F)$、$V\in M_{n,m}(\mathbb F)$，则
\begin{equation}
 \det(I_m+UV)=\det(I_n+VU).
 \label{eq:matrix-sylvester-determinant}
\end{equation}
若 $A\in M_m$ 可逆，$u,v\in\mathbb F^m$，则
\begin{equation}
 \det(A+uv^{\mathsf T})=\det A\,(1+v^{\mathsf T}A^{-1}u).
 \label{eq:matrix-determinant-lemma}
\end{equation}
\end{corollary}
\begin{proof}
对 $\begin{pmatrix}I_m&U\\-V&I_n\end{pmatrix}$ 分别按两个对角块消元，得到第一式。再分解 $A+uv^{\mathsf T}=A(I_m+A^{-1}uv^{\mathsf T})$，并将第一式用于中间维数为 $1$ 的情形，得到第二式。
\end{proof}

式~\eqref{eq:matrix-sylvester-determinant} 能把高阶行列式降为低阶行列式：例如 $D$ 为可逆对角矩阵、$U\in M_{m,r}$、$V\in M_{r,m}$ 时，$\det(D+UV)=\det D\det(I_r+VD^{-1}U)$，当 $r\ll m$ 时只须计算一个较小的行列式。特别地，$\det(I_n+aa^{\mathsf T})=1+a^{\mathsf T}a$；若 $a$ 为实向量，此值严格为正。对含参数问题，保留降阶后的多项式形式，通常比展开原行列式更易判断零点。

\begin{example}
\label{ex:matrix-bordered-determinant}
设 $A\in M_n(\mathbb F)$、$u,v\in\mathbb F^n$、$b\in\mathbb F$，证明
\begin{equation}
 \det\begin{pmatrix}A&u\\v^{\mathsf T}&b\end{pmatrix}
 =b\det A-v^{\mathsf T}A^*u.
 \label{eq:matrix-bordered-determinant}
\end{equation}
若 $A$ 可逆，给出这个加边矩阵的可逆条件。
\end{example}
\begin{solve}
$A$ 可逆时，由式~\eqref{eq:matrix-schur-a}，左侧为 $\det A\,(b-v^{\mathsf T}A^{-1}u)$，即所求表达式；因此加边矩阵可逆当且仅当 $b-v^{\mathsf T}A^{-1}u\ne0$。奇异时，式~\eqref{eq:matrix-bordered-determinant} 两侧均为 $A$ 元素的多项式，可按第~\ref{sec:matrix-perturbation} 节令 $A$ 替换为 $A+tI_n$ 后取极限；所得恒等式对所有 $A$ 成立，但含 $A^{-1}$ 的可逆条件仍只适用于 $A$ 可逆时。
\end{solve}

\begin{proposition}
\label{prop:matrix-symmetric-block-determinants}
对任意同阶方阵 $A,B$，无需交换假设即有
\begin{equation}
 \det\begin{pmatrix}A&B\\B&A\end{pmatrix}
 =\det(A+B)\det(A-B).
 \label{eq:matrix-symmetric-block-determinant}
\end{equation}
对复方阵还有
\begin{equation}
 \det\begin{pmatrix}A&B\\-B&A\end{pmatrix}
 =\det(A+\mathrm iB)\det(A-\mathrm iB).
 \label{eq:matrix-complex-block-determinant}
\end{equation}
若 $AB=BA$，第二式可进一步写为 $\det(A^2+B^2)$；若 $A,B$ 为实矩阵，则第二式等于 $|\det(A+\mathrm iB)|^2\ge0$。
\end{proposition}
\begin{proof}
第一式先把第二分块行加到第一分块行，再把第一分块列从第二分块列中减去，得到 $\begin{pmatrix}A+B&0\\B&A-B\end{pmatrix}$。第二式把第二分块行的 $\mathrm i$ 倍加到第一分块行，再把第一分块列的 $-\mathrm i$ 倍加到第二分块列，得到 $\begin{pmatrix}A-\mathrm iB&0\\-B&A+\mathrm iB\end{pmatrix}$。所用变换均不改变行列式。只有在 $AB=BA$ 时，$(A+\mathrm iB)(A-\mathrm iB)=A^2+B^2$；实数情形两个行列式互为共轭，得到非负性。
\end{proof}

\begin{proposition}
\label{prop:matrix-block-inverse}
若 $A,D$ 可逆，则
\begin{equation}
 \begin{pmatrix}A&B\\0&D\end{pmatrix}^{-1}
 =\begin{pmatrix}A^{-1}&-A^{-1}BD^{-1}\\0&D^{-1}\end{pmatrix}.
 \label{eq:matrix-block-triangular-inverse}
\end{equation}
更一般地，若 $A$ 与 $S=D-CA^{-1}B$ 均可逆，则
\begin{equation}
 \begin{pmatrix}A&B\\C&D\end{pmatrix}^{-1}
 =\begin{pmatrix}
 A^{-1}+A^{-1}BS^{-1}CA^{-1}&-A^{-1}BS^{-1}\\
 -S^{-1}CA^{-1}&S^{-1}
 \end{pmatrix}.
 \label{eq:matrix-schur-inverse}
\end{equation}
\end{proposition}
\begin{proof}
分块消元给出分解
\[
 \begin{pmatrix}A&B\\C&D\end{pmatrix}
 =\begin{pmatrix}I_m&0\\CA^{-1}&I_n\end{pmatrix}
 \begin{pmatrix}A&B\\0&S\end{pmatrix}.
\]
上三角块矩阵的逆由逐块求解乘积等于单位矩阵得到，正是式~\eqref{eq:matrix-block-triangular-inverse}；再按反次序乘两个因子的逆即得一般公式。由定理~\ref{thm:matrix-schur-determinant}，在 $A$ 可逆的前提下，整个矩阵可逆当且仅当 $S$ 可逆。
\end{proof}

\section{摄动法及其应用}
\label{sec:matrix-perturbation}

\begin{lemma}
\label{lem:matrix-perturbation}
对任意 $A\in M_n(\mathbb F)$，存在 $\varepsilon>0$，使每个满足 $0<t<\varepsilon$ 的实数 $t$ 都有 $\det(A+tI_n)\ne0$。可以选取有理数列 $t_k\to0$，使 $A+t_kI_n$ 均为可逆矩阵且元素仍在 $\mathbb F$ 中；对有限多个给定方阵，可以使用同一数列。
\end{lemma}
\begin{proof}
$\det(A+tI_n)$ 是关于 $t$ 的首一 $n$ 次多项式，故只有有限个复根。若存在非零根，取 $\varepsilon$ 小于这些根的模的最小值；若不存在非零根，任取 $\varepsilon>0$。于是所述区间没有根，区间内趋于零的有理数列即可使用。对有限多个方阵避开它们全部非零根的模，即可同时保证可逆。
\end{proof}

\emph{摄动法（perturbation method）}先在可逆情形证明结论，再以 $A+t_kI_n$ 代入，最后令 $t_k\to0$。这要求摄动后仍满足题设，而且所取极限中的表达式在目标矩阵处连续。行列式、伴随矩阵与有限矩阵乘积都是元素的多项式，因此连续；逆矩阵则只在可逆矩阵集合内连续，接近奇异矩阵时可能发散，不能直接把含逆的表达式带到极限中。交换关系 $AC=CA$ 在 $A\mapsto A+tI_n$ 下保持，但反对称、幂零或幂等条件通常不保持，须逐项检查。

例如，对可逆 $A,B$，由逆矩阵公式可直接证明 $(AB)^*=B^*A^*$；对一般 $A,B$，选同一数列使 $A+t_kI_n$ 与 $B+t_kI_n$ 都可逆，得到
\[
 \bigl((A+t_kI_n)(B+t_kI_n)\bigr)^*
 =(B+t_kI_n)^*(A+t_kI_n)^*.
\]
两侧均连续，取极限便得到一般的乘积伴随公式。同样，$\det A^*=(\det A)^{n-1}$（$n\ge2$）、二次伴随公式及分块对角伴随公式，都可先用逆阵证明，再用这一方法推广；它提供了定理~\ref{thm:matrix-adjugate-rules} 和命题~\ref{prop:matrix-block-adjugate} 的另一条证明路线。

\begin{proposition}
\label{prop:matrix-commuting-block-determinant}
设 $A,B,C,D\in M_n(\mathbb F)$ 且 $AC=CA$，则
\begin{equation}
 \det\begin{pmatrix}A&B\\C&D\end{pmatrix}=\det(AD-CB).
 \label{eq:matrix-commuting-block-determinant}
\end{equation}
\end{proposition}
\begin{proof}
若 $A$ 可逆，由 $AC=CA$ 得 $ACA^{-1}=C$，于是
\[
 \det\begin{pmatrix}A&B\\C&D\end{pmatrix}
 =\det A\det(D-CA^{-1}B)
 =\det(AD-ACA^{-1}B)=\det(AD-CB).
\]
一般情况下，$A+t_kI_n$ 仍与 $C$ 可交换，且可选为可逆。将其代入已证等式，两侧都是关于 $t_k$ 的多项式，令 $t_k\to0$ 即得结论。这里的条件指定的是 $A,C$ 可交换；不能只凭块的形式就将分块行列式写为 $\det(AD-CB)$。
\end{proof}

\begin{example}
\label{ex:matrix-adjugate-commutation}
若 $AB=BA$，证明 $AB^*=B^*A$。
\end{example}
\begin{solve}
$B$ 可逆时，$AB^{-1}=B^{-1}A$，乘 $\det B$ 得到结论。一般情况下，$A$ 仍与 $B+t_kI_n$ 可交换，故 $A(B+t_kI_n)^*=(B+t_kI_n)^*A$；利用伴随矩阵的连续性取极限即可。这个证明没有把奇异矩阵的伴随误写成行列式乘逆矩阵。
\end{solve}

\section{基础训练与方法整合}
\label{sec:matrix-training}

基础题首先检验运算能否进行、乘法次序是否保留、是否偷偷约去了奇异因子。由 $AB$ 的第 $j$ 列等于 $A(Be_j)$ 可知，$B$ 的零列必传递给乘积的对应列；但零行的传递方向不同。初等矩阵都可逆，乘积也可逆，却未必仍是一次初等变换的矩阵；初等变换一般改变行列式值，只保持其是否为零。另由式~\eqref{eq:matrix-adjugate-identity} 应区分 $AA^*=(\det A)I_n$ 与 $\det(AA^*)=(\det A)^n$，不能把矩阵等式与标量等式混用。以下将教材训练题中能联结多个方法的题目集中展开。

\begin{example}
\label{ex:matrix-training-polynomial}
证明以下结论：若 $A^2+A+I_n=0$，则 $A$ 可逆；若 $P^2=P$，则 $I_n-2P$ 可逆；若 $A$ 幂零且 $AB+BA=B$，则 $B=0$。
\end{example}
\begin{solve}
第一问由 $A(-A-I_n)=I_n$ 得 $A^{-1}=-A-I_n$。第二问由 $(I_n-2P)^2=I_n$，得到 $(I_n-2P)^{-1}=I_n-2P$。第三问把条件写为 $AB=B(I_n-A)$，归纳得 $A^kB=B(I_n-A)^k$。取 $A^s=0$ 得 $B(I_n-A)^s=0$；$I_n-A$ 的逆是有限几何级数，故右乘其 $s$ 次逆幂即可得到 $B=0$。这里不能仅从 $A^sB=0$ 约去 $A^s$。
\end{solve}

\begin{example}
\label{ex:matrix-training-companion-inverse}
设 $n\ge2$、$a_n\ne0$，求命题~\ref{prop:matrix-companion} 中友阵 $C(f)$ 的逆。
\end{example}
\begin{solve}
设 $C(f)x=y$。第一行给出 $x_n=-y_1/a_n$，其余各行给出 $x_{i-1}=y_i-(a_{n-i+1}/a_n)y_1$（$2\le i\le n$），因此
\[
 C(f)^{-1}=\begin{pmatrix}
 -a_{n-1}/a_n&1&0&\cdots&0\\
 -a_{n-2}/a_n&0&1&\cdots&0\\
 \vdots&\vdots&\vdots&\ddots&\vdots\\
 -a_1/a_n&0&0&\cdots&1\\
 -1/a_n&0&0&\cdots&0
 \end{pmatrix}.
\]
这相当于对增广矩阵进行消元，但直接逐行求解更能显示哪些未知量由哪一行确定。
\end{solve}

\begin{example}
\label{ex:matrix-training-trace-inequality}
设 $A,B$ 为实对称矩阵，证明 $\operatorname{tr}((AB)^2)\le\operatorname{tr}(A^2B^2)$，并确定等号条件。
\end{example}
\begin{solve}
将式~\eqref{eq:matrix-trace-square-bound} 应用于 $X=AB$。由 $X^{\mathsf T}=BA$ 与迹的循环性质，
\[
 \operatorname{tr}((AB)^2)\le\operatorname{tr}(ABBA)=\operatorname{tr}(A^2B^2).
\]
等号当且仅当 $AB$ 对称，即 $AB=(AB)^{\mathsf T}=BA$。这说明迹的不等式可以检测两个对称矩阵是否可交换。
\end{solve}

\begin{example}
\label{ex:matrix-training-orthogonal}
证明不存在实正交矩阵 $A,B$ 满足 $A^2=cAB+B^2$，其中 $c\ne0$ 为实数。
\end{example}
\begin{solve}
假设等式成立，左乘 $A^{\mathsf T}$、右乘 $B^{\mathsf T}$ 得 $AB^{\mathsf T}=cI_n+A^{\mathsf T}B$。但是
\[
 \operatorname{tr}(AB^{\mathsf T})
 =\operatorname{tr}((AB^{\mathsf T})^{\mathsf T})
 =\operatorname{tr}(BA^{\mathsf T})
 =\operatorname{tr}(A^{\mathsf T}B),
\]
因此两边取迹得到 $nc=0$，矛盾。正交性用于消去逆因子，迹用于消去剩余的两个乘积项。
\end{solve}

\begin{example}
\label{ex:matrix-training-b-circulant}
设 $b\ne0$，令 $J_b$ 的上副对角元均为 $1$、左下角为 $b$、其余为 $0$。把
\[
 A=a_1I_n+a_2J_b+\cdots+a_nJ_b^{n-1}
\]
称为 $b$-循环矩阵。证明同阶 $b$-循环矩阵对乘法封闭，并求其行列式。
\end{example}
\begin{solve}
由移位作用知 $J_b^n=bI_n$，所以两个关于 $J_b$ 的多项式相乘后可以约化为次数不超过 $n-1$ 的多项式，得到封闭性。设 $f(t)=a_1+a_2t+\cdots+a_nt^{n-1}$，在复数域取 $z^n=b$ 的全部根 $\eta_1,\ldots,\eta_n$。因为 $b\ne0$，这些根互异；向量 $(1,\eta_j,\ldots,\eta_j^{n-1})^{\mathsf T}$ 满足 $J_bv_j=\eta_jv_j$。与例~\ref{ex:matrix-circulant-determinant} 相同，Vandermonde 矩阵可逆，故
\[
 \det A=\prod_{j=1}^n f(\eta_j).
\]
$b=1$ 即退化为普通循环矩阵。这里定义的 $b$-循环矩阵第一行为 $(a_1,\ldots,a_n)$，每次右移时越过右端的元素须乘 $b$。
\end{solve}

\begin{example}
\label{ex:matrix-training-cofactors-equal}
设 $n\ge2$，$A\in M_n(\mathbb F)$ 的每行与每列元素之和都为零，证明所有代数余子式相等。
\end{example}
\begin{solve}
条件给出 $A\mathbf1=0$ 与 $\mathbf1^{\mathsf T}A=0$，故 $A$ 奇异。若秩不超过 $n-2$，则所有 $(n-1)$ 阶子式为零，结论成立。若秩为 $n-1$，其核为一维，且由 $A\mathbf1=0$ 知其核由 $\mathbf1$ 张成。由 $AA^*=0$，$A^*$ 每列都是 $\mathbf1$ 的倍数，故 $A^*=\mathbf1w^{\mathsf T}$。同理左核由 $\mathbf1^{\mathsf T}$ 张成，$A^*A=0$ 迫使 $w^{\mathsf T}=c\mathbf1^{\mathsf T}$。于是 $A^*=c\mathbf1\mathbf1^{\mathsf T}$，所有代数余子式都为 $c$。这里核的维数由相抵标准型中的 $n-r$ 个自由分量得到。
\end{solve}

\begin{example}
\label{ex:matrix-training-norm-similarity}
设 $P\in M_n(\mathbb R)$ 可逆，记 $F(X)=\sum_{i,j}x_{ij}^2$。证明：对每个实方阵 $X$ 均有 $F(PXP^{-1})=F(X)$，当且仅当 $P^{\mathsf T}P=cI_n$，其中 $c>0$。
\end{example}
\begin{solve}
记 $S=P^{\mathsf T}P$，它正定且可逆。由 $PE_{ij}P^{-1}=(Pe_i)(P^{-\mathsf T}e_j)^{\mathsf T}$，外积的元素平方和等于两向量平方和之积，故题设给出 $S_{ii}(S^{-1})_{jj}=1$。当 $i\ne k$ 时，将同一条件用于 $E_{ij}+E_{kj}$ 并减去两项各自的等式，得到 $2S_{ik}(S^{-1})_{jj}=0$；由于 $(S^{-1})_{jj}>0$，故 $S_{ik}=0$。因此 $S$ 对角，而 $S_{ii}(S^{-1})_{jj}=1$ 化为 $S_{ii}/S_{jj}=1$，所有对角元相等且为正，即 $S=cI_n$。$n=1$ 时结论直接成立。

反之，若 $S=cI_n$，则由迹的循环性质与式~\eqref{eq:matrix-trace-positive}，
\[
 F(PXP^{-1})
 =\operatorname{tr}(SXS^{-1}X^{\mathsf T})
 =\operatorname{tr}(XX^{\mathsf T})=F(X).
\]
所以保持全部矩阵元素平方和的相似变换，恰由正交矩阵的非零实数倍实施，因为 $P/\sqrt c$ 为正交矩阵。
\end{solve}

本章方法的共同点是先识别可以保留的信息：单位向量与基础矩阵提取行列及元素，初等变换保留秩和可逆性，迹提取乘积中能循环比较的标量，矩阵分解暴露中间维数，分块消元把高维问题降到舒尔补，而摄动只把满足连续性与条件保持要求的结论传到奇异情形。选择方法时应先检查这些条件，再进行具体运算。

\chapter{线性空间与线性方程组}
\label{chap:linear-spaces-systems}

% 研读来源：高等代数/textbooks/高等代数习题（谢启鸿）/高代习题-3.pdf，
% 第3章，书页116--191；按原书十节组织，整合基本结论、典型方法与训练题。
线性方程组的消元过程不仅给出未知数的取值，还揭示了向量之间有多少独立关系、一个空间需要多少个自由参数以及矩阵在乘法中损失了多少维信息。本章把这些问题统一在线性空间的语言中：先由线性表示建立基与维数，再用坐标把抽象问题化成矩阵计算，最后以秩和解空间解释方程组的可解性及其应用。学习时应把握两条相互联系的路线：具体计算依靠消元，结构证明依靠生成空间、线性无关性及维数。

除特别说明外，$\mathbb F$ 表示实数域或复数域的子域，各线性空间均在同一数域 $\mathbb F$ 上，讨论基、维数和维数公式时默认空间有限维。坐标向量统一写成列向量，$M_{m\times n}(\mathbb F)=M_{m,n}(\mathbb F)$ 表示 $m\times n$ 矩阵全体，$M_n(\mathbb F)=M_{n\times n}(\mathbb F)$，$I_n$ 表示 $n$ 阶单位矩阵。$A^{\mathsf T}$ 表示转置，$A^{\mathsf H}=\overline A^{\mathsf T}$ 表示共轭转置，$\operatorname{adj}(A)$ 表示伴随矩阵；三种记号所表示的运算不同。

\section{基本概念}
\label{sec:laex-basic}

\subsection{线性空间、线性表示与线性无关}

\begin{definition}
\label{def:laex-vector-space}
设 $V$ 是非空集合，在 $V$ 上定义了加法，在 $\mathbb F\times V$ 上定义了数乘。若任意 $u,v,w\in V$ 与 $a,b\in\mathbb F$ 满足
\[
\begin{gathered}
u+v=v+u,\qquad (u+v)+w=u+(v+w),\\
\text{存在 }0\in V\text{ 使 }v+0=v,\qquad
\text{对每个 }v\text{ 存在 }-v\text{ 使 }v+(-v)=0,\\
1v=v,\qquad a(bv)=(ab)v,\\
a(u+v)=au+av,\qquad (a+b)v=av+bv,
\end{gathered}
\]
则称 $V$ 为 $\mathbb F$ 上的\emph{线性空间（vector space）}，其元素称为向量。
\end{definition}

线性空间由集合、加法、数乘和标量域共同确定，不能只看元素的外观。通常运算下的 $\mathbb F^n$、矩阵空间、多项式空间以及实连续函数空间都是典型例子；在这些例子中，向量可以是有序数组、矩阵、多项式或函数。次数恰为 $n$ 的多项式全体则不是线性空间，因为它不含零多项式，而且两个这样的多项式相加可能降低次数。

\begin{definition}
\label{def:laex-span-independence}
给定向量 $v_1,\ldots,v_m\in V$，形如 $a_1v_1+\cdots+a_mv_m$ 的向量称为它们的\emph{线性组合（linear combination）}，其中 $a_i\in\mathbb F$。所有这些组合组成的集合记为
\[
\operatorname{span}(v_1,\ldots,v_m)
=\left\{\sum_{i=1}^m a_iv_i:a_i\in\mathbb F\right\},
\]
称为该向量组的生成空间。若等式 $a_1v_1+\cdots+a_mv_m=0$ 只能由 $a_1=\cdots=a_m=0$ 得到，则称向量组\emph{线性无关（linearly independent）}；否则称其线性相关。约定空向量组线性无关，且其生成空间为 $\{0\}$。
\end{definition}

这里的等式是空间中的向量等式。例如，函数的线性关系必须对定义域中的每一点成立，多项式的线性关系必须是恒等式，不能只在未经说明的若干个点上验证。一个非空向量组只要含零向量就线性相关，但不含零向量远不能保证线性无关。

\begin{proposition}
\label{prop:laex-dependence-basic}
有限向量组线性相关，当且仅当其中至少一个向量可以由其余向量线性表示。若 $v_1,\ldots,v_m$ 线性无关，则对任意 $v\in V$，加入 $v$ 后的向量组线性无关，当且仅当 $v\notin\operatorname{span}(v_1,\ldots,v_m)$。一个向量在给定向量组中的线性表示一旦存在，其表示系数唯一，当且仅当该向量组线性无关。
\end{proposition}

\begin{proof}
非平凡关系 $\sum_i a_iv_i=0$ 中至少有一个 $a_j\ne0$，移项后除以 $a_j$ 就得到 $v_j$ 的表示；反过来，将这样的表示移到等号一侧就得到非平凡关系。若 $\sum_i a_iv_i+bv=0$，当 $b=0$ 时，原向量组无关性迫使所有 $a_i=0$；因此加入新向量后出现非平凡关系，当且仅当可以从 $b\ne0$ 的关系中解出 $v$。最后，两种表示相减得到齐次关系；若存在非平凡齐次关系，将其加在已有表示上就会产生另一组系数。
\end{proof}

\subsection{基、维数、子空间与矩阵的秩}

\begin{definition}
\label{def:laex-basis-rank}
一个向量组的\emph{极大线性无关组（maximal linearly independent subset）}是其中既线性无关、又能生成原向量组中全部向量的部分组。其所含向量个数称为原向量组的\emph{秩（rank）}。若线性无关组能生成整个空间 $V$，则称其为 $V$ 的一组\emph{基（basis）}；有限基所含向量个数称为 $V$ 的\emph{维数（dimension）}，记为 $\dim_{\mathbb F}V$，标量域明确时简记为 $\dim V$。零空间的基为空组，维数为零。
\end{definition}

这些定义中的个数不依赖所选的极大无关组，依据是如下基本比较原理：若 $r$ 个线性无关向量均可由 $s$ 个向量生成，则 $r\le s$。事实上，把这 $r$ 个向量的表示系数排成 $s\times r$ 矩阵；若 $r>s$，对相应的齐次方程组消元，主元至多有 $s$ 个，必有自由未知数，因而存在非零系数使这 $r$ 个向量的组合为零，产生矛盾。将此原理分别用于两组基，就得到它们具有相同个数。

\begin{definition}
\label{def:laex-subspace}
若 $U\subseteq V$ 在原有加法和数乘下也是线性空间，则称 $U$ 为 $V$ 的\emph{子空间（subspace）}。判定时只需检查 $U\ne\varnothing$，且对任意 $u,v\in U$ 与 $a,b\in\mathbb F$ 有 $au+bv\in U$。两个子空间的交与和分别为
\[
U\cap W=\{v:v\in U,\ v\in W\},\qquad
U+W=\{u+w:u\in U,\ w\in W\}.
\]
两者都是子空间。若 $U\cap W=\{0\}$，则称这个和为\emph{直和（direct sum）}，记为 $U\oplus W$。
\end{definition}

矩阵 $A\in M_{m\times n}(\mathbb F)$ 的列空间是其列向量的生成空间，记为 $\operatorname{im}A=\{Ax:x\in\mathbb F^n\}$；核空间记为 $\ker A=\{x\in\mathbb F^n:Ax=0\}$。矩阵的秩定义为列空间的维数，记为 $\operatorname{rank}A$。下一节将说明它也等于行向量组的秩，以及行阶梯形矩阵中主元的个数。由消元可知，若 $\operatorname{rank}A=r$，齐次方程组就有 $n-r$ 个自由未知数，依次令各自由未知数取标准单位向量得到其解空间的一组基，因此
\begin{equation}
\label{eq:laex-rank-nullity}
\dim\ker A+\operatorname{rank}A=n.
\end{equation}
这个等式称为\emph{秩与零度公式（rank--nullity formula）}；$n$ 是未知数个数，不是方程个数。

\section{向量的线性关系}
\label{sec:laex-linear-relations}

\subsection{消元、秩与极大无关组}

具体向量的线性关系可以转化为矩阵问题。将 $v_1,\ldots,v_s\in\mathbb F^n$ 按列组成 $A=(v_1,\ldots,v_s)$，则 $\sum_j c_jv_j=0$ 等价于 $Ac=0$。若向量原来写成行向量，仍可先转置再按列排列，以便统一使用下面的方法。

\begin{proposition}
\label{prop:laex-pivot-columns}
设 $R=PA$ 是 $A$ 经初等行变换得到的行阶梯形矩阵，其中 $P$ 可逆。则 $A$ 与 $R$ 的列向量之间具有完全相同的线性关系。若 $R$ 有 $r$ 个非零行，主元列的列指标为 $j_1,\ldots,j_r$，则 $\operatorname{rank}A=r$，原矩阵中第 $j_1,\ldots,j_r$ 列组成列向量组的一个极大无关组。矩阵的行秩与列秩相等。
\end{proposition}

\begin{proof}
由 $P$ 可逆，$Ac=0$ 当且仅当 $Rc=0$。$R$ 的主元列在前 $r$ 行构成一个对角元非零的上三角矩阵，故这 $r$ 列无关；而所有列的后续分量均为零，其生成空间维数至多为 $r$，所以主元列是极大无关组。$R$ 的 $r$ 个非零行也线性无关：依次考察各主元位置，任何为零的行向量线性组合都迫使系数逐个为零。初等行变换保持行空间不变，故 $A$ 的行秩也为 $r$。
\end{proof}

\begin{notealert}[type=caution,title={返回原矩阵取列向量}]
求原列向量组的极大无关组时，应使用行阶梯形矩阵确定列指标，再返回原矩阵取相应列。行变换保持列之间的关系，却通常改变各列本身；任意列变换只保证矩阵的秩不变，不能保证原列指标仍有上述意义。
\end{notealert}

\begin{example}
\label{ex:laex-maximal-subset}
求向量组
\[
\begin{aligned}
v_1&=(1,1,2,2,1)^{\mathsf T},&v_2&=(0,2,1,5,-1)^{\mathsf T},\\
v_3&=(2,0,3,-1,3)^{\mathsf T},&v_4&=(1,1,0,4,-1)^{\mathsf T}
\end{aligned}
\]
的秩与一个极大无关组。
\end{example}

\begin{solve}
先观察到 $v_3=2v_1-v_2$，所以原向量组可由 $v_1,v_2,v_4$ 生成。取这三列的前三个分量得到子矩阵
\[
\begin{pmatrix}1&0&1\\1&2&1\\2&1&0\end{pmatrix},
\]
其行列式为 $-4\ne0$，故 $v_1,v_2,v_4$ 无关。因此秩为 $3$，它们是一个极大无关组。又因为 $v_3=2v_1-v_2$ 中两个系数均非零，也可以用 $v_3$ 替换 $v_1$ 或 $v_2$，得到其他极大无关组；向量的选择可以不同，个数则始终为 $3$。
\end{solve}

判定 $b$ 是否可由 $v_1,\ldots,v_s$ 线性表示，就是判断 $Ax=b$ 是否有解，等价条件为
\begin{equation}
\label{eq:laex-linear-representation-rank}
\operatorname{rank}(A\mid b)=\operatorname{rank}A.
\end{equation}
若还需要具体系数，就应对增广矩阵消元，而不能只比较两个秩。

\begin{example}
\label{ex:laex-representation}
设 $v_1=(1,0,-1,3)^{\mathsf T}$，$v_2=(2,1,0,1)^{\mathsf T}$，$v_3=(0,-2,1,1)^{\mathsf T}$，判断 $b=(5,4,-2,4)^{\mathsf T}$ 能否由它们表示，并说明表示是否唯一。
\end{example}

\begin{solve}
直接代入可得 $b=v_1+2v_2-v_3$，故表示存在。前三个分量组成的系数矩阵行列式为 $5\ne0$，所以三列线性无关，由命题~\ref{prop:laex-dependence-basic} 可知这组表示系数唯一。
\end{solve}

\subsection{由定义证明线性关系}

抽象向量没有现成分量时，通常先设一个线性组合为零，再利用已知的线性无关性逐步消去系数。一个常用推论是：若 $b$ 能由 $v_1,\ldots,v_m$ 表示，却不能由其中任何更少的向量表示，则这 $m$ 个向量无关。否则其中一个向量可由其余向量表示，将它代回 $b$ 的表达式就减少了所需向量个数。

\begin{proposition}
\label{prop:laex-generating-rank}
若向量组 $B$ 中每个向量均可由向量组 $A$ 线性表示，则
\[
\operatorname{span}B\subseteq\operatorname{span}A,
\qquad \operatorname{rank}B\le\operatorname{rank}A.
\]
两组向量互相线性表示，称为\emph{等价向量组}；这等价于它们生成同一个子空间。若已知其中一个生成空间包含于另一个，则两组秩相等也足以推出等价。
\end{proposition}

\begin{proof}
把 $B$ 中各向量的表达式代入任意线性组合，就得到生成空间的包含。取 $B$ 的一个极大无关组，用 $A$ 的极大无关组表示它，由基本比较原理得到秩的不等式。最后，两个有限维子空间若有包含关系且维数相同，就必相等：否则可从较大空间取一个不在较小空间中的向量，加入后增加一维，产生矛盾。
\end{proof}

\begin{example}
\label{ex:laex-coefficient-matrix}
已知 $v_1,\ldots,v_m$ 线性无关，令
\[
w_j=\sum_{i=1}^m c_{ij}v_i\quad(1\le j\le s),
\qquad C=(c_{ij})\in M_{m\times s}(\mathbb F).
\]
证明 $\operatorname{rank}(w_1,\ldots,w_s)=\operatorname{rank}C$。
\end{example}

\begin{solve}
对任意 $t\in\mathbb F^s$，有
\[
\sum_{j=1}^s t_jw_j
=\sum_{i=1}^m\left(\sum_{j=1}^s c_{ij}t_j\right)v_i=0
\quad\Longleftrightarrow\quad Ct=0.
\]
因此两个向量组有完全相同的系数关系，极大无关组的指标也相同，秩自然相等。特别地，当 $s=m$ 时，新向量组无关当且仅当 $\det C\ne0$。若不假设原向量组无关，一般只能断言新向量组的秩不超过 $\operatorname{rank}C$，不能直接取等号。
\end{solve}

\begin{example}
\label{ex:laex-cyclic-vectors}
设 $v_1,\ldots,v_s$ 是一组基，$s>1$，并令 $w_i=v_i+v_{i+1}$，其中 $v_{s+1}=v_1$。讨论 $w_1,\ldots,w_s$ 的线性关系。
\end{example}

\begin{solve}
关系 $\sum_i c_iw_i=0$ 在原基下等价于 $c_i+c_{i-1}=0$，下标按模 $s$ 解释。因此系数沿环交替变号。若 $s$ 为奇数，绕一周得到 $c_1=-c_1$，因 $\mathbb F$ 的特征为零，只能所有系数为零，新向量组是一组基；若 $s$ 为偶数，则关系空间恰为一维，关系可取 $w_1-w_2+\cdots-w_s=0$，秩为 $s-1$。
\end{solve}

\begin{notealert}[type=caution,title={区分不同层次的线性关系}]
两两无关不能推出整个向量组无关，例如 $e_1,e_2,e_1+e_2$。两组向量的秩相同也不能推出它们等价，例如 $\{e_1\}$ 与 $\{e_2\}$。矩阵相抵只要求同型矩阵的秩相同，而向量组等价要求生成空间相同，不能混用这两种“等价”。
\end{notealert}

\section{线性空间及其基}
\label{sec:laex-spaces-bases}

\subsection{基的判定与扩张}

\begin{theorem}
\label{thm:laex-basis-extension}
设 $\dim V=n$。$V$ 中任意线性无关组都可以扩张为一组基；任意有限生成组都可以删减为一组基。因此，恰含 $n$ 个向量的向量组只要满足“线性无关”或“生成 $V$”中的任一条件，就是 $V$ 的一组基。
\end{theorem}

\begin{proof}
从无关组出发，若尚未生成 $V$，就在其生成空间外选取一个向量，由命题~\ref{prop:laex-dependence-basic} 得到更大的无关组。基本比较原理保证无关向量个数不超过 $n$，故有限步后得到基。从有限生成组出发，若其相关，就删除一个可由其余向量表示的向量，生成空间不变，直至得到无关组。若原向量组已含 $n$ 个向量，无关组不能继续扩张，而生成组也不能删去任何向量，否则基的个数会小于 $n$。
\end{proof}

\begin{example}
\label{ex:laex-polynomial-bases}
记 $\mathcal P_n(\mathbb F)=\{f\in\mathbb F[x]:\deg f\le n\}\cup\{0\}$，其中 $n\ge0$。说明它的维数，并在 $n+1$ 个互异节点 $a_0,\ldots,a_n\in\mathbb F$ 上构造一组便于求坐标的基。
\end{example}

\begin{solve}
按系数比较可知 $1,x,\ldots,x^n$ 既无关又生成该空间，所以维数为 $n+1$。定义
\[
\ell_j(x)=\prod_{\substack{0\le k\le n\\k\ne j}}
\frac{x-a_k}{a_j-a_k}\quad(0\le j\le n),
\]
则 $\ell_j(a_i)=\delta_{ij}$。若 $\sum_j c_j\ell_j=0$，依次代入各节点得到所有 $c_j=0$，故这 $n+1$ 个多项式是一组基。任意 $f\in\mathcal P_n(\mathbb F)$ 因而满足
\[
f(x)=\sum_{j=0}^n f(a_j)\ell_j(x).
\]
这就是拉格朗日插值公式（Lagrange interpolation formula），其坐标是节点处的函数值。次数不超过 $n$ 的限制保证这些节点值足以唯一确定多项式。
\end{solve}

\begin{example}
\label{ex:laex-matrix-bases}
求全体 $m\times n$ 矩阵、全体 $n$ 阶上三角矩阵、对称矩阵和反对称矩阵空间的维数。
\end{example}

\begin{solve}
令 $E_{ij}$ 是仅在第 $(i,j)$ 个位置为 $1$、其余位置均为零的矩阵。全体 $E_{ij}$ 构成 $M_{m\times n}(\mathbb F)$ 的基，维数为 $mn$。上三角矩阵空间的基为 $\{E_{ij}:i\le j\}$，维数为 $n(n+1)/2$；对称矩阵空间的基为
\[
\{E_{ii}:1\le i\le n\}\cup
\{E_{ij}+E_{ji}:1\le i<j\le n\},
\]
维数也为 $n(n+1)/2$；反对称矩阵空间的基为 $\{E_{ij}-E_{ji}:i<j\}$，维数为 $n(n-1)/2$。这些组的生成性由逐项展开得到，无关性则由比较各独立位置的元素得到。
\end{solve}

\subsection{标量域与无限维现象}

\begin{example}
\label{ex:laex-hermitian}
令 $\mathcal H_n=\{A\in M_n(\mathbb C):A^{\mathsf H}=A\}$，$\mathcal K_n=\{A\in M_n(\mathbb C):A^{\mathsf H}=-A\}$，分别为埃尔米特（Hermitian）矩阵与斜埃尔米特（skew-Hermitian）矩阵全体。说明它们作为实线性空间的维数。
\end{example}

\begin{solve}
这两个集合对加法和实数数乘封闭。$\mathcal H_n$ 的一组实基为
\[
\{E_{ii}\}\cup\{E_{ij}+E_{ji}:i<j\}
\cup\{\mathrm i(E_{ij}-E_{ji}):i<j\},
\]
所以 $\dim_{\mathbb R}\mathcal H_n=n+2\binom n2=n^2$。映射 $A\mapsto\mathrm iA$ 将此基送到 $\mathcal K_n$ 的一组实基，故后者维数也为 $n^2$。两者通常都不是复线性空间：非零 Hermite 矩阵乘 $\mathrm i$ 后成为斜 Hermite 矩阵，已不在原集合中。
\end{solve}

同一个集合在不同标量域上的维数可能不同。例如 $\mathbb C$ 在自身上以 $1$ 为基，维数为 $1$；在 $\mathbb R$ 上则以 $1,\mathrm i$ 为基，维数为 $2$。遇到数域扩张时，必须把允许使用的系数写清楚。

\begin{example}
\label{ex:laex-cubic-field}
设 $\theta=\sqrt[3]{2}$，证明 $W=\{a+b\theta+c\theta^2:a,b,c\in\mathbb Q\}$ 是数域，并求其在 $\mathbb Q$ 上的维数。
\end{example}

\begin{solve}
多项式 $x^3-2$ 在 $\mathbb Q$ 上无有理根，因其次数为三，所以不可约。若一个次数至多为二的非零有理系数多项式也以 $\theta$ 为根，它与 $x^3-2$ 互素，由多项式的 Bézout 等式就会推出 $1=0$，故不可能。因此 $1,\theta,\theta^2$ 无关，$\dim_{\mathbb Q}W=3$。由 $\theta^3=2$ 可知 $W$ 对乘法封闭。对任意 $0\ne u\in W$，映射 $v\mapsto uv$ 是 $W$ 到自身的单射线性映射，因为 $W\subset\mathbb R$ 中无零因子；有限维且维数相同使它也是满射。于是存在 $v\in W$ 使 $uv=1$，证明非零元的逆仍在 $W$ 中，故 $W$ 为数域。
\end{solve}

\begin{proposition}
\label{prop:laex-scalar-restriction}
设 $\mathbb F\subseteq\mathbb K$，$\mathbb K$ 在 $\mathbb F$ 上有基 $a_1,\ldots,a_m$，而 $V$ 在 $\mathbb K$ 上有基 $e_1,\ldots,e_n$。则 $\{a_ie_j:1\le i\le m,\ 1\le j\le n\}$ 是 $V$ 在 $\mathbb F$ 上的一组基，特别地
\[
\dim_{\mathbb F}V=m\dim_{\mathbb K}V=mn.
\]
\end{proposition}

\begin{proof}
把任意 $v=\sum_j b_je_j$ 中的每个 $b_j\in\mathbb K$ 展开为 $\sum_i c_{ij}a_i$，就得到这些乘积向量的生成性。若 $\sum_{i,j}c_{ij}a_ie_j=0$，先用 $e_j$ 的 $\mathbb K$ 线性无关性得到 $\sum_i c_{ij}a_i=0$，再用 $a_i$ 的 $\mathbb F$ 线性无关性得到所有 $c_{ij}=0$。
\end{proof}

无限维不能仅理解为“元素有无穷多个”，因为任何非零有限维数域上线性空间都已经有无穷多个元素。要证明空间不是有限维，应构造任意多个线性无关向量。例如 $C(\mathbb R)$ 包含任意长的无关组 $1,x,\ldots,x^N$；趋于零的实数列空间 $c_0$ 包含标准单位数列 $e_1,\ldots,e_N$，它们也无关。无限向量族的线性无关是指其每个有限子族都无关，代数意义的基只允许\emph[caution]{有限线性组合}；因此不能把任意无限级数自动当作基展开。后文使用的维数相减与有限步扩张论证均限于已声明的有限维情形。

\section{线性同构和几何问题代数化}
\label{sec:laex-isomorphism}

\begin{definition}
\label{def:laex-isomorphism}
设 $V,W$ 是同一数域 $\mathbb F$ 上的线性空间。若映射 $T:V\to W$ 为双射，且对任意 $u,v\in V$、$a,b\in\mathbb F$ 满足 $T(au+bv)=aT(u)+bT(v)$，则称 $T$ 为\emph{线性同构（linear isomorphism）}，记为 $V\cong W$。
\end{definition}

同构保持全部有限线性关系，因为 $\sum_i c_iT(v_i)=T(\sum_i c_iv_i)$，再由单射性可知左端为零当且仅当原组合为零。因此同构保持线性无关性、生成关系、向量组的秩及空间维数。同一数域上两个有限维线性空间同构，当且仅当它们维数相同：必要性由同构送基到基得到；充分性则由把一组基逐个对应到另一组基，并按线性组合延拓得到。

\begin{theorem}
\label{thm:laex-coordinate-isomorphism}
给定有序基 $E=(e_1,\ldots,e_n)$，对 $v=\sum_i x_ie_i$ 定义坐标列向量 $[v]_E=(x_1,\ldots,x_n)^{\mathsf T}$，则映射 $v\mapsto[v]_E$ 是 $V$ 到 $\mathbb F^n$ 的线性同构。因此，抽象向量组的秩等于其坐标列组成的矩阵的秩，且极大无关组的列指标相同。
\end{theorem}

\begin{proof}
基展开的存在性与唯一性分别保证映射为双射，两个基展开相加或同乘一个标量则保证线性。其余结论由同构保持线性关系得到。
\end{proof}

\emph[tip]{几何问题代数化}的关键是先固定同一组基，把每个向量转换成坐标列，再用矩阵运算求解，最后将答案解释回原空间。不同向量若使用不同基，所得坐标列就不能直接拼成同一个矩阵。例如函数或多项式的坐标既可以是通常系数，也可以是在适当基下的节点值；选择坐标的目的，是使原问题的线性关系更容易计算。

\begin{example}
\label{ex:laex-vandermonde-basis}
设 $E=(e_1,\ldots,e_n)$ 是 $V$ 的一组基，$a_1,\ldots,a_n\in\mathbb F$ 互不相同，令
\[
v_j=e_1+a_je_2+\cdots+a_j^{n-1}e_n\quad(1\le j\le n).
\]
证明 $v_1,\ldots,v_n$ 也是一组基。
\end{example}

\begin{solve}
坐标矩阵是范德蒙德矩阵（Vandermonde matrix），其行列式为
\[
\det\begin{pmatrix}
1&\cdots&1\\a_1&\cdots&a_n\\\vdots&&\vdots\\a_1^{n-1}&\cdots&a_n^{n-1}
\end{pmatrix}
=\prod_{1\le i<j\le n}(a_j-a_i)\ne0.
\]
因此坐标列无关，原向量也无关，再由定理~\ref{thm:laex-basis-extension} 得到结论。这种构造把任意 $n$ 个互异标量变成一组基，后面可用于避开有限多个真子空间。
\end{solve}

\begin{example}
\label{ex:laex-unusual-operations}
在 $\mathbb R_{>0}$ 上定义 $u\oplus v=uv$、$a\odot u=u^a$。另在 $\mathbb R^2$ 上定义
\[
(u,v)\oplus(s,t)=(u+s,v+t+us),\qquad
a\odot(u,v)=\left(au,av+\frac{a(a-1)}2u^2\right).
\]
说明这两组运算都使相应集合成为实线性空间。
\end{example}

\begin{solve}
第一种情形的映射 $T:\mathbb R\to\mathbb R_{>0}$，$T(x)=e^x$ 为双射，满足 $T(x+y)=T(x)\oplus T(y)$、$T(ax)=a\odot T(x)$。因此新运算由通常实线性空间的运算传递而来，八条公理均成立；其零向量是数 $1$，加法逆元是通常乘法的倒数。

第二种情形使用双射 $S(x,y)=(x,y+x^2/2)$，其逆为 $S^{-1}(u,v)=(u,v-u^2/2)$。展开平方可验证 $S(z_1+z_2)=S(z_1)\oplus S(z_2)$ 与 $S(az)=a\odot S(z)$，因而同样构成线性空间。所谓“线性”是相对于所定义的加法和数乘而言，不能仅凭公式中出现平方或乘方就否定线性空间结构。
\end{solve}

\section{基变换与过渡矩阵}
\label{sec:laex-change-basis}

\subsection{列的含义与坐标变换方向}

\begin{definition}
\label{def:laex-transition}
设 $E=(e_1,\ldots,e_n)$ 与 $F=(f_1,\ldots,f_n)$ 是 $V$ 的两组有序基，令
\[
f_j=\sum_{i=1}^n p_{ij}e_i,\qquad P=(p_{ij}).
\]
称 $P$ 为从基 $E$ 到基 $F$ 的\emph{过渡矩阵（transition matrix）}。其第 $j$ 列就是 $[f_j]_E$，形式上记为 $F=EP$。
\end{definition}

若 $x=[v]_E$、$y=[v]_F$，则 $v=Ex=Fy=EPy$，由基展开唯一性得到
\begin{equation}
\label{eq:laex-coordinate-change}
x=Py,\qquad y=P^{-1}x.
\end{equation}
因此“从旧基到新基的过渡矩阵”将新坐标变成旧坐标；要把旧坐标变成新坐标，应使用其逆矩阵。$P$ 可逆，因为它的列是另一组基在 $E$ 下的坐标。若第三组基满足 $G=FQ$，则 $G=E(PQ)$，故过渡矩阵沿基变换次序相乘。

具体计算时，若基 $E,F$ 在某一共同参考基下的坐标矩阵分别为 $A,B$，则 $B=AP$，从而 $P=A^{-1}B$。可以对 $(A\mid B)$ 进行初等行变换，将左半块化为 $I_n$，右半块便同时化为 $P$，无须分别求出每个新基向量的系数。这里 $A,B$ 必须按列排列基向量。

\subsection{平移幂基与多项式的 Taylor 展开}

\begin{example}
\label{ex:laex-taylor-basis}
在 $\mathcal P_n(\mathbb F)$ 中，取 $E=(1,x,\ldots,x^n)$ 与 $F=(1,x-a,\ldots,(x-a)^n)$。求过渡矩阵，并解释后一组基下的坐标。
\end{example}

\begin{solve}
为简化公式，以 $0,1,\ldots,n$ 为行列指标。二项式定理给出
\[
p_{ij}=\begin{cases}\binom ji(-a)^{j-i},&i\le j,\\0,&i>j.\end{cases}
\]
矩阵 $P$ 为对角元全为 $1$ 的上三角矩阵，因而 $F$ 确为一组基。逆变换来自 $x=(x-a)+a$，故 $P^{-1}$ 的对应元素为 $\binom ji a^{j-i}$，$i\le j$。

设 $f(x)=\sum_{j=0}^n b_j(x-a)^j$。对等式求 $k$ 次导数，再令 $x=a$，只剩 $k!b_k$，得到
\begin{equation}
\label{eq:laex-polynomial-taylor}
f(x)=\sum_{k=0}^n\frac{f^{(k)}(a)}{k!}(x-a)^k.
\end{equation}
所以 Taylor 公式在这里就是一次基变换：新坐标为各阶导数值除以相应阶乘。这是有限次数多项式的恒等式，不涉及无穷级数、收敛性或余项估计。
\end{solve}

\section{子空间与商空间}
\label{sec:laex-subspace-quotient}

\subsection{交、和及其计算}

\begin{theorem}
\label{thm:laex-dimension-sum}
设 $U,W$ 是有限维线性空间 $V$ 的子空间，则
\begin{equation}
\label{eq:laex-dimension-sum}
\dim(U+W)=\dim U+\dim W-\dim(U\cap W).
\end{equation}
\end{theorem}

\begin{proof}
取 $U\cap W$ 的基 $e_1,\ldots,e_k$，分别扩张成 $U$ 的基 $e_1,\ldots,e_k,u_1,\ldots,u_p$ 和 $W$ 的基 $e_1,\ldots,e_k,w_1,\ldots,w_q$。合并且只保留一份公共基，所得向量组生成 $U+W$。若
\[
\sum_i a_ie_i+\sum_j b_ju_j+\sum_l c_lw_l=0,
\]
则 $\sum_l c_lw_l\in U\cap W$，于是 $W$ 中基展开唯一性迫使所有 $c_l=0$，再由 $U$ 的基无关性得到所有 $a_i,b_j=0$。所以合并组是一组基，维数为 $k+p+q$，恰等于右端。
\end{proof}

若 $U=\operatorname{im}A$、$W=\operatorname{im}B$，则 $U+W$ 的基可从 $(A\mid B)$ 的原列中选择。求交空间时解
\begin{equation}
\label{eq:laex-intersection-system}
As=Bt\quad\Longleftrightarrow\quad (A\mid-B)\binom st=0,
\end{equation}
再把解参数映射为公共向量 $As$。当 $A,B$ 的列分别为 $U,W$ 的基时，这个映射在上述系数解空间上是双射，因此基础解系直接对应交空间的基；若两侧只是含冗余的生成组，则不同系数可能给出同一向量，还应剔除映射后的零向量和相关向量。

\begin{example}
\label{ex:laex-intersection}
设
\[
\begin{aligned}
a_1&=(1,0,-1,0)^{\mathsf T},&a_2&=(0,1,2,1)^{\mathsf T},&a_3&=(2,1,0,1)^{\mathsf T},\\
b_1&=(-1,1,1,1)^{\mathsf T},&b_2&=(1,-1,-3,-1)^{\mathsf T},&b_3&=(-1,1,-1,1)^{\mathsf T}.
\end{aligned}
\]
令 $U=\operatorname{span}(a_1,a_2,a_3)$，$W=\operatorname{span}(b_1,b_2,b_3)$，求交与和的基。
\end{example}

\begin{solve}
因为 $a_3=2a_1+a_2$、$b_3=2b_1+b_2$，且 $a_1,a_2$ 和 $b_1,b_2$ 分别无关，故两空间维数都为 $2$。又 $b_2=a_1-a_2\in U$，而 $b_1\notin U$：若用前两个分量确定其系数为 $-1,1$，第三分量应为 $3$，实际却为 $1$。因此
\[
U+W=\operatorname{span}(a_1,a_2,b_1),\qquad
\dim(U\cap W)=2+2-3=1.
\]
所以 $a_1,a_2,b_1$ 是和空间的一组基，$b_2$ 是交空间的一组基。
\end{solve}

\subsection{直和、补空间与有限个真子空间}

\begin{proposition}
\label{prop:laex-direct-sum}
设 $V_1,\ldots,V_k$ 是有限维子空间，$S=V_1+\cdots+V_k$。下列条件等价：每个 $v\in S$ 的分解 $v=v_1+\cdots+v_k$ 唯一；$\sum_i v_i=0$ 只能由所有 $v_i=0$ 得到；对每个 $j\ge2$ 有 $V_j\cap(V_1+\cdots+V_{j-1})=\{0\}$；各空间的基拼接后是 $S$ 的一组基；$\dim S=\sum_i\dim V_i$。成立时记为 $S=\bigoplus_{i=1}^k V_i$。
\end{proposition}

\begin{proof}
两种分解相减就是零向量的分解，故前两个条件等价。在一个非平凡零分解中取最后一个非零分量，便得到相应的交中有非零向量，反向移项也成立。把每个分量按其空间的基展开，零分解的唯一性正好等价于拼接组无关；该组总能生成 $S$，所以又等价于其个数等于 $\dim S$。
\end{proof}

\begin{example}
\label{ex:laex-direct-decompositions}
全体方阵可以唯一写成对称矩阵与反对称矩阵之和；$\mathbb F^n$ 也可以唯一写成常数分量向量与分量和为零的向量之和。
\end{example}

\begin{solve}
对方阵有
\[
A=\frac{A+A^{\mathsf T}}2+\frac{A-A^{\mathsf T}}2.
\]
两项分别对称与反对称，且同时具有两种性质的矩阵只能为零，故为直和。对 $x=(x_1,\ldots,x_n)^{\mathsf T}$，令 $\bar x=(x_1+\cdots+x_n)/n$、$\mathbf1=(1,\ldots,1)^{\mathsf T}$，则
\[
x=\bar x\mathbf1+(x-\bar x\mathbf1).
\]
后一项的分量和为零，常数分量向量若分量和为零也只能是零。因此两空间构成直和，维数分别为 $1$ 和 $n-1$。这里除以 $2$ 和 $n$ 使用了数域的特征为零这一条件。
\end{solve}

\begin{notealert}[type=caution,title={多于两个子空间时不能只检验两两相交为零}]
在 $\mathbb R^2$ 中，$\operatorname{span}(e_1)$、$\operatorname{span}(e_2)$、$\operatorname{span}(e_1+e_2)$ 两两交为零，但三者之和不是直和，因为 $e_1+e_2-(e_1+e_2)=0$ 给出非平凡分解。应检验某一子空间与前面各空间之和的交，或检验维数是否相加。
\end{notealert}

任意子空间 $U\subseteq V$ 都有\emph{补空间（complementary subspace）} $W$，使 $V=U\oplus W$：取 $U$ 的基并扩张为 $V$ 的基，由新增基向量生成 $W$ 即可。补空间通常不唯一，例如 $\mathbb R^2$ 中 $\operatorname{span}(e_1)$ 的补空间可以是任意 $\operatorname{span}(e_2+ae_1)$，$a\in\mathbb R$。补空间与集合补集不同，二者交为 $\{0\}$，而且 $V=U+W$ 表示每个向量都能相加分解，并非 $V=U\cup W$。外直和 $U\times W$ 以逐分量运算定义，当 $U,W$ 已在某空间中构成内直和时，$(u,w)\mapsto u+w$ 就给出两种构造之间的同构。

\begin{proposition}
\label{prop:laex-finite-union}
设 $\dim V=n\ge1$，标量域 $\mathbb F$ 无限。有限个真子空间不能覆盖 $V$；更强地，存在 $V$ 的一组基，其中每个基向量都避开这些真子空间。
\end{proposition}

\begin{proof}
固定基 $e_1,\ldots,e_n$，对每个 $t\in\mathbb F$ 构造 $v(t)=e_1+te_2+\cdots+t^{n-1}e_n$。由例~\ref{ex:laex-vandermonde-basis}，任意 $n$ 个互异参数对应的向量组成一组基，故一个真子空间至多包含其中 $n-1$ 个不同参数对应的向量。有限个真子空间只排除了有限多个参数，从无限域中另选 $n$ 个互异参数就得到所需基。有限域上这一步失效，例如 $\mathbb F_2^2$ 是其三个一维子空间的并。
\end{proof}

\subsection{商空间及其维数}

\begin{definition}
\label{def:laex-quotient}
设 $U$ 是 $V$ 的子空间，$v+U=\{v+u:u\in U\}$ 称为 $U$ 的陪集（coset）。全体陪集构成的集合记为 $V/U$，并定义
\[
(v+U)+(w+U)=(v+w)+U,\qquad
a(v+U)=av+U.
\]
所得线性空间称为 $V$ 关于 $U$ 的\emph{商空间（quotient space）}。
\end{definition}

首先要验证运算不依赖代表元。由 $v+U=w+U$ 当且仅当 $v-w\in U$，若 $v-v'\in U$、$w-w'\in U$，就有 $(v+w)-(v'+w')\in U$ 及 $av-av'\in U$，所以运算良定义。其余线性空间公理由 $V$ 继承，零向量为陪集 $U=0+U$。任意两个陪集或者相等，或者不相交；一个陪集本身是 $V$ 的子空间，当且仅当它就是 $U$。

\begin{theorem}
\label{thm:laex-quotient-complement}
若 $V=U\oplus W$，则 $T:W\to V/U$，$T(w)=w+U$ 是线性同构。特别地，
\[
\dim(V/U)=\dim V-\dim U.
\]
\end{theorem}

\begin{proof}
陪集运算保证 $T$ 线性。若 $T(w)=U$，则 $w\in U\cap W=\{0\}$，故为单射。任意 $v\in V$ 唯一分解为 $u+w$，于是 $v+U=w+U$，故为满射。补空间维数为 $\dim V-\dim U$，从而得到维数公式。若 $e_1,\ldots,e_k$ 是 $U$ 的基，扩张为 $e_1,\ldots,e_n$ 后，$e_{k+1}+U,\ldots,e_n+U$ 就是商空间的一组基。
\end{proof}

商空间把相差一个 $U$ 中向量的元素视为同一个元素，补空间则从每个陪集中选出唯一代表。前者由 $U\subseteq V$ 自然确定，后者依赖选择；因此它们虽然同构，却不是同一个集合。




\section{矩阵的秩}
\label{sec:laex-matrix-rank}

\subsection{基本估计与乘积秩不等式}

秩既可以由主元或非零子式确定，也可以理解为列空间的维数或“未知数个数减去解空间维数”。三种观点对应不同的证明方法：初等变换适合显式分块结构，解空间包含适合比较核，生成空间包含适合比较像。选用哪一种，取决于题目给出的条件最直接地控制了什么。

\begin{proposition}
\label{prop:laex-rank-basic}
在矩阵尺寸相容的条件下，有
\begin{align}
\operatorname{rank}(AB)&\le\min\{\operatorname{rank}A,\operatorname{rank}B\},
\label{eq:laex-product-upper}\\
\operatorname{rank}(A+B)&\le\operatorname{rank}A+\operatorname{rank}B,
\label{eq:laex-sum-upper}\\
\bigl|\operatorname{rank}A-\operatorname{rank}B\bigr|
&\le\operatorname{rank}(A-B).
\label{eq:laex-difference-lower}
\end{align}
此外，可逆矩阵左乘、右乘及非零标量数乘均不改变秩，且
\begin{equation}
\label{eq:laex-block-diagonal-rank}
\operatorname{rank}\begin{pmatrix}A&0\\0&B\end{pmatrix}
=\operatorname{rank}A+\operatorname{rank}B,
\qquad
\operatorname{rank}\begin{pmatrix}A&C\\0&B\end{pmatrix}
\ge\operatorname{rank}A+\operatorname{rank}B.
\end{equation}
\end{proposition}

\begin{proof}
$AB$ 的列是 $A$ 的列的线性组合，故其秩不超过 $\operatorname{rank}A$；把 $B$ 的列用其一个极大无关组表示后，左乘 $A$ 得到 $AB$ 的生成组，故也不超过 $\operatorname{rank}B$。又 $\operatorname{im}(A+B)\subseteq\operatorname{im}A+\operatorname{im}B$，得到和的上界。对 $A=(A-B)+B$ 与 $B=(B-A)+A$ 分别使用该上界，得到差的下界。可逆乘法的结论可由乘积上界在正反两个方向应用得到。

分块对角阵的两组列占据互不重叠的分量，秩相加。对分块上三角阵，分别从 $A,B$ 中选取极大无关列，取它们在整个分块矩阵中的对应列。一个为零的线性组合先看下半部，$B$ 的无关性迫使第二组系数为零，再看上半部，$A$ 的无关性迫使其余系数为零。因此至少有 $\operatorname{rank}A+\operatorname{rank}B$ 个无关列。
\end{proof}

\begin{theorem}
\label{thm:laex-sylvester-frobenius}
若 $A\in M_{m\times n}(\mathbb F)$、$B\in M_{n\times p}(\mathbb F)$，则西尔维斯特（Sylvester）秩不等式为
\begin{equation}
\label{eq:laex-sylvester}
\operatorname{rank}(AB)\ge\operatorname{rank}A+\operatorname{rank}B-n.
\end{equation}
若另有 $C\in M_{p\times q}(\mathbb F)$，则弗罗贝尼乌斯（Frobenius）秩不等式为
\begin{equation}
\label{eq:laex-frobenius}
\operatorname{rank}(ABC)\ge
\operatorname{rank}(AB)+\operatorname{rank}(BC)-\operatorname{rank}B.
\end{equation}
\end{theorem}

\begin{proof}
先证明一个更准确的等式：
\begin{equation}
\label{eq:laex-product-rank-loss}
\operatorname{rank}(AB)
=\operatorname{rank}B-\dim(\operatorname{im}B\cap\ker A).
\end{equation}
取 $\operatorname{im}B\cap\ker A$ 的基并将其扩张为 $\operatorname{im}B$ 的基。交空间的基被 $A$ 送到零，新增基向量的像则无关：否则其某个非零组合落在 $\ker A$ 中，又落在 $\operatorname{im}B$ 中，与扩张后的无关性矛盾。这些像生成 $\operatorname{im}(AB)$，故得到等式。由交空间维数至多为 $\dim\ker A=n-\operatorname{rank}A$，即得 Sylvester 不等式。

分别对 $B$ 和 $BC$ 使用式~\eqref{eq:laex-product-rank-loss}，注意到 $\operatorname{im}(BC)\subseteq\operatorname{im}B$，因此
\[
\begin{aligned}
\operatorname{rank}(AB)-\operatorname{rank}(ABC)
&=\operatorname{rank}B-\operatorname{rank}(BC)\\
&\quad-\dim(\operatorname{im}B\cap\ker A)
+\dim(\operatorname{im}(BC)\cap\ker A)\\
&\le\operatorname{rank}B-\operatorname{rank}(BC).
\end{aligned}
\]
移项就得到 Frobenius 不等式。
\end{proof}

式~\eqref{eq:laex-product-rank-loss} 表明，左乘 $A$ 所损失的维数正是 $B$ 的像中被 $A$ 消去的部分。因此
\begin{equation}
\label{eq:laex-product-rank-equality}
\operatorname{rank}(AB)=\operatorname{rank}B
\quad\Longleftrightarrow\quad
\operatorname{im}B\cap\ker A=\{0\}
\quad\Longleftrightarrow\quad
\ker(AB)=\ker B.
\end{equation}
最后一个等价性也可由 $\ker B\subseteq\ker(AB)$ 与秩与零度公式得到。此时对任何尺寸相容的 $C$，都有 $\operatorname{rank}(ABC)=\operatorname{rank}(BC)$。特别地，$AB=0$ 蕴含 $\operatorname{rank}A+\operatorname{rank}B\le n$。反复应用 Sylvester 不等式还得到，对 $n$ 阶方阵 $A_1,\ldots,A_k$，
\[
\sum_{i=1}^k\operatorname{rank}A_i
\le(k-1)n+\operatorname{rank}(A_1\cdots A_k).
\]

\begin{example}
\label{ex:laex-commuting-rank}
设 $A,B\in M_n(\mathbb F)$ 且 $AB=BA$，证明
\[
\operatorname{rank}(A+B)+\operatorname{rank}(AB)
\le\operatorname{rank}A+\operatorname{rank}B.
\]
\end{example}

\begin{solve}
从列空间看，$\operatorname{im}(A+B)\subseteq\operatorname{im}A+\operatorname{im}B$，而交换条件保证 $\operatorname{im}(AB)\subseteq\operatorname{im}A\cap\operatorname{im}B$。将两个维数上界相加，再用式~\eqref{eq:laex-dimension-sum} 就得到结论。

从解空间看，$\ker A\cap\ker B\subseteq\ker(A+B)$，而 $\ker A+\ker B\subseteq\ker(AB)$，后一包含同样使用 $AB=BA$。将包含关系的维数不等式相加，再用交和维数公式与式~\eqref{eq:laex-rank-nullity}，得到同一结论。两个证明分别刻画列空间与核空间，不能省去交换假设。
\end{solve}

\subsection{分块消元与特殊矩阵}

\begin{proposition}
\label{prop:laex-schur-rank}
设 $M=\begin{pmatrix}A&B\\C&D\end{pmatrix}$，其中 $A$ 为 $r$ 阶可逆方阵，则
\begin{equation}
\label{eq:laex-schur-rank}
\operatorname{rank}M=r+\operatorname{rank}(D-CA^{-1}B).
\end{equation}
矩阵 $D-CA^{-1}B$ 称为关于 $A$ 的舒尔补（Schur complement，亦写作 Schur 补）。若 $D$ 可逆，则也有 $\operatorname{rank}M=\operatorname{rank}D+\operatorname{rank}(A-BD^{-1}C)$。
\end{proposition}

\begin{proof}
以下左右乘矩阵均可逆：
\[
\begin{pmatrix}I&0\\-CA^{-1}&I\end{pmatrix}
\begin{pmatrix}A&B\\C&D\end{pmatrix}
\begin{pmatrix}I&-A^{-1}B\\0&I\end{pmatrix}
=\begin{pmatrix}A&0\\0&D-CA^{-1}B\end{pmatrix}.
\]
再用分块对角阵秩相加公式即可。关于 $D$ 的公式先交换分块行列再同样处理。
\end{proof}

分块初等变换中，乘法顺序由矩阵尺寸决定：左乘系数对应行操作，右乘系数对应列操作。不能像标量一样随意交换各因子。Schur 补公式中的逆也必须由可逆条件保证，不能直接用于奇异分块。

\begin{proposition}
\label{prop:laex-idempotent-involution-rank}
对任意 $A\in M_n(\mathbb F)$，有
\begin{align}
\operatorname{rank}A+\operatorname{rank}(I_n-A)
&=n+\operatorname{rank}(A-A^2),
\label{eq:laex-idempotent-rank}\\
\operatorname{rank}(I_n+A)+\operatorname{rank}(I_n-A)
&=n+\operatorname{rank}(I_n-A^2).
\label{eq:laex-involution-rank}
\end{align}
因此 $A$ 为幂等矩阵（idempotent matrix），即 $A^2=A$，当且仅当第一式左端为 $n$；$A$ 为对合矩阵（involutory matrix），即 $A^2=I_n$，当且仅当第二式左端为 $n$。
\end{proposition}

\begin{proof}
用 $\sim$ 表示相抵变换，有
\[
\begin{pmatrix}A&0\\0&I_n-A\end{pmatrix}
\sim\begin{pmatrix}A&I_n-A\\0&I_n-A\end{pmatrix}
\sim\begin{pmatrix}I_n&I_n-A\\I_n-A&I_n-A\end{pmatrix}
\sim\begin{pmatrix}I_n&0\\0&A-A^2\end{pmatrix}.
\]
最后一步先以第一分块行消去左下块，再消去右上块。比较两侧秩即得第一式。将其中 $A$ 换成 $(I_n+A)/2$，并使用非零数乘不改变秩，得到第二式。
\end{proof}

\subsection{由核空间证明秩与非奇异性}

\begin{proposition}
\label{prop:laex-gram-rank}
若 $A$ 为实矩阵，则
\[
\operatorname{rank}(A^{\mathsf T}A)
=\operatorname{rank}(AA^{\mathsf T})=\operatorname{rank}A.
\]
若 $A$ 为复矩阵，则相应结论应写成 $\operatorname{rank}(A^{\mathsf H}A)=\operatorname{rank}(AA^{\mathsf H})=\operatorname{rank}A$。
\end{proposition}

\begin{proof}
实数情形中，$Ax=0$ 蕴含 $A^{\mathsf T}Ax=0$。反过来，若后者成立，则
\[
0=x^{\mathsf T}A^{\mathsf T}Ax
=(Ax)^{\mathsf T}(Ax)=\sum_i(Ax)_i^2,
\]
每一项非负，所以 $Ax=0$。两矩阵的核相同、列数相同，故秩相等；对 $A^{\mathsf T}$ 使用同一结论得到另一等式。复数情形改用 $x^{\mathsf H}A^{\mathsf H}Ax=\sum_i|(Ax)_i|^2$ 即可。若复数情形仍用转置，结论会失效，例如 $A=(1,\mathrm i)^{\mathsf T}$ 的秩为 $1$，但 $A^{\mathsf T}A=0$。
\end{proof}

\begin{example}
\label{ex:laex-diagonal-dominance}
设 $A=(a_{ij})\in M_n(\mathbb C)$ 满足严格行对角占优条件
\[
|a_{ii}|>\sum_{j\ne i}|a_{ij}|\quad(1\le i\le n).
\]
证明 $A$ 可逆。若 $A$ 为实矩阵且进一步有 $a_{ii}>\sum_{j\ne i}|a_{ij}|$，证明 $\det A>0$。
\end{example}

\begin{solve}
若 $Ax=0$ 有非零解，取 $|x_k|=\max_j|x_j|>0$。第 $k$ 个方程给出
\[
|a_{kk}|\,|x_k|
\le\sum_{j\ne k}|a_{kj}|\,|x_j|
\le\left(\sum_{j\ne k}|a_{kj}|\right)|x_k|,
\]
与严格不等式矛盾，故只有零解，$A$ 可逆。实数正对角元情形令 $D=\operatorname{diag}(a_{11},\ldots,a_{nn})$，考虑 $A(t)=D+t(A-D)$，$0\le t\le1$。路径上始终严格对角占优，行列式连续且处处非零，而 $\det A(0)=\prod_i a_{ii}>0$，由介值定理知其不变号，因此 $\det A>0$。
\end{solve}

\begin{example}
\label{ex:laex-skew-positive}
设 $A$ 为实反对称矩阵，$D$ 为对角元全正的实对角矩阵，证明 $\det(D+A)>0$。
\end{example}

\begin{solve}
若 $(D+tA)x=0$，其中 $t\in\mathbb R$、$x\in\mathbb R^n$，则 $x^{\mathsf T}Ax=0$，故 $0=x^{\mathsf T}Dx=\sum_i d_ix_i^2$，只能 $x=0$。于是实矩阵 $D+tA$ 对所有实数 $t$ 可逆，其行列式是处处非零的连续函数，在 $t=0$ 处为正，因而在 $t=1$ 处也为正。特别地，$\det(I_n\pm A)>0$。
\end{solve}

\subsection{子式、加边子式与伴随矩阵}

\begin{proposition}
\label{prop:laex-bordered-minors}
矩阵 $A$ 的秩为 $r\ge1$，当且仅当存在一个非零的 $r$ 阶子式，且所有 $r+1$ 阶子式都为零。更强地，固定一个非零 $r$ 阶子式后，只需检查包含它的全部 $r+1$ 阶\emph{加边子式（bordered minor）}为零。
\end{proposition}

\begin{proof}
将所选子式移到左上角，写成 $A=\begin{pmatrix}D&B\\C&E\end{pmatrix}$，其中 $D$ 可逆。每个相应加边子式都是
\[
\det\begin{pmatrix}D&b_j\\c_i&e_{ij}\end{pmatrix}
=\det D\,(e_{ij}-c_iD^{-1}b_j).
\]
因此这些加边子式全为零，当且仅当 Schur 补 $E-CD^{-1}B=0$，由式~\eqref{eq:laex-schur-rank} 即得秩为 $r$。必要性来自：任何阶数大于秩的方形子矩阵，其列不可能无关。若秩大于 $r$，从 $r$ 个独立列和相应独立行中可取得更大的非零子式，或直接由上述 Schur 补的非零元素得到非零加边子式。$r=0$ 的情形单独等价于 $A=0$。
\end{proof}

如果分别选取一个秩为 $r$ 的矩阵的 $r$ 个独立行和 $r$ 个独立列，它们交叉处的 $r$ 阶子矩阵必可逆。证明时只需注意所有行都是所选行的线性组合：如果所选列在这些行上的缩短列相关，同一关系就会在全部行上成立，与所选列无关矛盾。对称矩阵或反对称矩阵中，一组独立行的同指标列也独立，因此必有一个阶数等于秩的非零主子式。反对称矩阵的奇数阶行列式满足 $\det D=\det(-D^{\mathsf T})=-\det D$，只能为零，于是\emph{反对称矩阵的秩必为偶数}。

\begin{proposition}
\label{prop:laex-principal-border}
设对称矩阵 $A$ 有一个非零 $r$ 阶主子式 $\det D$，且所有包含它的 $r+1$ 阶及 $r+2$ 阶加边主子式均为零，则 $\operatorname{rank}A=r$。若 $A$ 为反对称矩阵，则只需检查包含该主子式的全部 $r+2$ 阶加边主子式为零。
\end{proposition}

\begin{proof}
将 $D$ 同时置于左上角，并记 Schur 补为 $S$。对称情形下 $S$ 仍对称，$r+1$ 阶加边主子式为零使 $s_{ii}=0$，$r+2$ 阶加边主子式为零再给出
\[
0=\det\begin{pmatrix}0&s_{ij}\\s_{ij}&0\end{pmatrix}=-s_{ij}^2,
\]
故 $S=0$。反对称情形下 $S$ 仍反对称，已自动有 $s_{ii}=0$，相应二阶主子式等于 $s_{ij}^2$，同样迫使 $S=0$。两种情形都由 Schur 补秩公式得到结论；若剩余维数小于二，只需应用存在的相应条件。
\end{proof}

\begin{proposition}
\label{prop:laex-adjugate-rank}
设 $A\in M_n(\mathbb F)$ 且 $n\ge2$，则
\[
\operatorname{rank}\operatorname{adj}(A)=
\begin{cases}
n,&\operatorname{rank}A=n,\\
1,&\operatorname{rank}A=n-1,\\
0,&\operatorname{rank}A\le n-2.
\end{cases}
\]
\end{proposition}

\begin{proof}
若 $A$ 可逆，$\operatorname{adj}(A)=(\det A)A^{-1}$，故满秩。若 $\operatorname{rank}A=n-1$，至少有一个 $n-1$ 阶子式非零，故伴随矩阵非零；又 $A\operatorname{adj}(A)=0$，使其全部列都落在一维核空间中，故秩为 $1$。若 $\operatorname{rank}A\le n-2$，全部 $n-1$ 阶子式为零，伴随矩阵即为零。
\end{proof}

当 $\operatorname{rank}A=n-1$ 时，伴随矩阵的任意非零列就构成 $Ax=0$ 的基础解系。还可用这一点证明：若矩阵每行、每列的元素和都为零，则所有代数余子式相等。秩低于 $n-1$ 时它们全为零；秩为 $n-1$ 时，右核与左核都由 $\mathbf1$ 生成，伴随矩阵的各列与各行都相应成比例，所以它只能是 $c\mathbf1\mathbf1^{\mathsf T}$ 的形式。

\section{相抵标准型及其应用}
\label{sec:laex-equivalence-form}

\subsection{标准型、单侧逆与满秩分解}

\begin{theorem}
\label{thm:laex-equivalence-form}
若 $A\in M_{m\times n}(\mathbb F)$ 的秩为 $r$，则存在可逆矩阵 $P\in M_m(\mathbb F)$、$Q\in M_n(\mathbb F)$，使
\begin{equation}
\label{eq:laex-equivalence-form}
PAQ=J_r=\begin{pmatrix}I_r&0\\0&0\end{pmatrix}.
\end{equation}
同型矩阵相抵，当且仅当秩相同。
\end{theorem}

\begin{proof}
先用行变换得到阶梯形，再用列交换把主元列移到前面，以可逆行列变换把对应可逆分块化为 $I_r$ 并消去其余块，即得到 $J_r$。反过来，左右乘可逆矩阵保持秩，所以同型矩阵有相同标准型当且仅当它们的秩相同。
\end{proof}

相抵允许独立选取左右两个可逆矩阵；相似变换则要求它们互为逆矩阵。因此相抵标准型用于秩问题，不能据此断言矩阵相似于 $J_r$，也不能直接从相抵标准型读出原矩阵的迹或特征值。

\begin{proposition}
\label{prop:laex-one-sided-inverse}
$A\in M_{m\times n}(\mathbb F)$ 存在左逆 $L\in M_{n\times m}(\mathbb F)$、$LA=I_n$，当且仅当 $\operatorname{rank}A=n$；存在右逆 $R\in M_{n\times m}(\mathbb F)$、$AR=I_m$，当且仅当 $\operatorname{rank}A=m$。列满秩矩阵适合左消去律，行满秩矩阵适合右消去律。
\end{proposition}

\begin{proof}
存在单侧逆时由乘积秩上界得到相应满秩性。若列满秩，则标准型为 $PAQ=\binom{I_n}{0}$，取 $L=Q(I_n\mid0)P$ 即有 $LA=I_n$。若行满秩，则 $PAQ=(I_m\mid0)$，取 $R=Q\binom{I_m}{0}P$ 即可。由 $AX=AY$ 左乘 $L$ 得 $X=Y$，右消去律同理。
\end{proof}

\begin{theorem}
\label{thm:laex-rank-factorization}
设 $A\in M_{m\times n}(\mathbb F)$ 的秩为 $r>0$，则存在列满秩的 $B\in M_{m\times r}(\mathbb F)$ 与行满秩的 $C\in M_{r\times n}(\mathbb F)$，使 $A=BC$，称为\emph{满秩分解（rank factorization）}。若还有 $A=B_1C_1=B_2C_2$ 两个满秩分解，则存在 $r$ 阶可逆矩阵 $S$，使
\[
B_2=B_1S,\qquad C_2=S^{-1}C_1.
\]
\end{theorem}

\begin{proof}
从 $A$ 的列中选一个极大无关组组成 $B$，把每一列在该组下的表示系数组成 $C$，便得到 $A=BC$。$C$ 在所选列的对应位置含 $I_r$，所以行满秩。对两个分解，$\operatorname{im}B_1=\operatorname{im}A=\operatorname{im}B_2$，故 $B_1,B_2$ 的列是同一空间的两组基，存在可逆过渡矩阵 $S$ 使 $B_2=B_1S$。代回 $B_1C_1=B_1SC_2$，左消去 $B_1$，得到 $C_2=S^{-1}C_1$。
\end{proof}

满秩分解还表明，秩为 $r$ 的矩阵恰可写成 $r$ 个秩一矩阵之和，且不能用更少个秩一矩阵相加得到。若 $B=(b_1,\ldots,b_r)$，$C$ 的各行为 $c_i^{\mathsf T}$，则 $A=\sum_i b_ic_i^{\mathsf T}$，每项的两因子均非零，秩为 $1$；若少于 $r$ 项，则秩的次可加性会给出矛盾。零矩阵对应空和。

\begin{example}
\label{ex:laex-generalized-inverse}
对任意 $A\in M_{m\times n}(\mathbb F)$，证明存在 $G\in M_{n\times m}(\mathbb F)$，使 $AGA=A$。
\end{example}

\begin{solve}
当 $A=0$ 时取 $G=0$。否则取满秩分解 $A=BC$，由命题~\ref{prop:laex-one-sided-inverse} 选取 $B$ 的左逆 $L$ 与 $C$ 的右逆 $R$，令 $G=RL$，则
\[
AGA=BCRLBC=B(CR)(LB)C=BC=A.
\]
这里依次使用 $CR=I_r$ 与 $LB=I_r$。这个条件称为一种广义逆条件（generalized inverse condition），不要求 $AG$ 或 $GA$ 为单位矩阵，也不保证 $G$ 唯一。
\end{solve}

\subsection{幂等矩阵与矩阵方程}

\begin{proposition}
\label{prop:laex-idempotent-factorization}
设 $A=BC$ 是秩为 $r>0$ 的方阵 $A$ 的满秩分解，则 $A^2=A$ 当且仅当 $CB=I_r$。此时 $\operatorname{tr}A=r$，且 $\mathbb F^n=\operatorname{im}A\oplus\ker A$。
\end{proposition}

\begin{proof}
由 $A^2=A$ 得 $B(CB-I_r)C=0$，用 $B$ 的左逆和 $C$ 的右逆消去两侧即得 $CB=I_r$；反向代入即可。逐元素求和给出 $\operatorname{tr}(BC)=\operatorname{tr}(CB)=r$。任意 $x$ 分解为 $Ax+(x-Ax)$，两项分别属于像和核；若 $v=Au$ 且 $Av=0$，又有 $Av=A^2u=Au=v$，所以交只有零向量。$A=0$ 时迹和直和结论也成立。
\end{proof}

\begin{example}
\label{ex:laex-rectangular-product}
设 $A$ 为 $3\times2$ 矩阵，$B$ 为 $2\times3$ 矩阵，且
\[
AB=C=\begin{pmatrix}8&2&-2\\2&5&4\\-2&4&5\end{pmatrix}.
\]
求 $BA$。
\end{example}

\begin{solve}
直接计算得到 $\operatorname{rank}C=2$ 及 $C^2=9C$。由乘积秩上界与尺寸限制知 $A$ 列满秩、$B$ 行满秩。于是 $A(BA-9I_2)B=0$，用单侧逆消去 $A,B$，得到 $BA=9I_2$。矩形因子的消去依赖满秩条件，不能仅从 $A,B$ 非零推断可消去。
\end{solve}

\begin{proposition}
\label{prop:laex-two-sided-equation}
设 $A\in M_{m\times n}(\mathbb F)$、$B\in M_{p\times q}(\mathbb F)$、$C\in M_{m\times q}(\mathbb F)$。矩阵方程
\[
AX+YB=C,\qquad
X\in M_{n\times q}(\mathbb F),\quad Y\in M_{m\times p}(\mathbb F),
\]
有解，当且仅当
\[
\operatorname{rank}\begin{pmatrix}A&C\\0&B\end{pmatrix}
=\operatorname{rank}A+\operatorname{rank}B.
\]
\end{proposition}

\begin{proof}
若有解，以第一分块列右乘 $-X$ 加到第二分块列，再以第二分块行左乘 $-Y$ 加到第一分块行，即得到分块对角阵，故秩相加。

反过来，将 $A,B$ 分别作相抵变换化为 $J_r,J_s$。若 $A'=PAQ$、$B'=RBS$、$C'=PCS$，则 $X'=Q^{-1}XS$、$Y'=PYR^{-1}$ 使新旧方程等价，相关分块矩阵的秩也不变。因此只需处理标准型。按前 $r$ 行和前 $s$ 列将 $C'$ 分块为 $\begin{pmatrix}C_{11}&C_{12}\\C_{21}&C_{22}\end{pmatrix}$，消去 $J_r,J_s$ 对应的块后，整个矩阵秩为 $r+s+\operatorname{rank}C_{22}$。假设使 $C_{22}=0$，而 $J_rX'$ 可以任取前 $r$ 行，$Y'J_s$ 可以任取前 $s$ 列，故可令前者覆盖 $C_{11},C_{12}$，后者覆盖 $C_{21}$，构造出解。
\end{proof}

\section{线性方程组的解及其应用}
\label{sec:laex-linear-systems}

\subsection{可解性、解的结构与参数讨论}

\begin{theorem}
\label{thm:laex-system-structure}
设 $A\in M_{m\times n}(\mathbb F)$、$b\in\mathbb F^m$、$\operatorname{rank}A=r$。方程组 $Ax=b$ 有解，当且仅当 $\operatorname{rank}(A\mid b)=r$。有解时，若 $x_0$ 是一个特解，$\eta_1,\ldots,\eta_{n-r}$ 是 $\ker A$ 的一组基，则全部解为
\begin{equation}
\label{eq:laex-affine-solution}
x=x_0+\sum_{i=1}^{n-r}t_i\eta_i,\qquad t_i\in\mathbb F.
\end{equation}
其中 $\eta_1,\ldots,\eta_{n-r}$ 称为齐次方程组 $Ax=0$ 的\emph{基础解系（fundamental system of solutions）}。当 $r=n$ 时解唯一，当 $r<n$ 时解有无穷多个。
\end{theorem}

\begin{proof}
可解性等价于 $b$ 属于 $A$ 的列空间，即增添 $b$ 不增加列秩。若 $x$ 与 $x_0$ 都为解，则 $A(x-x_0)=0$；反过来，特解加上任意核向量仍为解。由式~\eqref{eq:laex-rank-nullity}，核空间维数为 $n-r$，因此得到全部解的表达式。唯一性等价于核为零；核非零时，由数域无限可得到无穷多个解。
\end{proof}

非齐次解集 $x_0+\ker A$ 是一个\emph{仿射子空间（affine subspace）}，当 $b\ne0$ 时不含零向量，因而不是线性子空间。若 $x^{(1)},\ldots,x^{(k)}$ 都是解，则
\[
A\left(\sum_i c_ix^{(i)}\right)=\left(\sum_i c_i\right)b.
\]
当 $b\ne0$ 时，该组合仍为解当且仅当 $\sum_i c_i=1$；齐次情形则对系数没有这一限制。基础解系也不是任取 $n-r$ 个解向量，而必须是 $n-r$ 个线性无关的齐次解。两个基础解系互相通过一个可逆的 $n-r$ 阶系数矩阵联系，本质上就是同一核空间的基变换。

带参数的消元尤其要区分\emph[caution]{可能为零的主元}。若某一步拟除以 $g(\lambda)$，必须先分成 $g(\lambda)=0$ 与 $g(\lambda)\ne0$ 两种情况，零点处回到原矩阵重新计算；不能把一般参数情形的阶梯形直接代入所有特殊值。另一种办法是利用结构找出参数导致秩变化的位置，再逐一比较系数矩阵和增广矩阵的秩。

\begin{example}
\label{ex:laex-parameter-consistency}
讨论含参数 $\lambda\in\mathbb F$ 的方程组
\[
\begin{cases}
2x_1+3x_2+x_3+x_4=1,\\
x_1+2x_2-x_3+4x_4=2,\\
x_1+3x_2-4x_3+11x_4=\lambda.
\end{cases}
\]
\end{example}

\begin{solve}
第三个系数行等于第二行的三倍减去第一行，因此有解必须且只需 $\lambda=3\cdot2-1=5$。前两行独立，系数矩阵秩为 $2$。当 $\lambda=5$ 时令 $x_3=s$、$x_4=t$，得到
\[
x=\begin{pmatrix}-4\\3\\0\\0\end{pmatrix}
+s\begin{pmatrix}-5\\3\\1\\0\end{pmatrix}
+t\begin{pmatrix}10\\-7\\0\\1\end{pmatrix},\qquad s,t\in\mathbb F.
\]
当 $\lambda\ne5$ 时，增广矩阵秩为 $3$ 而系数矩阵秩为 $2$，所以无解。
\end{solve}

\begin{example}
\label{ex:laex-parameter-symmetric-system}
讨论参数 $k\in\mathbb F$ 下的方程组
\[
\begin{cases}
kx_1+x_2+x_3=-2,\\
x_1+kx_2+x_3=-2,\\
x_1+x_2+kx_3=-2.
\end{cases}
\]
\end{example}

\begin{solve}
任意两式相减得到 $(k-1)(x_i-x_j)=0$，三式相加得到 $(k+2)(x_1+x_2+x_3)=-6$。若 $k\ne1,-2$，则 $x_1=x_2=x_3=-2/(k+2)$，为唯一解。若 $k=1$，三个方程相同，通解为
\[
x=\begin{pmatrix}-2\\0\\0\end{pmatrix}
+s\begin{pmatrix}-1\\1\\0\end{pmatrix}
+t\begin{pmatrix}-1\\0\\1\end{pmatrix},\qquad s,t\in\mathbb F.
\]
若 $k=-2$，三式相加得到 $0=-6$，故无解。这里两个特殊值分别对应失去差分约束与出现不相容的总和约束。
\end{solve}

\subsection{解空间、行空间与可解性的另一种判据}

\begin{proposition}
\label{prop:laex-kernel-rowspace}
设 $A\in M_{m\times n}(\mathbb F)$，$c\in\mathbb F^n$。若 $Ax=0$ 的每个解都满足 $c^{\mathsf T}x=0$，则 $c^{\mathsf T}$ 是 $A$ 的行向量的线性组合；反之亦然。因此，两个含有相同未知数的齐次方程组同解，当且仅当其系数矩阵的行空间相同。
\end{proposition}

\begin{proof}
向 $A$ 添加行 $c^{\mathsf T}$ 得到矩阵 $D$。题设等价于 $\ker D=\ker A$，由秩与零度公式得到 $\operatorname{rank}D=\operatorname{rank}A$。添加的一行不能增加行空间维数，所以它已在原行空间中。反向结论只需对已有各方程取相应线性组合。同解的两个方程组逐行使用上述结论，即得到行空间互相包含。
\end{proof}

\begin{theorem}
\label{thm:laex-solvability-alternative}
设 $A\in M_{m\times n}(\mathbb F)$、$b\in\mathbb F^m$。方程组 $Ax=b$ 有解，当且仅当
\begin{equation}
\label{eq:laex-solvability-dual}
A^{\mathsf T}y=0\quad\Longrightarrow\quad b^{\mathsf T}y=0.
\end{equation}
等价地，下面两个方程组恰有一个可解：
\[
Ax=b;\qquad
A^{\mathsf T}y=0,\quad b^{\mathsf T}y=1.
\]
\end{theorem}

\begin{proof}
若 $Ax=b$，则对任意 $A^{\mathsf T}y=0$ 有 $b^{\mathsf T}y=x^{\mathsf T}A^{\mathsf T}y=0$。反过来，将命题~\ref{prop:laex-kernel-rowspace} 用于矩阵 $A^{\mathsf T}$，题设就使 $b^{\mathsf T}$ 属于其行空间，即 $b\in\operatorname{im}A$，故有解。若 $Ax=b$ 无解，上述蕴含失效，存在 $A^{\mathsf T}y=0$ 且 $b^{\mathsf T}y\ne0$，将 $y$ 除以这个非零标量得到第二组的解；两组同时有解则会推出 $0=1$。
\end{proof}

这个判据使用的是代数配对 $b^{\mathsf T}y$，在任意数域上都成立，不依赖正定性。在实数域可解释为 $b$ 与左核中的每个向量正交；在复数域若采用通常的 Hermite 内积，则应相应将条件写成 $A^{\mathsf H}y=0\Rightarrow y^{\mathsf H}b=0$，不能把两种记号随意混合。

任意子空间 $U\subseteq\mathbb F^n$ 都可以表示为某个齐次方程组的解空间。若基矩阵 $B\in M_{n\times r}(\mathbb F)$ 满足 $U=\operatorname{im}B$，先求 $B^{\mathsf T}y=0$ 的基础解系，并把其转置排成矩阵 $A$ 的各行，则 $AB=0$、$\operatorname{rank}A=n-r$。于是 $U\subseteq\ker A$ 且两空间维数同为 $r$，所以 $U=\ker A$。$r=n$ 时可取 $A$ 为零矩阵；$r=0$ 时可取 $A=I_n$。这一构造把“生成式”与“方程式”两种子空间描述联系起来。

\subsection{矩阵方程与公共解}

\begin{proposition}
\label{prop:laex-matrix-system}
给定 $A\in M_{m\times n}(\mathbb F)$、$B\in M_{m\times p}(\mathbb F)$，矩阵方程 $AX=B$ 有解，当且仅当 $\operatorname{rank}(A\mid B)=\operatorname{rank}A$。若 $X_0$ 是一个解，$N$ 的列组成 $\ker A$ 的一组基，则全部解为
\[
X=X_0+NZ,\qquad Z\in M_{(n-r)\times p}(\mathbb F),\quad r=\operatorname{rank}A.
\]
\end{proposition}

\begin{proof}
按列写成 $Ax_j=b_j$，可解性等价于 $B$ 的每一列都落在 $A$ 的列空间中。每一列的齐次部分可唯一写成 $Nz_j$，将这些系数列拼成 $Z$ 即得通解。若 $r=n$，齐次部分为空，解唯一。
\end{proof}

例如，对尺寸相容的 $A,B$，存在 $C$ 使 $ABC=A$，当且仅当 $\operatorname{rank}(AB)=\operatorname{rank}A$。必要性由乘积秩上界得到；充分性由 $\operatorname{im}(AB)\subseteq\operatorname{im}A$ 且维数相同得到两像空间相等，再逐列求解 $(AB)C=A$。

齐次方程组 $Ax=0$、$Bx=0$ 的公共解空间为
\[
\ker A\cap\ker B=\ker\binom AB,
\]
所以公共解空间维数为 $n-\operatorname{rank}\binom AB$。若只给出两个非齐次方程组的通解
\[
x=p+Nu,\qquad x=q+Mv,
\]
则求公共解等价于对参数求解
\begin{equation}
\label{eq:laex-common-affine}
(N\mid-M)\binom uv=q-p.
\end{equation}
解出后代回任一通解即可。若存在一个公共解 $x_*$，全部公共解构成 $x_*+(\operatorname{im}N\cap\operatorname{im}M)$；否则交集为空。两个非齐次解集即使各自非空，也可能没有公共解，因此必须先检查参数方程的相容性。

\subsection{解析几何中的秩判定}

设 $p_1,\ldots,p_s\in\mathbb R^d$，其仿射包是
\[
p_1+\operatorname{span}(p_2-p_1,\ldots,p_s-p_1),
\]
所以仿射维数就是差向量矩阵的秩。特别地，平面点共线当且仅当差向量的秩至多为 $1$，空间点共面当且仅当差向量的秩至多为 $2$。这里使用“至多”，允许点重合或落在更低维的集合中。若把点坐标作为列并在底部补一行 $1$，则其矩阵秩等于仿射维数加 $1$：将其余列减去第一列后，差向量的末分量为零，而第一列末分量为 $1$，恰多出一维。

\begin{example}
\label{ex:laex-circle}
设平面上三点 $(x_i,y_i)$，$1\le i\le3$，不共线，求经过它们的圆的方程，并给出第四点与它们共圆的判据。
\end{example}

\begin{solve}
圆可设为 $x^2+y^2+ux+vy+w=0$，代入三点得到关于 $u,v,w$ 的线性方程组。不共线保证系数矩阵 $(x_i,y_i,1)_{i=1}^3$ 可逆，故唯一确定圆，等价的行列式形式为
\begin{equation}
\label{eq:laex-circle-determinant}
\det\begin{pmatrix}
x^2+y^2&x&y&1\\
x_1^2+y_1^2&x_1&y_1&1\\
x_2^2+y_2^2&x_2&y_2&1\\
x_3^2+y_3^2&x_3&y_3&1
\end{pmatrix}=0.
\end{equation}
第一项的系数不为零，所以不会退化为直线；配方后半径也必为正，否则不能经过三个不共线实点。第四点共圆当且仅当将其坐标代入式~\eqref{eq:laex-circle-determinant} 后行列式为零。若原来的点全共线，这个行列式条件可能描述退化情形，不能无条件当作圆的判据。
\end{solve}

\begin{example}
\label{ex:laex-conic-pencil}
设平面四个不同点中任意三点不共线，$q_1,q_2$ 是经过它们的两个线性无关的次数至多为二的多项式。证明任意经过这四点的次数至多为二的多项式都可写为 $\lambda q_1+\mu q_2$。
\end{example}

\begin{solve}
全体次数至多为二的二元多项式以 $x^2,xy,y^2,x,y,1$ 为基，维数为 $6$。四点取值给出四个线性条件。为证明这些条件独立，对每个点 $p_i$，选另外三点中的两点确定一条直线，再取经过剩余一点且不经过 $p_i$ 的另一条直线。由于任意三点不共线，两条直线方程的乘积在另外三点为零而在 $p_i$ 非零。适当缩放后，四点的取值为第 $i$ 个标准单位向量，故取值映射的秩为 $4$，核维数为 $2$。$q_1,q_2$ 是此核中的两个无关向量，因而构成基，得到结论。

几何上，$(\lambda,\mu)\ne(0,0)$ 给出过四点的二次曲线束，其中可能含退化曲线。上述点位条件是保证本论证成立的充分条件；若四点全部共线，二次多项式在该直线上的限制只有三个自由系数，四个取值条件便不独立。仅说“两条不同二次曲线有四个公共点”也不足以排除共享直线分支，故应明确这里的点位与无关性条件。
\end{solve}

\section{基础训练：综合方法与易错辨析}
\label{sec:laex-training}

本节将原书基础训练中的代表性问题按方法整合。计算题首先检查秩与尺寸是否匹配，证明题首先寻找可以控制的生成空间、核空间或可逆分块；选择题中的错误选项则常由一个低维反例排除。例如行列式为零只保证行向量组相关，并不保证有零行或有两行成比例；方程比未知数多也并不必然无解；非齐次方程组行满秩则对每个常数向量都有解，但是否唯一还取决于列数。

\begin{example}
\label{ex:laex-independent-functions}
设 $a_1,\ldots,a_n$ 是互异实数，证明 $e^{a_1x},\ldots,e^{a_nx}$ 在实函数空间中线性无关，并说明三角函数组 $1,\sin x,\cos x,\ldots,\sin nx,\cos nx$ 为何也无关。
\end{example}

\begin{solve}
若 $\sum_j c_je^{a_jx}=0$ 对所有实数 $x$ 成立，分别求 $k=0,\ldots,n-1$ 阶导数并令 $x=0$，得到 $\sum_j a_j^kc_j=0$。系数矩阵为非奇异范德蒙德矩阵，因此所有 $c_j=0$。这里只对有限和求导，没有交换无穷过程。

对三角函数，若它们的线性组合恒为零，依次与 $1,\sin kx,\cos kx$ 相乘并在 $[-\pi,\pi]$ 上积分。由积化和差公式可得不同函数乘积的积分为零，而 $1$ 的平方积分为 $2\pi$，其余各函数的平方积分为 $\pi$，于是各系数也逐个为零。所用正交关系在此区间上成立，不能只凭不同函数的外观认定无关。
\end{solve}

\begin{example}
\label{ex:laex-outer-products}
设 $a_1,\ldots,a_p\in\mathbb F^m$ 无关，$b_1,\ldots,b_q\in\mathbb F^n$ 无关，证明 $pq$ 个矩阵 $a_ib_j^{\mathsf T}$ 线性无关。
\end{example}

\begin{solve}
令 $A=(a_1,\ldots,a_p)$、$B=(b_1,\ldots,b_q)$。若 $\sum_{i,j}c_{ij}a_ib_j^{\mathsf T}=0$，写成矩阵即 $ACB^{\mathsf T}=0$，其中 $C=(c_{ij})$。$A$ 列满秩，$B^{\mathsf T}$ 行满秩，故可用单侧逆先后消去它们，得到 $C=0$，所以全部矩阵无关。若缺少任一原向量组的无关性，就不能实施相应消去。
\end{solve}

\begin{example}
\label{ex:laex-biquadratic-field}
求 $\mathbb Q(\sqrt2,\sqrt3)=\{a+b\sqrt2+c\sqrt3+d\sqrt6:a,b,c,d\in\mathbb Q\}$ 在 $\mathbb Q$ 上的维数。
\end{example}

\begin{solve}
先令 $K=\mathbb Q(\sqrt2)$。它以 $1,\sqrt2$ 为有理基，且 $\sqrt3\notin K$：若 $\sqrt3=u+v\sqrt2$，平方并比较有理部分与 $\sqrt2$ 部分得到 $uv=0$；若 $v=0$，则 $u^2=3$，若 $u=0$，则 $v^2=3/2$，两者都不可能由有理数满足，依据是有理数平方中每个素因子的指数均为偶数。因此 $1,\sqrt3$ 在 $K$ 上无关。集合 $K+K\sqrt3$ 对乘法封闭，而且非零元的逆为
\[
(u+v\sqrt3)^{-1}=\frac{u-v\sqrt3}{u^2-3v^2}\quad(u,v\in K),
\]
分母非零，否则会推出 $\sqrt3\in K$ 或 $u=v=0$。它是 $K$ 上维数为 $2$ 的数域，由命题~\ref{prop:laex-scalar-restriction} 得到有理维数为 $4$，基为 $1,\sqrt2,\sqrt3,\sqrt6$。
\end{solve}

\begin{example}
\label{ex:laex-telescoping-rank}
证明 $n$ 阶方阵满足
\[
\operatorname{rank}(AB-I_n)
\le\operatorname{rank}(A-I_n)+\operatorname{rank}(B-I_n).
\]
若 $B_1,\ldots,B_k$ 都为 $n$ 阶幂等矩阵，且 $A=B_1\cdots B_k$，进一步证明
\[
\operatorname{rank}(I_n-A)\le k\bigl(n-\operatorname{rank}A\bigr).
\]
\end{example}

\begin{solve}
第一式由 $AB-I_n=(A-I_n)B+(B-I_n)$ 及秩的乘积上界、和的上界直接得到。第二式使用有限乘积的望远镜分解
\[
I_n-A=\sum_{j=1}^k B_1\cdots B_{j-1}(I_n-B_j),
\]
其中 $j=1$ 时前面的空乘积为 $I_n$。因为 $B_j$ 幂等，式~\eqref{eq:laex-idempotent-rank} 给出 $\operatorname{rank}(I_n-B_j)=n-\operatorname{rank}B_j$，而乘积秩上界给出 $\operatorname{rank}B_j\ge\operatorname{rank}A$。于是各项秩至多为 $n-\operatorname{rank}A$，相加得到所需估计。
\end{solve}

\begin{example}
\label{ex:laex-sum-idempotents}
设 $P=A+B$ 为幂等矩阵，且 $\operatorname{rank}P=\operatorname{rank}A+\operatorname{rank}B$，证明 $A^2=A$、$B^2=B$、$AB=BA=0$。
\end{example}

\begin{solve}
由 $\operatorname{im}P\subseteq\operatorname{im}A+\operatorname{im}B$ 与维数上界知
\[
\operatorname{im}P=\operatorname{im}A\oplus\operatorname{im}B.
\]
幂等矩阵 $P$ 在其像上作用为恒等，所以 $PA=A$、$PB=B$。由前式得到 $(A^2-A)x=-BAx$，左端属于 $\operatorname{im}A$，右端属于 $\operatorname{im}B$，两像交为零，故 $A^2=A$、$BA=0$；对后式同理得到 $B^2=B$、$AB=0$。仅有 $A+B$ 幂等而没有秩相加条件时，不能推出这些结论。
\end{solve}

\begin{example}
\label{ex:laex-normal-rank-stability}
设实方阵 $A$ 满足 $AA^{\mathsf T}=a^2I_n$，其中 $a\in\mathbb R$，证明
\[
\operatorname{rank}\bigl((aI_n-A)^2\bigr)=\operatorname{rank}(aI_n-A).
\]
\end{example}

\begin{solve}
若 $a=0$，由命题~\ref{prop:laex-gram-rank} 知 $A=0$，结论成立。若 $a\ne0$，则 $A$ 可逆，从 $AA^{\mathsf T}=a^2I_n$ 得到 $A^{\mathsf T}=a^2A^{-1}$，因而 $A^{\mathsf T}A=a^2I_n$。令 $M=aI_n-A$，可得 $MM^{\mathsf T}=M^{\mathsf T}M$。以 $\|x\|^2=x^{\mathsf T}x$ 表示实向量的欧几里得范数的平方，则 $\|Mx\|^2=\|M^{\mathsf T}x\|^2$，故 $\ker M=\ker M^{\mathsf T}$。若 $M^2x=0$，令 $v=Mx$，则 $v\in\ker M=\ker M^{\mathsf T}$，从而
\[
\|v\|^2=(Mx)^{\mathsf T}v=x^{\mathsf T}M^{\mathsf T}v=0.
\]
因此 $Mx=0$，证明 $\ker M^2=\ker M$，再由秩与零度公式得到结论。
\end{solve}

\begin{example}
\label{ex:laex-zero-coordinate}
设 $Ax=b$ 有解，记增广矩阵 $\widetilde A=(A\mid b)$，从中删除第 $k$ 个系数列得到 $D$。证明每个解的第 $k$ 个分量都为零，当且仅当 $\operatorname{rank}D<\operatorname{rank}A$。
\end{example}

\begin{solve}
记 $A=(a_1,\ldots,a_n)$。若所有解都满足 $x_k=0$，则 $b$ 可由其余系数列生成；而 $a_k$ 不能由这些列生成，否则将该齐次关系加到一个已有特解上，就得到 $x_k\ne0$ 的解。因此添加 $a_k$ 恰增加一维，$\operatorname{rank}D=\operatorname{rank}A-1$。

反过来，若有一个解满足 $x_k\ne0$，从 $\sum_jx_ja_j=b$ 可解出 $a_k$ 为其余列与 $b$ 的线性组合，所以删去 $a_k$ 不改变增广矩阵的秩。因方程组有解，$\operatorname{rank}\widetilde A=\operatorname{rank}A$，这就与 $\operatorname{rank}D<\operatorname{rank}A$ 矛盾。这里删的是增广矩阵中的系数列，并保留常数列；若仅考察 $A$ 删列，不能反映特解的分量限制。
\end{solve}

\begin{example}
\label{ex:laex-three-lines}
设实平面上三条不同的直线为
\[
ax+by+c=0,\qquad bx+cy+a=0,\qquad cx+ay+b=0.
\]
证明它们交于同一点，当且仅当 $a+b+c=0$。
\end{example}

\begin{solve}
将三式看成关于 $x,y$ 的方程组，其系数矩阵为 $K=\begin{pmatrix}a&b\\b&c\\c&a\end{pmatrix}$。若其秩小于 $2$，所有二阶子式为零，得到 $a^2=bc$、$b^2=ac$、$c^2=ab$，故
\[
(a-b)^2+(b-c)^2+(c-a)^2=0.
\]
实数条件迫使 $a=b=c$，这会使三式相同或不能定义直线，与假设矛盾。因此 $\operatorname{rank}K=2$，三条线有公共点当且仅当增广矩阵秩为 $2$，也即
\[
\det\begin{pmatrix}a&b&c\\b&c&a\\c&a&b\end{pmatrix}
=-(a+b+c)(a^2+b^2+c^2-ab-bc-ca)=0.
\]
第二因子是三个差的平方和的一半，在三条线不同的条件下严格为正，所以条件恰为 $a+b+c=0$。系数秩为 $2$ 还保证公共点唯一。
\end{solve}

\chapter{线性映射}
\label{chap:xie-ex-linear-maps}

% 依据《高等代数习题》（谢启鸿）第 4 章整理：高代习题-4.pdf，
% PDF 第 1--42 页（原书第 192--233 页），覆盖 §§4.1--4.8。
% 将重复例题的不同解法合并；原书题号仅在注释或正文中说明，不作为本笔记编号。
线性映射把向量空间的运算与矩阵的运算联系起来。本章围绕两条相互配合的思路展开：一条是在空间中选基、取补空间、研究像与核，另一条是把映射写成矩阵，利用秩、相似与分块运算。做题时先明确定义域、陪域和数域，再判断适合在哪一种语言中推导。以下 $K$ 表示数域，即 $\mathbb C$ 的子域；没有另作说明时，所讨论的空间均为 $K$ 上的有限维线性空间，$\mathcal L(V,U)$ 表示 $V\to U$ 的全体线性映射，$\mathcal L(V)=\mathcal L(V,V)$。无限维情形将单独指出。$K^n$ 统一采用列向量，$M_{m\times n}(K)$ 表示 $m\times n$ 矩阵空间，$\operatorname{GL}_n(K)$ 表示可逆 $n$ 阶矩阵全体。

\section{基本概念}
\label{sec:xie-ex-linear-basics}

\subsection{线性映射及其运算}

\begin{definition}
\label{def:xie-ex-linear-map}
映射 $T:V\to U$ 若满足
\[
    T(ax+by)=aT(x)+bT(y)
    \qquad(x,y\in V,\ a,b\in K),
\]
称为\emph{线性映射}（linear map）。$V$ 到自身的线性映射称为\emph{线性变换}（linear transformation）；双射的线性映射称为\emph{线性同构}（linear isomorphism），$V$ 到自身的线性同构称为\emph{自同构}（automorphism）。
\end{definition}

线性性等价于保持加法与数乘，因而线性映射保持零向量和任意有限线性组合。特别地，$T(0)=0$ 是必要条件，但不是充分条件。对 $S,T\in\mathcal L(V,U)$、$c\in K$，规定 $(S+T)(x)=S(x)+T(x)$、$(cT)(x)=cT(x)$，便使 $\mathcal L(V,U)$ 成为线性空间。若 $T:V\to U$、$R:U\to W$，规定 $RT=R\circ T$，即 $(RT)(x)=R(T(x))$；复合保持线性，满足结合律、分配律和与数乘的相容性。于是 $\mathcal L(V)$ 以复合为乘法成为含幺结合代数，单位元是恒等变换 $I_V$，但乘法一般不交换。\emph[caution]{复合时先作用右侧映射，只有定义域与陪域匹配时才能相乘。}

\subsection{表示矩阵与换基}

取 $V$ 的有序基 $\mathcal E=(e_1,\ldots,e_n)$ 和 $U$ 的有序基 $\mathcal F=(f_1,\ldots,f_m)$。若
\[
    T(e_j)=\sum_{i=1}^{m}a_{ij}f_i\qquad(1\leq j\leq n),
\]
则 $A=(a_{ij})\in M_{m\times n}(K)$ 称为 $T$ 的\emph{表示矩阵}（representing matrix），记为 $[T]_{\mathcal F\leftarrow\mathcal E}$；第 $j$ 列是 $T(e_j)$ 在陪域基下的坐标。若 $[x]_{\mathcal E}$ 表示坐标列向量，则
\begin{equation}
\label{eq:xie-ex-coordinate-map}
    [T(x)]_{\mathcal F}=A[x]_{\mathcal E}.
\end{equation}
因此取定两组基后，映射 $T\mapsto A$ 是 $\mathcal L(V,U)\to M_{m\times n}(K)$ 的线性同构；复合映射的矩阵是相应矩阵的乘积。对于 $V$ 上的变换，讨论相似时定义域与陪域使用同一组基。

若新基满足 $\mathcal E'=\mathcal E P$、$\mathcal F'=\mathcal F Q$，其中 $P,Q$ 可逆，则坐标关系为 $[x]_{\mathcal E}=P[x]_{\mathcal E'}$、$[y]_{\mathcal F}=Q[y]_{\mathcal F'}$。代入式~\eqref{eq:xie-ex-coordinate-map} 得
\begin{equation}
\label{eq:xie-ex-general-change-basis}
    [T]_{\mathcal F'\leftarrow\mathcal E'}=Q^{-1}AP.
\end{equation}
特别地，同一变换在两组基下的矩阵满足 $B=P^{-1}AP$，称 $A,B$ \emph{相似}（similar）。一般映射的两侧换基相互独立，对应矩阵相抵；同一变换的同基换基对应相似，不能混用这两种关系。

\subsection{像、核与维数公式}

\begin{definition}
\label{def:xie-ex-image-kernel}
线性映射 $T:V\to U$ 的\emph{像空间}（image）和\emph{核空间}（kernel）分别是
\[
    \operatorname{Im}T=\{T(x):x\in V\}\subseteq U,
    \qquad \ker T=\{x\in V:T(x)=0\}\subseteq V.
\]
它们都是子空间。$\operatorname{rank}T=\dim\operatorname{Im}T$ 称为秩，$\dim\ker T$ 称为零度（nullity）。
\end{definition}

\begin{theorem}
\label{thm:xie-ex-rank-nullity}
若 $\dim V=n$，则
\begin{equation}
\label{eq:xie-ex-rank-nullity}
    \dim\ker T+\dim\operatorname{Im}T=n.
\end{equation}
若 $A$ 是 $T$ 的表示矩阵，则 $\operatorname{rank}T=\operatorname{rank}A$。因此 $T$ 单射当且仅当 $\operatorname{rank}A=n$，$T$ 满射当且仅当 $\operatorname{rank}A=\dim U$；当 $\dim U=n$ 时，单射、满射、同构和 $A$ 可逆相互等价。
\end{theorem}

\begin{proof}
取 $\ker T$ 的基 $k_1,\ldots,k_s$，扩充为 $V$ 的基 $k_1,\ldots,k_s,v_1,\ldots,v_{n-s}$。任意向量的像均为 $T(v_1),\ldots,T(v_{n-s})$ 的线性组合。若 $\sum_j c_jT(v_j)=0$，则 $\sum_jc_jv_j\in\ker T$，由上述基的线性无关性知所有 $c_j=0$，故这些像组成 $\operatorname{Im}T$ 的基，得到维数公式。坐标同构把 $\ker T$ 对应到 $Ax=0$ 的解空间，把 $\operatorname{Im}T$ 对应到 $A$ 的列空间，故两种秩相等。最后，线性映射单射当且仅当核为零，满射当且仅当像等于陪域，其余结论由维数公式得到。
\end{proof}

\subsection{不变子空间与分块形式}

\begin{definition}
\label{def:xie-ex-invariant-subspace}
对 $T\in\mathcal L(V)$，若子空间 $W\subseteq V$ 满足 $T(W)\subseteq W$，称 $W$ 为 $T$ 的\emph{不变子空间}（invariant subspace）。$\{0\}$ 和 $V$ 称为平凡不变子空间。
\end{definition}

若 $e_1,\ldots,e_r$ 是 $W$ 的基，并扩充为 $V$ 的基，则 $W$ 不变当且仅当 $T$ 的矩阵具有形式
\begin{equation}
\label{eq:xie-ex-invariant-block}
    [T]=\begin{pmatrix}A_{11}&A_{12}\\0&A_{22}\end{pmatrix},
\end{equation}
其中 $A_{11}$ 是限制变换 $T|_W$ 的矩阵。理由是前 $r$ 个基向量的像在 $W$ 中，其后 $n-r$ 个坐标均为零。若 $V=W_1\oplus\cdots\oplus W_s$ 且每个 $W_i$ 都不变，拼接各子空间的基便得到分块对角矩阵。反过来，分块对角矩阵中各坐标块张成的子空间都是不变子空间。

\section{线性映射及其运算}
\label{sec:xie-ex-linear-construction}

\subsection{在线性无关组与直和上构造映射}

% 原书例 4.1、4.2、4.4。
\begin{theorem}
\label{thm:xie-ex-linear-extension}
给定 $V$ 的一组基 $e_1,\ldots,e_n$ 和 $U$ 中任意向量 $u_1,\ldots,u_n$，存在唯一线性映射 $T:V\to U$ 使 $T(e_i)=u_i$。若 $V=V_1\oplus\cdots\oplus V_s$，给定 $T_i\in\mathcal L(V_i,U)$，也存在唯一线性映射 $T:V\to U$，使 $T|_{V_i}=T_i$。
\end{theorem}

\begin{proof}
对 $x=\sum_i a_ie_i$ 定义 $T(x)=\sum_i a_iu_i$。坐标的唯一性保证定义良好，线性组合的运算保证线性；任何满足基向量取值要求的线性映射都必须遵守此公式，故唯一。对直和情形，将唯一分解 $x=x_1+\cdots+x_s$ 映为 $\sum_iT_i(x_i)$，同样验证即可。
\end{proof}

此结论称为\emph{线性扩张定理}（linear extension theorem）。给定子空间 $W\subseteq V$ 上的线性映射时，选补空间 $V=W\oplus W'$，在 $W'$ 上任意指定线性映射即可扩张到 $V$；常取零映射以简化构造。两个映射只要在一组基上的值相同，便在整个空间上相等。若给定的是线性相关的生成组，则其像不能任意指定，而必须保持原有的全部线性关系，详见第~\ref{sec:xie-ex-linear-training} 节。

\subsection{单侧逆、广义逆与因子分解}

% 原书例 4.3、4.5--4.7；合并几何与矩阵两种证明。
\begin{proposition}
\label{prop:xie-ex-splitting-inverses}
设 $T:V\to U$ 线性。$T$ 单射当且仅当存在 $L:U\to V$ 使 $LT=I_V$；$T$ 满射当且仅当存在 $R:U\to V$ 使 $TR=I_U$。此外，总存在 $G:U\to V$ 使 $TGT=T$。
\end{proposition}

\begin{proof}
$LT=I_V$ 直接推出 $T$ 单射，$TR=I_U$ 直接推出 $T$ 满射。反过来，若 $T$ 单射，在 $\operatorname{Im}T$ 上定义 $L(Tx)=x$，再把 $L$ 扩张到 $U$。若 $T$ 满射，取陪域的一组基并逐一选取原像，线性扩张得到 $R$。一般地，取分解 $V=\ker T\oplus W$ 和 $U=\operatorname{Im}T\oplus Z$，则 $T|_W:W\to\operatorname{Im}T$ 是同构；定义 $G$ 在像空间上等于此同构的逆，在 $Z$ 上为零，便有 $TGT=T$。
\end{proof}

矩阵形式分别为 $BA=I_n$、$AC=I_m$ 和 $AGA=A$；其中 $A\in M_{m\times n}(K)$，$B,C,G$ 均为 $n\times m$ 矩阵。最后一种 $G$ 称为一种\emph{广义逆}（generalized inverse），通常不唯一，也不要求来自内积结构。

\begin{proposition}
\label{prop:xie-ex-factorization}
设 $T,S:V\to U$ 为线性映射，则
\[
    S=HT\text{ 对某个 }H\in\mathcal L(U)\text{ 成立}
    \quad\Longleftrightarrow\quad \ker T\subseteq\ker S,
\]
\[
    S=TF\text{ 对某个 }F\in\mathcal L(V)\text{ 成立}
    \quad\Longleftrightarrow\quad \operatorname{Im}S\subseteq\operatorname{Im}T.
\]
\end{proposition}

\begin{proof}
两个必要性分别由复合的定义得到。若 $\ker T\subseteq\ker S$，在 $\operatorname{Im}T$ 上定义 $H(Tx)=Sx$；当 $Tx=Ty$ 时，$x-y\in\ker T\subseteq\ker S$，故 $Sx=Sy$，定义良好且线性，再扩张到 $U$。若 $\operatorname{Im}S\subseteq\operatorname{Im}T$，取 $V$ 的基 $e_i$，对每个 $S(e_i)$ 选取原像 $v_i$ 使 $T(v_i)=S(e_i)$，再由 $F(e_i)=v_i$ 线性扩张，即得 $TF=S$。
\end{proof}

\subsection{幂零链与选基}

% 原书例 4.8、4.9。
\begin{lemma}
\label{lem:xie-ex-nilpotent-chain}
设 $T\in\mathcal L(V)$、$v\in V$，$q\geq1$ 为整数。若 $T^{q}(v)=0$ 而 $T^{q-1}(v)\ne0$，则 $v,Tv,\ldots,T^{q-1}v$ 线性无关，因而 $q\leq\dim V$。
\end{lemma}

\begin{proof}
若 $\sum_{j=0}^{q-1}c_jT^jv=0$ 且存在非零系数，取最小的非零系数下标 $s$，对等式作用 $T^{q-1-s}$。高于 $s$ 的项全部为零，低于 $s$ 的系数全部为零，得到 $c_sT^{q-1}v=0$，与两个因子均非零矛盾。
\end{proof}

称满足某个 $T^q=0$ 的变换为\emph{幂零变换}（nilpotent transformation）。若 $\dim V=n\geq1$，$T$ 幂零且 $\operatorname{rank}T=n-1$，则存在 $v$ 使
\[
    (v,Tv,\ldots,T^{n-1}v)
\]
成为 $V$ 的一组基，$T$ 在此基下的矩阵仅在紧邻主对角线下方有 $n-1$ 个 $1$，其余元素均为零。为证明这一点，设 $q$ 为最小幂零指数，则引理给出 $q\leq n$；另一方面，由 $\dim\ker T=1$ 以及 $T:\ker T^{j+1}\to\ker T^j$ 的核为 $\ker T$，得到 $\dim\ker T^{j+1}\leq\dim\ker T^j+1$，从而 $n=\dim\ker T^q\leq q$。因此 $q=n$，再取 $T^{n-1}v\ne0$ 即得所需基。

\section{线性同构}
\label{sec:xie-ex-linear-isomorphisms}

\subsection{插值与同构的构造}

% 原书例 4.10、4.11。
记 $K[x]_{\leq n}$ 为次数不超过 $n$ 的多项式组成的空间，包含零多项式。取互不相同的 $a_0,\ldots,a_n\in K$，定义取值映射
\[
    E:K[x]_{\leq n}\to K^{n+1},\qquad
    E(f)=(f(a_0),\ldots,f(a_n))^{\mathsf T}.
\]
这是线性映射。若 $E(f)=0$，则 $f$ 有 $n+1$ 个不同的零点；由非零的 $d$ 次多项式至多有 $d$ 个不同零点可得 $f=0$。定义域与陪域均为 $n+1$ 维，故 $E$ 为同构。于是任意给定取值 $b_0,\ldots,b_n$，都存在唯一 $f\in K[x]_{\leq n}$ 满足 $f(a_i)=b_i$。

\begin{theorem}
\label{thm:xie-ex-lagrange-interpolation}
上述唯一多项式由\emph{Lagrange 插值公式}（Lagrange interpolation formula）给出：
\begin{equation}
\label{eq:xie-ex-lagrange-interpolation}
    f(x)=\sum_{i=0}^{n}b_i\ell_i(x),\qquad
    \ell_i(x)=\prod_{\substack{0\leq j\leq n\\j\ne i}}
    \frac{x-a_j}{a_i-a_j}.
\end{equation}
\end{theorem}

\begin{proof}
由于节点互异，各分母非零；$\ell_i(a_j)$ 在 $i=j$ 时为 $1$，否则为 $0$，且 $\deg\ell_i\leq n$，故右侧多项式具有指定的全部取值。唯一性已由取值映射的单射性证明。
\end{proof}

% 原书例 4.13、4.14。
有限维等维空间之间的线性映射，只需证明单射或满射之一即可证明同构。因此一个变换可逆，当且仅当它把某一组基变成一组基；此时它也把每一组基变成一组基。进一步，若 $U,W\subseteq V$ 且 $\dim U=\dim W$，分别扩充两者的一组基到 $V$ 的基，再逐一对应基向量，就得到自同构 $A$ 使 $A(U)=W$。这里两个子空间可以相交，不要求构成直和。

\subsection{由代数关系判定可逆}

% 原书例 4.12、4.15、4.16。
\begin{proposition}
\label{prop:xie-ex-polynomial-inverse}
设 $T\in\mathcal L(V)$，$q\geq1$ 为整数，$c_0,\ldots,c_q\in K$。若
\[
    \sum_{j=0}^{q}c_jT^j=0,\qquad c_0\ne0,
\]
则
\[
    T^{-1}=-c_0^{-1}\sum_{j=1}^{q}c_jT^{j-1}.
\]
该结论也适用于无限维空间。
\end{proposition}

\begin{proof}
公式右侧是 $T$ 的多项式，与 $T$ 可交换；将已知等式的常数项移至右侧，再分别从左侧和右侧提出 $T$，便得到所给变换与 $T$ 的两个乘积均为 $I_V$。
\end{proof}

若有限维空间上的 $T$ 满足：对每个 $v$ 都存在依赖于 $v$ 的正整数 $q_v$ 使 $T^{q_v}v=0$，则取一组基并令 $q$ 为各基向量对应指数的最大值，得到统一的 $T^q=0$。于是
\begin{equation}
\label{eq:xie-ex-nilpotent-geometric-series}
    (I_V-T)^{-1}=I_V+T+\cdots+T^{q-1}.
\end{equation}
这说明局部幂零条件在有限维时推出幂零。无限维时未必有统一指数，但仍可对每个向量分别使用有限和构造 $I_V-T$ 的逆：任取使 $T^qv=0$ 的 $q$，令 $R(v)=\sum_{j=0}^{q-1}T^jv$；增大 $q$ 不改变结果，使用共同的充分大指数可验证线性以及 $(I_V-T)R=R(I_V-T)=I_V$。

\begin{example}
\label{ex:xie-ex-two-sided-multiplication}
设 $A,B\in M_n(K)$，在 $M_n(K)$ 上定义 $T(X)=AXB$。证明：$T$ 可逆当且仅当 $A,B$ 都可逆。
\end{example}

\begin{solve}
若 $A,B$ 可逆，则 $T^{-1}(Y)=A^{-1}YB^{-1}$。反之，若 $T$ 可逆，则存在 $X$ 使 $AXB=I_n$。因此
\[
    n=\operatorname{rank}(AXB)
    \leq\min\{\operatorname{rank}A,\operatorname{rank}B\}\leq n,
\]
从而 $A,B$ 均满秩且可逆。一般矩形空间上的秩与零度将在例~\ref{ex:xie-ex-rectangular-multiplication} 中计算。
\end{solve}

\subsection{无限维空间中的单侧逆与交换子}

% 原书例 4.17--4.19。
在 $K[x]$ 上定义求导 $D(f)=f'$、积分 $J(f)(x)=\int_0^x f(t)\,dt$ 和乘法变换 $M(f)=xf$。这里积分按多项式系数逐项定义，数域的特征为零保证分母可逆。直接计算得到
\[
    DJ=I,\qquad JD(f)=f-f(0),\qquad DM-MD=I.
\]
$D$ 满而不单，$J$ 单而不满，所以无限维空间中不能由一个单侧逆推出可逆。$D$ 的右逆还可取 $J_\lambda(f)=J(f)+\lambda(f)\cdot1$，其中 $\lambda:K[x]\to K$ 为任意\emph{线性泛函}（linear functional），即取值于标量域的线性映射，因而右逆一般不唯一。

\begin{proposition}
\label{prop:xie-ex-infinite-inverse-criteria}
对任意线性空间上的变换 $S,T$，两者都可逆当且仅当 $ST,TS$ 都可逆；$T$ 可逆当且仅当它恰有一个左逆，也当且仅当它恰有一个右逆。
\end{proposition}

\begin{proof}
若 $ST,TS$ 可逆，则 $(ST)^{-1}S$ 是 $T$ 的左逆，$S(TS)^{-1}$ 是 $T$ 的右逆；左逆 $L$ 与右逆 $R$ 满足 $L=L(TR)=(LT)R=R$，故 $T$ 可逆，同理 $S$ 可逆。反向由逆的乘积公式得到。再设 $LT=I$ 且左逆唯一。因 $[L+(I-TL)]T=I$，唯一性迫使 $I-TL=0$，故 $L$ 也是右逆。若 $TR=I$ 且右逆唯一，则 $T[R+(I-RT)]=I$ 同样迫使 $RT=I$。可逆变换的任何左逆或右逆都等于其逆，得到另一方向。
\end{proof}

若 $0<\dim V=n<\infty$，不可能有 $ST-TS=I_V$：在同一基下取矩阵并取迹，由 $\operatorname{tr}(AB)=\sum_{i,j}a_{ij}b_{ji}=\operatorname{tr}(BA)$，左侧迹为 $0$，右侧迹为 $n\ne0$。因此多项式空间上的 $DM-MD=I$ 也提供了有限维结论失效的具体实例。

\section{线性映射与矩阵}
\label{sec:xie-ex-linear-matrices}

\subsection{坐标同构、秩标准型与秩一分解}

% 原书例 4.20--4.25、4.28。
坐标映射 $c_{\mathcal E}:V\to K^n$ 与 $c_{\mathcal F}:U\to K^m$ 都是同构，式~\eqref{eq:xie-ex-coordinate-map} 可以写成
\[
    c_{\mathcal F}T=L_Ac_{\mathcal E},\qquad L_A(x)=Ax.
\]
因此空间问题和矩阵问题可以双向转化。相似关系也有完整的空间解释：若 $B=P^{-1}AP$，在 $K^n$ 上取 $T(x)=Ax$，标准基记为 $\mathcal E$，新基取 $\mathcal E P$，则 $T$ 在两组基下分别为 $A,B$。

\begin{proposition}
\label{prop:xie-ex-rank-normal-form}
若 $T:V\to U$ 的秩为 $r$，则可以分别选择定义域与陪域的基，使
\begin{equation}
\label{eq:xie-ex-rank-normal-form}
    [T]=\begin{pmatrix}I_r&0\\0&0\end{pmatrix}_{m\times n}.
\end{equation}
此外，$T$ 可表示为 $r$ 个秩一映射之和；$r=0$ 时理解为空和。
\end{proposition}

\begin{proof}
取 $V=\ker T\oplus W$，$W$ 的基记为 $v_1,\ldots,v_r$，则 $T(v_1),\ldots,T(v_r)$ 是像空间的基。定义域按 $W$ 的基、核的基排列，陪域按这组像、其扩充基排列，即得所述矩阵。若 $x=\sum_{i=1}^{r}a_iv_i+k$，其中 $k\in\ker T$，令 $T_i(x)=a_iT(v_i)$，则每个 $T_i$ 的像为一维，且 $T=\sum_iT_i$。
\end{proof}

\begin{notealert}[type=caution,title={秩标准型使用两侧独立换基}]
式~\eqref{eq:xie-ex-rank-normal-form} 对任意线性映射成立，即使 $V=U$，也允许定义域和陪域使用不同的基。不能据此断言每个方阵都相似于 $\operatorname{diag}(I_r,0)$；后者要求变换本身幂等，见第~\ref{sec:xie-ex-idempotents} 节。
\end{notealert}

\subsection{核相同、矩阵相抵与秩不等式}

% 原书例 4.27 与重新讨论的 Frobenius 不等式。
\begin{theorem}
\label{thm:xie-ex-same-kernel}
设 $T,S:V\to U$，则 $\ker T=\ker S$ 当且仅当存在 $U$ 上的自同构 $H$ 使 $S=HT$。因此对 $A,B\in M_{m\times n}(K)$，齐次方程组 $Ax=0$ 与 $Bx=0$ 同解，当且仅当存在 $P\in\operatorname{GL}_m(K)$ 使 $B=PA$。
\end{theorem}

\begin{proof}
若 $S=HT$ 且 $H$ 可逆，则两者的核相同。反之，在 $\operatorname{Im}T$ 上定义 $H_0(Tx)=Sx$，共同的核保证定义良好、单射且满射到 $\operatorname{Im}S$。像空间维数相同，因此可分别把这两个像空间的基扩充成 $U$ 的基，并对应扩充的基向量，把 $H_0$ 扩张为 $U$ 的自同构。矩阵结论由坐标同构得到，也说明“同解”对应初等行变换等价。
\end{proof}

若矩阵 $A\in M_{p\times q}(K)$、$B\in M_{q\times r}(K)$、$C\in M_{r\times s}(K)$，把它们视为相应空间之间的线性映射。在 $\operatorname{Im}B$ 和 $\operatorname{Im}(BC)$ 上分别限制 $A$，由维数公式有
\[
\begin{aligned}
    \operatorname{rank}(AB)
    &=\operatorname{rank}B-\dim(\ker A\cap\operatorname{Im}B),\\
    \operatorname{rank}(ABC)
    &=\operatorname{rank}(BC)-\dim(\ker A\cap\operatorname{Im}(BC)).
\end{aligned}
\]
由于 $\operatorname{Im}(BC)\subseteq\operatorname{Im}B$，相减得\emph{Frobenius 秩不等式}
\begin{equation}
\label{eq:xie-ex-frobenius-rank}
    \operatorname{rank}(ABC)\geq
    \operatorname{rank}(AB)+\operatorname{rank}(BC)-\operatorname{rank}B.
\end{equation}
证明的关键是限制映射的核，而不是只比较三个乘积的像。

\subsection{根据空间结构选择基}

% 原书例 4.26、4.29、4.30。
\begin{proposition}
\label{prop:xie-ex-basis-independent-matrix}
有限维空间上的变换 $T$ 在每一组基下的表示矩阵都相同，当且仅当 $T=cI_V$，其中 $c\in K$。
\end{proposition}

\begin{proof}
纯量变换在每组基下的矩阵均为 $cI_n$。反过来，设共同矩阵为 $A$。任取非零 $v$ 并补成基，比较该基与只把第一项变为 $2v$ 的基下的第一列，得到 $T(v)$ 在其余基方向上的分量全为零，因此每条直线都不变。固定基 $e_i$，写 $T(e_i)=c_ie_i$，再利用 $K(e_i+e_j)$ 不变可得 $c_i=c_j$，故 $T=cI_V$。零空间的情形结论亦成立。
\end{proof}

在 $M_n(K)$ 上，转置变换 $\tau(X)=X^{\mathsf T}$ 的结构可由
\[
    X=\frac{X+X^{\mathsf T}}2+\frac{X-X^{\mathsf T}}2
\]
直接看出：对称矩阵空间与反对称矩阵空间构成直和，$\tau$ 在前者上为 $I$，在后者上为 $-I$。因此取基
\[
    E_{ii},\quad E_{ij}+E_{ji}\ (i<j),\quad E_{ij}-E_{ji}\ (i<j),
\]
其中 $E_{ij}$ 为第 $(i,j)$ 元素为 $1$、其余为零的矩阵，表示矩阵便为对角阵，$1$ 与 $-1$ 分别出现 $n(n+1)/2$ 和 $n(n-1)/2$ 次。

\begin{example}
\label{ex:xie-ex-paired-square-zero}
设 $S,T\in\mathcal L(V)$ 满足 $S^2=T^2=0$、$ST+TS=I_V$。证明 $\dim V$ 为偶数，并给出两者在同一组基下的矩阵。
\end{example}

\begin{solve}
任意 $x$ 有分解 $x=STx+TSx$，其中 $STx\in\ker S$、$TSx\in\ker T$；若 $x$ 同时属于两核，则 $x=(ST+TS)x=0$。故 $V=\ker S\oplus\ker T$。对 $u\in\ker S$，有 $STu=u$，故 $T:\ker S\to\ker T$ 单射；对 $v\in\ker T$，有 $TSv=v$，故它又满射，逆为 $S|_{\ker T}$。取 $\ker S$ 的基 $u_1,\ldots,u_k$，令 $v_i=Tu_i$，则 $(u_1,v_1,\ldots,u_k,v_k)$ 是 $V$ 的基，且每一对上
\[
    [T]_i=\begin{pmatrix}0&0\\1&0\end{pmatrix},\qquad
    [S]_i=\begin{pmatrix}0&1\\0&0\end{pmatrix}.
\]
两变换的整体矩阵由这些 $2\times2$ 块组成，因而 $\dim V=2k$。
\end{solve}

\section{像空间和核空间}
\label{sec:xie-ex-images-kernels}

\subsection{从坐标还原空间中的向量}

% 原书例 4.31。
\begin{example}
\label{ex:xie-ex-compute-image-kernel}
设 $T$ 在基 $\mathcal E=(e_1,e_2,e_3,e_4)$ 下的矩阵为
\[
    A=\begin{pmatrix}
    1&0&2&1\\-1&2&1&3\\1&2&5&5\\2&-2&1&-2
    \end{pmatrix}.
\]
求 $\operatorname{Im}T$ 与 $\ker T$ 的基。
\end{example}

\begin{solve}
设列向量为 $a_1,\ldots,a_4$，则 $a_3=2a_1+\tfrac32a_2$、$a_4=a_1+2a_2$，而前两列线性无关，故秩为 $2$。将前两列还原为原空间的向量得
\[
    \operatorname{Im}T=
    \operatorname{span}\{e_1-e_2+e_3+2e_4,\ 2e_2+2e_3-2e_4\}.
\]
上述列关系给出 $Ax=0$ 的两个无关解 $(-4,-3,2,0)^{\mathsf T}$ 和 $(-1,-2,0,1)^{\mathsf T}$，由零度为 $2$ 得
\[
    \ker T=\operatorname{span}
    \{-4e_1-3e_2+2e_3,\ -e_1-2e_2+e_4\}.\qedhere
\]
\end{solve}

求像空间时，行化简可以确定主元列，但必须在\emph[caution]{原矩阵}中取对应列；行变换通常改变列空间本身。求核空间时，先求坐标解，再以原来的定义域基还原，不能把坐标列向量直接当成抽象空间中的向量。

% 原书例 4.32、4.33。
\begin{proposition}
\label{prop:xie-ex-separate-maps}
若 $T_1,\ldots,T_s:V\to U$ 都非零，则存在 $v\in V$ 使每个 $T_i(v)\ne0$。若 $S_1,\ldots,S_s:V\to U$ 两两不同，则存在 $v$ 使所有 $S_i(v)$ 两两不同。
\end{proposition}

\begin{proof}
证明第一结论只需说明无限域上的空间不能被有限个真子空间覆盖。对真子空间数目归纳：取 $a$ 不属于前 $s-1$ 个子空间，若它也不属于最后一个则已得证；否则 $a$ 属于最后一个，取 $b$ 不属于最后一个，考察直线 $a+tb$。它与前 $s-1$ 个子空间各至多交于一点，否则可由两交点相减并线性组合推出 $a$ 属于该子空间；与最后一个仅可能在 $t=0$ 处相交。数域无限，故可避开有限多个参数值。将此结论用于 $\ker T_i$ 即得第一结论，再用于全部非零差映射 $S_i-S_j$ 即得第二结论。
\end{proof}

\subsection{幂次的稳定与空间分解}

% 原书例 4.34--4.36。
对任意 $T\in\mathcal L(V)$，有子空间链
\[
    V=\operatorname{Im}T^0\supseteq\operatorname{Im}T
    \supseteq\operatorname{Im}T^2\supseteq\cdots,
    \qquad
    0=\ker T^0\subseteq\ker T\subseteq\ker T^2\subseteq\cdots.
\]

\begin{theorem}
\label{thm:xie-ex-fitting}
设 $\dim V=n$，则存在 $0\leq m\leq n$，使所有 $j\geq m$ 都满足
\[
    \operatorname{Im}T^j=\operatorname{Im}T^m,\qquad
    \ker T^j=\ker T^m,
\]
并有
\begin{equation}
\label{eq:xie-ex-fitting}
    V=\ker T^m\oplus\operatorname{Im}T^m.
\end{equation}
两部分均为 $T$ 的不变子空间，$T$ 在前者上幂零，在后者上可逆。
\end{theorem}

\begin{proof}
像的维数每次严格下降至少 $1$，故不迟于 $m=n$ 就有 $\operatorname{Im}T^m=\operatorname{Im}T^{m+1}$；对等式继续作用 $T$，所有后续像均相同。维数公式与核的包含关系使后续核也相同。若 $x=T^my\in\ker T^m$，则 $T^{2m}y=0$，由核的稳定知 $y\in\ker T^m$，故 $x=0$。两部分维数之和为 $n$，得到直和。它们的不变性由幂次可交换得到；前者上 $T^m=0$，后者上 $T$ 满射且有限维，故可逆。$m=0$ 时直和即 $V=0\oplus V$。
\end{proof}

这也称为\emph{Fitting 分解}（Fitting decomposition），并推出任意 $n$ 阶矩阵满足 $\operatorname{rank}A^n=\operatorname{rank}A^{n+1}=\cdots$。不过，通常不能把指数 $m$ 直接取成 $1$。

\begin{proposition}
\label{prop:xie-ex-first-power-splitting}
对有限维空间上的 $T$，以下条件等价：$V=\ker T\oplus\operatorname{Im}T$；$V=\ker T+\operatorname{Im}T$；$\ker T\cap\operatorname{Im}T=0$；$\ker T=\ker T^2$；$\operatorname{Im}T=\operatorname{Im}T^2$；$\ker T$ 有 $T$-不变补空间；$\operatorname{Im}T$ 有 $T$-不变补空间。核或像在一次幂与二次幂相等时，其后全部幂次也相等。
\end{proposition}

\begin{proof}
维数公式保证前两子空间的维数之和为 $n$，故“和为全空间”“交为零”“构成直和”相互等价。若 $T^2x=0$，则 $Tx\in\ker T\cap\operatorname{Im}T$；反过来，若 $Tx\in\ker T$，由两核相等可得 $x\in\ker T$。因此交为零等价于两核相等，维数公式又将其转化为两像相等，后续稳定性同定理~\ref{thm:xie-ex-fitting} 的论证。直和本身提供不变补空间。若 $V=\ker T\oplus W$ 且 $W$ 不变，则 $\operatorname{Im}T=T(W)\subseteq W$，比较维数知 $W=\operatorname{Im}T$；若 $V=\operatorname{Im}T\oplus Z$ 且 $Z$ 不变，则 $T(Z)\subseteq Z\cap\operatorname{Im}T=0$，比较维数知 $Z=\ker T$。
\end{proof}

\subsection{商空间、预先指定像核与维数估计}

% 原书例 4.37--4.44。
若 $N\subseteq V$ 为子空间，\emph{商空间}（quotient space）$V/N$ 的元素是陪集 $x+N$，运算为 $(x+N)+(y+N)=(x+y)+N$、$a(x+N)=ax+N$；若 $N$ 的基扩充为 $V$ 的基，新增基向量的陪集构成商空间的基，故 $\dim(V/N)=\dim V-\dim N$。

\begin{theorem}
\label{thm:xie-ex-first-isomorphism}
线性映射 $T:V\to U$ 诱导同构
\[
    \overline T:V/\ker T\longrightarrow\operatorname{Im}T,
    \qquad \overline T(x+\ker T)=T(x).
\]
\end{theorem}

\begin{proof}
若 $x+\ker T=y+\ker T$，则 $x-y\in\ker T$，故 $Tx=Ty$，所以定义不依赖代表元。陪集运算与 $T$ 的线性性保证 $\overline T$ 线性；每个像均有原像，而 $\overline T(x+\ker T)=0$ 当且仅当 $x\in\ker T$，即该陪集为商空间零元。因此它既满又单。
\end{proof}

取维数又得到式~\eqref{eq:xie-ex-rank-nullity}。若 $V=N\oplus W$，则 $T|_W:W\to\operatorname{Im}T$ 是同构，所以核的任意补空间均可用来选择像向量的唯一原像；特别地，满射在这样的补空间上给出到整个陪域的同构。给定 $V$ 的子空间 $N,R$，存在 $T\in\mathcal L(V)$ 使 $\ker T=N$、$\operatorname{Im}T=R$，当且仅当 $\dim N+\dim R=\dim V$：必要性来自维数公式，充分性则取补空间 $V=N\oplus W$，在 $N$ 上令 $T=0$，在 $W$ 上取任意到 $R$ 的同构。\emph[caution]{这里并不要求 $N\cap R=0$。}

作为维数公式的应用，迹映射 $\operatorname{tr}:M_n(K)\to K$ 在 $n\geq1$ 时为满射，因为 $\operatorname{tr}E_{11}=1$，故无迹矩阵空间的维数为 $n^2-1$，一组基为
\[
    \{E_{ij}:i\ne j\}\ \cup
    \{E_{ii}-E_{i+1,i+1}:1\leq i<n\}.
\]
这些矩阵的非对角元素与对角元素分别可独立比较，因而线性无关；其数目恰为所需维数。

对 $T:V\to U$、子空间 $W\subseteq V$ 与 $Z\subseteq U$，限制映射的维数公式分别给出
\begin{equation}
\label{eq:xie-ex-restriction-dimension}
    \dim T(W)=\dim W-\dim(W\cap\ker T),
\end{equation}
\begin{equation}
\label{eq:xie-ex-preimage-dimension}
    \dim T^{-1}(Z)=\dim\ker T+\dim(Z\cap\operatorname{Im}T).
\end{equation}
第二式中的 $T^{-1}(Z)$ 是集合的原像，即使 $T$ 不可逆也有意义；限制 $T:T^{-1}(Z)\to Z\cap\operatorname{Im}T$ 为满射且核仍是 $\ker T$，从而该式成立。于是
\[
    \dim W-\dim\ker T\leq\dim T(W)\leq\dim W,
    \qquad \dim T^{-1}(Z)\leq\dim\ker T+\dim Z,
\]
且当 $Z\subseteq\operatorname{Im}T$ 时第二个不等式取等号。若 $\dim V>\dim U$，维数公式还立即给出 $\ker T\ne0$。对 $A\in M_{m\times n}(K)$、$B\in M_{n\times p}(K)$，将第一式用于 $A|_{\operatorname{Im}B}$，得到\emph{Sylvester 秩不等式}
\begin{equation}
\label{eq:xie-ex-sylvester-rank}
    \operatorname{rank}A+\operatorname{rank}B-n
    \leq\operatorname{rank}(AB)
    \leq\min\{\operatorname{rank}A,\operatorname{rank}B\}.
\end{equation}

\section{不变子空间}
\label{sec:xie-ex-invariant-subspaces}

\subsection{验证不变性与可交换变换}

% 原书例 4.45--4.50。
若 $W=\operatorname{span}\{w_1,\ldots,w_s\}$，验证 $W$ 不变只需逐一验证 $T(w_i)\in W$，不要求生成向量线性无关。两个不变子空间的交与和仍不变，因为像既保持各自归属，又保持向量的加法分解；相同论证也适用于有限个交与和。每条一维子空间都不变，当且仅当 $T$ 为纯量变换，证明见命题~\ref{prop:xie-ex-basis-independent-matrix} 中对 $e_i$ 与 $e_i+e_j$ 的比较。

更一般地，对固定的 $1\leq r\leq n-1$，如果每个 $r$ 维子空间都不变，那么每条直线也不变。事实上，给定直线 $L$，所有包含 $L$ 的 $r$ 维子空间之交就是 $L$：对任何 $x\notin L$，扩充 $L$ 的基并把 $x$ 作为第二个基向量，可以选出包含 $L$ 而避开 $x$ 的 $r$ 维坐标子空间。因此 $T$ 保持上述交，从而是纯量变换。反过来纯量变换保持所有子空间。这就得到“任意固定非平凡维数的全部子空间不变”与“纯量变换”等价。

\begin{proposition}
\label{prop:xie-ex-commuting-invariance}
若 $ST=TS$，则 $\ker S$ 和 $\operatorname{Im}S$ 都是 $T$ 的不变子空间。
\end{proposition}

\begin{proof}
若 $Sx=0$，则 $S(Tx)=T(Sx)=0$；若 $x=Sy$，则 $Tx=TSy=STy\in\operatorname{Im}S$。
\end{proof}

例如，设 $A$ 为幂零矩阵、$AB=BA$，且 $\operatorname{rank}(AB)=\operatorname{rank}B$。像空间 $\operatorname{Im}B$ 对 $A$ 不变，且 $A$ 在其上的限制满射；若 $A^q=0$，这个满射的 $q$ 次复合仍满射，却同时是零映射，只能有 $\operatorname{Im}B=0$，故 $B=0$。

若 $T$ 可逆且 $W$ 是有限维 $T$-不变子空间，则 $T|_W$ 单射，因而满射，故 $T(W)=W$、$T^{-1}(W)=W$。有限维条件不能随意去掉：在 Laurent 多项式空间 $K[x,x^{-1}]$ 上，乘以 $x$ 为自同构，$K[x]$ 对它不变，但乘以 $x^{-1}$ 并不保持 $K[x]$。

\subsection{求导变换与互异对角元的全部不变子空间}

% 原书例 4.51、4.52。
\begin{example}
\label{ex:xie-ex-derivative-invariant}
设 $n\geq1$。在 $K[x]_{<n}=\{f:\deg f<n\}\cup\{0\}$ 上，令 $D(f)=f'$。求 $D$ 的全部不变子空间。
\end{example}

\begin{solve}
除零空间外，取不变子空间 $W$ 中次数最大的多项式 $f$，记次数为 $d$。$f,Df,\ldots,D^df$ 均在 $W$ 中，次数分别为 $d,d-1,\ldots,0$，最高项系数均非零，故线性无关并张成 $K[x]_{\leq d}$；按 $d$ 的最大性又有 $W\subseteq K[x]_{\leq d}$。因此 $W=K[x]_{\leq d}$。全部不变子空间恰为
\[
    0,\quad K[x]_{\leq0},\quad K[x]_{\leq1},\quad\ldots,
    \quad K[x]_{\leq n-1}.
\]
特别地，每个非零不变子空间都包含常数多项式 $1$。
\end{solve}

\begin{example}
\label{ex:xie-ex-distinct-diagonal-invariant}
设 $T(e_i)=d_ie_i$，其中 $d_1,\ldots,d_n\in K$ 两两不同。求 $T$ 的全部不变子空间。
\end{example}

\begin{solve}
定义多项式
\[
    p_i(t)=\prod_{j\ne i}\frac{t-d_j}{d_i-d_j}.
\]
于是 $p_i(T)e_j=0$（$j\ne i$），$p_i(T)e_i=e_i$，即 $p_i(T)$ 提取第 $i$ 个坐标分量。若 $W$ 对 $T$ 不变，则它对 $T$ 的所有多项式也不变；只要 $w=\sum_jc_je_j\in W$ 且 $c_i\ne0$，就有 $c_ie_i=p_i(T)w\in W$，从而 $e_i\in W$。故 $W$ 必由某个基向量子集张成，反过来每个这样的子空间都不变，共有 $2^n$ 个，包含零空间与全空间。互异条件保证分母非零；若所有 $d_i$ 相同，则每个子空间均不变，结论会发生根本变化。
\end{solve}

\subsection{限制变换与商空间上的诱导变换}

% 原书例 4.53。
若 $W$ 为 $T$-不变子空间，则不仅能研究 $T|_W$，还可在商空间上定义
\[
    \widetilde T:V/W\to V/W,\qquad
    \widetilde T(x+W)=T(x)+W.
\]
当 $x-y\in W$ 时，$T(x)-T(y)\in W$，所以定义良好；陪集运算直接保证其线性。对式~\eqref{eq:xie-ex-invariant-block} 中的适配基，$T|_W$ 的矩阵为 $A_{11}$，$\widetilde T$ 在商空间基 $e_{r+1}+W,\ldots,e_n+W$ 下的矩阵为 $A_{22}$。非对角块 $A_{12}$ 描述所选补空间中的向量映入 $W$ 的分量；只有补空间也不变时，它才为零，不能因已有一个不变子空间就直接写成分块对角阵。

\section{幂等变换}
\label{sec:xie-ex-idempotents}

\subsection{幂等变换就是投影}

% 原书例 4.54--4.56。
\begin{definition}
满足 $P^2=P$ 的线性变换称为\emph{幂等变换}（idempotent transformation）。若 $V=U\oplus W$，将 $u+w$ 映为 $u$ 的变换称为沿 $W$ 到 $U$ 的\emph{投影}（projection）。这里不涉及内积，因而不要求正交。
\end{definition}

\begin{theorem}
\label{thm:xie-ex-idempotent-projection}
$P^2=P$ 当且仅当 $P$ 是沿其核到其像的投影。此时
\[
    V=\operatorname{Im}P\oplus\ker P,\qquad
    \operatorname{Im}P=\ker(I-P),\qquad
    \ker P=\operatorname{Im}(I-P).
\]
若 $\operatorname{rank}P=r$，则存在一组基使 $[P]=\operatorname{diag}(I_r,0)$。
\end{theorem}

\begin{proof}
若 $P^2=P$，则 $x=Px+(I-P)x$，其中 $P(Px)=Px$、$P((I-P)x)=0$。若 $y=Px\in\ker P$，则 $y=P^2x=Py=0$，因此构成直和。$P$ 在其像上为恒等，在核上为零，正是所述投影；两个像核等式可由上述分解逐项验证。反过来，投影作用一次后已经在目标子空间中，再作用不变，故 $P^2=P$。拼接像与核的基即得矩阵形式。
\end{proof}

因此幂等矩阵的秩等于其迹，两个同阶幂等矩阵相似当且仅当秩相同。特别地，若 $A,B$ 幂等且同秩，则存在可逆矩阵 $C$ 使 $B=C^{-1}AC$，即 $CB=AC$。这些结论来自投影的空间分解，无须预先使用特征值理论。

\subsection{两个投影的像、核与加减}

% 原书例 4.57、4.58。
\begin{proposition}
\label{prop:xie-ex-two-projections}
设 $P^2=P$、$Q^2=Q$，则
\[
\begin{aligned}
    \operatorname{Im}P=\operatorname{Im}Q
    &\iff PQ=Q,\quad QP=P,\\
    \ker P=\ker Q
    &\iff PQ=P,\quad QP=Q,\\
    (P+Q)^2=P+Q
    &\iff PQ=QP=0,\\
    (P-Q)^2=P-Q
    &\iff PQ=QP=Q.
\end{aligned}
\]
\end{proposition}

\begin{proof}
投影在自身像上为恒等，因此两像相同推出第一组乘积等式，反过来乘积等式给出两个方向的像包含。若两核相同，利用 $x-Qx\in\ker Q=\ker P$ 得 $Px=PQx$，同理 $Q=QP$；反向则由 $Px=0\Rightarrow Qx=QPx=0$ 及对称论证得到。

对和的幂等性，展开得 $PQ+QP=0$。分别从左、右乘 $P$，得到 $PQ+PQP=0$、$PQP+QP=0$，所以 $PQ=QP$，再由 $2PQ=0$ 得二者为零。对差，展开得 $PQ+QP=2Q$。左右乘 $P$ 分别得到 $PQP=PQ$ 和 $PQP=QP$，故 $PQ=QP$，再代回得 $PQ=Q$。各充分性均由展开验证；这里使用了数域中特征不为 $2$ 的事实。
\end{proof}

\subsection{多个投影、直和与秩的可加性}

% 原书例 4.59、4.60，以及训练题三 10、11。
\begin{theorem}
\label{thm:xie-ex-family-projections}
设 $P_1,\ldots,P_s\in\mathcal L(V)$ 满足 $P_i^2=P_i$、$P_iP_j=0$（$i\ne j$）。令 $P=\sum_iP_i$，则 $P^2=P$，并且
\begin{equation}
\label{eq:xie-ex-family-projections}
    V=\operatorname{Im}P_1\oplus\cdots\oplus\operatorname{Im}P_s
    \oplus\bigcap_{i=1}^{s}\ker P_i.
\end{equation}
因此 $\bigcap_i\ker P_i=0$ 当且仅当 $P=I_V$，此时 $P_i$ 正是该直和分解的各分量投影。
\end{theorem}

\begin{proof}
展开乘积得 $P^2=P$，且 $P_iP=PP_i=P_i$。若 $x_i\in\operatorname{Im}P_i$ 且 $\sum_ix_i=0$，作用 $P_j$ 得 $x_j=0$，故各像的和为直和，其和等于 $\operatorname{Im}P$。又 $Px=0$ 推出 $P_ix=P_iPx=0$，反向由求和得到，所以 $\ker P=\bigcap_i\ker P_i$。应用定理~\ref{thm:xie-ex-idempotent-projection} 即得分解及其余结论。
\end{proof}

\begin{theorem}
\label{thm:xie-ex-idempotent-rank-additivity}
设 $P^2=P$ 且 $P=T_1+\cdots+T_s$，则
\begin{equation}
\label{eq:xie-ex-idempotent-rank-additivity}
    \operatorname{rank}P=\sum_{i=1}^{s}\operatorname{rank}T_i
    \quad\Longleftrightarrow\quad
    T_i^2=T_i,\quad T_iT_j=0\ (i\ne j).
\end{equation}
\end{theorem}

\begin{proof}
若右侧成立，定理~\ref{thm:xie-ex-family-projections} 给出像空间的直和，因而秩可加。反过来，设左侧成立，记 $R_i=\operatorname{Im}T_i$。由
\[
    \operatorname{Im}P\subseteq\sum_iR_i,\qquad
    \dim\operatorname{Im}P=\sum_i\dim R_i
\]
以及和空间维数不超过各维数之和，得 $\operatorname{Im}P=\bigoplus_iR_i$。对任意 $y\in R_i$，因 $P$ 在其像上为恒等，$y=Py=\sum_jT_jy$；由直和表示的唯一性，$T_iy=y$、$T_jy=0$（$j\ne i$）。令 $y=T_ix$ 即得 $T_i^2=T_i$、$T_jT_i=0$。
\end{proof}

特别地，若 $\sum_iT_i=I_V$ 且 $\sum_i\operatorname{rank}T_i=\dim V$，则各 $T_i$ 自动成为互相消去的投影。几何上，秩的可加性排除了各像空间之间的重叠，而和为恒等又保证这些像覆盖整个空间。

\section{基础训练：方法与综合应用}
\label{sec:xie-ex-linear-training}

\subsection{基本判断与换基计算}

基础判断先归结为定义与维数，而不直接套选项。含非零常数平移的映射不保持零向量；坐标中出现平方的映射一般不保持数乘，例如 $(x_1,x_2,x_3)\mapsto(x_1^2,x_2,x_3)$ 在 $x_1=1$ 处与数乘 $2$ 不相容。一维空间上的每个线性变换都是 $cI$，$n$ 维空间的变换空间维数为 $n^2$。两个变换在一组基上的作用相同便相等，但仅有相同的秩、像或核均不能推出相等；例如 $I$ 与 $2I$ 有相同的像和核。线性相关组的像仍相关，因为原有非零系数关系在作用线性映射后仍成立。

维数公式并不自动给出像与核的直和。例如在 $K[x]_{<2}$ 上，求导 $D$ 的像与核均为 $K\cdot1$，其维数之和虽为 $2$，交却非零。幂等变换保证像核直和；幂零变换则可能有非零的像核交。非零有限维变换若没有非平凡不变子空间，则其核只能为零，从而可逆。对实二维空间，旋转 $\pi/2$ 的变换 $R(x,y)=(-y,x)$ 没有不变直线：若 $R(v)=cv$ 且 $v\ne0$，由 $R^2=-I$ 得 $c^2=-1$，与 $c\in\mathbb R$ 矛盾。这也说明不变子空间问题必须明确数域。

计算时坚持使用“基向量的像为列”的约定。例如 $D$ 在 $1,x,\ldots,x^{n-1}$ 下的矩阵满足第 $(i,i+1)$ 元素为 $i$（$1\leq i<n$），其余为零。若 $S,T$ 在同一组基下的矩阵为 $A,B$，则 $ST+2S^2$ 的矩阵为 $AB+2A^2$；若 $S$ 可逆，基的过渡矩阵为 $P$，则 $TS^{-2}+2S+I$ 在新基下的矩阵为
\[
    P^{-1}(BA^{-2}+2A+I_n)P.
\]
若恰有 $T(e_i)=f_i$ 且 $(f_1,\ldots,f_n)=(e_1,\ldots,e_n)P$，则 $T$ 在旧基下为 $P$，在新基下仍为 $P^{-1}PP=P$；这是该映射的特殊关系，不能推广为一般换基不改变矩阵。

\subsection{构造、幂零链与复合映射}

% 对应训练题三 1--7；已在正文完整证明的结论不重复展开。
\begin{example}
\label{ex:xie-ex-map-space-basis}
取 $V$ 的基 $e_1,\ldots,e_n$ 和 $U$ 的基 $f_1,\ldots,f_m$，定义 $T_{ij}(e_j)=f_i$、$T_{ij}(e_k)=0$（$k\ne j$）。证明这些映射构成 $\mathcal L(V,U)$ 的一组基。
\end{example}

\begin{solve}
若 $T(e_j)=\sum_i a_{ij}f_i$，则在线性扩张定理下 $T=\sum_{i,j}a_{ij}T_{ij}$，故这些映射张成整个空间。若 $\sum_{i,j}c_{ij}T_{ij}=0$，将等式作用于 $e_j$ 得 $\sum_ic_{ij}f_i=0$，由 $f_i$ 线性无关可知全部 $c_{ij}=0$。因此共有 $mn$ 个基向量，且对应 $T_{ij}\mapsto E_{ij}$ 正是表示矩阵同构。
\end{solve}

有限支撑数列空间 $K^{(\mathbb N_0)}$ 是所有仅有有限多个非零项的数列组成的空间，映射 $(a_0,a_1,\ldots)\mapsto\sum_{j\geq0}a_jx^j$ 是到 $K[x]$ 的线性同构：该和实际有限，系数的唯一性保证单射，每个多项式又都来自其系数数列。对 $K[x]$ 上的 $D(f)=f'$、$M(f)=xf$，乘积求导对任意正整数 $q$ 进一步给出
\[
    DM^q-M^qD=qM^{q-1}\qquad(q\geq1),
\]
因为 $D(x^qf)-x^qD(f)=qx^{q-1}f$。这推广了前面的交换子等式。

若 $V=U_1\oplus U_2$，且 $T(U_1)\subseteq U_2$、$T(U_2)\subseteq U_1$，则 $\operatorname{Im}T=T(U_1)\oplus T(U_2)$，因为两像分别落在交为零的两个子空间中。因此 $\operatorname{rank}T=\dim T(U_1)+\dim T(U_2)$，该结论不要求 $U_1,U_2$ 维数相等。

\begin{example}
\label{ex:xie-ex-cyclic-nilpotent-invariant}
设 $T(e_i)=e_{i+1}$（$i<n$）、$T(e_n)=0$。证明：包含 $e_1$ 的不变子空间只能是 $V$；每个非零不变子空间都包含 $e_n$；$V$ 不能分解为两个非零不变子空间的直和。
\end{example}

\begin{solve}
包含 $e_1$ 且不变就必须包含 $Te_1=e_2$ 及其后所有基向量，因此等于 $V$。若 $W\ne0$ 不变，取 $w=\sum_i a_ie_i\ne0$，令 $k$ 为最小非零系数下标，则 $T^{n-k}w=a_ke_n\in W$，故 $e_n\in W$。若有两个非零不变子空间构成直和，它们都包含 $e_n$，交非零，产生矛盾。
\end{solve}

若 $A$ 可逆，则 $BA=A^{-1}(AB)A$，所以 $AB$ 与 $BA$ 相似；这里 $B$ 无须可逆。对一般复合映射 $ST$，把 $T$ 限制到 $\ker(ST)$，其像包含于 $\ker S$，核正是 $\ker T$，由维数公式有
\begin{equation}
\label{eq:xie-ex-product-nullity}
    \dim\ker(ST)\leq\dim\ker S+\dim\ker T.
\end{equation}
对同一 $n$ 维空间上的变换，这也等价于式~\eqref{eq:xie-ex-sylvester-rank} 的下界。

\subsection{限制映射、对合与矩形矩阵空间}

% 对应训练题三 8--13；三 10、11 已归入投影族定理。
\begin{example}
\label{ex:xie-ex-restricted-matrix-image}
设 $A\in M_{m\times n}(K)$、$B\in M_{\ell\times n}(K)$，$W=\ker B$，$T(x)=Ax$。证明
\[
    \dim T(W)=\operatorname{rank}\begin{pmatrix}A\\B\end{pmatrix}
    -\operatorname{rank}B.
\]
\end{example}

\begin{solve}
限制 $T|_W$ 的核为 $\ker A\cap\ker B$，等于纵向拼接矩阵的核。又 $\dim W=n-\operatorname{rank}B$，故由式~\eqref{eq:xie-ex-restriction-dimension} 得
\[
    \dim T(W)=(n-\operatorname{rank}B)
    -\left(n-\operatorname{rank}\begin{pmatrix}A\\B\end{pmatrix}\right),
\]
化简即得结论。
\end{solve}

\begin{example}
\label{ex:xie-ex-involution}
设 $T^2=I_V$，称 $T$ 为\emph{对合变换}（involution）。证明
\[
    V=\ker(T-I_V)\oplus\ker(T+I_V).
\]
\end{example}

\begin{solve}
任意 $v$ 可写为
\[
    v=\frac{v+Tv}{2}+\frac{v-Tv}{2}.
\]
由 $T^2=I_V$，第一项被 $T$ 保持，第二项被 $T$ 映为其相反向量。若 $Tv=v=-Tv$，则 $2v=0$，故交为零。相应分量投影是 $(I_V+T)/2$ 和 $(I_V-T)/2$。条件 $T\ne I_V$ 只保证第二个子空间非零，不保证第一个也非零，例如 $T=-I_V$。
\end{solve}

\begin{example}
\label{ex:xie-ex-rational-cyclic}
设 $V$ 是有理数域上的三维空间，$T(\alpha)=\beta$、$T(\beta)=\gamma$、$T(\gamma)=\alpha+\beta$，且 $\alpha\ne0$。证明 $\alpha,\beta,\gamma$ 线性无关。
\end{example}

\begin{solve}
先说明两个将用到的算术事实：多项式 $t^3-t-1$ 和 $t^3-t+1$ 都没有有理根。事实上，若约分后的 $a/b$ 是其中某个多项式的根，由整系数首一方程可知 $b=1$ 且 $a$ 整除常数项，所以只需检验 $a=1,-1$，代入均不为零。

若 $\alpha,\beta$ 相关，由 $\alpha\ne0$ 可设 $\beta=c\alpha$，其中 $c\in\mathbb Q$。于是 $\gamma=c^2\alpha$，再由 $T\gamma=\alpha+\beta$ 得 $c^3=c+1$，与第一个算术事实矛盾。因此 $\alpha,\beta$ 无关。若三向量相关，便可设 $\gamma=u\alpha+v\beta$，其中 $u,v\in\mathbb Q$；作用 $T$ 并与 $T\gamma=\alpha+\beta$ 比较，得到
\[
    uv=1,\qquad u+v^2=1.
\]
因此 $v\ne0$，消去 $u$ 得 $v^3-v+1=0$，与第二个算术事实矛盾。故三向量线性无关；有理数域条件在排除这些方程的根时使用，不能省略。
\end{solve}

\begin{example}
\label{ex:xie-ex-rectangular-multiplication}
设 $A\in M_m(K)$、$B\in M_n(K)$，定义 $T:M_{m\times n}(K)\to M_{m\times n}(K)$ 为 $T(X)=AXB$。求 $T$ 的秩、零度及核的一组基。
\end{example}

\begin{solve}
记 $r=\operatorname{rank}A$、$s=\operatorname{rank}B$，取可逆矩阵 $L_A,R_A,L_B,R_B$ 使
\[
    L_AAR_A=\operatorname{diag}(I_r,0),\qquad
    L_BBR_B=\operatorname{diag}(I_s,0).
\]
对 $X=R_AYL_B$，有
\[
    L_AT(X)R_B=
    \operatorname{diag}(I_r,0)\,Y\,\operatorname{diag}(I_s,0).
\]
左右两侧的坐标变换都可逆，而右侧仅保留 $Y$ 的前 $r$ 行、前 $s$ 列组成的块。因此
\[
    \operatorname{rank}T=rs,\qquad \dim\ker T=mn-rs,
\]
且核的一组基为
\[
    \{R_AE_{ij}L_B:1\leq i\leq m,\ 1\leq j\leq n,
    \ i>r\text{ 或 }j>s\}.
\]
这里 $E_{ij}$ 是 $m\times n$ 基础矩阵；上述条件恰好去除左上角 $r\times s$ 块。基向量的独立性与完备性由可逆变换 $Y\mapsto R_AYL_B$ 保证。
\end{solve}

\subsection{线性关系的保持与秩可加的同时标准型}

% 对应训练题三 14、15。
\begin{example}
\label{ex:xie-ex-relations-preservation}
给定 $n$ 维空间 $V$ 中的两个向量组 $S=(v_1,\ldots,v_m)$、$R=(w_1,\ldots,w_m)$，定义关系空间
\[
    \mathcal R_S=\{(a_1,\ldots,a_m)^{\mathsf T}\in K^m:
    \textstyle\sum_i a_iv_i=0\},
\]
并类似定义 $\mathcal R_R$。证明：存在 $T\in\mathcal L(V)$ 使 $T(v_i)=w_i$ 当且仅当 $\mathcal R_S\subseteq\mathcal R_R$；存在满足此条件的自同构，当且仅当 $\mathcal R_S=\mathcal R_R$。
\end{example}

\begin{solve}
关系空间是线性映射 $(a_i)\mapsto\sum_i a_iv_i$ 的核，所以是子空间。若存在 $T$，对任何线性关系作用 $T$，立即得到关系空间包含；若 $T$ 可逆，对逆映射同样论证，得到相等。反过来，若包含成立，在 $\operatorname{span}S$ 上定义
\[
    T_0\left(\sum_i a_iv_i\right)=\sum_i a_iw_i.
\]
两种表示之差属于 $\mathcal R_S$，从而也属于 $\mathcal R_R$，所以定义良好且线性。扩张 $T_0$ 到 $V$ 即得所需变换。若两关系空间相等，则此定义还给出 $\operatorname{span}S\to\operatorname{span}R$ 的同构；分别扩充两边的基为 $V$ 的基，再对应补充的基向量，即扩张为自同构。
\end{solve}

\begin{example}
\label{ex:xie-ex-additive-rank-normal-form}
设 $A,B\in M_{m\times n}(K)$，且 $\operatorname{rank}(A+B)=\operatorname{rank}A+\operatorname{rank}B$。记 $r=\operatorname{rank}A$、$s=\operatorname{rank}B$。证明存在可逆矩阵 $P\in\operatorname{GL}_m(K)$、$Q\in\operatorname{GL}_n(K)$，使
\[
    PAQ=\begin{pmatrix}I_r&0&0\\0&0&0\\0&0&0\end{pmatrix},
    \qquad
    PBQ=\begin{pmatrix}0&0&0\\0&I_s&0\\0&0&0\end{pmatrix},
\]
两矩阵的行分块大小均为 $r,s,m-r-s$，列分块大小均为 $r,s,n-r-s$。
\end{example}

\begin{solve}
将 $A,B$ 看作 $K^n\to K^m$ 的映射。由
\[
    \operatorname{Im}(A+B)\subseteq\operatorname{Im}A+\operatorname{Im}B,
    \qquad \operatorname{rank}(A+B)=r+s
\]
及和空间维数上界，得到 $\operatorname{Im}(A+B)=\operatorname{Im}A\oplus\operatorname{Im}B$。另一方面，纵向拼接矩阵满足
\[
    r+s=\operatorname{rank}(A+B)
    \leq\operatorname{rank}\begin{pmatrix}A\\B\end{pmatrix}
    \leq r+s,
\]
故 $N=\ker A\cap\ker B$ 的维数为 $n-r-s$，且由交和维数公式有 $\ker A+\ker B=K^n$。

取 $N$ 在 $\ker B$ 中的补空间 $X$、在 $\ker A$ 中的补空间 $Y$，则 $\dim X=r$、$\dim Y=s$，并有 $K^n=X\oplus Y\oplus N$。取基 $x_1,\ldots,x_r$、$y_1,\ldots,y_s$、$z_1,\ldots,z_{n-r-s}$，则 $Ax_i$ 是 $\operatorname{Im}A$ 的基，$By_j$ 是 $\operatorname{Im}B$ 的基；两组像共同线性无关，故可扩充为 $K^m$ 的基。在这些定义域与陪域基下，两映射恰有题述矩阵形式。若 $Q$ 的列为定义域新基、$F$ 的列为陪域新基，则取 $P=F^{-1}$ 即得结论。
\end{solve}

% 教材第 4 章（高代习题-4.pdf）整理内容至此结束。

% BEGIN textbook-xie-exercises-ch5
\chapter{多项式}
\label{chap:polynomials}

% 依据《高等代数习题（谢启鸿）》第 5 章，原书第 234--286 页整理。
% 第 5.1 节的概念分别融入相应专题；第 5.12 节按主题选取典型训练，不逐题誊录。
多项式把数的运算、函数的取值和线性代数中的变换联系起来。本章依照原书第 5 章的主题，从一元多项式的整除出发，建立最大公因式、互素性和唯一分解，再讨论根、系数域、对称多项式以及结式与判别式，最后把 Bézout 恒等式用于矩阵、线性变换和代数数。学习时应注意两条线索：\emph[tip]{整除关系可转化为带余除法或线性组合，根的条件可转化为因式及其重数的条件}。本章始终把数域 $\mathbb F$ 理解为复数域的一个子域，因而 $\mathbb F$ 是无限域，且其特征为零；涉及有理系数、实系数时再明确指定 $\mathbb Q$、$\mathbb R$。

\section{基本概念与运算}
\label{sec:poly-basic}

\begin{definition}
\label{def:poly-formal}
设 $x$ 是未定元，$a_0,\ldots,a_n\in\mathbb F$。形式表达式
\[
    f(x)=a_nx^n+a_{n-1}x^{n-1}+\cdots+a_1x+a_0
\]
称为 $\mathbb F$ 上的\emph{一元多项式}（univariate polynomial）。两个多项式相等，是指同次幂的系数分别相等。所有这类多项式构成的集合记为 $\mathbb F[x]$。若 $a_n\ne0$，则 $n$ 是 $f$ 的\emph{次数}（degree），记为 $\deg f=n$，$a_n$ 是首项系数；首项系数为 $1$ 的多项式称为\emph{首一多项式}（monic polynomial）。所有系数都为零的多项式称为零多项式，约定其次数为 $-\infty$。
\end{definition}

多项式是由系数确定的代数对象，未定元 $x$ 不预先表示某个数。把 $x$ 换为数域中的数，才得到相应的多项式函数。加法按同次幂的系数相加；数乘按每个系数同时乘以同一个数定义；乘法由分配律确定。具体地，若 $f=\sum_i a_ix^i$、$g=\sum_j b_jx^j$，则
\[
    fg=\sum_k\left(\sum_{i+j=k}a_ib_j\right)x^k.
\]
这些运算使 $\mathbb F[x]$ 成为交换代数（commutative algebra）：它既是以 $1,x,x^2,\ldots$ 为一组基的 $\mathbb F$ 上向量空间，又具有满足结合律、交换律和分配律的乘法，乘法与数乘相容。这里的乘法不是逐项系数相乘，而是系数的卷积。

\begin{proposition}
\label{prop:poly-degree}
若 $f,g\in\mathbb F[x]$ 均非零，$c\in\mathbb F\setminus\{0\}$，则
\[
    \deg(fg)=\deg f+\deg g,\qquad \deg(cf)=\deg f,
    \qquad \deg(f+g)\leq\max\{\deg f,\deg g\}.
\]
最后一个不等式也适用于零多项式。当 $\deg f\ne\deg g$ 时，最后一式取等号。
\end{proposition}
\begin{proof}
若 $f,g$ 的首项分别为 $a_nx^n,b_mx^m$，则 $fg$ 的最高次项是 $a_nb_mx^{n+m}$，而数域中 $a_nb_m\ne0$，故乘积的次数为 $n+m$。数乘的结论同理。相加不会产生超过原来最高次数的项；只有两者次数相等且首项系数互为相反数时，最高次项才可能消去。
\end{proof}

由次数公式，两个非零多项式的乘积不为零，因此 $\mathbb F[x]$ 是\emph{整环}（integral domain），可以约去等式两边相同的非零因子。另一个后果是：多项式有多项式乘法逆元，当且仅当它是非零常数。因此讨论因式分解时，非零常数是可忽略的可逆因子，但不是不可约因式。

\section{整除与带余除法}
\label{sec:poly-division}

\begin{definition}
若 $f,g\in\mathbb F[x]$，存在 $h\in\mathbb F[x]$ 使 $f=gh$，就称 $g$ \emph{整除}（divide）$f$，记为 $g\mid f$，并称 $g$ 是 $f$ 的因式。两个非零多项式若只相差一个非零常数倍，称为相伴多项式（associate polynomials）。
\end{definition}

整除关系具有传递性；若 $g\mid f_1$、$g\mid f_2$，则 $g$ 整除任意多项式线性组合 $u_1f_1+u_2f_2$。非零多项式相互整除，当且仅当它们相伴。任何多项式都整除零，而 $0\mid f$ 只在 $f=0$ 时成立；实际使用除法时必须另加除式非零的条件。

\begin{theorem}[带余除法]
\label{thm:poly-division}
设 $f,g\in\mathbb F[x]$ 且 $g\ne0$，则存在唯一的 $q,r\in\mathbb F[x]$，使
\begin{equation}
\label{eq:poly-division}
    f=gq+r,\qquad \deg r<\deg g.
\end{equation}
其中 $q$、$r$ 分别称为商和余式；$r=0$ 等价于 $g\mid f$。
\end{theorem}
\begin{proof}
若 $f=0$ 或 $\deg f<\deg g$，直接取 $q=0,r=f$。否则设 $f,g$ 的首项为 $ax^n,bx^m$，从 $f$ 中减去 $(a/b)x^{n-m}g$，使余下多项式的次数严格降低；重复这一过程，至多经过 $n-m+1$ 步便得到式~\eqref{eq:poly-division}。若另有 $f=gq_1+r_1$ 且 $\deg r_1<\deg g$，则 $g(q-q_1)=r_1-r$。当 $q\ne q_1$ 时，左边的次数至少为 $\deg g$，右边却小于 $\deg g$，矛盾，故商和余式都唯一。
\end{proof}

\begin{corollary}[余数定理与因式定理]
\label{cor:poly-remainder}
对 $a\in\mathbb F$，$f(x)$ 除以 $x-a$ 的余式是 $f(a)$，因此
\[
    (x-a)\mid f(x)\quad\Longleftrightarrow\quad f(a)=0.
\]
\end{corollary}
\begin{proof}
式~\eqref{eq:poly-division} 中的余式此时为常数，代入 $x=a$ 即得结论。
\end{proof}

\begin{example}
\label{ex:poly-divisibility-basic}
设 $g(x)=ax+b$，其中 $a\ne0$，证明 $g\mid f^2$ 当且仅当 $g\mid f$；再设 $g(x)=ax^2+bx+c$，其中 $abc\ne0$，且 $g\mid x^3+px^2+qx+r$，求系数必须满足的关系。
\end{example}
\begin{solve}
第一问中，$g$ 与 $x+b/a$ 相伴，因而 $g\mid f^2$ 等价于 $f(-b/a)^2=0$，也就等价于 $f(-b/a)=0$。第二问的商为一次多项式，设为 $mx+n$，比较系数得到
\[
    am=1,\qquad an+bm=p,\qquad bn+cm=q,\qquad cn=r.
\]
消去 $m,n$，得
\[
    \frac{ap-b}{a^2}=\frac{aq-c}{ab}=\frac rc.
\]
这说明低次整除问题往往可以直接由商的次数确定待定系数的个数，再比较系数求解。
\end{solve}

\begin{example}
\label{ex:poly-power-divisibility}
设 $a\in\mathbb F\setminus\{0\}$，$d,n$ 为正整数，证明
\[
    (x^d-a^d)\mid(x^n-a^n)\quad\Longleftrightarrow\quad d\mid n.
\]
并证明对非负整数 $m,n,p$，$x^2+x+1$ 整除 $x^{3m}+x^{3n+1}+x^{3p+2}$。
\end{example}
\begin{solve}
写 $n=qd+r$，其中 $0\leq r<d$，则
\[
    x^n-a^n=x^r(x^{qd}-a^{qd})+a^{qd}(x^r-a^r).
\]
第一项可被 $x^d-a^d$ 整除；若 $r>0$，第二项是非零的低次多项式，不能被它整除，所以整除成立当且仅当 $r=0$。这里 $a\ne0$ 必不可少，否则只需 $n\geq d$ 就有整除关系。第二问利用 $x^3-1=(x-1)(x^2+x+1)$，把待证多项式写为
\[
    (x^{3m}-1)+x(x^{3n}-1)+x^2(x^{3p}-1)+(1+x+x^2),
\]
每一项都含有所需因式。这种方法也可表述为模 $x^2+x+1$ 计算，并使用 $x^3\equiv1$。
\end{solve}

\section{最大公因式与互素多项式}
\label{sec:poly-gcd}

\begin{definition}
\label{def:poly-gcd}
对不全为零的 $f_1,\ldots,f_s\in\mathbb F[x]$，若首一多项式 $d$ 整除每个 $f_i$，且每个公因式都整除 $d$，则称 $d$ 为它们的\emph{最大公因式}（greatest common divisor），记为 $\gcd(f_1,\ldots,f_s)$。若 $d=1$，就称这些多项式互素（coprime）；若任意两个都互素，则称两两互素（pairwise coprime）。
\end{definition}

首一化保证最大公因式唯一，特别地，$\gcd(f,0)$ 是非零多项式 $f$ 的首一化。多个多项式互素不等于两两互素，例如 $x(x-1),x(x+1),(x-1)(x+1)$ 的共同最大公因式是 $1$，但任意两个都有非常数公因式。下文的最大公因式均采用首一约定，不再把相差非零常数倍的式子视为不同结果。

\begin{theorem}[Bézout 恒等式]
\label{thm:poly-bezout}
不全为零的 $f_1,\ldots,f_s\in\mathbb F[x]$ 的最大公因式存在，且有多项式 $u_1,\ldots,u_s$ 使
\begin{equation}
\label{eq:poly-bezout}
    f_1u_1+\cdots+f_su_s=\gcd(f_1,\ldots,f_s).
\end{equation}
特别地，它们互素当且仅当 $1$ 是它们的多项式线性组合。
\end{theorem}
\begin{proof}
在所有非零线性组合 $\sum_i f_iu_i$ 中，选取次数最小者并首一化，记为 $d$。将任一 $f_i$ 除以 $d$，其余式仍是各 $f_j$ 的线性组合；若余式非零，便与 $d$ 的次数最小性矛盾，故 $d\mid f_i$。而每个公因式必整除所有线性组合，特别整除 $d$，所以 $d$ 正是最大公因式。反之，若 $1$ 是线性组合，每个公因式均整除 $1$，因而只能是非零常数。
\end{proof}

实际计算两个多项式的最大公因式，可反复作带余除法：从 $f=q_1g+r_1$ 开始，继续将前一个除式除以当前非零余式，直至出现零余式。等式 $\gcd(f,g)=\gcd(g,r_1)$ 保证最后一个非零余式的首一化就是所求；从最后一步向前回代，还能得到 Bézout 系数。余式次数严格递减，保证此\emph{Euclid 算法}（Euclidean algorithm）经过有限步终止。

\begin{remark}
仅有 $d=fu+gv$ 不能说明 $d=\gcd(f,g)$。例如 $f=x,g=x^2$ 时，$d=x+x^2$ 是线性组合，却不是公因式。若同时知道 $d\mid f,d\mid g$，则每个公因式又整除 $d$，此时 $d$ 才与最大公因式相伴。Bézout 系数一般不唯一，因为可把 $(u,v)$ 替换为 $(u+tg,v-tf)$。
\end{remark}

\begin{proposition}
\label{prop:poly-coprime-rules}
若 $\gcd(f,g)=1$，则 $f\mid gh$ 蕴含 $f\mid h$。若 $\gcd(f_i,g_j)=1$ 对所有 $i,j$ 成立，则 $\prod_i f_i$ 与 $\prod_j g_j$ 互素。因此，对任意正整数 $m,n$，有
\[
    \gcd(f,g)=1\quad\Longleftrightarrow\quad \gcd(f^m,g^n)=1.
\]
此外，若 $f,g$ 互素，则 $fg$ 与 $f+g$ 互素；对任意 $h\in\mathbb F[x]$，$f(h(x))$ 与 $g(h(x))$ 也互素。
\end{proposition}
\begin{proof}
由 $fu+gv=1$ 乘以 $h$，可得 $h=fuh+gvh$，右边两项均可被 $f$ 整除，证明了第一条。若 $f$ 分别与 $g_1,g_2$ 互素，将两条 Bézout 恒等式相乘，得到 $1=fU+g_1g_2V$；反复应用这一结论便得到乘积的互素性。反向的幂次结论来自：$f,g$ 的公因式也整除它们的正整数次幂。$f$ 与 $f+g$ 的公因式恰是 $f,g$ 的公因式，对 $g$ 同理，再用乘积的结论即得 $fg$ 与 $f+g$ 互素。最后，在 Bézout 恒等式中代入 $x=h(x)$ 即得复合后的互素性。
\end{proof}

\begin{proposition}
\label{prop:poly-reduced-bezout}
若 $f,g$ 均为正次数多项式且互素，则存在唯一的 $u,v\in\mathbb F[x]$，使
\[
    fu+gv=1,\qquad \deg u<\deg g,\qquad \deg v<\deg f.
\]
\end{proposition}
\begin{proof}
先取任一 Bézout 表示 $fh+gk=1$，把 $h$ 除以 $g$，写为 $h=gq+u$，便有 $fu+g(k+fq)=1$。令 $v=k+fq$。若 $\deg v\geq\deg f$，则 $gv$ 的次数至少为 $\deg f+\deg g$，而 $fu$ 的次数严格小于这个数，不能相消得到常数，故 $\deg v<\deg f$。若还有另一组 $u_1,v_1$，相减得到 $f(u-u_1)=g(v_1-v)$。互素性给出 $g\mid u-u_1$，次数条件迫使 $u=u_1$，进而 $v=v_1$。
\end{proof}

\begin{example}
\label{ex:poly-gcd-powers}
设 $m,n$ 为正整数，$d=\gcd(m,n)$，证明 $\gcd(x^m-1,x^n-1)=x^d-1$。再设 $\gcd(f,g)=h$，证明对正整数 $N$，
\[
    \gcd(f^N,f^{N-1}g,\ldots,fg^{N-1},g^N)=h^N.
\]
\end{example}
\begin{solve}
若 $m=qn+r$，则
\[
    x^m-1=x^r(x^{qn}-1)+(x^r-1),
\]
而 $x^n-1$ 整除 $x^{qn}-1$，故最大公因式可化为 $\gcd(x^n-1,x^r-1)$。按整数的 Euclid 算法反复进行，最终得到 $x^d-1$。第二问中，$h^N$ 是全部所列多项式的公因式；又由 $fu+gv=h$，有
\[
    h^N=\sum_{j=0}^N\binom Nj f^{N-j}g^ju^{N-j}v^j.
\]
任何公因式都整除右边，因而整除 $h^N$，结论成立。
\end{solve}

\begin{theorem}[多项式的中国剩余定理]
\label{thm:poly-crt}
设正次数多项式 $g_1,\ldots,g_s$ 两两互素，任给 $r_1,\ldots,r_s\in\mathbb F[x]$，存在 $f\in\mathbb F[x]$，使 $f\equiv r_i\pmod{g_i}$ 对所有 $i$ 成立。令 $G=\prod_i g_i$，则所有这样的 $f$ 模 $G$ 同余，且有唯一一个满足 $\deg f<\deg G$ 的代表。
\end{theorem}
\begin{proof}
令 $G_i=G/g_i$，则 $\gcd(G_i,g_i)=1$，可取 $a_i,b_i$ 使 $a_iG_i+b_ig_i=1$。于是 $E_i=a_iG_i$ 模 $g_i$ 同余于 $1$，模其余各 $g_j$ 同余于零，故 $\sum_i r_iE_i$ 满足全部要求。两个解之差可被每个 $g_i$ 整除，由两两互素性可被 $G$ 整除。再对 $G$ 作带余除法，即得到唯一的低次代表。
\end{proof}

\section{不可约多项式与因式分解}
\label{sec:poly-factorization}

\begin{definition}
\label{def:poly-irreducible}
正次数多项式 $p\in\mathbb F[x]$ 若不能分解为两个正次数的 $\mathbb F$ 上多项式之积，称为 $\mathbb F$ 上的\emph{不可约多项式}（irreducible polynomial），否则称为可约多项式（reducible polynomial）。
\end{definition}

\emph[caution]{不可约性必须连同系数域一起说明}。例如 $x^2-2$ 在 $\mathbb Q$ 上不可约，在 $\mathbb R$ 上可约；$x^2+1$ 在 $\mathbb R$ 上不可约，在 $\mathbb C$ 上可约。扩大系数域可能使不可约多项式变得可约，但已有的分解始终保留。反之，互素性在数域扩张下仍然成立，因为原来的 Bézout 恒等式仍可使用。

\begin{theorem}[不可约多项式的素性]
\label{thm:poly-prime}
正次数多项式 $p$ 在 $\mathbb F$ 上不可约，当且仅当对任意 $f,g\in\mathbb F[x]$，$p\mid fg$ 蕴含 $p\mid f$ 或 $p\mid g$。
\end{theorem}
\begin{proof}
若 $p$ 不可约，则 $\gcd(p,f)$ 只能是 $1$ 或 $p$ 的首一化；在 $p\nmid f$ 时，命题~\ref{prop:poly-coprime-rules} 给出 $p\mid g$。反之，若 $p=p_1p_2$ 且两个因子的次数都小于 $\deg p$，则 $p$ 整除乘积，却不整除任一个因子，与所述性质矛盾。
\end{proof}

\begin{theorem}[唯一因式分解]
\label{thm:poly-unique-factorization}
任一非零的 $f\in\mathbb F[x]$ 都能唯一写为
\begin{equation}
\label{eq:poly-standard-factorization}
    f=c\prod_{i=1}^s p_i^{e_i},
\end{equation}
其中 $c\in\mathbb F\setminus\{0\}$，$p_i$ 是互不相同的首一不可约多项式，$e_i$ 是正整数；唯一性允许调换因子的次序。非零常数对应空乘积。
\end{theorem}
\begin{proof}
对次数归纳。正次数多项式若不可约，首一化即可；若可约，则分解成两个次数更低的因子，再分别归纳，得到存在性。对两个分解式，一边的不可约因子由素性必整除另一边的某个不可约因子，这两个首一因子只能相等。约去它们，重复此过程便确定所有因子及其重数，最后由首项系数确定 $c$。
\end{proof}

\begin{proposition}
\label{prop:poly-gcd-lcm}
对非零多项式，采用公共不可约因子表写成
\[
    f=c\prod_i p_i^{e_i},\qquad g=d\prod_i p_i^{h_i},\qquad e_i,h_i\geq0,
\]
则首一最大公因式和首一\emph{最小公倍式}（least common multiple）分别为
\begin{equation}
\label{eq:poly-gcd-lcm}
    \gcd(f,g)=\prod_i p_i^{\min\{e_i,h_i\}},\qquad
    \operatorname{lcm}(f,g)=\prod_i p_i^{\max\{e_i,h_i\}}.
\end{equation}
最小公倍式是可被 $f,g$ 整除、并整除每个公倍式的首一多项式。因而
\[
    fg=cd\,\gcd(f,g)\operatorname{lcm}(f,g),
\]
且对任意正整数 $N$，最大公因式、最小公倍式都与同时取 $N$ 次幂相容；此外 $g^N\mid f^N$ 当且仅当 $g\mid f$。
\end{proposition}
\begin{proof}
整除等价于每个不可约因子的指数都不超过被除式中相应的指数，故公因式取各指数的最小值，公倍式取最大值。利用 $\min\{e,h\}+\max\{e,h\}=e+h$，以及所有指数同时乘以 $N$ 不改变其大小关系，即得余下结论。
\end{proof}

\begin{example}
\label{ex:poly-primary-power}
证明：非零非常数多项式 $f$ 只含一种首一不可约因子，当且仅当对每个非常数多项式 $g$，或者 $f,g$ 互素，或者 $f$ 整除 $g$ 的某个正整数次幂。
\end{example}
\begin{solve}
若 $f=cp^e$，与 $f$ 不互素的 $g$ 必含因子 $p$，所以 $f\mid g^e$。反之，若 $f=cp^eh$，其中 $\deg h>0$ 且 $p\nmid h$，取 $g=p$，则 $f,g$ 不互素，却没有任何 $N$ 能使 $f\mid p^N$，矛盾。这里结论是 $f$ 为不可约多项式幂的非零常数倍，除非先假定 $f$ 首一，不能省去常数倍。
\end{solve}

\begin{definition}
\label{def:poly-minimal}
若 $\alpha\in\mathbb C$ 是某个非零有理系数多项式的根，则称 $\alpha$ 为\emph{代数数}（algebraic number）。使 $m_\alpha(\alpha)=0$ 的非零首一有理系数多项式中，次数最小者称为 $\alpha$ 在 $\mathbb Q$ 上的\emph{极小多项式}（minimal polynomial）。
\end{definition}

\begin{theorem}
\label{thm:poly-minimal}
代数数的极小多项式存在且唯一，并有
\[
    f(\alpha)=0\quad\Longleftrightarrow\quad m_\alpha\mid f
    \qquad(f\in\mathbb Q[x]).
\]
一个首一有理系数多项式是 $\alpha$ 的极小多项式，当且仅当它以 $\alpha$ 为根且在 $\mathbb Q$ 上不可约。
\end{theorem}
\begin{proof}
次数是非负整数，故在非空的消去多项式集合中可取最小次数者，再首一化。将任一 $f$ 除以它，得 $f=m_\alpha q+r$；若 $f(\alpha)=0$，则 $r(\alpha)=0$，最小性迫使 $r=0$。任意两个极小多项式因此相互整除，首一性保证它们相等。如果 $m_\alpha$ 可约，由乘积在 $\alpha$ 处为零，其中一个低次因子也在 $\alpha$ 处为零，矛盾。反之，以 $\alpha$ 为根的首一不可约多项式含有因子 $m_\alpha$，只能与之相等。
\end{proof}

同样的论证适用于任意数域 $\mathbb F$ 上的代数元。特别地，若 $p\in\mathbb F[x]$ 不可约，且它的某个复根 $\alpha$ 也是 $f\in\mathbb F[x]$ 的根，则 $p\mid f$：否则 $p,f$ 互素，将 Bézout 恒等式代入 $\alpha$ 就得到 $0=1$。因此，\emph[tip]{在原系数域中识别一个根的极小多项式，可以一次确定它必须带来的全部复根}。

\section{多项式函数、根与重数}
\label{sec:poly-roots}

\begin{definition}
对 $f\in\mathbb F[x]$，代入 $a\in\mathbb F$ 得到的映射 $a\mapsto f(a)$ 称为相应的\emph{多项式函数}（polynomial function）。若 $f(\alpha)=0$，就称 $\alpha$ 为 $f$ 的\emph{根}或\emph{零点}（root or zero）。根也可以在包含 $\mathbb F$ 的更大数域中讨论，此时应说明取值范围。
\end{definition}

\begin{theorem}
\label{thm:poly-root-bound}
次数为 $n$ 的非零一元多项式，在任意包含其系数域的数域中，至多有 $n$ 个互不相同的根。因此，若两个一元多项式在无限多个点处取值相同，它们就是同一多项式；若它们的次数均不超过 $n$，则在 $n+1$ 个不同点处取值相同已足够。
\end{theorem}
\begin{proof}
若 $\alpha$ 是根，因式定理给出 $f=(x-\alpha)g$。其他不同于 $\alpha$ 的根必须是 $g$ 的根，对次数归纳即得根数上界。再对两个多项式的差应用这一结论即可。
\end{proof}

\begin{proposition}
\label{prop:poly-values-field}
设 $\mathbb F\subseteq\mathbb K\subseteq\mathbb C$，$f\in\mathbb K[x]$ 且 $\deg f\leq n$。若在 $n+1$ 个不同的 $\mathbb F$ 中元素处，$f$ 的值都属于 $\mathbb F$，则 $f\in\mathbb F[x]$。
\end{proposition}
\begin{proof}
设这些点为 $a_0,\ldots,a_n$。多项式
\[
    P(x)=\sum_{i=0}^n f(a_i)
    \prod_{j\ne i}\frac{x-a_j}{a_i-a_j}
\]
的系数均在 $\mathbb F$ 中，次数不超过 $n$，且 $P(a_i)=f(a_i)$。由定理~\ref{thm:poly-root-bound}，$f=P$。
\end{proof}

这同时解释了两个常见结论：在所有实数处都取实值的复系数多项式必为实系数多项式；在所有有理数处都取有理值的多项式必为有理系数多项式。这里起决定作用的是足够多的取值条件，而不是连续性。相反，$\sin x$ 在实数域上不可能是多项式函数，因为它非恒为零，却有无限多个根。

\begin{definition}
对 $f(x)=\sum_{j=0}^n a_jx^j$，定义其\emph{形式导数}（formal derivative）为 $f'(x)=\sum_{j=1}^n ja_jx^{j-1}$。若非零多项式 $f$ 满足 $(x-\alpha)^k\mid f$ 而 $(x-\alpha)^{k+1}\nmid f$，则称 $\alpha$ 为 $f$ 的 $k$ \emph{重根}（root of multiplicity $k$）；$k=1$ 时称为单根。
\end{definition}

形式导数由系数直接定义，不依赖函数极限；它满足线性法则和乘积法则 $(fg)'=f'g+fg'$，这些法则可在单项式上验证后按线性展开。由于数域的特征为零，次数为 $n\geq1$ 的多项式的导数次数恰为 $n-1$。

\begin{theorem}[重根的导数判据]
\label{thm:poly-multiple-root}
非零多项式 $f$ 的根 $\alpha$ 恰为 $k$ 重根，当且仅当
\[
    f(\alpha)=f'(\alpha)=\cdots=f^{(k-1)}(\alpha)=0,
    \qquad f^{(k)}(\alpha)\ne0.
\]
\end{theorem}
\begin{proof}
把 $f$ 按 $x-\alpha$ 的幂展开为 $\sum_{j=0}^n c_j(x-\alpha)^j$，逐项求导可得 $c_j=f^{(j)}(\alpha)/j!$。因此最低的非零幂次恰为 $k$，等价于上述导数条件。特别地，若 $f=(x-\alpha)^kg$ 且 $g(\alpha)\ne0$，则 $f^{(k)}(\alpha)=k!g(\alpha)$。
\end{proof}

\begin{proposition}[重因式与无重因式部分]
\label{prop:poly-squarefree}
设 $f=c\prod_i p_i^{e_i}$ 是正次数多项式的标准分解，则
\[
    \gcd(f,f')=\prod_i p_i^{e_i-1}.
\]
因此 $f$ 没有重因式，当且仅当 $\gcd(f,f')=1$；$f/\gcd(f,f')$ 没有重因式，且与 $f$ 含有相同的不可约因子，只是不再保留重数。数域上的不可约多项式在复数域内也没有重根。
\end{proposition}
\begin{proof}
对固定的 $p=p_i$ 写 $f=p^eh$，其中 $p\nmid h$。由不可约性和 $\deg p'<\deg p$，有 $\gcd(p,p')=1$，从而
\[
    f'=p^{e-1}(ep'h+ph')
\]
中括号内的多项式不能被 $p$ 整除，故 $p$ 在 $f'$ 中的指数恰为 $e-1$，结论由式~\eqref{eq:poly-gcd-lcm} 得到。若 $p$ 在复数域有重根，则它与 $p'$ 在该根处同时为零，与原系数域上的 Bézout 恒等式矛盾。
\end{proof}

\begin{example}
\label{ex:poly-derivative-divides}
设 $\deg f=n\geq1$。证明 $f'\mid f$ 当且仅当存在 $a\in\mathbb F$ 和 $c\in\mathbb F\setminus\{0\}$，使 $f=c(x-a)^n$。
\end{example}
\begin{solve}
若 $f'\mid f$，比较次数和首项系数可写成 $f=(x-a)f'/n$，其中 $a\in\mathbb F$。将 $f$ 展开为 $\sum_{j=0}^n c_j(x-a)^j$，代入此式，得到 $c_j=(j/n)c_j$，所以 $j<n$ 时 $c_j=0$，只有 $c_n\ne0$。反向结论直接求导即可验证。
\end{solve}

\begin{example}
\label{ex:poly-finite-roots-applications}
若存在非零 $a\in\mathbb F$ 使 $f(x+a)=f(x)$，证明 $f$ 为常数；若非常数多项式 $f$ 满足 $f(x)\mid f(x^m)$，其中整数 $m>1$，证明它的每个复根不是零就是某个单位根。
\end{example}
\begin{solve}
第一问中，$f(ka)=f(0)$ 对所有非负整数 $k$ 成立。由于 $a\ne0$ 且特征为零，这些点互不相同，根数上界迫使 $f-f(0)=0$。第二问中，若 $\alpha$ 是根，则 $\alpha^m,\alpha^{m^2},\ldots$ 也都是根。不同根只有有限多个，所以存在 $r>s\geq0$ 使 $\alpha^{m^r}=\alpha^{m^s}$。若 $\alpha\ne0$，便有 $\alpha^{m^s(m^{r-s}-1)}=1$。这一结论只是必要条件，不能仅凭根都是单位根就反推原来的整除关系。
\end{solve}

\begin{example}
\label{ex:poly-interpolation-value}
设 $n\geq1$，$\deg f\leq n$，且 $f(k)=k/(k+1)$ 对 $k=0,1,\ldots,n$ 成立，求 $f(n+1)$。
\end{example}
\begin{solve}
令 $H(x)=(x+1)f(x)-x$，则 $\deg H\leq n+1$，且 $0,1,\ldots,n$ 都是它的根，故 $H=Cx(x-1)\cdots(x-n)$。代入 $x=-1$，由 $H(-1)=1$ 得 $C=(-1)^{n+1}/(n+1)!$，于是
\[
    f(n+1)=\frac{n+1+(-1)^{n+1}}{n+2}
    =\begin{cases}1,&n\text{ 为奇数},\\n/(n+2),&n\text{ 为偶数}.\end{cases}
\]
关键是先构造一个具有已知零点的低次多项式，而不直接求出 $f$ 的全部系数。
\end{solve}

\section{复系数多项式与根和系数的关系}
\label{sec:poly-complex}

\begin{theorem}[代数基本定理]
\label{thm:poly-fundamental}
每个正次数的复系数一元多项式至少有一个复根。因而，若 $f\in\mathbb C[x]$ 的次数为 $n\geq1$、首项系数为 $a_n$，则存在 $\alpha_1,\ldots,\alpha_n\in\mathbb C$，使
\[
    f(x)=a_n\prod_{i=1}^n(x-\alpha_i).
\]
根按重数计算恰有 $n$ 个；复数域上的不可约多项式恰为一次多项式。
\end{theorem}

这里采用代数基本定理本身，不在本章证明；其后的分解结论由因式定理反复降低次数得到。它保证可以在复数域内完整地描述任意数域上多项式的根，但并不说明这些根属于原来的系数域。

\begin{theorem}[Vieta 定理]
\label{thm:poly-vieta}
若首一多项式 $f(x)=x^n+c_1x^{n-1}+\cdots+c_n$ 的复根按重数记为 $\alpha_1,\ldots,\alpha_n$，则对 $1\leq k\leq n$，
\begin{equation}
\label{eq:poly-vieta}
    \sum_{1\leq i_1<\cdots<i_k\leq n}\alpha_{i_1}\cdots\alpha_{i_k}
    =(-1)^kc_k.
\end{equation}
\end{theorem}
\begin{proof}
展开 $\prod_i(x-\alpha_i)$ 时，$x^{n-k}$ 的系数来自恰选 $k$ 个根的乘积，符号为 $(-1)^k$，再与 $f$ 比较系数即可。
\end{proof}

\begin{example}
\label{ex:poly-cubic-progressions}
设 $f(x)=x^3+px^2+qx+r$。证明其三个复根可排列为等差数列，当且仅当 $2p^3-9pq+27r=0$；若 $r\ne0$，它们可排列为具有非零公比的等比数列，当且仅当 $rp^3=q^3$。
\end{example}
\begin{solve}
等差情形中，若根为 $c-d,c,c+d$，则 Vieta 定理给出 $c=-p/3$，所以 $f(-p/3)=0$，这恰是第一个条件。反之，若 $c=-p/3$ 是根，另外两根之和为 $-p-c=2c$，可以写为 $c-d,c+d$。

等比情形中，设三根为 $c/d,c,cd$，则 $r=-c^3$，且 $q=-pc$，故 $q^3=rp^3$。反之，若 $p\ne0$，令 $c=-q/p$，由条件得到 $c^3=-r$ 和 $f(c)=0$；另外两根的积为 $-r/c=c^2$，因它们均非零，可写为 $c/d,cd$。若 $p=0$，条件迫使 $q=0$，任取一个根 $c$，另外两根的积同样为 $c^2$，仍得到所需排列。这里讨论的是复数公比；若要求根与公比都是实数，还需检查相应的实数条件。
\end{solve}

\begin{example}
\label{ex:poly-transform-roots}
设三次方程 $x^3+px^2+qx+r=0$ 的根为 $\alpha,\beta,\gamma$。当 $r\ne0$ 时，求根的倒数的平方和；并求以 $\alpha^3,\beta^3,\gamma^3$ 为根的首一三次多项式。
\end{example}
\begin{solve}
由根与系数的关系，
\[
    \frac1{\alpha^2}+\frac1{\beta^2}+\frac1{\gamma^2}
    =\frac{(\alpha\beta+\alpha\gamma+\beta\gamma)^2
    -2\alpha\beta\gamma(\alpha+\beta+\gamma)}{(\alpha\beta\gamma)^2}
    =\frac{q^2-2pr}{r^2}.
\]
再利用 $u^3+v^3+w^3=(u+v+w)^3-3(u+v+w)(uv+uw+vw)+3uvw$，得到
\[
    \alpha^3+\beta^3+\gamma^3=-p^3+3pq-3r,
    \qquad (\alpha\beta)^3+(\alpha\gamma)^3+(\beta\gamma)^3=q^3-3pqr+3r^2.
\]
三根的立方的乘积为 $-r^3$，所以所求多项式为
\[
    t^3+(p^3-3pq+3r)t^2+(q^3-3pqr+3r^2)t+r^3.
\]
\end{solve}

\begin{proposition}[倒数多项式]
\label{prop:poly-reciprocal}
若 $f(x)=a_nx^n+\cdots+a_0$ 且 $a_na_0\ne0$，定义\emph{倒数多项式}（reciprocal polynomial）$f^*(x)=x^nf(1/x)=a_0x^n+\cdots+a_n$。其根恰为 $f$ 的各根的倒数，重数不变；而且 $f$ 在 $\mathbb F$ 上可约当且仅当 $f^*$ 在 $\mathbb F$ 上可约。
\end{proposition}
\begin{proof}
由 $f=a_n\prod_i(x-\alpha_i)$，有 $f^*=a_n\prod_i(1-\alpha_ix)$，且各 $\alpha_i\ne0$，故得到根的结论。若 $f=gh$ 是正次数因子的分解，则 $g,h$ 的常数项均非零，从而 $f^*=g^*h^*$ 仍是正次数分解。再用 $(f^*)^*=f$ 得到反向结论。
\end{proof}

\begin{example}
\label{ex:poly-unity-roots}
令 $\zeta=\cos(2\pi/n)+i\sin(2\pi/n)$，其中 $n\geq2$。证明 $\zeta^{mk}$（$0\leq k<n$）遍历 $x^n-1$ 的全部根，当且仅当 $\gcd(m,n)=1$；并求 $\prod_{k=1}^{n-1}(1-\zeta^k)$。若 $\deg f<n$，再求 $\sum_{k=0}^{n-1}f(\zeta^k)$。
\end{example}
\begin{solve}
若 $m,n$ 互素，$\zeta^{ms}=\zeta^{mt}$ 就蕴含 $n\mid m(s-t)$，进而 $n\mid s-t$，所以这些根互不相同。若 $d=\gcd(m,n)>1$，则 $\zeta^{m(n/d)}=1$，其中 $0<n/d<n$，出现重复，不能遍历。又由
\[
    \frac{x^n-1}{x-1}=1+x+\cdots+x^{n-1}
    =\prod_{k=1}^{n-1}(x-\zeta^k),
\]
令 $x=1$ 得乘积为 $n$。最后，$\sum_{k=0}^{n-1}\zeta^{kj}=0$ 对 $1\leq j<n$ 成立，这是有限等比数列求和；把 $f$ 按幂展开，便得 $\sum_{k=0}^{n-1}f(\zeta^k)=nf(0)$。
\end{solve}

\section{实系数多项式}
\label{sec:poly-real}

\begin{theorem}
\label{thm:poly-real-factorization}
实系数多项式的非实根按共轭成对出现，且重数相同。实数域上的首一不可约多项式恰为一次多项式和判别式小于零的二次多项式。因此每个实系数多项式都可以分解为实的一次因式与实的不可约二次因式的乘积；奇数次实系数多项式必有实根。
\end{theorem}
\begin{proof}
实系数给出 $f(\overline z)=\overline{f(z)}$，对各阶导数同样成立，由重根判据可知共轭根的重数相同。在复数分解中，将 $a\pm bi$（$b\ne0$）的因子合并为
\[
    (x-a-bi)(x-a+bi)=x^2-2ax+a^2+b^2,
\]
其判别式为 $-4b^2<0$，没有实根，故在实数域上不可约。分解中其余因子均为一次，便得全部结论。
\end{proof}

若实多项式的全部系数均为正，则代入任意 $x\geq0$ 都得到正数，故没有非负实根；全部系数均负时同理。若 $f(-x)$ 的全部系数同号，则 $f$ 没有非正实根。这类符号判据只能排除相应半轴上的根，不能给出所有复根的信息。若三次多项式 $x^3+px^2+qx+r$ 的三个根都是实数，则由
\[
    (\alpha-\beta)^2+(\alpha-\gamma)^2+(\beta-\gamma)^2=2(p^2-3q)
\]
得到 $p^2\geq3q$；仅此不等式并不足以保证三个根都为实数，例如 $x^3+1$ 满足它却有一对非实根。

\begin{example}
\label{ex:poly-real-two-squares}
设 $f\in\mathbb R[x]$ 首一且没有实根，证明存在实系数多项式 $u,v$，使 $f=u^2+v^2$。
\end{example}
\begin{solve}
从每对共轭根中选取位于上半平面的一个，按重数构造首一多项式 $h$。则 $f=h\overline h$，其中 $\overline h$ 表示将各系数取共轭。把系数实部、虚部分开写成 $h=u+iv$，有 $\overline h=u-iv$，因而 $f=u^2+v^2$。若 $f=1$，可取 $u=1,v=0$。
\end{solve}

\begin{example}
\label{ex:poly-cubic-stability}
设 $p,q,r\in\mathbb R$。证明 $x^3+px^2+qx+r$ 的全部根的实部均小于零，当且仅当
\[
    p>0,\qquad r>0,\qquad pq>r.
\]
\end{example}
\begin{solve}
先证必要性。若三根都是负实数 $\alpha,\beta,\gamma$，则 $p>0,r>0$，且
\[
    pq-r=-(\alpha+\beta)(\alpha+\gamma)(\beta+\gamma)>0.
\]
若三根为 $s,a+bi,a-bi$，其中 $s<0,a<0,b>0$，则 $p=-(s+2a)>0$、$r=-s(a^2+b^2)>0$，并且
\[
    pq-r=-2a\bigl((s+a)^2+b^2\bigr)>0.
\]
反之，由所给三个条件可知 $q>0$，所以四个系数全正，全部实根均为负数。如果还有一对非实根 $a\pm bi$，取其余实根 $s<0$，由同一恒等式及 $b>0$ 得 $a<0$。两种情形都得到全部根的实部为负。
\end{solve}

\begin{example}
对实数 $a,b,c$，证明 $a,b,c$ 全为正数，当且仅当 $a+b+c>0$、$ab+ac+bc>0$、$abc>0$。
\end{example}
\begin{solve}
必要性直接成立。反之，令 $P=a+b+c,Q=ab+ac+bc,R=abc$，则 $t^3+Pt^2+Qt+R=(t+a)(t+b)(t+c)$ 的系数均为正，没有非负实根，所以它的实根 $-a,-b,-c$ 全为负数，即 $a,b,c>0$。
\end{solve}

\section{有理系数多项式与不可约性判别}
\label{sec:poly-rational}

\begin{theorem}[有理根的必要条件]
\label{thm:poly-rational-root}
设 $f(x)=a_nx^n+\cdots+a_0\in\mathbb Z[x]$，$a_n\ne0$。若既约分数 $p/q$ 是它的根，其中 $q>0$，则 $p\mid a_0$、$q\mid a_n$。特别地，首一整系数多项式的有理根必为整数。
\end{theorem}
\begin{proof}
将 $f(p/q)=0$ 乘以 $q^n$，分别移项得 $q\mid a_np^n$ 和 $p\mid a_0q^n$。由 $p,q$ 互素，便得所需结论；$p=0$ 时原等式直接给出 $a_0=0$。
\end{proof}

\begin{definition}
非零整系数多项式各系数的正最大公因数称为其内容（content）；内容为 $1$ 的多项式称为\emph{本原多项式}（primitive polynomial）。
\end{definition}

\begin{lemma}[Gauss 引理]
\label{lem:poly-gauss}
两个本原整系数多项式的乘积仍为本原多项式。因此，一个整系数多项式在 $\mathbb Q$ 上可约，当且仅当它能写为两个正次数整系数多项式之积。
\end{lemma}
\begin{proof}
若乘积的全部系数都能被某素数 $\ell$ 整除，将系数模 $\ell$ 化简，就得到两个非零的 $\mathbb F_\ell$ 上多项式之积为零，与域上多项式的次数公式矛盾。故本原多项式的乘积本原。

对有理系数分解 $f=gh$，分别清除分母、除去内容，可写成 $g=cu,h=dv$，其中 $c,d\in\mathbb Q\setminus\{0\}$，$u,v\in\mathbb Z[x]$ 本原。于是 $f=cd\,uv$。由于 $uv$ 本原，若 $cd$ 的既约分母大于 $1$，便不可能使全部系数为整数，因此 $cd\in\mathbb Z$，从而 $f=((cd)u)v$ 是整系数分解。反向结论直接成立。这里 $\mathbb F_\ell$ 表示模素数 $\ell$ 的剩余类构成的域。
\end{proof}

\begin{theorem}[Eisenstein 判别法]
\label{thm:poly-eisenstein}
设 $f(x)=a_nx^n+\cdots+a_0\in\mathbb Z[x]$，$n\geq1$。若存在素数 $\ell$，使
\[
    \ell\nmid a_n,\qquad \ell\mid a_j\ (0\leq j<n),
    \qquad \ell^2\nmid a_0,
\]
则 $f$ 在 $\mathbb Q$ 上不可约。
\end{theorem}
\begin{proof}
若 $f=gh$ 是两个正次数整系数多项式的乘积，由 $\ell\nmid a_n$，它们的首项系数模 $\ell$ 均非零，降模不改变次数。又 $\overline g\,\overline h=\overline{a_n}x^n$，故两个因子降模后均为单项式：否则其最低非零次数之和严格小于其最高次数之和，乘积不可能只有一项。两者的次数均正，所以它们的常数项都能被 $\ell$ 整除，迫使 $\ell^2\mid a_0$，矛盾。由 Gauss 引理即得结论。
\end{proof}

\begin{example}
\label{ex:poly-eisenstein-examples}
设 $\ell_1,\ldots,\ell_s$ 为互不相同的素数，证明 $x^n-\ell_1\cdots\ell_s$ 对任意 $n\geq1$ 在 $\mathbb Q$ 上不可约；并证明 $x^8+1$ 在 $\mathbb Q$ 上不可约。
\end{example}
\begin{solve}
第一问任选一个 $\ell_i$ 应用 Eisenstein 判别法即可。第二问直接使用素数 $2$ 不满足常数项被 $2$ 整除的条件，但作平移 $x=y+1$ 后，有
\[
    (y+1)^8+1=y^8+8y^7+28y^6+56y^5+70y^4+56y^3+28y^2+8y+2,
\]
满足素数 $2$ 的全部条件。代换 $x=y+1$ 及其逆代换均保持正次数分解，因此原多项式也不可约。
\end{solve}

有理根判据只能排除一次因式。\emph[caution]{次数至少为四时，没有有理根并不意味着不可约}，例如 $(x^2-2)(x^2-3)$ 没有有理根却可约。对二次或三次多项式，可约必有一次因式，因此这两种次数下，无有理根才等价于在 $\mathbb Q$ 上不可约。另一个便于使用的排根方法是降模：若首一整系数多项式 $f$ 的 $f(0),f(1)$ 都为奇数，则整数代入后，偶数输入与 $0$ 同余、奇数输入与 $1$ 同余，所得值都为奇数，故没有整数根，也就没有有理根。

\begin{proposition}[双二次多项式的可约性]
\label{prop:poly-biquadratic}
设 $p,q\in\mathbb Q$。多项式 $x^4+px^2+q$ 在 $\mathbb Q$ 上可约，当且仅当下列两种条件至少有一种成立：$p^2-4q$ 是有理数的平方；或者存在 $s,t\in\mathbb Q$，使 $q=s^2$ 且 $2s-p=t^2$。
\end{proposition}
\begin{proof}
若有有理根 $u$，则 $u^2$ 是 $y^2+py+q$ 的有理根，故 $p^2-4q$ 为有理数的平方。其余可约情形必有分解
\[
    x^4+px^2+q=(x^2+ax+b)(x^2-ax+d),
\]
其中 $a,b,d\in\mathbb Q$。比较系数得 $b+d-a^2=p$、$a(d-b)=0$、$bd=q$。当 $a=0$ 时，$p^2-4q=(b-d)^2$；当 $d=b$ 时，取 $s=b,t=a$ 得第二种条件。反之，若 $p^2-4q=k^2$，则有分解
\[
    \left(x^2+\frac{p+k}{2}\right)\left(x^2+\frac{p-k}{2}\right);
\]
若 $q=s^2,2s-p=t^2$，则有分解 $(x^2+tx+s)(x^2-tx+s)$。
\end{proof}

\begin{example}
\label{ex:poly-conjugate-radicals}
求 $\sqrt2+\sqrt3$ 的极小多项式；并说明有理系数多项式以 $\sqrt[n]{2}$ 为根时，还必须有哪些复根。
\end{example}
\begin{solve}
令 $\alpha=\sqrt2+\sqrt3$，平方后得 $\alpha^2-5=2\sqrt6$，再平方得 $\alpha^4-10\alpha^2+1=0$。对 $p=-10,q=1$，命题~\ref{prop:poly-biquadratic} 中的第一种条件需要 $96$ 为有理数平方，第二种条件需要 $12$ 或 $8$ 为有理数平方，均不成立，故极小多项式为 $x^4-10x^2+1$。于是以 $\alpha$ 为根的有理系数多项式也必须以 $\pm\sqrt2\pm\sqrt3$ 的四种符号组合为根。

由于 $x^n-2$ 由 Eisenstein 判别法不可约，它是实数 $\sqrt[n]{2}$ 的极小多项式。因此任何以该数为根的有理系数多项式都可被 $x^n-2$ 整除，必须同时含有根 $\sqrt[n]{2}\zeta^k$（$0\leq k<n$），其中 $\zeta=\cos(2\pi/n)+i\sin(2\pi/n)$。特别地，次数小于 $n$ 的非零有理系数多项式不可能以 $\sqrt[n]{2}$ 为根。
\end{solve}

更一般地，若 $a,b,c\in\mathbb Q$ 且 $\sqrt c\notin\mathbb Q$，则有理系数多项式以 $a+b\sqrt c$ 为根时，也以 $a-b\sqrt c$ 为根。因为代入后可整理为 $A+B\sqrt c$，其中 $A,B\in\mathbb Q$；其值为零迫使 $A=B=0$，再改变根号的符号仍为零。若涉及两个独立的平方根，应核对所用的线性无关条件；尤其在 $a=0$ 或 $b=0$ 时，不能把 $(x^2-a^2c-b^2d)^2-4a^2b^2cd$ 一概称为次数为四的极小多项式。

\begin{example}
\label{ex:poly-odd-degree-root-sum}
设 $f\in\mathbb Q[x]$ 不可约，其次数是大于 $1$ 的奇数。证明它的两个不同复根之和不可能为有理数。
\end{example}
\begin{solve}
将 $f$ 首一化，设其次数为奇数 $n$。若两根之和为 $r\in\mathbb Q$，则 $f(x)$ 与 $f(r-x)$ 有公根，由不可约性可知前者整除后者。比较次数和首项系数，得到 $f(r-x)=-f(x)$。令 $x=r/2$，便有 $f(r/2)=0$，使次数大于 $1$ 的不可约多项式具有有理根，矛盾。
\end{solve}

\begin{example}
\label{ex:poly-integer-values-irreducible}
设 $a_1,\ldots,a_n$ 是互不相同的整数，证明下列两个多项式在 $\mathbb Q$ 上不可约：
\[
    P(x)-1,\qquad P(x)^2+1,\qquad P(x)=\prod_{i=1}^n(x-a_i).
\]
\end{example}
\begin{solve}
若 $P-1=uv$ 可约，由 Gauss 引理可令 $u,v$ 均为首一整系数多项式，且次数都小于 $n$。在 $a_i$ 处，$u(a_i)v(a_i)=-1$，所以 $u(a_i)+v(a_i)=0$。于是低于 $n$ 次的多项式 $u+v$ 有 $n$ 个不同根，只能恒为零，这与两个因子都首一矛盾。

若 $P^2+1=uv$ 可约，仍取首一整系数正次数因子。因为乘积在实数轴上恒正，两个因子均无实根；又各自首一，故在实数轴上都恒正。因此 $u(a_i)=v(a_i)=1$。不妨设 $\deg u\leq n$；若 $\deg u<n$，则 $u-1$ 有 $n$ 个不同根，迫使 $u=1$，矛盾。故两个因子的次数都为 $n$，而 $u-v$ 的次数小于 $n$，又在各 $a_i$ 处为零，故 $u=v$。由 $u^2=P^2+1$ 得 $(u-P)(u+P)=1$，两个因子都只能为常数，从而 $P$ 也是常数，与 $n\geq1$ 矛盾。
\end{solve}

\section{多元多项式与对称多项式}
\label{sec:poly-multivariable}

\begin{definition}
未定元 $x_1,\ldots,x_n$ 的有限个单项式 $x_1^{a_1}\cdots x_n^{a_n}$（各 $a_i$ 为非负整数）的 $\mathbb F$ 系数线性组合称为\emph{多元多项式}（multivariate polynomial），全体记为 $\mathbb F[x_1,\ldots,x_n]$。非零多项式的总次数是其非零单项式中 $a_1+\cdots+a_n$ 的最大值；若所有非零单项式的总次数相同，则称为齐次多项式（homogeneous polynomial）。
\end{definition}

\begin{proposition}
\label{prop:poly-multivariate-nonzero}
$\mathbb F[x_1,\ldots,x_n]$ 是整环。每个非零的多元多项式，都在 $\mathbb F^n$ 中至少有一个取值不为零的点。因此，多元多项式恒为零当且仅当它在全部 $\mathbb F^n$ 上取零值。
\end{proposition}
\begin{proof}
按字典序排列指数向量，两个非零多项式相乘时，其最高单项式的乘积就是乘积中的最高单项式，系数为两个非零最高系数的乘积，所以不会消失。非零取值的结论对元数归纳：把多项式视为关于 $x_n$ 的一元多项式，选取它的某个非零系数多项式，由归纳假设可给前 $n-1$ 个变量赋值，使该系数非零；于是得到非零的一元多项式。因为 $\mathbb F$ 无限，可以在其有限根集之外选取 $x_n$ 的值。
\end{proof}

多元多项式也满足唯一因式分解：每个非零非常数多项式可分解为有限个不可约因子的乘积，且分解在调换次序、改变非零常数倍的意义下唯一。此处采用这一结果而不证明。需要区分的是，一元情形的 Bézout 判据不能原样推广：例如 $x,y$ 在 $\mathbb F[x,y]$ 中没有非常数公因式，却不可能有 $xu+yv=1$，因为代入 $(x,y)=(0,0)$ 即产生矛盾。

\emph[caution]{多元情形不能把“在全部点上相等”弱化为“在无限多个点上相等”}。例如 $f(x,y)=x$ 与零多项式在无限点集 $\{(0,y):y\in\mathbb F\}$ 上相等，却不是同一多项式。与此相对，下面的结论允许在排除若干代数退化条件后先作计算，再恢复完整的多项式恒等式。

\begin{proposition}[非退化取值上的恒等式]
\label{prop:poly-generic-identity}
设非零多项式 $H_1,\ldots,H_s\in\mathbb F[x_1,\ldots,x_n]$。若 $P,Q$ 在所有满足 $H_1\cdots H_s\ne0$ 的点上取值相同，则 $P=Q$。
\end{proposition}
\begin{proof}
若 $P-Q\ne0$，整性保证 $(P-Q)H_1\cdots H_s\ne0$。由命题~\ref{prop:poly-multivariate-nonzero}，存在一个点使这个乘积不为零；在该点，全部非退化条件都成立，却有 $P\ne Q$，矛盾。
\end{proof}

例如，文字行列式的元素为参数的多项式时，行列式本身也是多项式。若在若干非零多项式条件下算出它等于某个多项式表达式，命题~\ref{prop:poly-generic-identity} 保证这一恒等式对特殊取值也成立。若中间答案是有理式，应先清除分母得到多项式恒等式；不能在分母为零处直接代入原有理式。

\begin{proposition}[以多项式为根的分解]
\label{prop:poly-polynomial-roots}
设 $R=\mathbb F[x_2,\ldots,x_n]$，$P\in R[x_1]$ 关于 $x_1$ 的次数为 $d$，首项系数为 $c\in R\setminus\{0\}$。若 $g_1,\ldots,g_d\in R$ 互不相同，且代入 $x_1=g_i$ 后 $P$ 恒为零，则
\[
    P=c\prod_{i=1}^d(x_1-g_i).
\]
\end{proposition}
\begin{proof}
由恒等式 $x_1^k-g_1^k=(x_1-g_1)\sum_{j=0}^{k-1}x_1^{k-1-j}g_1^j$，$P(g_1)=0$ 蕴含 $P=(x_1-g_1)P_1$。再代入 $x_1=g_2$，得到 $(g_2-g_1)P_1(g_2)=0$。因 $R$ 是整环且 $g_2-g_1\ne0$，有 $P_1(g_2)=0$，可继续约去下一因子。反复 $d$ 次，再比较首项系数，即得结论。这里不要求 $g_i-g_j$ 是可逆元，只需它是非零多项式。
\end{proof}

\begin{definition}
\label{def:poly-symmetric}
若交换任意两个变量后多项式不变，则称为\emph{对称多项式}（symmetric polynomial）。定义 $n$ 元\emph{初等对称多项式}（elementary symmetric polynomials）
\[
    e_k=\sum_{1\leq i_1<\cdots<i_k\leq n}x_{i_1}\cdots x_{i_k}
    \quad(1\leq k\leq n),\qquad e_0=1,
\]
并记 $s_k=x_1^k+\cdots+x_n^k$（$k\geq1$）为 $k$ 次幂和（power sum）。特别地，$e_1=\sum_i x_i$，$e_n=\prod_i x_i$。
\end{definition}

\begin{theorem}[对称多项式基本定理]
\label{thm:poly-symmetric}
每个 $\mathbb F$ 上的 $n$ 元对称多项式，都能唯一表示为 $e_1,\ldots,e_n$ 的多项式，即存在唯一的 $G\in\mathbb F[y_1,\ldots,y_n]$，使 $P(x_1,\ldots,x_n)=G(e_1,\ldots,e_n)$。
\end{theorem}
\begin{proof}
先对单项式按总次数、同次时按 $x_1>\cdots>x_n$ 的字典序排列。对称性保证 $P$ 的最高单项式 $ax_1^{b_1}\cdots x_n^{b_n}$ 满足 $b_1\geq\cdots\geq b_n$，否则交换某两个变量便产生更高的单项式。多项式
\[
    ae_1^{b_1-b_2}e_2^{b_2-b_3}\cdots e_{n-1}^{b_{n-1}-b_n}e_n^{b_n}
\]
恰以同一单项式为最高项，减去它以后最高项严格降低，而且总次数不增。在有限的次数范围内只有有限多个单项式，所以反复相减必终止，得到所需表示。

为证唯一性，注意 $e_1^{c_1}\cdots e_n^{c_n}$ 的最高单项式的指数依次为 $c_1+\cdots+c_n,c_2+\cdots+c_n,\ldots,c_n$，这个对应是一一的。故不同的 $e$ 单项式具有不同的最高单项式，任意非零线性组合中的最高项都不能抵消，不能表示零多项式。因此两种表示的差只能为零。
\end{proof}

对称有理函数也能表示为初等对称多项式的有理函数。事实上，若 $Q=f/g$ 对变量置换不变，令 $D=\prod_{\sigma\in S_n}\sigma(g)$，其中 $S_n$ 是全部变量置换组成的对称群，则 $D$ 是非零对称多项式，$QD=f\prod_{\sigma\ne\mathrm{id}}\sigma(g)$ 也是多项式且对称。分别应用定理~\ref{thm:poly-symmetric}，再相除即可。这里的相等是在有理函数域内成立，具体代入时仍须避开相应分母的零点。

\begin{theorem}[Newton 恒等式]
\label{thm:poly-newton}
对 $1\leq k\leq n$，有
\begin{equation}
\label{eq:poly-newton-low}
    s_k-e_1s_{k-1}+e_2s_{k-2}-\cdots+(-1)^{k-1}e_{k-1}s_1+(-1)^kke_k=0.
\end{equation}
对 $k>n$，有
\begin{equation}
\label{eq:poly-newton-high}
    s_k-e_1s_{k-1}+\cdots+(-1)^ne_ns_{k-n}=0.
\end{equation}
\end{theorem}
\begin{proof}
在形式幂级数中令 $E(t)=\prod_i(1-x_it)=\sum_{j=0}^n(-1)^je_jt^j$。因 $E(0)=1$，它有形式乘法逆元，乘积求导给出
\[
    -\frac{tE'(t)}{E(t)}=\sum_{i=1}^n\frac{x_it}{1-x_it}
    =\sum_{k\geq1}s_kt^k.
\]
两边乘以 $E(t)$，比较各次系数，即得两组公式。这里使用形式级数的系数运算，不涉及解析收敛或交换极限。
\end{proof}

这些公式可以从系数计算幂和，也可以由前 $n$ 个幂和反求全部初等对称多项式。例如
\[
    s_1=e_1,\qquad s_2=e_1^2-2e_2,\qquad
    s_3=e_1^3-3e_1e_2+3e_3.
\]
对 $1\leq k\leq n$，递推还给出行列式表达式
\begin{equation}
\label{eq:poly-newton-determinants}
    e_k=\frac1{k!}
    \begin{vmatrix}
    s_1&1&0&\cdots&0\\
    s_2&s_1&2&\cdots&0\\
    \vdots&\vdots&\vdots&\ddots&\vdots\\
    s_{k-1}&s_{k-2}&s_{k-3}&\cdots&k-1\\
    s_k&s_{k-1}&s_{k-2}&\cdots&s_1
    \end{vmatrix},\qquad
    s_k=
    \begin{vmatrix}
    e_1&1&0&\cdots&0\\
    2e_2&e_1&1&\cdots&0\\
    \vdots&\vdots&\vdots&\ddots&\vdots\\
    (k-1)e_{k-1}&e_{k-2}&e_{k-3}&\cdots&1\\
    ke_k&e_{k-1}&e_{k-2}&\cdots&e_1
    \end{vmatrix}.
\end{equation}
在 $k=1$ 时，两个行列式都按一阶理解。第一式可由式~\eqref{eq:poly-newton-low} 中关于 $e_1,\ldots,e_k$ 的下三角线性方程组使用 Cramer 法则得到；第二式按第一列展开，递推式恰为 $s_k=e_1s_{k-1}-\cdots+(-1)^{k-1}ke_k$，再对 $k$ 归纳即可。

\begin{example}
\label{ex:poly-symmetric-root-condition}
设 $\alpha,\beta,\gamma$ 为 $x^3+px^2+qx+r$ 的根，证明其中某一根的平方等于另外两根的平方和，当且仅当
\[
    p^4(p^2-2q)=2(p^3-2pq+2r)^2.
\]
\end{example}
\begin{solve}
所述条件等价于 $T=(\alpha^2+\beta^2-\gamma^2)(\alpha^2+\gamma^2-\beta^2)(\beta^2+\gamma^2-\alpha^2)=0$。令 $S=\alpha^2+\beta^2+\gamma^2=p^2-2q$，则由 Vieta 定理，
\[
    \alpha^2\beta^2+\alpha^2\gamma^2+\beta^2\gamma^2=q^2-2pr,
    \qquad \alpha^2\beta^2\gamma^2=r^2.
\]
因此
\[
    T=-S^3+4(q^2-2pr)S-8r^2
    =p^4(p^2-2q)-2(p^3-2pq+2r)^2,
\]
即得结论。先把“某一根满足”写成对称乘积为零，可以避免事先指定哪一个根。
\end{solve}

\begin{example}
\label{ex:poly-newton-system}
在复数范围内求解方程组 $x_1^k+\cdots+x_n^k=n$（$k=1,\ldots,n$）。
\end{example}
\begin{solve}
这组方程给出前 $n$ 个幂和 $s_1=\cdots=s_n=n$，与 $n$ 元组 $(1,\ldots,1)$ 的幂和相同。由 Newton 恒等式递次唯一确定 $e_1,\ldots,e_n$，故两组数的初等对称多项式相同，即 $e_k=\binom nk$。于是以 $x_1,\ldots,x_n$ 为根的首一多项式为 $(t-1)^n$，只能有 $x_1=\cdots=x_n=1$。这一方法在复数范围内同样成立，不依赖平方和的非负性。
\end{solve}

\section{结式与判别式}
\label{sec:poly-resultant}

\begin{definition}
\label{def:poly-resultant}
设 $f(x)=a_0x^n+a_1x^{n-1}+\cdots+a_n$、$g(x)=b_0x^m+b_1x^{m-1}+\cdots+b_m$，其中 $a_0b_0\ne0$，$n,m\geq1$。将
\[
    x^{m-1}f,\ldots,f,\quad x^{n-1}g,\ldots,g
\]
依次按 $x^{n+m-1},x^{n+m-2},\ldots,1$ 展开，把它们的系数向量作为行组成 $(n+m)$ 阶矩阵，称为 Sylvester 矩阵（Sylvester matrix）。其行列式称为 $f,g$ 的\emph{结式}（resultant），记为 $R(f,g)$。例如 $n=m=2$ 时，
\[
    R(f,g)=\begin{vmatrix}
    a_0&a_1&a_2&0\\0&a_0&a_1&a_2\\
    b_0&b_1&b_2&0\\0&b_0&b_1&b_2
    \end{vmatrix}.
\]
\end{definition}

本节采用降幂系数 $a_0,\ldots,a_n$，特别地 $a_0$ 是首项系数，这与前面某些按升幂索引的式子不同。固定上述行次序后，结式的符号也固定；使用其他矩阵排列时，应先核对符号约定。

\begin{theorem}
\label{thm:poly-resultant-criterion}
对定义~\ref{def:poly-resultant} 中的两个多项式，$R(f,g)=0$ 当且仅当它们不互素，也当且仅当它们在复数域内有公共根。
\end{theorem}
\begin{proof}
Sylvester 矩阵奇异，当且仅当存在不全为零的 $u,v$，满足 $\deg u<m,\deg v<n$ 且 $uf+vg=0$。若 $f,g$ 互素，则 $f\mid v$，次数条件迫使 $v=0$，再得 $u=0$，矛盾。若它们有正次数公因式 $d$，取 $u=g/d,v=-f/d$ 就得到非零的行间线性关系。最后，正次数公因式在复数域有根；反之，若 $f,g$ 有公根，却在原数域互素，将 Bézout 恒等式代入公根会得到矛盾。
\end{proof}

\begin{theorem}[结式的根表达式]
\label{thm:poly-resultant-product}
若 $f$ 的复根为 $\alpha_1,\ldots,\alpha_n$，$g$ 的复根为 $\beta_1,\ldots,\beta_m$，均按重数排列，则
\begin{equation}
\label{eq:poly-resultant-product}
    R(f,g)=a_0^m\prod_{i=1}^n g(\alpha_i)
    =a_0^mb_0^n\prod_{i=1}^n\prod_{j=1}^m(\alpha_i-\beta_j).
\end{equation}
\end{theorem}
\begin{proof}
先设 $f$ 的根互异。把结式看成 $g$ 的系数 $b_0,\ldots,b_m$ 的多项式，由行列式定义它是 $n$ 次齐次多项式。每当线性式 $g(\alpha_i)$ 为零，结式就为零，故该线性式整除结式；这一点可用线性代换把它作为一个新变量，再令该变量为零来验证。不同根对应的线性式互不相伴，所以结式等于它们乘积的某个常数倍。将 $g$ 的系数专门取为常数多项式 $1$ 的系数（仍保留定义矩阵的大小），矩阵行列式为 $a_0^m$，而各 $g(\alpha_i)=1$，因而常数倍为 $a_0^m$。

再把 $f=a_0\prod_i(x-\alpha_i)$ 代入两边。两边都是 $a_0,\alpha_i,b_j$ 的多项式，已知它们在 $a_0\prod_{i<j}(\alpha_i-\alpha_j)\ne0$ 时相等，由命题~\ref{prop:poly-generic-identity}，在根有重复时也相等。最后分解每个 $g(\alpha_i)$ 即得第二式。
\end{proof}

\begin{corollary}
\label{cor:poly-resultant-properties}
对正次数的 $f,g,g_1,g_2$ 及非零常数 $c,d$，有
\[
    R(f,g)=(-1)^{nm}R(g,f),\qquad
    R(cf,dg)=c^md^nR(f,g),\qquad
    R(f,g_1g_2)=R(f,g_1)R(f,g_2).
\]
若将非零常数 $c$ 也纳入讨论，约定 $R(f,c)=c^{\deg f}$、$R(c,g)=c^{\deg g}$，则这些乘积公式仍然成立。
\end{corollary}
\begin{proof}
在式~\eqref{eq:poly-resultant-product} 中交换两组根，每个差改变一次符号，总共 $nm$ 次；数乘只改变首项系数；乘积 $g_1g_2$ 则使每个代入值分成两个因子。常数情形由约定直接验证。
\end{proof}

\begin{definition}
\label{def:poly-discriminant}
设 $f(x)=a_0x^n+\cdots+a_n$，$a_0\ne0$，$n\geq2$，复根为 $\alpha_1,\ldots,\alpha_n$。定义其\emph{判别式}（discriminant）为
\begin{equation}
\label{eq:poly-discriminant}
    \Delta(f)=a_0^{2n-2}\prod_{1\leq i<j\leq n}(\alpha_i-\alpha_j)^2
    =\frac{(-1)^{n(n-1)/2}}{a_0}R(f,f').
\end{equation}
为统一乘积公式，约定一次多项式的判别式为 $1$。
\end{definition}

第二个表达式由 $f'(\alpha_i)=a_0\prod_{j\ne i}(\alpha_i-\alpha_j)$ 与结式的根表达式得到；若有重根，两边同时为零。根表达式直接说明：$\Delta(f)=0$ 当且仅当 $f$ 有重根。判别式虽然按根写出，却是系数的多项式：$R(f,f')$ 本来就是系数的多项式，而其 Sylvester 矩阵第一列的非零候选项都含有因子 $a_0$，故该行列式可被 $a_0$ 整除，式~\eqref{eq:poly-discriminant} 中的除法仍得到多项式。

\begin{proposition}
\label{prop:poly-discriminant-products}
若 $f,g$ 均为正次数多项式，则
\begin{equation}
\label{eq:poly-discriminant-product}
    \Delta(fg)=\Delta(f)\Delta(g)R(f,g)^2.
\end{equation}
特别地，$\Delta((x-a)g)=g(a)^2\Delta(g)$。若 $g,h$ 均首一，$\deg g=n,\deg h=m$，且 $g$ 的根为 $\alpha_1,\ldots,\alpha_n$，则
\begin{equation}
\label{eq:poly-discriminant-composite}
    \Delta(g\circ h)=\Delta(g)^m\prod_{i=1}^n\Delta(h-\alpha_i).
\end{equation}
\end{proposition}
\begin{proof}
在 $fg$ 的根表达式中，根差分成 $f$ 内部、$g$ 内部以及两者之间的三类。前两类分别产生 $\Delta(f),\Delta(g)$，第三类连同所需首项系数产生 $R(f,g)^2$，得到式~\eqref{eq:poly-discriminant-product}，一次因子情形也包含在内。

再写 $g(h(x))=\prod_i(h(x)-\alpha_i)$。由结式的根表达式，$R(h-\alpha_i,h-\alpha_j)=(\alpha_i-\alpha_j)^m$。反复应用乘积公式，所有不同因子之间的结式平方的乘积为 $\prod_{i<j}(\alpha_i-\alpha_j)^{2m}=\Delta(g)^m$，即得式~\eqref{eq:poly-discriminant-composite}。
\end{proof}

\begin{example}
\label{ex:poly-discriminant-substitution}
设 $f(x)=a_0x^n+\cdots+a_n$，$a_0\ne0$，$n\geq1$，求 $\Delta(f(x^2))$。
\end{example}
\begin{solve}
设 $f$ 的根为 $\alpha_i$，则 $f(x^2)=a_0\prod_i(x^2-\alpha_i)$。每个二次因子的判别式为 $4\alpha_i$，不同因子之间的结式为 $(\alpha_i-\alpha_j)^2$。此外，次数为 $N$ 的多项式乘以常数 $c$ 后，判别式乘以 $c^{2N-2}$。因此
\[
    \Delta(f(x^2))=a_0^{4n-2}4^n\prod_i\alpha_i
    \prod_{i<j}(\alpha_i-\alpha_j)^4
    =(-4)^na_0a_n\Delta(f)^2.
\]
若 $a_n=0$，则 $f(x^2)$ 含因子 $x^2$，两边均为零；有重根时同样由乘积公式成立。
\end{solve}

\begin{example}
\label{ex:poly-trinomial-discriminant}
求 $F(x)=x^n+px+q$（$n\geq2$）的判别式，并给出四次情形的重根条件。
\end{example}
\begin{solve}
导数为 $F'=nx^{n-1}+p$。设其根为 $\beta_1,\ldots,\beta_{n-1}$，则 $\beta_j^{n-1}=-p/n$，且 $F(\beta_j)=q+(n-1)p\beta_j/n$。因为 $n(n-1)$ 为偶数，
\[
    R(F,F')=R(F',F)
    =n^n\prod_{j=1}^{n-1}\left(q+\frac{n-1}{n}p\beta_j\right)
    =n^nq^{n-1}+(-1)^{n-1}(n-1)^{n-1}p^n.
\]
最后一步由 $\prod_j(t-\beta_j)=t^{n-1}+p/n$ 代入并整理得到；当 $p=0$ 时也直接成立。故
\begin{equation}
\label{eq:poly-trinomial-discriminant}
    \Delta(x^n+px+q)=(-1)^{n(n-1)/2}
    \left(n^nq^{n-1}+(-1)^{n-1}(n-1)^{n-1}p^n\right).
\end{equation}
当 $n=4$ 时，$\Delta(x^4+px+q)=256q^3-27p^4$，因此有重根当且仅当 $27p^4=256q^3$。作为另一应用，令 $p=0,q=-1$，再利用 $x^n-1=(x-1)(1+x+\cdots+x^{n-1})$，可得
\[
    \Delta(1+x+\cdots+x^{n-1})
    =(-1)^{(n-1)(n-2)/2}n^{n-2}\qquad(n\geq2).
\]
\end{solve}

\begin{proposition}[实多项式判别式的符号]
\label{prop:poly-real-discriminant-sign}
若实系数多项式 $f$ 没有重根，且有 $s$ 对非实共轭根，则 $\Delta(f)$ 的符号为 $(-1)^s$。
\end{proposition}
\begin{proof}
首项系数的因子 $a_0^{2n-2}$ 为正。每对非实共轭根 $z,\overline z$ 之间的因子 $(z-\overline z)^2$ 为负；实根之间的根差平方为正；一个实根与一对共轭根产生 $|r-z|^4>0$；两对不同共轭根之间的全部根差平方的乘积为 $|z-w|^4|z-\overline w|^4>0$。故唯一的负号来自每对共轭根内部，共有 $s$ 个。
\end{proof}

\begin{example}
\label{ex:poly-real-cubic-discriminant}
对实系数方程 $x^3+px+q=0$，令 $D=q^2/4+p^3/27$，用 $D$ 判别其实根情况。
\end{example}
\begin{solve}
式~\eqref{eq:poly-trinomial-discriminant} 给出 $\Delta=-4p^3-27q^2=-108D$。若 $D>0$，判别式负，三次方程恰有一对非实共轭根和一个实根；若 $D<0$，判别式正，非实根对数为偶数，而三次方程最多有一对，故三个根均为互不相同的实根。若 $D=0$，有重根，但非实重根会连同它的共轭根共占至少四个重数，所以所有根都为实数；其中 $p=q=0$ 时零为三重根，其余情形是一个二重实根和一个不同的单实根。这里 $D$ 与本节定义的判别式相差负常数，不能混用两者的符号结论。
\end{solve}

\begin{example}
\label{ex:poly-resultant-elimination}
对实参数曲线
\[
    x=\frac{2(t+1)}{t^2+1},\qquad y=\frac{t^2}{2t-1},\qquad t\in\mathbb R\setminus\{1/2\},
\]
用结式消去参数，并说明所得方程的适用范围。
\end{example}
\begin{solve}
清除分母得到 $f(t)=xt^2-2t+x-2=0$、$g(t)=t^2-2yt+y=0$。计算对应的四阶系数行列式得
\[
    R(f,g)=5x^2y^2-2x^2y-12xy^2+x^2-4x+12y+4,
\]
所以参数曲线上的点都满足右边等于零。当 $x=0$ 时，$f$ 的实际次数降低，但上述固定大小行列式仍是各参数的多项式，必要条件依旧成立；事实上此时原参数式给出点 $(0,-1/3)$，也满足该式。

消元方程不自动保证存在实参数。例如 $(x,y)=(3,1/3)$ 满足该方程，但此时 $f=3g=3t^2-2t+1$，没有实根，不能来自原实参数曲线。因此把结式为零解释为原实曲线的充要条件时，还必须核查公共根是否为实数、分母是否非零以及多项式是否降次。
\end{solve}

\section{互素多项式在线性代数与数域中的应用}
\label{sec:poly-applications}

多项式恒等式可以代入同一个方阵或同一个线性变换。若 $T$ 是 $\mathbb F$ 上向量空间 $V$ 的线性变换，定义 $f(T)=\sum_i a_iT^i$，其中 $T^0=I_V$。由于 $T$ 的各次幂相互交换，有 $(fg)(T)=f(T)g(T)=g(T)f(T)$。\emph[caution]{这里的乘法相容性依赖于都代入同一个变换}；不能任意把涉及多个交换变量的恒等式用于不交换的矩阵。

\begin{proposition}
\label{prop:poly-matrix-inverse}
若 $f,g\in\mathbb F[x]$ 互素，且 $f(T)=0$，则 $g(T)$ 可逆。若 $fu+gv=1$，其逆就是 $v(T)$。
\end{proposition}
\begin{proof}
代入 $T$ 后有 $f(T)u(T)+g(T)v(T)=I_V$，故 $g(T)v(T)=I_V$。因为两者都是 $T$ 的多项式，也有 $v(T)g(T)=I_V$。
\end{proof}

\begin{theorem}[互素因子诱导的核分解]
\label{thm:poly-kernel-decomposition}
设 $V$ 为有限维向量空间，$T$ 是其线性变换，$f,g\in\mathbb F[x]$ 互素，则
\begin{equation}
\label{eq:poly-kernel-decomposition}
    \ker\bigl(f(T)g(T)\bigr)=\ker f(T)\oplus\ker g(T).
\end{equation}
各核空间均为 $T$ 的不变子空间。若 $\dim V=n$，还有
\[
    \operatorname{rank}f(T)+\operatorname{rank}g(T)
    =n+\operatorname{rank}\bigl(f(T)g(T)\bigr).
\]
因此 $f(T)g(T)=0$ 等价于两者秩的和为 $n$，此时 $V=\ker f(T)\oplus\ker g(T)$。
\end{theorem}
\begin{proof}
取 $fu+gv=1$。若 $w\in\ker(f(T)g(T))$，则
\[
    w=u(T)f(T)w+v(T)g(T)w.
\]
第一项在 $\ker g(T)$ 中，第二项在 $\ker f(T)$ 中，因为各多项式变换相互交换。反过来，这两个核中的向量之和均被 $f(T)g(T)$ 消去。若 $w$ 同时属于两个核，仍由上式得 $w=0$，所以和是直和。对任一多项式 $h$，$h(T)Tw=Th(T)w$ 表明 $\ker h(T)$ 在 $T$ 下不变。最后对式~\eqref{eq:poly-kernel-decomposition} 取维数，使用秩与零度之和为 $n$，即得秩恒等式及其推论。
\end{proof}

当 $fg$ 消去 $T$ 时，$v(T)g(T)$ 是到 $\ker f(T)$ 的投影，$u(T)f(T)$ 是到 $\ker g(T)$ 的投影：它们的和是恒等变换，在各自对应的子空间上为恒等、在另一个子空间上为零。因此，若两两互素的 $f_1,\ldots,f_s$ 满足 $(\prod_i f_i)(T)=0$，反复应用上述定理得到 $V=\bigoplus_i\ker f_i(T)$。这与中国剩余定理中构造 $E_i$ 的方法相通，是以后研究相似标准型时分解空间的基础。

\begin{theorem}[代数元生成的数域]
\label{thm:poly-generated-field}
设 $\alpha\in\mathbb C$ 在数域 $\mathbb F$ 上代数，其首一极小多项式的次数为 $d$。则
\[
    \mathbb F[\alpha]=\{a_0+a_1\alpha+\cdots+a_{d-1}\alpha^{d-1}:a_i\in\mathbb F\}
\]
是数域，等于包含 $\mathbb F$ 和 $\alpha$ 的最小数域 $\mathbb F(\alpha)$，且 $1,\alpha,\ldots,\alpha^{d-1}$ 是它在 $\mathbb F$ 上的一组基。因此其扩张次数 $[\mathbb F(\alpha):\mathbb F]$，即作为 $\mathbb F$ 向量空间的维数，等于 $d$。
\end{theorem}
\begin{proof}
对极小多项式 $m$ 作带余除法，任意多项式在 $\alpha$ 处的值都可化为上述低次表达式，因此该集合对加、减、乘封闭。若非零元素为 $g(\alpha)$，$\deg g<d$，则 $m$ 不可约保证 $\gcd(m,g)=1$。取 $mu+gv=1$ 并把 $v$ 模 $m$ 化简，代入 $\alpha$ 得 $g(\alpha)v(\alpha)=1$，且 $\deg v<d$，故非零元素的倒数仍在集合中。这便是一个包含 $\mathbb F,\alpha$ 的数域，而任何这样的数域必包含这些多项式值，故它最小。若 $\sum_{j=0}^{d-1}a_j\alpha^j=0$，极小性迫使所有 $a_j=0$，从而所列向量线性无关，同时按定义生成整个空间，故构成一组基。
\end{proof}

\begin{example}
\label{ex:poly-radical-field}
设 $\alpha=\sqrt[n]{2}$，求 $\mathbb Q(\alpha)$ 的一组有理数域上的基和维数。
\end{example}
\begin{solve}
由 Eisenstein 判别法，$x^n-2$ 是 $\alpha$ 的极小多项式，所以 $1,\alpha,\ldots,\alpha^{n-1}$ 是一组基，$[\mathbb Q(\alpha):\mathbb Q]=n$。例如 $\alpha^n=2$ 用于降低乘积的幂次，非零元素的倒数则由它所对应多项式与 $x^n-2$ 的 Bézout 恒等式计算。二者分别解决乘法封闭和除法封闭，不能仅凭可降低幂次就断言所得集合是域。
\end{solve}

\begin{proposition}[扩张次数的乘法公式]
\label{prop:poly-tower-degree}
若 $\mathbb F\subseteq\mathbb K\subseteq\mathbb L$，且两个扩张次数均有限，则
\[
    [\mathbb L:\mathbb F]=[\mathbb L:\mathbb K][\mathbb K:\mathbb F].
\]
\end{proposition}
\begin{proof}
取 $u_1,\ldots,u_r$ 为 $\mathbb K$ 在 $\mathbb F$ 上的基，$v_1,\ldots,v_s$ 为 $\mathbb L$ 在 $\mathbb K$ 上的基。先按 $v_j$、再把相应系数按 $u_i$ 展开，说明 $u_iv_j$ 生成 $\mathbb L$。若这些向量的 $\mathbb F$ 线性组合为零，按 $v_j$ 分组，由 $v_j$ 的 $\mathbb K$ 线性无关性，每组系数都为零；再由 $u_i$ 的线性无关性，全部系数都为零。因此 $u_iv_j$ 构成一组含 $rs$ 个元素的基。
\end{proof}

\begin{example}
\label{ex:poly-mixed-radical-minimal}
求 $\alpha=\sqrt2+\sqrt[3]{3}$ 的极小多项式。
\end{example}
\begin{solve}
从 $(\alpha-\sqrt2)^3=3$ 得
\[
    \alpha^3+6\alpha-3=(3\alpha^2+2)\sqrt2.
\]
平方并移项，得到 $F(\alpha)=0$，其中
\[
    F(x)=x^6-6x^4-6x^3+12x^2-36x+1.
\]
得到一个消去多项式后，还必须证明它的次数最小。令 $\beta=\sqrt[3]{3}$，由 Eisenstein 判别法，$[\mathbb Q(\beta):\mathbb Q]=3$。若 $\sqrt2\in\mathbb Q(\beta)$，扩张次数的乘法公式会使 $2$ 整除 $3$，不可能；故 $x^2-2$ 在 $\mathbb Q(\beta)$ 上仍不可约，$[\mathbb Q(\beta,\sqrt2):\mathbb Q]=6$。

另一方面，$3\alpha^2+2>0$，上面的恒等式给出 $\sqrt2\in\mathbb Q(\alpha)$，继而 $\beta=\alpha-\sqrt2\in\mathbb Q(\alpha)$。所以 $\mathbb Q(\alpha)=\mathbb Q(\beta,\sqrt2)$，其扩张次数为 $6$，极小多项式也为六次。它整除首一六次多项式 $F$，故只能等于 $F$。
\end{solve}

\begin{example}
\label{ex:poly-cyclic-basis}
设 $V$ 是 $n$ 维 $\mathbb F$ 向量空间，$T$ 是其线性变换，$f\in\mathbb F[x]$ 是 $n$ 次不可约多项式。若 $v\ne0$ 且 $f(T)v=0$，证明 $v,Tv,\ldots,T^{n-1}v$ 是 $V$ 的一组基。
\end{example}
\begin{solve}
若有非平凡线性关系，就存在非零 $g\in\mathbb F[x]$，使 $\deg g<n$ 且 $g(T)v=0$。因为 $f$ 不可约，$f,g$ 互素，故有 $fu+gv_1=1$。代入 $T$ 再作用于向量 $v$，得到 $v=u(T)f(T)v+v_1(T)g(T)v=0$，矛盾。因此这些 $n$ 个向量线性无关，构成一组基。

例如，若 $f(x)=x^n+c_1x^{n-1}+\cdots+c_n$，向量 $v_1\ne0,v_2,\ldots,v_n$ 满足 $Tv_i=v_{i+1}$（$i<n$）及 $Tv_n=-c_nv_1-\cdots-c_1v_n$，则这些关系正好给出 $f(T)v_1=0$，可直接应用上述结论。
\end{solve}

\section{基础训练中的典型方法}
\label{sec:poly-training}

本节将原书基础训练中可反复使用的方法整理为例题。判断题或选择题中的断言，应先检查系数域、次数与非零条件，再决定使用定理还是构造反例。例如，三个多项式共同互素不保证两两互素；导数在某点为零不保证原多项式也在该点为零；可约的实多项式可以没有实根，如 $(x^2+1)(x^2+2)$。这些区别分别涉及公因式、重根和系数域，不能相互替代。

\begin{example}
\label{ex:poly-training-remainders}
设 $m,n$ 是正整数，求 $x^{3m}+x^{3n}$ 除以 $x^2+x+1$ 的余式；再求首一三次多项式 $f$，使 $f(1)=1,f(2)=2,f(3)=3$。
\end{example}
\begin{solve}
因为模 $x^2+x+1$ 有 $x^3\equiv1$，第一问的余式为 $2$。第二问中，$f-x$ 是首一三次多项式，以 $1,2,3$ 为根，故
\[
    f(x)=(x-1)(x-2)(x-3)+x=x^3-6x^2+12x-6.
\]
已知多个取值时，先减去一个具有相同取值的简单多项式，通常比逐一解系数方程更直接。
\end{solve}

\begin{example}
\label{ex:poly-training-triple-root}
设 $n\geq1$，证明 $1$ 是 $P_n(x)=x^{2n+1}-(2n+1)x^{n+1}+(2n+1)x^n-1$ 的三重根。
\end{example}
\begin{solve}
逐项求导并代入 $x=1$，得到 $P_n(1)=P_n'(1)=P_n''(1)=0$，而
\[
    P_n'''(1)=n(n+1)(2n+1)\ne0.
\]
由定理~\ref{thm:poly-multiple-root}，重数恰为 $3$。若只验证前三个等式而不检查三阶导数，最多只能得到重数至少为 $3$。
\end{solve}

\begin{example}
\label{ex:poly-training-trinomial-multiplicity}
设 $n>m\geq1$，$a,b\in\mathbb F$ 且 $b\ne0$。证明 $P(x)=x^n+ax^{n-m}+b$ 不可能有重数大于 $2$ 的复根。
\end{example}
\begin{solve}
因常数项非零，任何根 $\alpha$ 均非零。导数为
\[
    P'(x)=x^{n-m-1}\bigl(nx^m+(n-m)a\bigr).
\]
若 $a=0$，任何非零点都不是导数的根，故 $P$ 无重根。若 $a\ne0$，括号内多项式在非零根 $\alpha$ 处的导数为 $nm\alpha^{m-1}\ne0$，故 $P'$ 在 $\alpha$ 处至多是单根。原多项式若在该点有 $k$ 重根，导数恰有 $k-1$ 重根，因此 $k\leq2$。
\end{solve}

\begin{example}
\label{ex:poly-training-euler-operator}
设 $f(x)=\sum_{j=0}^n a_jx^j$，$k\geq0$。证明 $(x-1)^{k+1}\mid f$ 当且仅当
\[
    \sum_{j=0}^n j^\ell a_j=0\qquad(\ell=0,1,\ldots,k),
\]
其中 $\ell=0$ 时每个 $j^0$ 都按 $1$ 理解。
\end{example}
\begin{solve}
令 $\Theta f=xf'$，则 $\Theta^\ell f=\sum_j j^\ell a_jx^j$。另一方面，对 $\ell$ 归纳并使用乘积求导法则，$\Theta^\ell f$ 可写为 $x^\ell f^{(\ell)}$ 加上 $x^rf^{(r)}$（$r<\ell$）的常数线性组合。在 $x=1$ 处，这给出两组值 $\{\Theta^\ell f(1)\}_{\ell=0}^k$ 与 $\{f^{(\ell)}(1)\}_{\ell=0}^k$ 之间的下三角变换，且对角元均为 $1$，因而它们同时全为零。再由重根判据得到结论；$f=0$ 的情形也直接成立。
\end{solve}

\begin{example}
\label{ex:poly-training-additive}
设 $f\in\mathbb F[x]$，对任意 $a,b\in\mathbb F$ 都有 $f(a+b)=f(a)+f(b)$，证明 $f(x)=cx$，其中 $c\in\mathbb F$。
\end{example}
\begin{solve}
令 $a=b=0$ 得 $f(0)=0$。取 $c=f(1)$，反复相加得到 $f(N)=Nc$ 对所有正整数 $N$ 成立。于是 $f(x)-cx$ 有无限多个根，只能为零多项式。这里利用了多项式这一限制；不能把同样的论证直接用于任意函数。
\end{solve}

\begin{example}
\label{ex:poly-training-derivation}
设 $D$ 是 $\mathbb F[x]$ 上的 $\mathbb F$ 线性变换，满足 $D(x)=1$ 及 $D(fg)=gD(f)+fD(g)$。证明 $D$ 就是形式求导变换。
\end{example}
\begin{solve}
由 $D(1)=D(1\cdot1)=2D(1)$，得 $D(1)=0$。若 $D(x^j)=jx^{j-1}$，则 $D(x^{j+1})=xD(x^j)+x^jD(x)=(j+1)x^j$，因此对次数归纳得到全部单项式上的结论。再用 $\mathbb F$ 线性性，$D(\sum_j a_jx^j)=\sum_j ja_jx^{j-1}$。
\end{solve}

\begin{example}
\label{ex:poly-training-minimal-common-roots}
求同时以 $1+\sqrt3$ 和 $3+i\sqrt2$ 为根的次数最小的首一有理系数多项式。
\end{example}
\begin{solve}
两数的极小多项式分别为 $p=x^2-2x-2$ 和 $q=x^2-6x+11$，因为它们分别以相应数为根，且二次判别式 $12,-8$ 都不是有理数平方。两个不同的首一不可约多项式互素，所以任何符合要求的多项式都可被 $pq$ 整除；$pq$ 本身又符合条件。因此所求为
\[
    pq=x^4-8x^3+21x^2-10x-22.
\]
\end{solve}

\begin{example}
\label{ex:poly-training-odd-irreducible}
设奇数 $n\geq1$，$a_1,\ldots,a_n$ 为互不相同的整数，证明 $F(x)=\prod_{i=1}^n(x-a_i)+1$ 在 $\mathbb Q$ 上不可约。
\end{example}
\begin{solve}
若 $F=uv$ 是两个首一正次数整系数因子的乘积，则 $u(a_i)v(a_i)=1$，整数性给出 $u(a_i)=v(a_i)=1$ 或 $u(a_i)=v(a_i)=-1$。故 $u-v$ 在 $n$ 个不同点处为零，且次数小于 $n$，只能恒为零。于是 $F=u^2$，次数必须为偶数，与 $n$ 为奇数矛盾。这个证明与例~\ref{ex:poly-integer-values-irreducible} 中处理常数项修正为 $-1$ 的方法相似，但两处的符号与次数条件不可互换。
\end{solve}

\begin{example}
\label{ex:poly-training-integer-root}
设 $f\in\mathbb Z[x]$，存在四个不同整数 $a_1,\ldots,a_4$，使 $f(a_i)=5$。证明 $f$ 没有整数根。
\end{example}
\begin{solve}
因 $f-5$ 在这四点为零，逐次用首一因式 $x-a_i$ 相除，得到 $f(x)-5=\prod_{i=1}^4(x-a_i)h(x)$，其中 $h\in\mathbb Z[x]$。若 $c$ 是整数根，则
\[
    -5=\prod_{i=1}^4(c-a_i)h(c).
\]
四个互不相同的非零整数 $c-a_i$ 都必须整除 $5$，因而只能恰为 $1,-1,5,-5$；它们的乘积为 $25$，不能整除 $-5$，矛盾。
\end{solve}

\begin{example}
\label{ex:poly-training-parameter-common-root}
求使 $x^3+cx+1$ 与 $x^2+cx+1$ 有公共复根的全部 $c\in\mathbb C$。
\end{example}
\begin{solve}
若公共根为 $\alpha$，两式相减得 $\alpha^2(\alpha-1)=0$。因两式常数项均为 $1$，$\alpha\ne0$，只能 $\alpha=1$，代入得 $c=-2$。反之，$c=-2$ 时两式确实以 $1$ 为根。这说明结式虽然给出统一判据，但具体题目也可先用带余除法或直接相减降低次数。
\end{solve}

\begin{example}
\label{ex:poly-training-power-sums}
设 $a,b\in\mathbb C$，$h_j=\sum_{r=0}^j a^rb^{j-r}$，$h_0=1$。若 $\alpha_1,\ldots,\alpha_n$ 是 $x^n+h_1x^{n-1}+\cdots+h_n$ 的根，求 $s_k=\sum_i\alpha_i^k$（$1\leq k\leq n$）。
\end{example}
\begin{solve}
由 Vieta 定理，$e_j=(-1)^jh_j$。在形式级数中，$\sum_{j\geq0}h_jt^j=((1-at)(1-bt))^{-1}$。因此本题的 $E(t)=\sum_{j=0}^nh_jt^j$ 与该逆元的系数直到 $t^n$ 都相同。使用 Newton 恒等式证明中的 $\sum_{k\geq1}s_kt^k=-tE'(t)/E(t)$，比较前 $n$ 次系数，得到
\[
    s_k=-(a^k+b^k)\qquad(1\leq k\leq n).
\]
只比较有限阶系数即可，不需要知道所有根的具体形式。
\end{solve}
% END textbook-xie-exercises-ch5

% 教材来源：高代习题-6.pdf，第6章，印刷页287--353；按主题整合，不逐题转录。
\chapter{特征值}
\label{chap:xie-ex-eig-eigenvalues}

特征值研究线性变换在某些方向上只产生数乘的现象。特征向量给出这些方向，特征子空间把同一数乘规律的向量组织起来，特征多项式与极小多项式则分别记录特征值的重数信息和矩阵所满足的多项式关系。本章先讨论特征值的计算与估计，再研究乘法交换性、相似与可对角化问题，最后用多项式方法和 Kronecker 积统一处理矩阵方程及矩阵空间上的线性变换。

全章以数域 $\mathbb F\subseteq\mathbb C$ 上的有限维线性空间和矩阵为对象，矩阵阶数均为正整数。全部不同的复特征值组成的集合称为矩阵的\emph{谱（spectrum）}，记为 $\sigma(A)$。除非明确讨论在 $\mathbb F$ 上的性质，计算全部特征值时允许扩充到复数域，并将特征值按代数重数计数。记 $I_n$ 为 $n$ 阶单位矩阵，$\operatorname{tr}A$、$\operatorname{rank}A$、$\operatorname{adj}A$ 分别表示迹、秩和伴随矩阵；$A^{\mathsf T}$ 表示转置，复矩阵的共轭转置写作 $A^{\mathsf H}$，以免将转置与伴随矩阵混淆。

\section{基本概念}
\label{sec:xie-ex-eig-eigen-basic}

\subsection{特征值、特征子空间与两种重数}

\begin{definition}
设 $T\colon V\to V$ 为 $\mathbb F$ 上的线性变换。若存在 $\lambda\in\mathbb F$ 和 $0\ne v\in V$，使 $Tv=\lambda v$，则称 $\lambda$ 为 $T$ 的\emph{特征值（eigenvalue）}，$v$ 为属于 $\lambda$ 的\emph{特征向量（eigenvector）}。矩阵 $A\in\mathbb F^{n\times n}$ 的定义相应为 $Av=\lambda v$，其中 $0\ne v\in\mathbb F^n$。属于 $\lambda$ 的\emph{特征子空间（eigenspace）}为
\[
E_\lambda(T)=\ker(T-\lambda I_V),\qquad
E_\lambda(A)=\ker(A-\lambda I_n).
\]
特征子空间包含零向量，而特征向量本身必须非零。
\end{definition}

\begin{definition}
矩阵 $A$ 的\emph{特征多项式（characteristic polynomial）}定义为
\begin{equation}
\chi_A(t)=\det(tI_n-A)=t^n+c_1t^{n-1}+\cdots+c_n.
\label{eq:xie-ex-eig-characteristic-polynomial}
\end{equation}
若 $\lambda$ 是其根，则根的重数 $a_A(\lambda)$ 称为 $\lambda$ 的\emph{代数重数（algebraic multiplicity）}，而
\[
g_A(\lambda)=\dim E_\lambda(A)=n-\operatorname{rank}(\lambda I_n-A)
\]
称为\emph{几何重数（geometric multiplicity）}。若两者相等，教材称该特征值有完全的特征向量系；若每个特征值都如此，则称矩阵有完全的特征向量系。
\end{definition}

对 $\lambda\in\mathbb F$，齐次方程组有非零解当且仅当系数矩阵奇异，故 $\lambda$ 是特征值当且仅当 $\chi_A(\lambda)=0$。这一说法必须连同数域理解：例如
\[
R=\begin{pmatrix}0&-1\\1&0\end{pmatrix}
\]
在 $\mathbb R$ 上没有特征值，在 $\mathbb C$ 上却有特征值 $i,-i$。复数域上的 $n$ 次特征多项式总有 $n$ 个根，不能据此断言任一实矩阵都有实特征向量。

\begin{proposition}
\label{prop:xie-ex-eig-geometric-algebraic}
对任一特征值 $\lambda$，有 $1\le g_A(\lambda)\le a_A(\lambda)$。
\end{proposition}
\begin{proof}
设 $g=g_A(\lambda)$，取 $E_\lambda(A)$ 的一组基并扩充为全空间的基。相应表示矩阵为
\[
\begin{pmatrix}\lambda I_g&C\\0&B\end{pmatrix},
\]
其特征多项式含因子 $(t-\lambda)^g$，所以代数重数至少为 $g$；特征子空间非零又保证 $g\ge1$。
\end{proof}

\subsection{相似与上三角化}

\begin{definition}
若存在 $P\in\mathbb F^{n\times n}$ 可逆，使 $B=P^{-1}AP$，则称 $A,B$ 在 $\mathbb F$ 上\emph{相似（similar）}，记作 $A\sim B$。这表示同一个线性变换在不同基下的矩阵；其中 $P$ 的列是新基在旧基下的坐标。
\end{definition}

由
\[
\det(tI_n-P^{-1}AP)=\det\bigl(P^{-1}(tI_n-A)P\bigr)=\chi_A(t)
\]
可知特征多项式是相似不变量，因此线性变换的特征多项式可以用任意一组基下的表示矩阵定义。相似还保持 $\operatorname{rank}f(A)$、$\operatorname{tr}A$ 和 $\det A$，因为 $f(P^{-1}AP)=P^{-1}f(A)P$。

若矩阵相似于一个上三角矩阵，就称它可以\emph{上三角化（triangularization）}。

\begin{theorem}
\label{thm:xie-ex-eig-triangularization}
$A$ 在 $\mathbb F$ 上可相似上三角化，当且仅当 $\chi_A(t)$ 在 $\mathbb F$ 上分裂为一次因式的乘积。特别地，每个复矩阵都可在复数域上相似上三角化。
\end{theorem}
\begin{proof}
上三角矩阵的特征多项式是其对角元对应的一次因式之积，故必要性成立。充分性对阶数归纳。选取特征向量 $v_1\in\mathbb F^n$ 并扩充为一组基，得到
\[
P^{-1}AP=\begin{pmatrix}\lambda_1&*\\0&A_1\end{pmatrix}.
\]
由于 $\chi_{A_1}(t)=\chi_A(t)/(t-\lambda_1)$ 仍在 $\mathbb F$ 上分裂，可在后 $n-1$ 个坐标中继续上三角化；与第一步合成即得所需变换。
\end{proof}

\subsection{可对角化的基本判据}

\begin{definition}
若存在 $\mathbb F$ 上的可逆矩阵 $P$，使 $P^{-1}AP$ 为对角矩阵，则称 $A$ 在 $\mathbb F$ 上\emph{可对角化（diagonalizable）}。线性变换可对角化，是指存在一组由特征向量组成的基。
\end{definition}

\begin{lemma}
\label{lem:xie-ex-eig-eigenspace-direct-sum}
属于互异特征值 $\lambda_1,\ldots,\lambda_s$ 的特征子空间之和是直和。因此，从这些子空间各取一组线性无关向量，其合并仍线性无关。
\end{lemma}
\begin{proof}
设 $v_i\in E_{\lambda_i}(T)$ 且 $\sum_i v_i=0$。固定 $i$，将 $\prod_{j\ne i}(T-\lambda_jI_V)$ 作用于等式，得到
\[
\left(\prod_{j\ne i}(\lambda_i-\lambda_j)\right)v_i=0.
\]
括号内的数非零，所以每个 $v_i$ 都为零，这正是直和的判据。
\end{proof}

\begin{theorem}
\label{thm:xie-ex-eig-diagonalization-basic}
设 $\chi_A(t)$ 在 $\mathbb F$ 上分裂，不同特征值为 $\lambda_1,\ldots,\lambda_s$。下列条件等价：$A$ 在 $\mathbb F$ 上可对角化；$A$ 有 $n$ 个线性无关的特征向量；$\mathbb F^n=\bigoplus_{i=1}^sE_{\lambda_i}(A)$；对每个 $i$ 均有 $g_A(\lambda_i)=a_A(\lambda_i)$。特别地，$n$ 个不同特征值足以保证可对角化。
\end{theorem}
\begin{proof}
设 $P=(v_1,\ldots,v_n)$，则 $AP=P\operatorname{diag}(\mu_1,\ldots,\mu_n)$ 当且仅当 $Av_j=\mu_jv_j$，所以前两个条件等价。由引理~\ref{lem:xie-ex-eig-eigenspace-direct-sum}，所有特征子空间的基合并后线性无关，故它们构成全空间的基当且仅当几何重数之和等于 $n$。又由命题~\ref{prop:xie-ex-eig-geometric-algebraic}，每个几何重数不超过对应代数重数，而代数重数之和为 $n$，因此总和相等当且仅当逐个相等。
\end{proof}

\emph[caution]{重特征值本身不妨碍对角化，缺少足够的线性无关特征向量才是障碍。}例如 $I_n$ 只有一个不同的特征值，却已经是对角矩阵；$\begin{pmatrix}1&1\\0&1\end{pmatrix}$ 的特征值 $1$ 的代数重数为 $2$、几何重数为 $1$，因而不可对角化。

\subsection{极小多项式与特征值估计的入口}

\begin{definition}
使 $m_A(A)=0$ 的非零多项式中，次数最低且首项系数为 $1$ 的多项式，称为 $A$ 的\emph{极小多项式（minimal polynomial）}。满足 $f(A)=0$ 时，也说 $A$ 适合多项式 $f$；非零的这样的 $f$ 称为 $A$ 的\emph{零化多项式（annihilating polynomial）}。
\end{definition}

特征多项式由行列式定义，极小多项式由矩阵所满足的关系定义。二者并不总相同，例如 $I_n$ 的特征多项式为 $(t-1)^n$，极小多项式为 $t-1$。极小多项式的存在唯一性及 Cayley--Hamilton 定理将在第~\ref{sec:xie-ex-eig-minimal-polynomial} 节证明；特征值不易精确求出时，则可用第~\ref{subsec:xie-ex-eig-gershgorin} 小节的圆盘定理估计其位置。

\section{特征值和特征向量}
\label{sec:xie-ex-eig-eigen-computation}

\subsection{从定义和不变子空间出发}

\begin{example}
% 教材例6.1、例6.3、例6.4。
设 $L_A\colon\mathbb F^{n\times n}\to\mathbb F^{n\times n}$ 定义为 $L_A(X)=AX$。说明 $L_A$ 与 $A$ 的特征值有何关系。
\end{example}
\begin{solve}
若 $Av=\lambda v$ 且 $v\ne0$，则 $X=(v,0,\ldots,0)$ 是 $L_A$ 的特征向量。反之，$AX=\lambda X$ 且 $X\ne0$ 时，$X$ 至少有一个非零列，该列就是 $A$ 的特征向量。因此二者的不同特征值相同。进一步，按列排列矩阵空间的基，$L_A$ 的表示矩阵为 $n$ 个 $A$ 组成的分块对角矩阵，所以
\[
\chi_{L_A}(t)=\chi_A(t)^n,\qquad
\dim E_\lambda(L_A)=n\dim E_\lambda(A).
\]
可见特征值相同并不意味着重数相同。
\end{solve}

一般地，若 $V=V_1\oplus\cdots\oplus V_r$，每个 $V_i$ 均为 $T$ 的不变子空间，则拼接各子空间的基可得分块对角表示，从而
\begin{equation}
\chi_T(t)=\prod_{i=1}^r\chi_{T|_{V_i}}(t),\qquad
E_\lambda(T)=\bigoplus_{i=1}^rE_\lambda(T|_{V_i}).
\label{eq:xie-ex-eig-block-eigenspaces}
\end{equation}
第二个等式来自直和分量的唯一性：将 $Tv=\lambda v$ 分解到各 $V_i$ 中，即逐项得到相同的特征向量方程。因此分块对角矩阵的代数重数与几何重数都按块相加，分块的特征向量只需在其余坐标处补零。

\begin{example}
% 教材例6.5。
整数矩阵不可能有非整数的有理特征值。
\end{example}
\begin{solve}
整数矩阵的特征多项式为首一整系数多项式。若互素整数 $p,q$ 满足 $q>0$ 且 $p/q$ 是其根，将等式乘以 $q^n$ 可得 $q\mid p^n$，因而 $q=1$。所以 $q>1$ 时，方程 $Ax=(p/q)x$ 只有零解。
\end{solve}

\subsection{正向使用矩阵多项式}

\begin{theorem}
\label{thm:xie-ex-eig-spectral-mapping}
设 $A$ 的全部复特征值为 $\lambda_1,\ldots,\lambda_n$，$f\in\mathbb F[t]$，则 $f(A)$ 的全部特征值为 $f(\lambda_1),\ldots,f(\lambda_n)$，按此列表计重。
\end{theorem}
\begin{proof}
由定理~\ref{thm:xie-ex-eig-triangularization}，可取 $P$ 使 $U=P^{-1}AP$ 为上三角矩阵，其对角元依次为 $\lambda_i$。于是 $P^{-1}f(A)P=f(U)$ 仍上三角，对角元依次为 $f(\lambda_i)$。
\end{proof}

对单个特征向量，$Av=\lambda v$ 可直接推出 $f(A)v=f(\lambda)v$；上述证明进一步保证列出了全部特征值及其代数重数。若不同的 $\lambda_i$ 经 $f$ 映到同一点，相应重数需要合并；也不能反过来断言 $f(A)$ 的每个特征向量都是 $A$ 的特征向量，例如 $f$ 为常数时便不成立。

若 $A$ 可逆，则 $A^{-1}$ 的特征值为 $\lambda_i^{-1}$，因为 $U^{-1}$ 的对角元为 $\lambda_i^{-1}$。无论 $A$ 是否可逆，$\operatorname{adj}A$ 的特征值均为
\[
\prod_{j\ne1}\lambda_j,\quad\ldots,\quad\prod_{j\ne n}\lambda_j,
\]
其中 $n=1$ 时空积为 $1$。其理由是上三角矩阵的伴随矩阵仍上三角，且相似变换保持伴随运算；后者对可逆矩阵由 $\operatorname{adj}A=(\det A)A^{-1}$ 得到，对一般矩阵则用 $A+sI_n$ 的多项式恒等式延拓到 $s=0$。

\begin{example}
% 教材例6.8。
若 $A$ 的特征值为 $\lambda_1,\ldots,\lambda_n$，求
\[
M=\begin{pmatrix}A&A^2\\A^2&A\end{pmatrix}
\]
的特征值。
\end{example}
\begin{solve}
取 $S=\begin{pmatrix}I_n&I_n\\I_n&-I_n\end{pmatrix}$，直接相乘可得
\[
S^{-1}MS=\operatorname{diag}(A+A^2,A-A^2).
\]
故其特征值为 $\lambda_i+\lambda_i^2$ 与 $\lambda_i-\lambda_i^2$，$1\le i\le n$。这一结论无需假设 $A$ 可对角化。
\end{solve}

\begin{example}
% 教材例6.10、例6.11；统一采用下次对角线为1的友阵。
设 $p(t)=t^n+a_{n-1}t^{n-1}+\cdots+a_0$，其友阵为
\[
C_p=\begin{pmatrix}
0&0&\cdots&0&-a_0\\
1&0&\cdots&0&-a_1\\
0&1&\cdots&0&-a_2\\
\vdots&\vdots&\ddots&\vdots&\vdots\\
0&0&\cdots&1&-a_{n-1}
\end{pmatrix}.
\]
如何构造以 $f(\lambda_1),\ldots,f(\lambda_n)$ 为根的首一多项式，其中 $\lambda_i$ 是 $p$ 的根？
\end{example}
\begin{solve}
沿 $\det(tI_n-C_p)$ 的第一行展开，最后一项的余子式为对角元均为 $-1$ 的三角矩阵，连同展开符号后贡献 $a_0$；第一项则为 $t$ 乘以系数为 $a_1,\ldots,a_{n-1}$ 的低一阶友阵的特征行列式。因此对阶数递推得到 $\chi_{C_p}(t)=t(t^{n-1}+a_{n-1}t^{n-2}+\cdots+a_1)+a_0=p(t)$；一阶时友阵为 $(-a_0)$，结论直接成立。由定理~\ref{thm:xie-ex-eig-spectral-mapping}，所求多项式为
\[
h(z)=\det\bigl(zI_n-f(C_p)\bigr).
\]
同时，$\det f(C_p)=\prod_i f(\lambda_i)$；按首一多项式的结式约定，此积正是 $\operatorname{Res}(p,f)$。一般的 $A$ 也满足 $\det f(A)=\operatorname{Res}(\chi_A,f)$。
\end{solve}

\subsection{反向使用矩阵适合的多项式}

若 $q(A)=0$ 且 $Av=\lambda v$、$v\ne0$，则 $0=q(A)v=q(\lambda)v$，所以 $q(\lambda)=0$。这给出特征值的候选集合，\emph[caution]{并不保证 $q$ 的每个根都实际出现}。例如 $I_n$ 适合 $t^2-1$，但没有特征值 $-1$。

\begin{proposition}
\label{prop:xie-ex-eig-nilpotent-spectrum}
$A$ 为\emph{幂零矩阵（nilpotent matrix）}，即存在正整数 $r$ 使 $A^r=0$，当且仅当它的全部复特征值都为零。此时一定有 $A^n=0$。
\end{proposition}
\begin{proof}
$A^r=0$ 推出每个特征值满足 $\lambda^r=0$。反之，全部特征值为零时，$A$ 可复相似于严格上三角矩阵 $U$；$U$ 的每次相乘都使可能非零的位置至少向右移动一列，故 $U^n=0$，于是 $A^n=0$。
\end{proof}

若 $A^m=I_n$，则其特征值均为 $m$ 次单位根。每行每列恰有一个非零元、且非零元为 $\pm1$ 的矩阵也具有有限阶：这样的矩阵仅有 $2^nn!$ 个，乘法封闭且均可逆，所以某两个正整数次幂相等，约去可逆因子便得到 $A^m=I_n$。对合矩阵 $A^2=I_n$ 的特征值只能为 $\pm1$，幂等矩阵 $A^2=A$ 的特征值只能为 $0,1$；它们的可对角化性还需要第~\ref{sec:xie-ex-eig-diagonalization-tests} 节的无重根多项式判据。

\begin{example}
% 教材例6.16。
设 $A$ 为实矩阵，$I_n-A$ 的所有特征值模长均小于 $1$，则 $0<\det A<2^n$。
\end{example}
\begin{solve}
$A$ 的特征值满足 $|1-\lambda|<1$。实特征值因此属于 $(0,2)$；非实特征值按共轭成对出现，每对之积为 $|\lambda|^2\in(0,4)$。将这些正因子相乘，就得到所述严格不等式。
\end{solve}

\subsection{降阶公式、低秩结构与交织关系}

\begin{theorem}
\label{thm:xie-ex-eig-ab-ba}
设 $A\in\mathbb F^{m\times n}$、$B\in\mathbb F^{n\times m}$ 且 $m\ge n$，则
\begin{equation}
\det(tI_m-AB)=t^{m-n}\det(tI_n-BA).
\label{eq:xie-ex-eig-ab-ba-characteristic}
\end{equation}
因而 $AB,BA$ 的非零特征值连同代数重数相同；若 $m=n$，则特征多项式完全相同。
\end{theorem}
\begin{proof}
当 $t\ne0$ 时，对分块矩阵 $\begin{pmatrix}tI_m&A\\B&I_n\end{pmatrix}$ 分别消去右上块和左下块，得到
\[
\det(tI_m-AB)=t^m\det(I_n-t^{-1}BA)
=t^{m-n}\det(tI_n-BA).
\]
两端都是 $t$ 的多项式，且在无穷多个非零点相等，故恒等，也包括 $t=0$。
\end{proof}

\begin{example}
\label{ex:xie-ex-eig-rank-one-spectrum}
% 教材例6.20、例6.21。
设 $u,v\in\mathbb F^n$，$R=uv^{\mathsf T}$。当 $n\ge2$ 时，降阶公式给出
\[
\chi_R(t)=t^{n-1}(t-v^{\mathsf T}u),\qquad
R^2=(v^{\mathsf T}u)R.
\]
一阶时直接有 $R=(v^{\mathsf T}u)$。特别地，实单位列向量 $u$ 对应的矩阵 $H=I_n-2uu^{\mathsf T}$ 在 $\mathbb Ru$ 上为 $-I$，在 $\ker u^{\mathsf T}$ 上为 $I$，所以其特征值为 $-1$ 一重、$1$ 共 $n-1$ 重。
\end{example}

\begin{example}
% 教材例6.22；说明退化时的重数合并。
设 $n\ge2$，$a=(a_1,\ldots,a_n)^{\mathsf T}\in\mathbb R^n$ 满足 $\sum_i a_i=0$，$e=(1,\ldots,1)^{\mathsf T}$。求 $M=-I_n+aa^{\mathsf T}+ee^{\mathsf T}$ 的特征值。
\end{example}
\begin{solve}
令 $U=(a,e)$，则 $M=-I_n+UU^{\mathsf T}$，而
\[
U^{\mathsf T}U=\operatorname{diag}(\|a\|^2,n).
\]
故
\[
\chi_M(t)=(t+1)^{n-2}(t+1-\|a\|^2)(t+1-n).
\]
根重合时应合并重数。若 $a\ne0$，则 $a,e$ 正交，分别给出特征值 $\|a\|^2-1,n-1$，它们的公共正交补上为 $-I$；若 $a=0$，只保留 $e$ 及其正交补。因此各种退化情形下 $M$ 都可在实数域上对角化。
\end{solve}

\begin{proposition}
\label{prop:xie-ex-eig-intertwining}
若 $A\in\mathbb F^{m\times m}$、$B\in\mathbb F^{n\times n}$、$X\in\mathbb F^{m\times n}$ 满足 $AX=XB$，且 $\operatorname{rank}X=r$，则 $\chi_A,\chi_B$ 有一个次数为 $r$ 的公共首一因式。特别地，它们至少有 $r$ 个公共特征值，按重数计。
\end{proposition}
\begin{proof}
由交织等式可知 $\operatorname{Im}X$ 为 $A$ 的不变子空间，$\ker X$ 为 $B$ 的不变子空间。映射
\[
\overline X\colon\mathbb F^n/\ker X\longrightarrow\operatorname{Im}X,
\qquad [v]\longmapsto Xv
\]
为线性同构，并将商空间上由 $B$ 诱导的变换与 $A|_{\operatorname{Im}X}$ 联系起来。因此二者有相同的 $r$ 次特征多项式；各自在扩充基下都是全矩阵的对角分块，故该多项式同时整除 $\chi_A,\chi_B$。
\end{proof}

\subsection{特征多项式系数、主子式与迹}

对 $J\subseteq\{1,\ldots,n\}$，记 $A[J,J]$ 为同时保留 $J$ 中行列得到的主子矩阵。将行列式 $\det(tI_n-A)$ 按多重线性展开：选择 $k$ 个位置使用 $-A$ 的元素，其余位置使用 $tI_n$ 的对角元，便得到
\begin{equation}
c_k=(-1)^k\sum_{|J|=k}\det A[J,J]
=(-1)^k\sum_{i_1<\cdots<i_k}\lambda_{i_1}\cdots\lambda_{i_k}.
\label{eq:xie-ex-eig-principal-minors}
\end{equation}
这里 $c_k$ 沿用式~\eqref{eq:xie-ex-eig-characteristic-polynomial}。特别地，$\operatorname{tr}A=\sum_i\lambda_i$，$\det A=\prod_i\lambda_i$，而 $A$ 可逆当且仅当零不是其特征值。

记 $p_k=\operatorname{tr}(A^k)=\sum_i\lambda_i^k$，并令 $c_0=1$。Newton 恒等式为
\begin{equation}
kc_k+\sum_{j=1}^k c_{k-j}p_j=0,\qquad 1\le k\le n.
\label{eq:xie-ex-eig-newton-identities}
\end{equation}
它可由形式幂级数恒等式
\[
\frac{\mathrm d}{\mathrm dz}\prod_i(1-\lambda_i z)
=-\prod_i(1-\lambda_i z)\sum_{j\ge1}p_jz^{j-1}
\]
比较系数得到，因而不涉及级数的解析收敛问题。由于数域的特征为零，可依次除以 $1,\ldots,n$，从前 $n$ 个幂次的迹恢复特征多项式。

\begin{example}
% 教材例6.30、例6.31、例6.32。
若四阶矩阵满足 $\operatorname{tr}(A^k)=k$，$1\le k\le4$，则依次得到
\[
c_1=-1,\quad c_2=-\frac12,\quad c_3=-\frac16,\quad
\det A=c_4=\frac1{24}.
\]
一般地，$n$ 阶矩阵幂零当且仅当 $\operatorname{tr}(A^k)=0$ 对 $1\le k\le n$ 均成立。
\end{example}
\begin{solve}
四阶结论直接代入式~\eqref{eq:xie-ex-eig-newton-identities}。若前 $n$ 个迹均为零，递推给出 $c_1=\cdots=c_n=0$，于是 $\chi_A(t)=t^n$，由命题~\ref{prop:xie-ex-eig-nilpotent-spectrum} 得到幂零性；反向由全部特征值为零即得。
\end{solve}

\begin{example}
\label{ex:xie-ex-eig-commutator-nilpotent}
设 $C=AB-BA$。若 $AC=CA$，则 $C$ 幂零；假设 $BC=CB$ 也可得到相同结论。
\end{example}
\begin{solve}
当 $AC=CA$ 时，对任意 $k\ge1$，
\[
C^k=A(C^{k-1}B)-(C^{k-1}B)A.
\]
乘积交换次序不改变迹，因此 $\operatorname{tr}(C^k)=0$，再用前一例。另一种条件下交换 $A,B$ 并对 $-C$ 作相同论证即可。这里并不需要同时假设两个交换关系。
\end{solve}

\subsection{Gershgorin 圆盘与特征值的估计}
\label{subsec:xie-ex-eig-gershgorin}

\begin{theorem}
\label{thm:xie-ex-eig-gershgorin}
设 $A=(a_{ij})\in\mathbb C^{n\times n}$，令
\[
R_i=\sum_{j\ne i}|a_{ij}|,\qquad
D_i=\{z\in\mathbb C:|z-a_{ii}|\le R_i\}.
\]
则 $A$ 的全部特征值均在 $\bigcup_iD_i$ 内。若这一并集的某个连通分支由 $k$ 个圆盘组成，则该分支中恰有 $k$ 个特征值。圆盘及特征值都按重数计。这分别称为\emph{Gershgorin 第一、第二圆盘定理}。
\end{theorem}
\begin{proof}
若 $Ax=\lambda x$，选取 $|x_i|=\max_j|x_j|>0$，则
\[
|\lambda-a_{ii}|\,|x_i|
=\left|\sum_{j\ne i}a_{ij}x_j\right|
\le R_i|x_i|,
\]
从而 $\lambda\in D_i$。第二个结论用连续变形说明。令 $A(s)$ 保留 $A$ 的对角元，将非对角元乘以 $s\in[0,1]$。其圆盘半径为 $sR_i$，始终包含于原圆盘内。原并集的不同连通分支都是彼此分离的紧集，可在它们之间选取不经过任何圆盘的分隔边界。多项式的根作为计重的多重集合随系数连续变化，因此没有特征值能穿过该边界，分支内根的总数保持不变。$s=0$ 时矩阵为对角矩阵，该分支包含的对角元恰为对应的 $k$ 个圆心，故 $s=1$ 时仍有 $k$ 个特征值。
\end{proof}

对 $A^{\mathsf T}$ 应用同一定理，还能得到以列的非对角元绝对值和为半径的圆盘；全部特征值同时受两组圆盘并集约束。第二圆盘定理对互相连接的圆盘给出的是整个分支的计数，不能未经分离就把每个圆盘各分配一个特征值。

\begin{example}
% 教材§6.2中的例3.83及例6.33。
若 $|a_{ii}|>\sum_{j\ne i}|a_{ij}|$ 对所有 $i$ 成立，则 $A$ 可逆。若 $A$ 为实矩阵且进一步有 $a_{ii}>\sum_{j\ne i}|a_{ij}|$，则 $\det A>0$。
\end{example}
\begin{solve}
第一种条件使所有圆盘避开原点，故零不是特征值。第二种条件使所有圆盘位于右半平面，故实特征值均为正，非实特征值按共轭成对出现、每对之积为正，因此行列式为正。若一个独立连通分支由两个外切圆盘组成，则在 $s<1$ 时它们缩小为互不相交的圆盘，各含一个特征值；令 $s\to1$ 可知，一个圆盘除去切点后不可能容纳这两个特征值。
\end{solve}

\begin{example}
% 教材例6.34。
任给 $A\in\mathbb C^{n\times n}$，存在 $\delta>0$，使
\[
A(s)=A+\operatorname{diag}(s,s^2,\ldots,s^n)
\]
在每个实参数 $0<s<\delta$ 时都有 $n$ 个不同特征值。
\end{example}
\begin{solve}
当 $s$ 充分大时，不同圆心 $a_{ii}+s^i$ 之间的距离均大于相应两个固定半径之和，所以圆盘互不相交，各含一个特征值。$\chi_{A(s)}(t)$ 关于 $t$ 的判别式是 $s$ 的多项式；既然在充分大的 $s$ 处非零，它就不是零多项式，因而只有有限个零点。取 $\delta$ 小于其所有非零根模长的最小值；若没有非零根，任取 $\delta>0$。此时 $0<s<\delta$ 不会使判别式为零，于是特征值互异。这里先使用大参数证明判别式非恒零，再由零点有限性得到小参数结论。
\end{solve}

\section{乘法交换性诱导的同时性质}
\label{sec:xie-ex-eig-simultaneous}

\subsection{特征子空间的不变性与公共特征向量}

\begin{proposition}
\label{prop:xie-ex-eig-commuting-invariance}
若 $ST=TS$，则 $S$ 的每个特征子空间都是 $T$ 的不变子空间，反之亦然。这个结论只要求所讨论的特征值属于当前数域。
\end{proposition}
\begin{proof}
若 $Sv=\lambda v$，则 $S(Tv)=T(Sv)=\lambda Tv$，故 $Tv\in E_\lambda(S)$。
\end{proof}

\begin{theorem}
\label{thm:xie-ex-eig-common-eigenvector}
设 $T_1,\ldots,T_r$ 为非零有限维空间上的两两可交换线性变换，且各自的特征多项式均在 $\mathbb F$ 上分裂，则它们在 $\mathbb F$ 上有公共特征向量。
\end{theorem}
\begin{proof}
先取 $T_1$ 的一个非零特征子空间 $W$，它对其余变换都不变。$T_2|_W$ 的特征多项式整除 $T_2$ 的特征多项式：将 $W$ 的基扩充后，$T_2$ 的矩阵为分块上三角矩阵。因此该限制仍有 $\mathbb F$ 中的特征值。再取它在 $W$ 中的非零特征子空间，所得空间对所有后续变换仍不变。有限次重复后，所得非零空间中的任一非零向量就是所需公共特征向量。
\end{proof}

需要分裂条件的是公共特征向量的存在性，特征子空间的不变性本身不依赖它。例如 $I_2$ 与实旋转矩阵 $R$ 可交换，但 $R$ 没有实特征向量，所以二者没有公共实特征向量。

\begin{example}
% 教材例6.36。
设 $V$ 为非零有限维复线性空间，$\mathcal S$ 为一族线性变换，没有非平凡公共不变子空间。若 $T$ 与族内每个变换都可交换，则 $T$ 为纯量变换。
\end{example}
\begin{solve}
复数域上 $T$ 有非零特征子空间 $E_\lambda(T)$。由命题~\ref{prop:xie-ex-eig-commuting-invariance}，它是 $\mathcal S$ 的公共不变子空间，因而只能等于 $V$，于是 $T=\lambda I_V$。
\end{solve}

\subsection{同时上三角化与同时对角化}

\begin{theorem}
\label{thm:xie-ex-eig-simultaneous-triangular}
两两可交换且特征多项式均在 $\mathbb F$ 上分裂的有限个矩阵，可以在 $\mathbb F$ 上同时上三角化。
\end{theorem}
\begin{proof}
对空间维数归纳。由定理~\ref{thm:xie-ex-eig-common-eigenvector}，可取公共特征向量作为首个基向量，扩充后各矩阵均为
\[
\begin{pmatrix}\lambda_i&*\\0&B_i\end{pmatrix}.
\]
原矩阵两两可交换，故比较右下分块得到 $B_iB_j=B_jB_i$，且各 $\chi_{B_i}$ 仍分裂。对右下块应用归纳假设，就得到同一组基下的上三角表示。
\end{proof}

\begin{lemma}
\label{lem:xie-ex-eig-diagonal-restriction}
若 $T$ 在 $\mathbb F$ 上可对角化，$W$ 为 $T$ 的不变子空间，则 $T|_W$ 及 $V/W$ 上的诱导变换也在 $\mathbb F$ 上可对角化。
\end{lemma}
\begin{proof}
设 $\lambda_1,\ldots,\lambda_s$ 为不同特征值。多项式算子
\[
P_i=\prod_{j\ne i}\frac{T-\lambda_jI_V}{\lambda_i-\lambda_j}
\]
在 $E_{\lambda_i}(T)$ 上为恒等变换，在其余特征子空间上为零，故 $\sum_iP_i=I_V$。$W$ 对 $T$ 不变，因而对每个 $P_i$ 不变，于是
\[
W=\bigoplus_i\bigl(W\cap E_{\lambda_i}(T)\bigr).
\]
这给出 $W$ 的特征向量基。对各 $W\cap E_{\lambda_i}(T)$ 的基在对应 $E_{\lambda_i}(T)$ 内扩充，新增基向量的陪集构成 $V/W$ 的特征向量基，故商空间结论也成立。
\end{proof}

\begin{theorem}
\label{thm:xie-ex-eig-simultaneous-diagonal}
一族有限个在 $\mathbb F$ 上可对角化的矩阵可同时对角化，当且仅当它们两两可交换。
\end{theorem}
\begin{proof}
同一组基下的对角矩阵两两可交换，故必要性成立。反之，按第一个变换的特征子空间分解全空间，其余变换都保持这些子空间，并且由引理~\ref{lem:xie-ex-eig-diagonal-restriction}，在其上的限制仍可对角化。再在各分量内按第二个变换的特征子空间细分，依次进行，最终得到由公共特征向量组成的基。
\end{proof}

\emph[caution]{交换性保证同时上三角化还需要分裂条件，保证同时对角化还需要每个矩阵分别可对角化。}例如 $I_2$ 与非对角化的 Jordan 块可交换，却不能同时对角化；反过来，同时上三角化本身也不推出交换性。

\subsection{幂零扰动与交换子的应用}

\begin{proposition}
若 $NB=BN$ 且 $N$ 幂零，则 $\chi_{B+N}(t)=\chi_B(t)$，特别地 $\det(B+N)=\det B$。
\end{proposition}
\begin{proof}
扩充到复数域后同时上三角化 $N,B$。$N$ 的对角元全为零，因此 $B+N$ 与 $B$ 的上三角表示有完全相同的对角元，从而特征多项式相同。两端本来均在 $\mathbb F[t]$ 中，所以所得恒等式也在原数域上成立。
\end{proof}

\begin{proposition}
% 教材例6.32的延拓。
在复数域上，若 $C=AB-BA$ 同时满足 $AC=CA$、$BC=CB$，则 $A,B,C$ 可同时上三角化。
\end{proposition}
\begin{proof}
由例~\ref{ex:xie-ex-eig-commutator-nilpotent}，$C$ 幂零，所以 $W=\ker C\ne\{0\}$。交换关系使 $W$ 对 $A,B$ 都不变，而在 $W$ 上有 $A|_WB|_W=B|_WA|_W$。故这两个限制有公共特征向量，它同时也是 $C$ 的零特征向量。以此为首个基向量，在商空间上仍有同样的交换子等式与交换条件；对维数归纳便得到共同的上三角基。这里 $A,B$ 本身未必可交换，不能直接套用定理~\ref{thm:xie-ex-eig-simultaneous-triangular}。
\end{proof}

\section{矩阵相似和可对角化的计算}
\label{sec:xie-ex-eig-similarity-computation}

\subsection{相似初等变换及相似不变量}

初等行变换一般只能保证矩阵相抵，不能单独保证相似。若左乘初等矩阵 $E$，就必须同时右乘 $E^{-1}$，得到 $EAE^{-1}$。例如，交换第 $i,j$ 行后同时交换第 $i,j$ 列；第 $i$ 行乘以 $c\ne0$ 后，第 $i$ 列乘以 $c^{-1}$；第 $i$ 行加上第 $j$ 行的 $c$ 倍后，第 $j$ 列减去第 $i$ 列的 $c$ 倍。最后一种对应 $E=I_n+cE_{ij}$，$i\ne j$，其中 $E_{ij}$ 是只有第 $(i,j)$ 个元素为 $1$ 的矩阵。分块变换也遵循相同原则，尤其要保持乘法次序：分块行左乘的逆操作，对应分块列右乘。

因为每个可逆矩阵都是初等矩阵的乘积，每次相似变换都可以分解为若干次相似初等变换。计算中若要追踪总变换矩阵，必须按发生次序记录左右因子，最后直接检验 $AP=PB$，从而确认 $B=P^{-1}AP$。

\begin{example}
% 教材例6.46、例6.48及相似不变量的补充反例。
分块对角矩阵的方块可任意重排而保持相似；另一方面，相同的特征多项式不足以保证相似。
\end{example}
\begin{solve}
前一结论只需同时置换各块对应的行和列，即改变直和分量所用基的排列。对后一结论，$I_2$ 与 $J=\begin{pmatrix}1&1\\0&1\end{pmatrix}$ 都有特征多项式 $(t-1)^2$，但 $\operatorname{rank}(I_2-I_2)=0$，$\operatorname{rank}(J-I_2)=1$，因此不相似。类似地，$AB-BA$ 不可能相似于 $cI_n$，$c\ne0$，因为前者的迹为零，后者的迹为 $nc\ne0$。相似不变量不同时足以否定相似；只核对少量不变量相同，通常还不能肯定相似。
\end{solve}

\begin{example}
% 教材例6.47。
设 $A,B$ 同为 $n$ 阶矩阵，且 $\operatorname{rank}(ABA)=\operatorname{rank}B$，证明 $AB\sim BA$。
\end{example}
\begin{solve}
若 $A=0$ 或 $A$ 可逆，结论分别由两个乘积均为零或 $BA=A^{-1}(AB)A$ 得到。其余情形，取可逆 $P,Q$，使 $PAQ=\operatorname{diag}(I_r,0)$，同时将 $B$ 换成 $Q^{-1}BP^{-1}$。此变换分别把 $AB,BA$ 换成与它们相似的矩阵，并保持秩条件，所以可设
\[
A=\begin{pmatrix}I_r&0\\0&0\end{pmatrix},\qquad
B=\begin{pmatrix}B_{11}&B_{12}\\B_{21}&B_{22}\end{pmatrix}.
\]
此时 $\operatorname{rank}B_{11}=\operatorname{rank}B$。所以 $B_{12}$ 的列在 $B_{11}$ 的列空间内，$B_{21}$ 的行在 $B_{11}$ 的行空间内，可写成 $B_{12}=B_{11}N$、$B_{21}=MB_{11}$。于是
\[
AB=\begin{pmatrix}B_{11}&B_{11}N\\0&0\end{pmatrix},\qquad
BA=\begin{pmatrix}B_{11}&0\\MB_{11}&0\end{pmatrix}.
\]
分别用 $S=\begin{pmatrix}I_r&-N\\0&I_{n-r}\end{pmatrix}$ 与 $T=\begin{pmatrix}I_r&0\\M&I_{n-r}\end{pmatrix}$ 作相似变换，得到
\[
S^{-1}(AB)S=T^{-1}(BA)T=\operatorname{diag}(B_{11},0).
\]
这说明附加的秩条件足以把降阶公式中的同特征多项式加强为相似。
\end{solve}

\subsection{过渡矩阵与含参数矩阵}

若 $A$ 可对角化，计算可按以下顺序进行：先求 $\chi_A$ 并确定当前数域中的特征值及代数重数；对每个 $\lambda$ 解 $(A-\lambda I_n)x=0$，比较几何重数与代数重数；若全部相等，将各特征子空间的基拼成 $P$，并按列向量的对应次序写出对角矩阵 $D$，检验 $AP=PD$。特征向量可以数乘，也可以在同一特征子空间内换基，所以 $P$ 一般不唯一。

\begin{example}
% 教材例6.50。
设三阶矩阵 $A$ 的特征值为 $1,1,4$，对应特征向量为 $(2,1,0)^{\mathsf T}$、$(-1,0,1)^{\mathsf T}$、$(0,1,1)^{\mathsf T}$，求 $A$。
\end{example}
\begin{solve}
以所给向量为列组成
\[
P=\begin{pmatrix}2&-1&0\\1&0&1\\0&1&1\end{pmatrix},\qquad
D=\operatorname{diag}(1,1,4).
\]
由于 $\det P=-1\ne0$，它们构成特征向量基，故唯一确定
\[
A=PDP^{-1}=\begin{pmatrix}1&0&0\\-3&7&-3\\-3&6&-2\end{pmatrix}.
\]
将此矩阵乘以 $P$ 即得 $PD$，可直接核对全部给定条件。
\end{solve}

\begin{example}
% 补充例：说明教材例6.51、例6.52中的参数判定步骤。
判定 $A_s=\begin{pmatrix}2&s&0\\0&2&0\\0&0&3\end{pmatrix}$ 何时可对角化。
\end{example}
\begin{solve}
其特征多项式恒为 $(t-2)^2(t-3)$，参数不影响特征值，却影响特征子空间。特征值 $2$ 的几何重数为 $3-\operatorname{rank}(A_s-2I_3)$；当 $s=0$ 时秩为 $1$，当 $s\ne0$ 时秩为 $2$。特征值 $3$ 为单根，几何重数必为 $1$，所以 $A_s$ 可对角化当且仅当 $s=0$。含参数问题中，应先对重特征值施加秩条件，而不能只看特征多项式是否随参数改变。
\end{solve}

\subsection{矩阵幂与线性递推}

若 $A=PDP^{-1}$，则对整数 $k\ge0$ 有 $A^k=PD^kP^{-1}$；当 $A$ 可逆时也适用于负整数。设不同特征值为 $\lambda_1,\ldots,\lambda_s$，还可以使用谱投影
\begin{equation}
\Pi_i=\prod_{j\ne i}\frac{A-\lambda_jI_n}{\lambda_i-\lambda_j},\qquad
A^k=\sum_{i=1}^s\lambda_i^k\Pi_i.
\label{eq:xie-ex-eig-spectral-projections}
\end{equation}
该公式在 $k=0$ 时将每个系数 $\lambda_i^0$ 约定为 $1$；它来自各特征子空间上逐一比较作用，因此仅在 $A$ 可对角化时这样使用。它也说明，可对角化矩阵的任意多项式仍可对角化，即使多项式把若干特征值合并也不影响结论。

\begin{example}
% 教材例6.54。
设 Fibonacci 数列由 $F_0=0$、$F_1=1$、$F_{k+2}=F_{k+1}+F_k$ 定义，求其通项。
\end{example}
\begin{solve}
将二阶递推改写为一阶向量递推：
\[
\begin{pmatrix}F_{k+1}\\F_k\end{pmatrix}
=M^k\begin{pmatrix}1\\0\end{pmatrix},\qquad
M=\begin{pmatrix}1&1\\1&0\end{pmatrix}.
\]
特征多项式 $t^2-t-1$ 有两个不同实根 $\alpha=(1+\sqrt5)/2$、$\beta=(1-\sqrt5)/2$，对应特征向量 $(\alpha,1)^{\mathsf T}$、$(\beta,1)^{\mathsf T}$。据此对角化并代入初值，得到
\[
F_k=\frac{\alpha^k-\beta^k}{\sqrt5},\qquad k\ge0.
\]
一般的常系数线性递推也可先组成状态向量，再研究递推矩阵的幂；若该矩阵不可对角化，就需要其他求幂方法，不能强行使用对角分解。
\end{solve}

\section{可对角化的判定（一）}
\label{sec:xie-ex-eig-diagonalization-tests}

\subsection{无重根零化多项式与特征子空间分解}

\begin{theorem}
\label{thm:xie-ex-eig-squarefree-annihilator}
若 $A$ 适合多项式 $q(t)=\prod_{i=1}^s(t-\mu_i)$，其中 $\mu_i\in\mathbb F$ 两两不同，则 $A$ 在 $\mathbb F$ 上可对角化。反之，若 $A$ 在 $\mathbb F$ 上可对角化，则它适合以全部不同特征值为根的无重根多项式。
\end{theorem}
\begin{proof}
定义 Lagrange 多项式
\[
\ell_i(t)=\prod_{j\ne i}\frac{t-\mu_j}{\mu_i-\mu_j}.
\]
$\sum_i\ell_i(t)-1$ 次数不超过 $s-1$，却有 $s$ 个不同的根，所以为零多项式。又 $(t-\mu_i)\ell_i(t)=q(t)/\prod_{j\ne i}(\mu_i-\mu_j)$，故 $q(A)=0$ 推出
\[
\sum_i\ell_i(A)=I_n,\qquad
\operatorname{Im}\ell_i(A)\subseteq\ker(A-\mu_iI_n).
\]
这表明全空间为这些核的和；去掉零子空间后，由引理~\ref{lem:xie-ex-eig-eigenspace-direct-sum} 知其为特征子空间的直和，因而可对角化。反向在特征向量基下逐个对角元验证 $q(A)=0$ 即可。
\end{proof}

\emph[caution]{在 $\mathbb F$ 上无重根还不够，零化多项式必须在 $\mathbb F$ 上分裂。}例如 $R^2+I_2=0$，$t^2+1$ 在复数域中无重根，因此旋转矩阵 $R$ 可在复数域上对角化，但在实数域上不能。极小多项式形式的充要条件将在命题~\ref{prop:xie-ex-eig-minimal-diagonalization} 中统一给出；Jordan 块判据则需要后续的标准型理论。

\begin{example}
% 教材例6.14、例6.67、例6.69。
对合矩阵与幂等矩阵在任一数域上均可对角化。更一般地，若 $A^m=I_n$，则 $A$ 可在复数域上对角化；在原数域上还须所有实际出现的特征值都属于该数域。
\end{example}
\begin{solve}
对合矩阵适合 $(t-1)(t+1)$，幂等矩阵适合 $t(t-1)$，它们在数域上分裂且无重根。$t^m-1$ 与导数 $mt^{m-1}$ 互素，因而在复数域中无重根。作为矩阵空间的例子，设 $P$ 是次对角线上为 $1$ 的反序置换矩阵，$\Phi(X)=PX^{\mathsf T}P$，则 $\Phi^2=I$。写 $X=PY$，有 $\Phi(PY)=PY^{\mathsf T}$，故 $1$ 和 $-1$ 的特征子空间分别由 $Y^{\mathsf T}=Y$ 与 $Y^{\mathsf T}=-Y$ 给出，维数为 $n(n+1)/2$ 与 $n(n-1)/2$，合计 $n^2$。当 $n=1$ 时第二个空间为零，$-1$ 并非实际特征值。
\end{solve}

\subsection{构造完整的特征向量系}

\begin{example}
% 教材例6.9、例6.55；使用首行右移的循环约定。
设 $S$ 为满足 $Se_1=e_n$、$Se_j=e_{j-1}$（$2\le j\le n$）的循环移位矩阵，$C=f(S)$，其中 $f(t)=a_0+a_1t+\cdots+a_{n-1}t^{n-1}$。证明循环矩阵 $C$ 在复数域上可对角化。
\end{example}
\begin{solve}
记 $\omega_k=\exp(2\pi i k/n)$，$0\le k<n$，则
\[
v_k=(1,\omega_k,\ldots,\omega_k^{n-1})^{\mathsf T},\qquad
Sv_k=\omega_kv_k,\qquad Cv_k=f(\omega_k)v_k.
\]
矩阵 $P=(v_0,\ldots,v_{n-1})$ 的行列式为 Vandermonde 行列式 $\prod_{j<k}(\omega_k-\omega_j)\ne0$，故
\[
P^{-1}CP=\operatorname{diag}\bigl(f(\omega_0),\ldots,f(\omega_{n-1})\bigr).
\]
即使某些 $f(\omega_k)$ 相等，这 $n$ 个向量仍线性无关，故不影响对角化。
\end{solve}

\begin{example}
% 教材例6.57、例6.58。
若 $A=P\operatorname{diag}(\lambda_1,\ldots,\lambda_n)P^{-1}$，则矩阵空间上的 $X\mapsto AXA$ 与 $X\mapsto AX-XA$ 都可对角化。
\end{example}
\begin{solve}
取 $F_{ij}=PE_{ij}P^{-1}$。映射 $X\mapsto PXP^{-1}$ 为线性同构，所以这些 $F_{ij}$ 构成矩阵空间的一组基。直接计算得
\[
AF_{ij}A=\lambda_i\lambda_jF_{ij},\qquad
AF_{ij}-F_{ij}A=(\lambda_i-\lambda_j)F_{ij}.
\]
因此它们构成两个变换的公共特征向量基。这里通过完整的基保证了可对角化性，并非仅仅列出 $n^2$ 个可能的特征值。
\end{solve}

\subsection{单谱矩阵、交换矩阵与插值}

\begin{proposition}
\label{prop:xie-ex-eig-simple-spectrum-centralizer}
设 $A$ 有 $n$ 个不同的特征值 $\lambda_1,\ldots,\lambda_n\in\mathbb F$。若 $AB=BA$，则 $A,B$ 可同时在 $\mathbb F$ 上对角化，且存在唯一的次数小于 $n$ 的多项式 $f\in\mathbb F[t]$，使 $B=f(A)$。
\end{proposition}
\begin{proof}
每个 $E_{\lambda_i}(A)$ 都是一维且对 $B$ 不变，故 $B$ 在其上为数乘 $\mu_i$。取这些一维空间的非零向量为基，就同时对角化了 $A,B$。Lagrange 插值给出
\[
f(t)=\sum_{i=1}^n\mu_i\prod_{j\ne i}\frac{t-\lambda_j}{\lambda_i-\lambda_j},
\]
在该基下逐个对角元验证可得 $B=f(A)$。若另一个次数小于 $n$ 的多项式也适用，两者之差在 $n$ 个不同点为零，只能是零多项式。
\end{proof}

固定任一具有 $n$ 个不同复特征值的矩阵 $A$，上述插值还说明：任一可对角化的复矩阵 $B$ 都相似于某个 $f(A)$，$\deg f<n$。先将 $B$ 对角化，再把它的对角元插值到 $A$ 的不同特征值上即可；反向因为 $f(A)$ 可对角化而成立。取 $A=S$ 为循环移位矩阵，便得到每个可对角化复矩阵都相似于某个循环矩阵。

\begin{example}
% 教材例6.59、例6.65。
二阶实矩阵若行列式小于零，则可在实数域上对角化；三对角矩阵
\[
T_n=\begin{pmatrix}
a&b&&&\\c&a&b&&\\&c&a&\ddots&\\&&\ddots&\ddots&b\\&&&c&a
\end{pmatrix},\qquad bc\ne0,
\]
则可在复数域上对角化。
\end{example}
\begin{solve}
前者的特征多项式为 $t^2-(\operatorname{tr}A)t+\det A$，判别式 $(\operatorname{tr}A)^2-4\det A>0$，所以有两个不同实根。对后者选取 $r\ne0$ 使 $r^2=c/b$，令 $d=br$、$D=\operatorname{diag}(1,r,\ldots,r^{n-1})$，则 $D^{-1}T_nD$ 的上下次对角元同为 $d\ne0$。向量的第 $j$ 个分量取 $\sin(jk\pi/(n+1))$，利用正弦加法公式和两端的零边界可得其特征值为
\[
\lambda_k=a+2d\cos\frac{k\pi}{n+1},\qquad 1\le k\le n.
\]
这些值互异，因此对应向量线性无关，足以对角化。换取另一个平方根只会把这一组特征值反序排列。
\end{solve}

\subsection{分块矩阵、低秩乘积与反例}

\begin{proposition}
\label{prop:xie-ex-eig-block-diagonalization}
设 $M=\begin{pmatrix}A&C\\0&B\end{pmatrix}$。若 $M$ 可在 $\mathbb F$ 上对角化，则 $A,B$ 都可在 $\mathbb F$ 上对角化；若 $A,B$ 均可对角化且没有公共复特征值，则任意 $C$ 对应的 $M$ 都可在 $\mathbb F$ 上对角化。
\end{proposition}
\begin{proof}
第一结论对前一坐标子空间及其商空间应用引理~\ref{lem:xie-ex-eig-diagonal-restriction}。对第二结论，取 $A$ 的特征向量基 $u_i$ 和 $B$ 的特征向量基 $v_j$，分别对应 $\lambda_i,\mu_j$。由于 $A-\mu_jI$ 可逆，方程
\[
(A-\mu_jI)x_j=-Cv_j
\]
有唯一解。于是 $(u_i,0)^{\mathsf T}$ 和 $(x_j,v_j)^{\mathsf T}$ 构成 $M$ 的特征向量组；它们组成的矩阵为对角块均可逆的分块上三角矩阵，故全体线性无关。第~\ref{subsec:xie-ex-eig-sylvester} 小节还将用矩阵方程直接消去 $C$。
\end{proof}

\begin{proposition}
\label{prop:xie-ex-eig-ab-ba-diagonalization}
设 $A\in\mathbb F^{m\times n}$、$B\in\mathbb F^{n\times m}$，且 $BA$ 可逆，则 $m\ge n$，并且 $AB\sim\operatorname{diag}(BA,0_{m-n})$。因此 $AB$ 在 $\mathbb F$ 上可对角化当且仅当 $BA$ 如此。
\end{proposition}
\begin{proof}
$BA$ 可逆保证 $A$ 单射、$B$ 满射，故 $\dim\operatorname{Im}A=n$、$\dim\ker B=m-n$。若 $Ay\in\ker B$，则 $BAy=0$，从而 $y=0$，所以
\[
\mathbb F^m=\operatorname{Im}A\oplus\ker B.
\]
两分量都对 $AB$ 不变：第一分量上利用同构 $y\mapsto Ay$，有 $AB(Ay)=A(BAy)$；第二分量上 $AB$ 为零。选择相应基即得分块对角表示。
\end{proof}

\begin{example}
\label{ex:xie-ex-eig-rank-one-diagonalization}
% 教材例6.73、例6.74。
当 $n>1$ 且 $\operatorname{rank}R=1$ 时，$R$ 可对角化当且仅当 $\operatorname{tr}R\ne0$。命题~\ref{prop:xie-ex-eig-ab-ba-diagonalization} 中 $BA$ 可逆的假设不能直接删去。
\end{example}
\begin{solve}
写 $R=uv^{\mathsf T}$，则 $\operatorname{tr}R=v^{\mathsf T}u$。若该数非零，例~\ref{ex:xie-ex-eig-rank-one-spectrum} 给出的零化多项式 $t(t-\operatorname{tr}R)$ 无重根，故可对角化。若迹为零，则 $R^2=0$ 而 $R\ne0$；一个全部特征值为零的可对角化矩阵只能是零矩阵，因此矛盾。取 $A=(1,0)^{\mathsf T}$、$B=(0,1)$，便有 $BA=(0)$ 可对角化，$AB=\begin{pmatrix}0&1\\0&0\end{pmatrix}$ 却不可对角化。
\end{solve}

更一般地，若 $A$ 只有一个不同的特征值 $\lambda$，则 $A$ 可对角化当且仅当 $A=\lambda I_n$。有共同特征值的分块上三角矩阵，其非对角块可能阻碍对角化，例如 $\begin{pmatrix}1&1\\0&1\end{pmatrix}$；而没有共同特征值时，命题~\ref{prop:xie-ex-eig-block-diagonalization} 保证这种障碍可以消去。

\section{极小多项式与 Cayley--Hamilton 定理}
\label{sec:xie-ex-eig-minimal-polynomial}

\subsection{存在唯一性、整除性质与不变量}

\begin{theorem}
\label{thm:xie-ex-eig-minimal-properties}
任一 $n$ 阶矩阵 $A$ 的极小多项式 $m_A$ 存在且唯一，并且
\[
f(A)=0\quad\Longleftrightarrow\quad m_A(t)\mid f(t).
\]
若 $d=\deg m_A$，则代数 $\mathbb F[A]=\{f(A):f\in\mathbb F[t]\}$ 作为线性空间的维数为 $d$，一组基为 $I_n,A,\ldots,A^{d-1}$。
\end{theorem}
\begin{proof}
矩阵空间维数为 $n^2$，所以 $I_n,A,\ldots,A^{n^2}$ 线性相关，存在非零零化多项式。在其中取次数最低者并首一化。对任一 $f$ 作带余除法 $f=qm_A+r$，$\deg r<\deg m_A$。若 $f(A)=0$，则 $r(A)=0$，最低次数性质迫使 $r=0$；反向直接代入可得。两个最低次数的首一零化多项式因此相互整除，只能相同。带余除法还说明每个 $f(A)$ 均在所列 $d$ 个矩阵的张成空间内；若这 $d$ 个矩阵线性相关，就产生次数低于 $d$ 的非零零化多项式，故它们线性无关。
\end{proof}

相似矩阵有相同的极小多项式，$A$ 与 $A^{\mathsf T}$ 也有相同的极小多项式，因为两种变换都保持 $f(A)=0$。对分块对角矩阵则有
\begin{equation}
m_{\operatorname{diag}(A_1,\ldots,A_r)}
=\operatorname{lcm}(m_{A_1},\ldots,m_{A_r}),
\label{eq:xie-ex-eig-block-minimal}
\end{equation}
其中最小公倍式取首一形式：$f$ 零化全矩阵当且仅当零化每个分块，也就当且仅当被每个 $m_{A_i}$ 整除。这与特征多项式按块相乘的规律不同。

\begin{theorem}[Cayley--Hamilton 定理]
\label{thm:xie-ex-eig-cayley-hamilton}
矩阵适合自己的特征多项式，即 $\chi_A(A)=0$。
\end{theorem}
\begin{proof}
沿用 $\chi_A(t)=t^n+c_1t^{n-1}+\cdots+c_n$，写
\[
\operatorname{adj}(tI_n-A)=B_0t^{n-1}+B_1t^{n-2}+\cdots+B_{n-1}.
\]
由 $(tI_n-A)\operatorname{adj}(tI_n-A)=\chi_A(t)I_n$ 比较 $t$ 的系数，得到
\[
B_0=I_n,\qquad B_j=AB_{j-1}+c_jI_n\ (1\le j<n),\qquad
-AB_{n-1}=c_nI_n.
\]
逐次代入可得 $B_{n-1}=A^{n-1}+c_1A^{n-2}+\cdots+c_{n-1}I_n$，代入最后一个等式便是 $\chi_A(A)=0$。此证明通过系数递推处理矩阵乘法，未将带矩阵系数的恒等式未经说明地直接代入 $t=A$。
\end{proof}

由定理~\ref{thm:xie-ex-eig-minimal-properties}，立即有
\begin{equation}
m_A\mid\chi_A,\qquad 1\le\deg m_A\le n.
\label{eq:xie-ex-eig-minimal-divides-characteristic}
\end{equation}
若 $\lambda$ 为特征值，取非零特征向量便得到 $m_A(\lambda)=0$；反之 $m_A$ 的根也是 $\chi_A$ 的根。因此二者在复数域中的不同根完全相同，区别在于重数，并且 $\chi_A\mid m_A^n$。特别地，$n$ 个不同特征值时有 $m_A=\chi_A$；幂零矩阵的极小多项式为 $t^r$，其中 $r$ 是使 $A^r=0$ 的最小正整数。

\begin{proposition}
\label{prop:xie-ex-eig-minimal-diagonalization}
$A$ 在 $\mathbb F$ 上可对角化，当且仅当 $m_A$ 在 $\mathbb F$ 上分裂且无重根。此时
\[
m_A(t)=\prod_{\lambda\text{ 为不同特征值}}(t-\lambda).
\]
\end{proposition}
\begin{proof}
若可对角化，对角矩阵各纯量分块的极小多项式为 $t-\lambda$，由式~\eqref{eq:xie-ex-eig-block-minimal} 得到所列乘积。反向对 $m_A$ 应用定理~\ref{thm:xie-ex-eig-squarefree-annihilator}。
\end{proof}

当 $n>1$ 且 $\operatorname{rank}A=1$ 时，$m_A(t)=t^2-(\operatorname{tr}A)t$。因为 $A$ 适合此多项式而不能适合一次多项式：一次关系会使它成为纯量矩阵，而非零纯量矩阵的秩为 $n$。这也说明矩阵的极小多项式不一定不可约，不能照搬代数数极小多项式的不可约性。

\subsection{多项式可逆性、逆矩阵与伴随矩阵}

\begin{proposition}
\label{prop:xie-ex-eig-polynomial-inverse}
对 $g\in\mathbb F[t]$，下列条件等价：$g(A)$ 可逆；$\gcd(g,m_A)=1$；$\gcd(g,\chi_A)=1$；$g(\lambda)\ne0$ 对每个复特征值 $\lambda$ 成立。成立时 $g(A)^{-1}$ 可以表示为 $A$ 的次数小于 $\deg m_A$ 的多项式。
\end{proposition}
\begin{proof}
谱映射定理给出第一个条件与最后一个条件的等价性，而 $m_A,\chi_A$ 有相同的不同根，故两个互素条件也与之等价。由 Bézout 恒等式，存在 $u,v\in\mathbb F[t]$，使 $ug+vm_A=1$。代入 $A$ 得 $u(A)g(A)=I_n$，故 $g(A)^{-1}=u(A)$；再用 $m_A$ 对 $u$ 取余即可降低次数。
\end{proof}

取 $g(t)=t$ 可知 $A$ 可逆当且仅当 $m_A(0)\ne0$。用特征多项式也能直接写出
\begin{equation}
A^{-1}=-\frac1{c_n}\left(A^{n-1}+c_1A^{n-2}+\cdots+c_{n-1}I_n\right),
\qquad c_n\ne0.
\label{eq:xie-ex-eig-inverse-polynomial}
\end{equation}
无论是否可逆，伴随矩阵都满足
\begin{equation}
\operatorname{adj}A=(-1)^{n-1}
\left(A^{n-1}+c_1A^{n-2}+\cdots+c_{n-1}I_n\right).
\label{eq:xie-ex-eig-adjugate-polynomial}
\end{equation}
这是 Cayley--Hamilton 定理证明中 $\operatorname{adj}(tI_n-A)$ 的常数项：$\operatorname{adj}(-A)=(-1)^{n-1}\operatorname{adj}A$，而该常数项正是 $B_{n-1}$。因此这里无需对奇异矩阵使用并不存在的逆矩阵。

\subsection{Sylvester 矩阵方程及其应用}
\label{subsec:xie-ex-eig-sylvester}

\begin{theorem}
\label{thm:xie-ex-eig-sylvester}
设 $A\in\mathbb F^{m\times m}$、$B\in\mathbb F^{n\times n}$。对任意 $C\in\mathbb F^{m\times n}$，方程
\begin{equation}
AX-XB=C
\label{eq:xie-ex-eig-sylvester-equation}
\end{equation}
均有唯一的 $\mathbb F$ 上的解，当且仅当 $A,B$ 没有公共复特征值。
\end{theorem}
\begin{proof}
先设没有公共特征值，考虑齐次方程 $AX=XB$。对任意多项式 $f$，由幂次递推可得 $f(A)X=Xf(B)$。取 $f=\chi_A$，便有 $0=X\chi_A(B)$；两矩阵无公共特征值，故命题~\ref{prop:xie-ex-eig-polynomial-inverse} 保证 $\chi_A(B)$ 可逆，因此 $X=0$。映射 $\mathcal L\colon X\mapsto AX-XB$ 是有限维空间到自身的单射，因而也是满射，于是每个 $C$ 都有唯一解。

反之，若有公共特征值 $\lambda$，在复数域中取 $Au=\lambda u$、$B^{\mathsf T}v=\lambda v$，$u,v$ 均非零，则 $uv^{\mathsf T}\ne0$ 且 $Auv^{\mathsf T}-uv^{\mathsf T}B=0$。故 $\mathcal L$ 的复表示矩阵奇异。它的元素原本在 $\mathbb F$ 中，行列式为零也意味着在 $\mathbb F$ 上奇异，所以在原数域上同样不能对每个 $C$ 存在唯一解。
\end{proof}

存在公共特征值时，应对具体的右端 $C$ 另外判断相容性：方程可能无解；若有解，因齐次核非零，解便不唯一。定理给出的是对所有右端都存在唯一解的判据，并非对相交谱情形的完整相容性条件。

\begin{example}
% 教材例6.70的第二种证明。
若 $A,B$ 无公共复特征值，则任意分块上三角矩阵 $M=\begin{pmatrix}A&C\\0&B\end{pmatrix}$ 都相似于 $\operatorname{diag}(A,B)$。
\end{example}
\begin{solve}
解 $AX-XB=C$，取 $P=\begin{pmatrix}I_m&-X\\0&I_n\end{pmatrix}$，则
\[
P^{-1}MP
=\begin{pmatrix}A&-AX+C+XB\\0&B\end{pmatrix}
=\operatorname{diag}(A,B).
\]
这比命题~\ref{prop:xie-ex-eig-block-diagonalization} 的结论更强，因为消去非对角块本身不要求 $A,B$ 可对角化。
\end{solve}

\begin{example}
% 教材例6.89、例6.90、例6.92。
Sylvester 方程可用于证明矩阵平方根的唯一性、交换矩阵的分块结构及实对称解的存在唯一性。
\end{example}
\begin{solve}
若 $A^2=B^2$，且二者全部特征值均为正实数，令 $X=A-B$，则 $AX=X(-B)$。$A$ 与 $-B$ 无公共特征值，故 $X=0$，即 $A=B$，这一证明无需假设 $AB=BA$。若 $A=\operatorname{diag}(A_1,\ldots,A_r)$，各块的谱两两不交，且 $AB=BA$，按同样方式对 $B$ 分块，便有 $A_iB_{ij}=B_{ij}A_j$；当 $i\ne j$ 时只能 $B_{ij}=0$，所以 $B$ 也按同样方式分块对角。

再设 $A$ 为实矩阵，全部特征值的实部为正，$C$ 为实对称矩阵。$A^{\mathsf T}$ 与 $-A$ 的谱分别在左右半平面，所以 $A^{\mathsf T}X+XA=C$ 有唯一实矩阵解。转置后可知 $X^{\mathsf T}$ 也是解，由唯一性得到 $X=X^{\mathsf T}$。教材假设所有特征值为正实数；由上述谱分离论证可见，实部为正已经足够。
\end{solve}

\subsection{互素因式诱导的不变子空间直和}

\begin{theorem}
\label{thm:xie-ex-eig-coprime-decomposition}
设 $T\colon V\to V$ 为线性变换，$f_1,f_2\in\mathbb F[t]$ 为互素首一多项式，且 $f_1(T)f_2(T)=0$。令 $V_i=\ker f_i(T)$，则
\begin{equation}
V=V_1\oplus V_2,\qquad
V_1=\operatorname{Im}f_2(T),\qquad
V_2=\operatorname{Im}f_1(T),
\label{eq:xie-ex-eig-coprime-decomposition}
\end{equation}
各 $V_i$ 均为 $T$ 的不变子空间。若 $\chi_T=f_1f_2$，则 $\chi_{T|_{V_i}}=f_i$；若 $m_T=f_1f_2$，则 $m_{T|_{V_i}}=f_i$。零维分量的特征多项式与极小多项式约定为 $1$。
\end{theorem}
\begin{proof}
取 $u,v\in\mathbb F[t]$ 使 $uf_1+vf_2=1$。对任意 $x\in V$，有
\[
x=u(T)f_1(T)x+v(T)f_2(T)x.
\]
第一项属于 $V_2$，第二项属于 $V_1$，故二者之和为全空间。若 $x$ 同属两个核，上式两项均为零，所以交集为零。$f_i(T)$ 与 $T$ 可交换，因而各核不变。由 $f_1(T)f_2(T)=0$ 得到 $\operatorname{Im}f_2(T)\subseteq V_1$；而 $x\in V_1$ 时 $x=v(T)f_2(T)x$，又属于 $\operatorname{Im}f_2(T)$，从而相等，另一个等式同理。

记 $g_i=\chi_{T|_{V_i}}$，由直和表示可得 $g_1g_2=\chi_T$，而 $f_i$ 零化 $T|_{V_i}$，所以 $g_i$ 的每个复根只能是 $f_i$ 的根。若 $\chi_T=f_1f_2$，因 $f_1,f_2$ 互素，比较根及重数便有 $g_i=f_i$。对极小多项式，记 $d_i=m_{T|_{V_i}}$，则 $d_i\mid f_i$，且 $m_T=\operatorname{lcm}(d_1,d_2)=d_1d_2$。若 $m_T=f_1f_2$，再由整除性得到 $d_i=f_i$。
\end{proof}

若只要求 $f_1(T)f_2(T)=0$，不能断言限制变换的特征多项式就是 $f_i$，因为零化多项式可能带有多余因子。若没有全空间零化的假设，仍有 $\ker(f_1(T)f_2(T))=\ker f_1(T)\oplus\ker f_2(T)$，只需把上述证明限制在左侧的核中。复数域上将特征多项式分解为 $(t-\lambda_i)^{a_i}$ 后，逐次应用定理便得到全空间的广义特征子空间直和分解，这是后续 Jordan 标准型的起点。

\subsection{有限次幂生成的子空间与系数递推}

\begin{example}
% 教材例6.96。
设 $A\in\mathbb F^{n\times n}$、$B\in\mathbb F^{n\times p}$，令
\[
\mathcal U=\operatorname{Im}B+\operatorname{Im}(AB)+\cdots+
\operatorname{Im}(A^{n-1}B),\qquad \dim\mathcal U=r.
\]
则存在可逆矩阵 $P$，使
\[
P^{-1}AP=\begin{pmatrix}A_{11}&A_{12}\\0&A_{22}\end{pmatrix},\qquad
P^{-1}B=\begin{pmatrix}B_1\\0\end{pmatrix},
\]
其中 $A_{11}$ 为 $r$ 阶矩阵。
\end{example}
\begin{solve}
由 Cayley--Hamilton 定理，$A^nB$ 是 $B,AB,\ldots,A^{n-1}B$ 的线性组合，故 $A\mathcal U\subseteq\mathcal U$。取 $\mathcal U$ 的一组基并扩充为全空间的基，即得到所述形式。事实上 $\mathcal U$ 是包含 $\operatorname{Im}B$ 的最小 $A$ 不变子空间，因为任一这样的子空间都必须包含所有 $A^jB$ 的列空间。
\end{solve}

\begin{example}
% 教材例6.95；使用等价的核与不变子空间证明。
设 $C\in\mathbb C^{p\times n}$。矩阵
\[
\mathcal O=\begin{pmatrix}C\\CA\\\vdots\\CA^{n-1}\end{pmatrix}
\]
列满秩，当且仅当对每个 $\lambda\in\mathbb C$，$\begin{pmatrix}A-\lambda I_n\\C\end{pmatrix}$ 列满秩。此外，将 $A$ 换成 $A-\mu I_n$ 后得到的堆叠矩阵与 $\mathcal O$ 秩相同，$\mu\in\mathbb C$ 任意。
\end{example}
\begin{solve}
令 $K=\bigcap_{j=0}^{n-1}\ker(CA^j)=\ker\mathcal O$。Cayley--Hamilton 定理保证 $K$ 对 $A$ 不变。若 $K\ne0$，则 $A|_K$ 有特征向量 $x\ne0$，满足 $(A-\lambda I_n)x=0$、$Cx=0$，所以对应堆叠矩阵不满秩。反之，若存在这样的非零 $x$，则 $CA^jx=\lambda^jCx=0$，于是 $x\in K$。最后，由二项式展开，$C(A-\mu I_n)^j$ 是 $C,CA,\ldots,CA^j$ 的线性组合，且 $CA^j$ 的系数为 $1$；全部行块之间的变换是对角块为单位矩阵的下三角变换，因此可逆并保持秩。
\end{solve}

\begin{example}
% 教材例6.97；等价的 Faddeev--LeVerrier 递推。
特征多项式系数可由下述有限递推计算：置 $B_0=I_n$，对 $k=1,\ldots,n$ 依次令
\begin{equation}
c_k=-\frac1k\operatorname{tr}(AB_{k-1}),\qquad
B_k=AB_{k-1}+c_kI_n.
\label{eq:xie-ex-eig-faddeev-leverrier}
\end{equation}
输出的 $c_1,\ldots,c_n$ 正是 $\chi_A(t)$ 的系数，并有 $B_n=0$。
\end{example}
\begin{solve}
归纳得到 $B_k=A^k+c_1A^{k-1}+\cdots+c_kI_n$，故
\[
\operatorname{tr}(AB_{k-1})=p_k+c_1p_{k-1}+\cdots+c_{k-1}p_1.
\]
式~\eqref{eq:xie-ex-eig-newton-identities} 恰好说明递推算出的系数等于特征多项式系数，而 $B_n=\chi_A(A)=0$。这个递推给出精确代数关系；使用浮点数计算时，舍入误差可能累积，不能把理论上的终止等式当作数值计算中必然精确出现的零矩阵。
\end{solve}

\section{矩阵的 Kronecker 积}
\label{sec:xie-ex-eig-kronecker}

\subsection{定义、维度与基本运算}

\begin{definition}
设 $A=(a_{ij})\in\mathbb F^{m\times n}$、$B\in\mathbb F^{p\times q}$，其\emph{Kronecker 积（Kronecker product）}定义为 $mp\times nq$ 矩阵
\[
A\otimes B=\begin{pmatrix}
a_{11}B&\cdots&a_{1n}B\\
\vdots&&\vdots\\
a_{m1}B&\cdots&a_{mn}B
\end{pmatrix}.
\]
对应单个元素的索引为 $(A\otimes B)_{(i-1)p+\alpha,(j-1)q+\beta}=a_{ij}b_{\alpha\beta}$。
\end{definition}

这一运算对每个变量分别线性，满足结合律，并有
\begin{equation}
(A\otimes B)(C\otimes D)=(AC)\otimes(BD),\qquad
(A\otimes B)^{\mathsf T}=A^{\mathsf T}\otimes B^{\mathsf T},
\label{eq:xie-ex-eig-kronecker-rules}
\end{equation}
其中各普通矩阵乘积须有意义。第一式称为混合乘积公式：其第 $(i,j)$ 个分块为 $\sum_k a_{ik}c_{kj}BD=(AC)_{ij}BD$，从而得到结论。结合律则由两种括号顺序都产生元素 $a_{ij}b_{\alpha\beta}c_{uv}$ 验证。若 $A,B$ 分别是可逆方阵，则
\[
(A\otimes B)^{-1}=A^{-1}\otimes B^{-1},\qquad
I_m\otimes I_n=I_{mn}.
\]
Kronecker 积通常不可交换，不能将 $A\otimes B$ 与 $B\otimes A$ 直接视为相同矩阵；它们交换的是索引次序，须由相应置换联系。

\begin{proposition}
\label{prop:xie-ex-eig-kronecker-rank}
对任意矩形矩阵 $A,B$，有 $\operatorname{rank}(A\otimes B)=\operatorname{rank}A\operatorname{rank}B$。若 $A$ 为 $m$ 阶、$B$ 为 $n$ 阶，则
\[
\operatorname{tr}(A\otimes B)=\operatorname{tr}A\operatorname{tr}B,\qquad
\det(A\otimes B)=(\det A)^n(\det B)^m.
\]
\end{proposition}
\begin{proof}
将 $A,B$ 分别用可逆的左右因子化为 $\operatorname{diag}(I_r,0)$、$\operatorname{diag}(I_s,0)$ 形式的矩形秩标准形。由混合乘积公式及逆矩阵公式，它们的 Kronecker 积与 $A\otimes B$ 相抵，恰有 $rs$ 个非零单位对角元，故秩为 $rs$。迹等式由对角元素相乘后求和得到。对行列式，先在复数域中分别相似上三角化 $A,B$，再由混合乘积公式同时变换其 Kronecker 积；对角元为所有 $\lambda_i\mu_j$，乘积就是 $(\prod_i\lambda_i)^n(\prod_j\mu_j)^m$。
\end{proof}

因此，在各维度均为正时，$A\otimes B$ 列满秩当且仅当 $A,B$ 均列满秩，行满秩也同理。对方阵，$A\otimes B$ 可逆当且仅当两个因子都可逆。

\subsection{谱与矩阵空间上线性变换的表示}

\begin{theorem}
\label{thm:xie-ex-eig-kronecker-spectrum}
设 $A\in\mathbb C^{m\times m}$、$B\in\mathbb C^{n\times n}$ 的特征值分别为 $\lambda_i$、$\mu_j$，则 $A\otimes B$ 的全部特征值为 $\lambda_i\mu_j$；$A\otimes I_n-I_m\otimes B$ 的全部特征值为 $\lambda_i-\mu_j$，其中 $1\le i\le m$、$1\le j\le n$，均计重。
\end{theorem}
\begin{proof}
分别取 $P,Q$ 上三角化 $A,B$。利用 $P\otimes Q$ 作相似变换，两种矩阵分别变为两个上三角矩阵的 Kronecker 积，以及相应上三角矩阵与单位矩阵的 Kronecker 积之差，所得对角元恰为所列的乘积和差。这里无需假设 $A,B$ 可对角化。
\end{proof}

为把矩阵方程写成普通线性方程，约定 $\operatorname{vec}(X)$ 将 $X\in\mathbb F^{m\times n}$ 的各列从左到右堆叠。对应的矩阵空间基为 $E_{11},E_{21},\ldots,E_{m1},E_{12},\ldots,E_{mn}$。按每一列展开乘积得到
\begin{equation}
\operatorname{vec}(AXB)=(B^{\mathsf T}\otimes A)\operatorname{vec}(X),
\label{eq:xie-ex-eig-vectorization-product}
\end{equation}
从而
\begin{equation}
\operatorname{vec}(AX-XB)
=(I_n\otimes A-B^{\mathsf T}\otimes I_m)\operatorname{vec}(X).
\label{eq:xie-ex-eig-vectorization-sylvester}
\end{equation}
教材按行排列基础矩阵时得到的表示分别是 $A\otimes B^{\mathsf T}$ 和 $A\otimes I_n-I_m\otimes B^{\mathsf T}$。两种表示来自不同的基次序，通过置换相似联系；使用时必须先固定向量化约定，不能将其中的转置和因子顺序混用。

\begin{proposition}
\label{prop:xie-ex-eig-matrix-space-spectrum}
在 $m\times n$ 矩阵空间上，$\Phi(X)=AXB$ 的特征值为全部 $\lambda_i\mu_j$，$\Psi(X)=AX-XB$ 的特征值为全部 $\lambda_i-\mu_j$。此外，$\operatorname{rank}\Phi=\operatorname{rank}A\operatorname{rank}B$，$\Phi$ 为自同构当且仅当 $A,B$ 均可逆；$\Psi$ 为自同构当且仅当两者的谱不交。
\end{proposition}
\begin{proof}
利用式~\eqref{eq:xie-ex-eig-vectorization-product}、式~\eqref{eq:xie-ex-eig-vectorization-sylvester}，分别应用命题~\ref{prop:xie-ex-eig-kronecker-rank} 与定理~\ref{thm:xie-ex-eig-kronecker-spectrum}，再注意 $B^{\mathsf T}$ 与 $B$ 的特征多项式相同。最后一条也给出了 Sylvester 方程定理的另一种证明。
\end{proof}

若 $A=P\Lambda P^{-1}$、$B=QMQ^{-1}$ 均可对角化，则 $PE_{ij}Q^{-1}$ 构成矩阵空间的一组基，且
\[
A(PE_{ij}Q^{-1})B=\lambda_i\mu_jPE_{ij}Q^{-1},\qquad
A(PE_{ij}Q^{-1})-(PE_{ij}Q^{-1})B
=(\lambda_i-\mu_j)PE_{ij}Q^{-1}.
\]
故这两个矩阵空间变换都可对角化。与单纯列出特征值不同，这里同时构造了完整的特征向量基。

\subsection{幂零性与 Kronecker 和}

\begin{proposition}
在复数域上，变换 $X\mapsto AXB$ 幂零，当且仅当 $A,B$ 至少有一个幂零。若 $A,B$ 均幂零，则 $X\mapsto AX-XB$ 也幂零。特别地，$X\mapsto AX-XA$ 幂零，当且仅当存在 $\lambda\in\mathbb C$ 使 $A-\lambda I_n$ 幂零。
\end{proposition}
\begin{proof}
若 $A^r=0$ 或 $B^r=0$，则 $\Phi^r(X)=A^rXB^r=0$。反之，若两者都不幂零，则各有一个非零特征值，乘积给出 $\Phi$ 的非零特征值，所以 $\Phi$ 不幂零。对差型变换，$A,B$ 均幂零时其全部特征值为零，故 $\Psi$ 幂零；也可以直接给出幂次界：若 $A^r=0$、$B^s=0$，左右乘法相互可交换，二项式展开表明 $\Psi^{r+s-1}=0$，因为每项中至少含 $A^r$ 或 $B^s$。最后，$AX-XA$ 的变换幂零当且仅当所有差 $\lambda_i-\lambda_j$ 都为零，即 $A$ 只有一个不同的特征值 $\lambda$，再由命题~\ref{prop:xie-ex-eig-nilpotent-spectrum} 即得结论。
\end{proof}

\begin{example}
% 教材例6.108。
令 $A$ 为 $n$ 阶矩阵，$g(t)=\det(tI_n+A)$，则 $K=A\otimes I_n+I_n\otimes A$ 可逆当且仅当 $g(A)$ 可逆。
\end{example}
\begin{solve}
将差谱公式中的第二个矩阵取为 $-A$，可知 $K$ 的特征值为 $\lambda_i+\lambda_j$。因此
\[
\det K=\prod_{i,j}(\lambda_i+\lambda_j)
=\prod_i g(\lambda_i)=\det g(A).
\]
行列式是否为零便同时决定两者的可逆性。这里 $K$ 称为相应的\emph{Kronecker 和（Kronecker sum）}，不能将它与矩阵乘积 $A\otimes A$ 混淆。
\end{solve}

\section{基础训练与方法整合}
\label{sec:xie-ex-eig-eigen-exercises}

教材基础训练可归结为几条反复使用的路线：由 $Av=\lambda v$ 提取参数约束，由矩阵适合的多项式限制特征值，由代数重数与核的维数判定可对角化，由不变子空间和互素多项式构造分解，再由谱分离判断矩阵方程的唯一可解性。以下选取能连接不同方法的训练题，给出完整推导。

\subsection{基本概念与实矩阵的特征值}

\begin{example}
% 教材基础训练：单选题5、7及填空题14。
证明：若每个非零向量都是 $A$ 的特征向量，则 $A$ 是纯量矩阵；若 $\operatorname{rank}A=r<n$，则零特征值的代数重数至少为 $n-r$。
\end{example}
\begin{solve}
由 $Ae_i=\lambda_i e_i$，$A$ 首先是对角矩阵。对任意 $i\ne j$，$e_i+e_j$ 也是特征向量，所以 $\lambda_ie_i+\lambda_je_j=\mu(e_i+e_j)$，比较系数得到 $\lambda_i=\lambda_j$，故 $A=\lambda I_n$；一阶情形直接成立。第二问由 $\dim\ker A=n-r$ 以及几何重数不超过代数重数得到。注意，同一特征值的多个特征向量既可能相关，也可能无关，例如 $I_n$ 的全部标准基向量都属于特征值 $1$。
\end{solve}

\begin{example}
% 教材基础训练：解答题3。
设实矩阵 $A$ 满足 $x^{\mathsf T}Ax>0$ 对每个非零实列向量 $x$ 成立，教材称其为亚正定矩阵。证明 $A$ 的全部复特征值实部均为正。
\end{example}
\begin{solve}
令 $H=(A+A^{\mathsf T})/2$，则 $H$ 为实对称矩阵，且 $x^{\mathsf T}Hx=x^{\mathsf T}Ax>0$。若 $Az=\lambda z$、$z=x+iy\ne0$，则
\[
\operatorname{Re}(z^{\mathsf H}Az)
=z^{\mathsf H}Hz=x^{\mathsf T}Hx+y^{\mathsf T}Hy>0.
\]
另一方面，$z^{\mathsf H}Az=\lambda z^{\mathsf H}z$，所以
\[
\operatorname{Re}\lambda
=\frac{x^{\mathsf T}Hx+y^{\mathsf T}Hy}{z^{\mathsf H}z}>0.
\]
条件只约束对称部分，不要求 $A$ 本身对称，故结论是实部为正，而不能加强为所有特征值均为正实数。
\end{solve}

\begin{example}
% 教材基础训练：解答题6。
设实矩阵 $A$ 的主对角元均为 $1$，特征值均为非负实数，证明 $\det A\le1$。
\end{example}
\begin{solve}
由 $\sum_i\lambda_i=\operatorname{tr}A=n$ 及算术平均不小于几何平均，得到
\[
\det A=\prod_i\lambda_i\le\left(\frac1n\sum_i\lambda_i\right)^n=1.
\]
等号当且仅当全部特征值为 $1$，但不必有 $A=I_n$，例如非平凡的上三角 Jordan 块也可满足这些条件。若另加 $A$ 可对角化，才可由单一特征值得到 $A=I_n$。
\end{solve}

\subsection{互素关系、统一逆多项式与循环向量}

\begin{example}
% 教材基础训练：解答题11；明确复数域，避免把复对角化误写成原数域对角化。
设 $A\in\mathbb C^{n\times n}$，$f,g\in\mathbb C[t]$ 互素且各自无重根。若 $\operatorname{rank}f(A)+\operatorname{rank}g(A)=n$，证明 $A$ 可对角化。
\end{example}
\begin{solve}
Bézout 恒等式给出 $\ker f(A)\cap\ker g(A)=\{0\}$。由秩条件，两核的维数之和为 $n$，所以全空间是它们的直和。对第一分量，$f(A)$ 为零，故 $f(A)g(A)$ 也为零；对第二分量，$g(A)$ 为零，故乘积同样为零。于是 $A$ 适合 $fg$。由于 $f,g$ 互素且各自无重根，$fg$ 也无重根，再用定理~\ref{thm:xie-ex-eig-squarefree-annihilator}。若把条件写在一般数域 $\mathbb F$ 上，这只保证复对角化；要得到 $\mathbb F$ 上的对角化，还需保证实际特征值属于 $\mathbb F$。
\end{solve}

\begin{example}
% 教材基础训练：解答题13。
设 $A_1,\ldots,A_r$ 是 $\mathbb F$ 上的方阵，$g\in\mathbb F[t]$，且每个 $g(A_i)$ 都可逆。证明存在同一个多项式 $h\in\mathbb F[t]$，使 $g(A_i)^{-1}=h(A_i)$ 对所有 $i$ 成立。
\end{example}
\begin{solve}
令 $m=\operatorname{lcm}(m_{A_1},\ldots,m_{A_r})$。由命题~\ref{prop:xie-ex-eig-polynomial-inverse}，$g$ 与每个 $m_{A_i}$ 互素，从而与 $m$ 互素。取 $h,u$ 使 $hg+um=1$，依次代入 $A_i$，便得到 $h(A_i)g(A_i)=I$。再将 $h$ 对 $m$ 取余，可令 $\deg h<\deg m$。用同一个 Bézout 恒等式即可统一处理所有矩阵，无需分别寻找不同的逆多项式。
\end{solve}

\begin{example}
% 教材基础训练：解答题15。
若 $v\in\mathbb F^n$ 使 $K=(A^{n-1}v,A^{n-2}v,\ldots,Av,v)$ 可逆，求 $\chi_A(t)=t^n+c_1t^{n-1}+\cdots+c_n$ 的系数。
\end{example}
\begin{solve}
由 Cayley--Hamilton 定理作用于 $v$ 得
\[
A^nv+K(c_1,\ldots,c_n)^{\mathsf T}=0,
\]
因此 $(c_1,\ldots,c_n)^{\mathsf T}=-K^{-1}A^nv$。这样的 $v$ 称为\emph{循环向量（cyclic vector）}，因为 $v,Av,\ldots,A^{n-1}v$ 构成全空间的基。若存在次数小于 $n$ 的非零多项式零化 $A$，作用于 $v$ 就会与这组基的线性无关性矛盾，所以 $\deg m_A=n$，再由式~\eqref{eq:xie-ex-eig-minimal-divides-characteristic} 得 $m_A=\chi_A$。
\end{solve}

\subsection{有限乘法结构与非交换矩阵的同时化}

\begin{example}
% 教材基础训练：解答题7。
设 $\mathcal S=\{A_1,\ldots,A_r\}$ 为有限个互异可逆矩阵组成的非空集合，且关于乘法封闭。证明 $\sum_iA_i=0$ 当且仅当 $\sum_i\operatorname{tr}A_i=0$。
\end{example}
\begin{solve}
必要性由迹的线性性得到。令 $M=\sum_iA_i$。右乘任一 $A_j$ 在有限集合 $\mathcal S$ 上是单射，故也是置换，所以 $MA_j=M$。求和得到 $M^2=rM$，于是 $M$ 的特征值只能是 $0,r$。若 $\operatorname{tr}M=0$，因 $r>0$，特征值 $r$ 的重数只能为零，故 $M$ 幂零。设 $M^k=0$，由 $M^k=r^{k-1}M$ 得 $M=0$，充分性成立。
\end{solve}

\begin{example}
% 教材基础训练：解答题8。
设 $A,B$ 为复矩阵，且 $AB-BA=aA+bB$，其中 $a,b\in\mathbb C$。证明 $A,B$ 可同时上三角化。
\end{example}
\begin{solve}
若 $a=b=0$，由定理~\ref{thm:xie-ex-eig-simultaneous-triangular} 得到结论。其余情形，必要时交换 $A,B$，可设 $a\ne0$。令 $D=aA+bB$、$E=B/a$，则 $DE-ED=D$，而 $A,B$ 均为 $D,E$ 的线性组合，所以只需同时上三角化 $D,E$。

将例~\ref{ex:xie-ex-eig-commutator-nilpotent} 用于交换子 $D=DE-ED$，因 $D$ 与自身可交换，得到 $D$ 幂零。故 $W=\ker D\ne0$；由 $DE=ED+D$ 可知 $W$ 对 $E$ 不变。取 $E|_W$ 的特征向量，它就是 $D,E$ 的公共特征向量。在对应一维不变子空间的商空间上，仍有 $\overline D\,\overline E-\overline E\,\overline D=\overline D$。对维数归纳，得到共同上三角基，从而也同时上三角化 $A,B$。
\end{solve}

\chapter{相似标准型}
\label{chap:xie-ex-similarity}

% 依据《高等代数习题（谢启鸿）》第7章整理。
% 来源：高等代数/textbooks/高等代数习题（谢启鸿）/高代习题-7.pdf，PDF第1--81页，书页354--434。
本章研究同一个线性变换在不同基下的矩阵如何分类。相抵允许独立地改变定义域和陪域的基，其全系不变量是秩；相似要求两边使用同一组基，因而必须保留更细的信息。基本路线是把数字矩阵 $A$ 转化为特征矩阵 $\lambda I-A$，通过多项式矩阵的相抵得到行列式因子、不变因子和初等因子，再分别构造有理标准型与 Jordan 标准型。循环子空间和广义特征向量说明这些代数不变量的几何意义，矩阵函数及矩阵方程则展示标准型的用途。以下 $\mathbb K$ 表示数域，即 $\mathbb C$ 的子域；除另有说明外，$V$ 为 $\mathbb K$ 上的 $n$ 维线性空间，$n\geq1$，$T$ 为其线性变换，$A$ 为 $T$ 在某组基下的矩阵。记 $\chi_A(\lambda)=\det(\lambda I_n-A)$，$m_A(\lambda)$ 为 $A$ 的首一极小多项式。

\section{基本概念}
\label{sec:xie-ex-sim-basic}

\subsection{多项式矩阵与 Smith 标准形}
\label{subsec:xie-ex-sim-smith}

\begin{definition}
元素属于多项式环 $\mathbb K[\lambda]$ 的矩阵称为\emph{多项式矩阵（polynomial matrix）}或 $\lambda$-矩阵。其初等行、列变换包括交换两行或两列，以非零常数乘一行或一列，以及把一行或一列的多项式倍加到另一行或另一列。能够通过有限次这些变换互相转化的同型矩阵称为\emph{相抵（equivalent）}。
\end{definition}

这些变换都在 $\mathbb K[\lambda]$ 内可逆；其中第二类只能乘非零常数，因为多项式环的可逆元恰是非零常数。例如，把一行乘以 $\lambda$ 一般不可逆，不能作为求标准形时的合法变换。方阵 $U(\lambda)$ 在多项式环上可逆，当且仅当 $\det U(\lambda)$ 是非零常数：必要性来自行列式的乘法公式，充分性来自伴随矩阵公式。这样的矩阵又称\emph{幺模矩阵（unimodular matrix）}。因此 $M,N$ 相抵可表为 $N=UMW$，其中 $U,W$ 为相应阶数的幺模矩阵；它们在每个标量 $\lambda_0$ 处的取值仍是可逆数字矩阵。

\begin{theorem}[Smith 标准形]
\label{thm:xie-ex-sim-smith}
设 $M(\lambda)$ 为 $a\times b$ 多项式矩阵，$r$ 为它在有理函数域 $\mathbb K(\lambda)$ 上的秩，则 $M$ 相抵于
\[
 \operatorname{diag}(d_1,\ldots,d_r,0,\ldots,0),
 \qquad d_1\mid d_2\mid\cdots\mid d_r,
\]
其中 $d_i$ 是非零首一多项式；这一对角形状按 $a\times b$ 理解，且 $d_i$ 由 $M$ 唯一确定，称为其\emph{不变因子（invariant factors）}。
\end{theorem}

存在性的构造思路是反复使用多项式带余除法。先把非零元素移到左上角，用行、列相减降低同一行、列中非零余式的次数，直到左上角元素能整除这一行和列的元素并将它们消去；若它不能整除右下子矩阵的某个元素，可把该元素引入首行，再作除法，使主元次数继续下降。次数不能无限下降，故最终得到能整除其余元素的首个主元，再对右下子矩阵递归处理。这里只说明构造的关键步骤；唯一性由下面的行列式因子公式直接给出。

\begin{definition}
令 $D_0=1$。对 $1\leq k\leq\min(a,b)$，若 $M$ 有非零 $k$ 阶子式，则所有 $k$ 阶子式的首一最大公因式称为 $k$ 阶\emph{行列式因子（determinantal divisor）}，记为 $D_k$；若这些子式全为零，则约定 $D_k=0$。
\end{definition}

\begin{proposition}
\label{prop:xie-ex-sim-divisors}
行列式因子在多项式矩阵的相抵变换下不变，且
\begin{equation}
 D_k=d_1\cdots d_k\quad(1\leq k\leq r),
 \qquad d_k=\frac{D_k}{D_{k-1}}.
 \label{eq:xie-ex-sim-divisors}
\end{equation}
\end{proposition}
\begin{proof}
由 Cauchy--Binet 公式，$UMW$ 的每个 $k$ 阶子式都是 $M$ 的 $k$ 阶子式的多项式线性组合，所以 $D_k(M)$ 整除这些新子式。对逆变换使用相同论证，即得两个首一最大公因式相等。在 Smith 形中，每个非零 $k$ 阶子式都是若干 $d_i$ 的乘积，由整除链知它是 $d_1\cdots d_k$ 的倍式，而前 $k$ 个对角元本身给出这个子式。因此其首一最大公因式恰为该乘积，逐项相除即得不变因子及其唯一性。
\end{proof}

\subsection{不变因子与有理标准型}

数字矩阵 $A$ 的行列式因子和不变因子，均指其特征矩阵 $\lambda I_n-A$ 的相应因子。由于该特征矩阵的行列式是 $n$ 次首一多项式，它有恰好 $n$ 个非零不变因子。把其中的 $1$ 省去后，记非常数不变因子为
\[
 q_1\mid q_2\mid\cdots\mid q_s,
 \qquad \sum_{i=1}^s\deg q_i=n.
\]
因而 $\lambda I_n-A$ 的 Smith 形是 $\operatorname{diag}(1,\ldots,1,q_1,\ldots,q_s)$。这里常数因子 $1$ 可以省写，但在使用第几个不变因子时应恢复完整的 $n$ 项。

\begin{definition}
对 $d$ 次首一多项式 $p(\lambda)=\lambda^d+a_{d-1}\lambda^{d-1}+\cdots+a_1\lambda+a_0$，定义其\emph{友阵（companion matrix）}
\[
 C(p)=\begin{pmatrix}
 0&0&\cdots&0&-a_0\\
 1&0&\cdots&0&-a_1\\
 0&1&\cdots&0&-a_2\\
 \vdots&\vdots&\ddots&\vdots&\vdots\\
 0&0&\cdots&1&-a_{d-1}
 \end{pmatrix},
\]
并记 Frobenius 块为 $F(p)=C(p)^{\mathsf T}$；当 $d=1$ 时，两者均为 $(-a_0)$。
\end{definition}

\begin{theorem}[有理标准型]
\label{thm:xie-ex-sim-rational}
矩阵 $A$ 在 $\mathbb K$ 上相似于
\[
 \operatorname{diag}\bigl(C(q_1),\ldots,C(q_s)\bigr),
\]
也相似于将每个 $C(q_i)$ 换成 $F(q_i)$ 的矩阵。它们称为\emph{有理标准型（rational canonical form）}，其中
\begin{equation}
 \chi_A=q_1\cdots q_s,
 \qquad m_A=q_s.
 \label{eq:xie-ex-sim-char-min}
\end{equation}
这里“有理”表示构造无需离开给定数域，并不要求 $\mathbb K=\mathbb Q$。
\end{theorem}

有理标准型的存在性来自特征矩阵的 Smith 形及相似判据，见第~\ref{sec:xie-ex-sim-invariants} 节；它与循环子空间的对应见第~\ref{sec:xie-ex-sim-cyclic} 节。每个友阵的特征多项式、极小多项式都是相应的 $q_i$；分块对角矩阵的特征多项式为各块的乘积，极小多项式为各块极小多项式的首一最小公倍式，故整除链给出式~\eqref{eq:xie-ex-sim-char-min}。

\subsection{初等因子与 Jordan 标准型}

\begin{definition}
把每个非常数不变因子在 $\mathbb K[\lambda]$ 中分解为互异首一不可约多项式的幂，其所有\emph{非平凡不可约幂因子}组成的多重集合称为 $A$ 的\emph{初等因子组（elementary divisors）}。例如，$q=(\lambda-1)^2(\lambda+1)^3$ 贡献两个初等因子 $(\lambda-1)^2$、$(\lambda+1)^3$；同一因子若在不同不变因子中出现，必须保留其重复次数。
\end{definition}

从初等因子恢复不变因子时，应对每个不可约多项式分别把正指数按升序排列，在左端补零，使各列等长，再逐列相乘；这保证所得多项式构成整除链。反过来，只需将每个不变因子作不可约分解。\emph[caution]{有理标准型的块数等于非常数不变因子的个数，Jordan 块数等于复数域上初等因子的个数，两者一般不同。}

\begin{theorem}[Jordan 标准型]
\label{thm:xie-ex-sim-jordan}
若 $A\in M_n(\mathbb C)$，则 $A$ 相似于若干 Jordan 块的直和，其中
\[
 J_r(\mu)=\mu I_r+N_r,
 \qquad
 N_r=\begin{pmatrix}
 0&1&&0\\
 &0&\ddots&\\
 &&\ddots&1\\
 0&&&0
 \end{pmatrix},
 \qquad N_r^r=0,
 \qquad N_1=(0).
\]
每个初等因子 $(\lambda-\mu)^r$ 对应一个 $J_r(\mu)$，因此 Jordan 标准型除各块的排列顺序外唯一。属于 $\mu$ 的各块阶数之和是其代数重数，块数是其几何重数，最大块阶数是 $\lambda-\mu$ 在 $m_A$ 中的指数。
\end{theorem}

Jordan 型要求特征多项式在当前数域上分裂；一般数域总存在有理标准型，却未必能在该数域上写出通常的 Jordan 型。第~\ref{sec:xie-ex-sim-general-field} 节将用不可约多项式对应的广义 Jordan 块处理这一差别。

\subsection{矩阵序列与矩阵幂级数}

矩阵序列 $A_k$ 收敛于 $B$，是指每个矩阵元分别收敛。对复幂级数 $f(z)=\sum_{j=0}^{\infty}c_jz^j$，若部分和矩阵 $\sum_{j=0}^{k}c_jA^j$ 收敛，则其极限记为 $f(A)$。设标量幂级数的收敛半径为 $R$，$\rho(A)=\max\{|\mu|:\mu\text{ 是 }A\text{ 的特征值}\}$ 为\emph{谱半径（spectral radius）}，则 $\rho(A)<R$ 保证 $f(A)$ 收敛，$\rho(A)>R$ 时不收敛；边界 $\rho(A)=R$ 需要另行判断。相似变换与分块对角运算都可逐项作用于部分和，并在收敛后保留下来。具体计算公式及收敛理由将在第~\ref{sec:xie-ex-sim-functions} 节展开。

\section{矩阵相似的全系不变量}
\label{sec:xie-ex-sim-invariants}

\begin{theorem}[相似的四种等价判据]
\label{thm:xie-ex-sim-criteria}
对同一数域 $\mathbb K$ 上的同阶矩阵 $A,B$，下列条件等价：$A$ 与 $B$ 在 $\mathbb K$ 上相似；$\lambda I-A$ 与 $\lambda I-B$ 在 $\mathbb K[\lambda]$ 上相抵；两者有相同的行列式因子组；两者有相同的不变因子组；两者有相同的初等因子组。
\end{theorem}

相似推出特征矩阵相抵，可直接写成 $\lambda I-B=P^{-1}(\lambda I-A)P$；反方向的常数过渡矩阵提取方法见命题~\ref{prop:xie-ex-sim-polynomial-transition}，其证明也适用于 $B$ 代替其中的 $J$。行列式因子与不变因子的等价由式~\eqref{eq:xie-ex-sim-divisors} 给出，不变因子与初等因子的等价由不可约分解及补零重组给出。这些判据是完整分类；只比较迹、行列式、秩、特征多项式或极小多项式，一般只能得到必要条件。

\begin{example}
\label{ex:xie-ex-sim-char-min-insufficient}
比较 $A=\operatorname{diag}(J_3(0),J_3(0))$ 与 $B=\operatorname{diag}(J_3(0),J_2(0),J_1(0))$，判断相同的特征多项式和极小多项式是否足以保证相似。
\end{example}
\begin{solve}
两者都有 $\chi(\lambda)=\lambda^6$、$m(\lambda)=\lambda^3$，但 $\operatorname{rank}A=4$、$\operatorname{rank}B=3$，故不相似。特征多项式记录各特征值的总维数，极小多项式记录最大块，二者没有完整记录其余块的分配。
\end{solve}

\begin{proposition}
\label{prop:xie-ex-sim-diagonal-gcd}
对首一非零多项式 $f,g$，令 $d=\gcd(f,g)$、$h=\operatorname{lcm}(f,g)$，则 $\operatorname{diag}(f,g)$ 相抵于 $\operatorname{diag}(d,h)$。特别地，将特征矩阵相抵化为任一非零对角矩阵后，直接收集各对角元的不可约幂因子，就得到初等因子组，无须先整理出整除链。
\end{proposition}
\begin{proof}
前一矩阵的一阶行列式因子为 $d$，二阶行列式因子为 $fg$，故其不变因子为 $d,fg/d=d,h$。对每个不可约多项式 $p$，这一操作将两个指数 $a,b$ 换成 $\min(a,b),\max(a,b)$，保留正指数的多重集合。依次作这样的操作直至得到整除链，即证第二个结论。
\end{proof}

\begin{example}
证明 $A$ 与 $A^{\mathsf T}$ 相似，并说明分块对角矩阵的块顺序为何不影响相似类。
\end{example}
\begin{solve}
$\lambda I-A^{\mathsf T}=(\lambda I-A)^{\mathsf T}$，转置把每个子式变成对应的转置子式，行列式值不变，故两者行列式因子相同。对 $\operatorname{diag}(A_1,\ldots,A_k)$，可分别将各块的特征矩阵对角化，再用命题~\ref{prop:xie-ex-sim-diagonal-gcd} 收集初等因子；所得多重集合是各块初等因子组的并，保留重复次数。交换块顺序不改变它，也可直接用置换基向量构造相似变换。
\end{solve}

\begin{proposition}[相似性在数域扩张下不变]
\label{prop:xie-ex-sim-field-extension}
若 $\mathbb K\subseteq\mathbb L\subseteq\mathbb C$，且 $A,B\in M_n(\mathbb K)$，则它们在 $\mathbb L$ 上相似当且仅当在 $\mathbb K$ 上相似。
\end{proposition}
\begin{proof}
系数扩张不会改变一组 $\mathbb K[\lambda]$ 多项式的首一最大公因式：原最大公因式仍整除各式，而它在 $\mathbb K[\lambda]$ 中可表为这些多项式的线性组合，故扩张后的公因式也整除它。因此行列式因子及不变因子不变，再用定理~\ref{thm:xie-ex-sim-criteria} 即得结论。初等因子的不可约分解则可能随数域改变，不能将它与不变因子混为一谈。
\end{proof}

\begin{example}
设 $A,C$ 为可逆同阶矩阵，证明存在可逆常数矩阵 $P,Q$ 使 $A=PCQ$、$B=PDQ$，当且仅当 $\lambda A-B$ 与 $\lambda C-D$ 相抵。
\end{example}
\begin{solve}
必要性由 $\lambda A-B=P(\lambda C-D)Q$ 给出。反之，分别左乘 $A^{-1}$、$C^{-1}$，可知 $\lambda I-A^{-1}B$ 与 $\lambda I-C^{-1}D$ 相抵，故存在可逆 $Q$ 使 $A^{-1}B=Q^{-1}C^{-1}DQ$。令 $P=AQ^{-1}C^{-1}$ 即得所求。$A,C$ 可逆使矩阵束能化为特征矩阵，是这一推导的必要前提。
\end{solve}

\section{有理标准型的几何与应用}
\label{sec:xie-ex-sim-cyclic}

\begin{definition}
对 $0\ne v\in V$，称
\[
 \mathcal C(T,v)=\operatorname{span}_{\mathbb K}\{v,Tv,T^2v,\ldots\}
\]
为 $v$ 生成的\emph{循环子空间（cyclic subspace）}，$v$ 为其循环向量。它是包含 $v$ 的最小 $T$-不变子空间；若它等于 $V$，则称 $V$ 是关于 $T$ 的循环空间。
\end{definition}

\begin{proposition}
\label{prop:xie-ex-sim-cyclic-basis}
若 $r=\dim\mathcal C(T,v)$，则 $v,Tv,\ldots,T^{r-1}v$ 构成该空间的一组基，$T$ 的限制在此基下的矩阵是某个 $r$ 次首一多项式的友阵。反之，表示矩阵为友阵的变换具有循环向量。
\end{proposition}
\begin{proof}
设首次出现依赖的幂次为 $r_0$，则有唯一关系 $T^{r_0}v=-a_0v-\cdots-a_{r_0-1}T^{r_0-1}v$。反复作用 $T$ 后，更高幂也属于前 $r_0$ 项的张成空间，故 $r_0=r$；各基向量的像给出 $C(p)$，其中 $p(\lambda)=\lambda^r+a_{r-1}\lambda^{r-1}+\cdots+a_0$。反过来，友阵满足 $C(p)^je_1=e_{j+1}$，$0\leq j<r$，故 $e_1$ 是循环向量。
\end{proof}

\begin{theorem}[循环分解与循环空间判据]
\label{thm:xie-ex-sim-cyclic}
若 $T$ 的非常数不变因子为 $q_1\mid\cdots\mid q_s$，则存在直和分解
\[
 V=\mathcal C(T,v_1)\oplus\cdots\oplus\mathcal C(T,v_s),
 \qquad \dim\mathcal C(T,v_i)=\deg q_i,
\]
各限制变换的特征多项式和极小多项式均为 $q_i$。特别地，$V$ 为循环空间，当且仅当只有一个非常数不变因子，当且仅当 $m_T=\chi_T$。
\end{theorem}
\begin{proof}
取有理标准型对应的基，并按每个友阵块分组，命题~\ref{prop:xie-ex-sim-cyclic-basis} 即给出循环分解。循环空间的矩阵是友阵，故特征多项式等于极小多项式；反之，若 $m_T=\chi_T$，式~\eqref{eq:xie-ex-sim-char-min} 表明除最大不变因子外不可能再有非常数因子，于是整个空间就是一个循环分量。
\end{proof}

若 $A$ 有 $n$ 个不同特征值，则 $\chi_A=m_A=\prod_i(\lambda-\mu_i)$，所以它是循环的。在分裂域中取相应特征向量 $u_i$，向量 $v=u_1+\cdots+u_n$ 是一个循环向量，因为 $v,Av,\ldots,A^{n-1}v$ 在基 $u_i$ 下的坐标矩阵是非奇异的 Vandermonde 矩阵。更一般地，若 $\chi_A$ 在 $\mathbb K[\lambda]$ 中无重不可约因子，则整除链迫使所有这些因子都出现在最大不变因子中，因此 $m_A=\chi_A$；由于数域的特征为零，不可约多项式在 $\mathbb C$ 中无重根，故此时 $A$ 在 $\mathbb C$ 上有互异特征值。

\begin{proposition}
$\chi_T$ 在 $\mathbb K[\lambda]$ 中不可约，当且仅当 $V$ 没有非零真 $T$-不变子空间，当且仅当每个非零向量都是循环向量。若只假设 $m_T=p$ 不可约，则每个非零不变子空间上的限制变换的极小多项式都是 $p$，并且该子空间可分解为若干个维数 $\deg p$ 的循环子空间的直和。
\end{proposition}
\begin{proof}
若 $U$ 是非零真不变子空间，把它的一组基扩张成 $V$ 的基，所得分块上三角矩阵使 $\chi_T$ 分解为两个正次数多项式的乘积，故不可约性排除了这样的 $U$。若不存在这样的 $U$，则每个非零 $v$ 的循环子空间只能等于 $V$。反之，若每个非零向量都循环，而 $\chi_T=gh$ 是非平凡分解，则 $g(T)h(T)=0$，至少一个因子变换奇异；设 $0\ne v\in\ker g(T)$，则它的循环空间维数至多 $\deg g<n$，产生矛盾。最后，非零限制变换的极小多项式是 $m_T=p$ 的非常数首一因子，只能为 $p$；其所有非常数不变因子也只能为 $p$，应用循环分解即可。
\end{proof}

\begin{example}
证明存在满足 $A^2+2A+5I_n=0$ 的实矩阵，当且仅当 $n$ 为偶数；并解释 $n\geq4$ 时为何有无穷多个不变子空间。
\end{example}
\begin{solve}
$p(\lambda)=\lambda^2+2\lambda+5$ 在 $\mathbb R$ 上不可约，故每个非常数不变因子都为 $p$，从而 $n=2k$。反之，取 $k$ 个 $C(p)=\left(\begin{smallmatrix}0&-5\\1&-2\end{smallmatrix}\right)$ 的直和即可。若 $k\geq2$，在前两个循环分量的基 $e_1,e_2$ 与 $e_3,e_4$ 中，对每个 $c\in\mathbb R$，空间 $\operatorname{span}(e_1+ce_3,e_2+ce_4)$ 都是不变子空间，并且不同 $c$ 给出不同子空间。一般的解与此标准型相似，故也有同样的性质。
\end{solve}

\begin{proposition}
\label{prop:xie-ex-sim-rank-bounds}
对任意 $A$，有 $\deg m_A\leq\operatorname{rank}A+1$。若 $q_1,\ldots,q_s$ 是其非常数不变因子，则对 $\mu\in\mathbb K$ 有
\begin{equation}
 \dim\ker(A-\mu I_n)=\#\{i:q_i(\mu)=0\}.
 \label{eq:xie-ex-sim-geometric-invariant}
\end{equation}
\end{proposition}
\begin{proof}
设 $d=\deg m_A$。有理标准型中包含 $d$ 阶块 $C(m_A)$，其秩至少为 $d-1$，因为其次对角线上的 $1$ 构成一个非零 $d-1$ 阶子式，故 $\operatorname{rank}A\geq d-1$；$d=1$ 时此不等式直接成立。再将特征矩阵的 Smith 等式在 $\lambda=\mu$ 处取值，左右幺模矩阵仍可逆，于是 $\mu I_n-A$ 的秩等于 $\operatorname{diag}(1,\ldots,1,q_1(\mu),\ldots,q_s(\mu))$ 的秩，即得所述零度公式。
\end{proof}

\begin{proposition}[零迹与交换子]
\label{prop:xie-ex-sim-commutator}
$C\in M_n(\mathbb K)$ 可写成 $C=AB-BA$，当且仅当 $\operatorname{tr}C=0$。
\end{proposition}
\begin{proof}
必要性由 $\operatorname{tr}(AB)=\operatorname{tr}(BA)$ 得出。充分性先证零迹矩阵相似于主对角元全为零的矩阵，对阶数归纳：纯量矩阵 $cI_n$ 若迹为零，则 $c=0$；若 $C$ 非纯量，则存在 $v$ 使 $v,Cv$ 线性无关，否则每个向量都是特征向量，比较一组基及两两之和便会推出 $C$ 是纯量。以 $v,Cv$ 开始选基，第一主对角元为零，右下 $n-1$ 阶块的迹仍为零。用归纳假设只在该块内换基，就使全部主对角元为零。设由此得到 $C_0=(c_{ij})$，取 $\mathbb K$ 中互异的 $a_1,\ldots,a_n$，令 $A_0=\operatorname{diag}(a_1,\ldots,a_n)$、$b_{ij}=c_{ij}/(a_i-a_j)$（$i\ne j$）、$b_{ii}=0$，则 $A_0B_0-B_0A_0=C_0$。最后把两矩阵同时变换回原基即可。
\end{proof}

\section{乘法交换性诱导的多项式表示}
\label{sec:xie-ex-sim-centralizer}

记 $\mathbb K[A]=\{f(A):f\in\mathbb K[\lambda]\}$，以及\emph{中心化子（centralizer）} $\mathcal Z(A)=\{B\in M_n(\mathbb K):AB=BA\}$。总有 $\mathbb K[A]\subseteq\mathcal Z(A)$，但相等需要条件，例如 $\mathbb K[I_n]$ 仅含纯量矩阵，而 $\mathcal Z(I_n)=M_n(\mathbb K)$。

\begin{theorem}
\label{thm:xie-ex-sim-centralizer}
下列条件等价：$A$ 是循环矩阵；$m_A=\chi_A$；$\mathcal Z(A)=\mathbb K[A]$。这些条件成立时，$\mathcal Z(A)$ 的一组基为 $I_n,A,\ldots,A^{n-1}$，每个交换矩阵都唯一表示为次数小于 $n$ 的 $A$ 的多项式。
\end{theorem}
\begin{proof}
设 $v$ 为循环向量，$B$ 与 $A$ 交换。唯一写成 $Bv=\sum_{j=0}^{n-1}c_jA^jv=f(A)v$ 后，交换性给出 $BA^jv=A^jBv=f(A)A^jv$。因为这些 $A^jv$ 构成一组基，故 $B=f(A)$。若次数小于 $n$ 的多项式满足 $f(A)=0$，则由 $\deg m_A=n$ 可知 $f=0$，从而基与唯一性也成立。

反之，若非常数不变因子有 $s\geq2$ 个，在有理标准型 $F=\operatorname{diag}(F(q_1),\ldots,F(q_s))$ 中取 $E=\operatorname{diag}(I,0,\ldots,0)$，则 $EF=FE$。若 $E=f(F)$，最后一块给出 $q_s\mid f$，而 $q_1\mid q_s$，便有 $f(F(q_1))=0$，与第一块为 $I$ 矛盾。将 $E$ 变换回原基即得一个不属于 $\mathbb K[A]$ 的交换矩阵。
\end{proof}

对单个 Jordan 块 $J=\mu I_n+N_n$，中心化子具体为
\[
 \mathcal Z(J)=\left\{c_0I_n+c_1N_n+\cdots+c_{n-1}N_n^{n-1}:c_j\in\mathbb K\right\},
\]
即各条上对角线为常数的上三角矩阵。更一般地，若 $A$ 是下 Hessenberg（lower Hessenberg）矩阵，即主对角线上方第二条及更高的对角线全为零，并且上次对角元全部非零，则特征矩阵的前 $n-1$ 行、后 $n-1$ 列给出非零常数子式，故 $D_{n-1}=1$，$A$ 也是循环矩阵。$n=1$ 时结论直接成立。

\begin{proposition}[分块多项式的拼接]
\label{prop:xie-ex-sim-crt}
设 $A=\operatorname{diag}(A_1,\ldots,A_k)$，每个 $A_i$ 适合非零多项式 $g_i$，且 $g_i$ 两两互素。若 $B=\operatorname{diag}(f_1(A_1),\ldots,f_k(A_k))$，则存在次数小于 $\deg m_A\leq n$ 的多项式 $f$，使 $B=f(A)$。
\end{proposition}
\begin{proof}
多项式中国剩余定理给出 $h\equiv f_i\pmod{g_i}$，于是 $h(A_i)=f_i(A_i)$，故 $h(A)=B$。再用 $m_A$ 除 $h$ 并取余式，即得所需次数界。互素性使各块上的要求彼此独立；任意分块上的多项式表示一般不能直接拼接。
\end{proof}

\begin{example}
设 $\operatorname{rank}A=n-1$ 且 $AB=BA=0$，证明 $B$ 可表示为 $A$ 的次数小于 $n$ 的多项式。
\end{example}
\begin{solve}
若 $B=0$ 无须证明。否则 $AB=0$ 使 $B$ 的列属于一维空间 $\ker A$，$BA=0$ 使其行属于一维左零空间，所以所有满足这两个条件的矩阵组成一个一维空间。因 $A$ 奇异，可写 $m_A(\lambda)=\lambda g(\lambda)$；由极小性知 $g(A)\ne0$，而 $Ag(A)=g(A)A=0$。故 $B$ 是 $g(A)$ 的常数倍，且 $\deg g=\deg m_A-1<n$。结合满秩时的 Cayley--Hamilton 公式与秩至多 $n-2$ 时伴随矩阵为零，可知任意方阵的伴随矩阵都能表示为 $A$ 的次数小于 $n$ 的多项式。
\end{solve}

\begin{theorem}[复数域上的 Jordan--Chevalley 分解]
\label{thm:xie-ex-sim-jc}
任意复矩阵 $A$ 唯一分解为 $A=D+N$，其中 $D$ 可对角化，$N$ 幂零，且 $DN=ND$。此外 $D,N$ 都是 $A$ 的多项式。
\end{theorem}
\begin{proof}
将属于同一特征值 $\mu_i$ 的 Jordan 块合并为 $J_i=\mu_iI+N_i$，令 $s_i$ 为这些块的最大阶数。多项式中国剩余定理给出 $f(\lambda)\equiv\mu_i\pmod{(\lambda-\mu_i)^{s_i}}$，故在 Jordan 基下，$f(A)$ 的各块为 $\mu_iI$。取 $D=f(A)$、$N=A-f(A)$，即满足全部条件。

若还有 $A=D'+N'$ 满足这些条件，则 $D',N'$ 都与 $A$ 交换，因而与 $D,N$ 交换。两个交换的可对角化矩阵 $D,D'$ 可同时对角化，因为其中一个的各特征子空间对另一个不变，限制变换仍可对角化，故 $D-D'$ 可对角化。另一方面，$N,N'$ 交换且幂零，由二项式公式知 $N'-N$ 仍幂零。等式 $D-D'=N'-N$ 表明这一矩阵同时可对角化且幂零，只能为零，于是分解唯一。
\end{proof}

\section{可对角化的判定}
\label{sec:xie-ex-sim-diagonalizable}

\begin{theorem}
\label{thm:xie-ex-sim-diagonalizable}
复矩阵 $A$ 可对角化，当且仅当 $m_A$ 无重根，当且仅当其初等因子全为一次多项式，当且仅当其 Jordan 块全为一阶。对一般数域 $\mathbb K$，在 $\mathbb K$ 上可对角化要求 $m_A$ 是\emph[caution]{互异一次因子的乘积}，仅在 $\mathbb K[\lambda]$ 中没有重不可约因子还不够。
\end{theorem}
\begin{proof}
在复数域上，$\lambda-\mu$ 在 $m_A$ 中的指数恰为属于 $\mu$ 的最大 Jordan 块阶数，故这些条件等价。一般数域上，若 $m_A=\prod_i(\lambda-\mu_i)$ 且各 $\mu_i\in\mathbb K$ 互异，则互素因子对应的核空间直和分解给出 $V=\bigoplus_i\ker(A-\mu_iI)$，因此存在特征向量基。反之，对角矩阵的极小多项式就是其互异对角元对应的一次因子之积。
\end{proof}

因此，只要找到一个无重根的零化多项式 $g$，由 $m_A\mid g$ 就能判定复可对角化。幂等矩阵 $A^2=A$ 的标准型是 $\operatorname{diag}(I_r,0)$，$r=\operatorname{rank}A$；满足 $A^k=I$ 的复矩阵也可对角化，因为 $z^k-1$ 与其导数没有公共根。相反，非零幂零矩阵一定不能对角化：若能对角化，其所有对角元只能为零，便会推出矩阵本身为零。实矩阵 $C(\lambda^2+1)$ 的极小多项式无重不可约因子，但没有实特征值，说明“复可对角化”与“实可对角化”必须区别。

\begin{proposition}[限制、商空间与不变补空间]
\label{prop:xie-ex-sim-complement}
若复线性变换 $T$ 可对角化，则它在任意不变子空间 $U$ 上的限制变换及在 $V/U$ 上的诱导变换都可对角化。并且，$T$ 可对角化，当且仅当它的每个不变子空间都有不变补空间。
\end{proposition}
\begin{proof}
$m_T$ 同时零化限制变换和商变换，故两者的极小多项式都整除无重根的 $m_T$，得到第一结论。设 $E_i$ 为 $T$ 的各特征子空间，多项式谱投影保持 $U$，所以 $U=\bigoplus_i(U\cap E_i)$；在每个 $E_i$ 内选取 $U\cap E_i$ 的补空间并求直和，就得到不变补空间。反之，设所有不变子空间都有不变补空间，令 $E$ 为全部特征向量张成的空间。若 $E\ne V$，其非零不变补空间上的限制变换在复数域中有一个非零特征向量，这又应属于 $E$，与直和矛盾。因此 $E=V$，$T$ 可对角化。
\end{proof}

\begin{proposition}[核、像与秩的判据]
\label{prop:xie-ex-sim-kernel-criterion}
复矩阵 $A$ 可对角化，当且仅当对每个特征值 $\mu$，以下任一条件成立：
\begin{align*}
 &\ker(A-\mu I)\cap\operatorname{Im}(A-\mu I)=\{0\};\\
 &V=\ker(A-\mu I)\oplus\operatorname{Im}(A-\mu I);\\
 &\ker(A-\mu I)=\ker((A-\mu I)^2);\\
 &\operatorname{rank}(A-\mu I)=\operatorname{rank}((A-\mu I)^2).
\end{align*}
上述条件均需对每个特征值分别检验。
\end{proposition}
\begin{proof}
记 $L=A-\mu I$。映射 $L:\ker L^2\to\ker L\cap\operatorname{Im}L$ 满射，核为 $\ker L$，故交空间为零等价于两级核相等；秩零度定理又给出其余等价条件。若 Jordan 型中有属于 $\mu$ 的长度至少为二的链 $Lv_2=v_1\ne0$、$Lv_1=0$，则 $v_1$ 属于上述交空间，所以这些条件迫使每个块均为一阶。反之，对角型下核与像由互补的基向量张成，条件成立。核或像连续两级相等后，高次核或像也随之稳定。
\end{proof}

\begin{example}
设 $g\in\mathbb C[z]$，且对 $A$ 的每个特征值 $\mu$ 都有 $g'(\mu)\ne0$，证明 $A$ 可对角化当且仅当 $g(A)$ 可对角化。
\end{example}
\begin{solve}
必要性由在特征向量基中计算多项式得到。对充分性，在 $J_r(\mu)=\mu I+N_r$ 上有
\[
 g(J_r(\mu))-g(\mu)I=N_rH(N_r),\qquad H(0)=g'(\mu)\ne0.
\]
$H(N_r)$ 可逆且与 $N_r$ 交换，所以左边各次幂的秩与 $N_r$ 的对应幂相同，它有且只有一个大小为 $r$ 的 Jordan 块。若 $A$ 有 $r>1$ 的块，则 $g(A)$ 也不能对角化。导数条件不可省，例如 $A=J_2(0)$、$g(z)=z^2$ 时 $g(A)=0$ 可对角化而 $A$ 不可对角化。
\end{solve}

\begin{example}
设 $A\in M_n(\mathbb C)$，在 $M_n(\mathbb C)$ 上定义 $\Phi(X)=AXA^{\mathsf T}$、$\Psi(X)=AX-XA^{\mathsf T}$。证明这两个线性变换分别可对角化，当且仅当 $A$ 可对角化。
\end{example}
\begin{solve}
若 $A=PDP^{-1}$，用坐标变换 $X=PYP^{\mathsf T}$，两个算子分别化为 $Y\mapsto DYD^{\mathsf T}$、$Y\mapsto DY-YD^{\mathsf T}$；基础矩阵 $E_{ij}$ 分别是特征值为 $d_id_j$、$d_i-d_j$ 的特征向量，故充分性成立。

反之，若 $A$ 有 $r\geq2$ 的块 $J_r(\mu)$，同样的坐标变换把问题化到 Jordan 型；只在这个块对应的 $r\times r$ 位置非零的矩阵组成一个共同不变子空间。在该子空间的基础矩阵基下，$\Phi$ 为三角形且对角元全为 $\mu^2$，但不是纯量算子：$\mu\ne0$ 时 $\Phi(E_{21})=\mu^2E_{21}+\mu E_{11}$，$\mu=0$ 时 $\Phi(E_{22})=E_{11}$。同理，$\Psi$ 为三角形且对角元全为零，而 $\Psi(E_{21})=E_{11}\ne0$。只有一个特征值的可对角化算子必为纯量算子，所以这两个限制均不可对角化，与命题~\ref{prop:xie-ex-sim-complement} 矛盾。这里右侧明确使用转置；具体计算必须与题设中的乘法次序一致。
\end{solve}

\begin{example}
设 $s=\deg m_A$，定义 $s$ 阶矩阵 $H=(\operatorname{tr}(A^{i+j-2}))_{i,j=1}^{s}$，其中 $\operatorname{tr}(A^0)=n$。证明 $A$ 在复数域上可对角化，当且仅当 $H$ 可逆。
\end{example}
\begin{solve}
设互异特征值为 $\mu_1,\ldots,\mu_k$，代数重数为 $a_1,\ldots,a_k$。令 $W_{ij}=\mu_j^{i-1}$，则 $H=W\operatorname{diag}(a_1,\ldots,a_k)W^{\mathsf T}$。总有 $s\geq k$；若 $A$ 不可对角化，$m_A$ 有重根，故 $s>k$，于是 $\operatorname{rank}H\leq k<s$。若 $A$ 可对角化，则 $s=k$，$W$ 为方阵，且
\[
 \det H=\left(\prod_{j=1}^{k}a_j\right)
 \prod_{i<j}(\mu_j-\mu_i)^2\ne0.
\]
这是 Vandermonde 分解的应用，不需要把复数域上的双线性矩阵 $W^{\mathsf T}$ 当作共轭转置。
\end{solve}

\section{Jordan 标准型的求法}
\label{sec:xie-ex-sim-jordan-computation}

计算 Jordan 型时，先用特征多项式确定每个特征值的总维数，再用极小多项式确定最大块，或用 $\dim\ker(A-\mu I)$ 确定块数。若信息仍不足，就继续计算 $(A-\mu I)^k$ 的秩。对于结构简单的特征矩阵，非零常数子式可以迅速确定若干行列式因子；但含参数时，除以参数之前必须分情形，不能把可能为零的量作为初等变换中的非零常数。

\begin{theorem}[秩序列确定 Jordan 块]
\label{thm:xie-ex-sim-rank-sequence}
固定特征值 $\mu$，记 $R_k=\operatorname{rank}((A-\mu I)^k)$，$R_0=n$。则属于 $\mu$ 且阶数至少为 $k$ 的 Jordan 块数为 $R_{k-1}-R_k$，阶数恰为 $k$ 的块数为
\begin{equation}
 b_k=R_{k-1}-2R_k+R_{k+1},\qquad k\geq1.
 \label{eq:xie-ex-sim-block-count}
\end{equation}
因此对 $A,B\in M_n(\mathbb C)$，两者相似当且仅当对其特征值集合之并中的每个 $\mu$，及 $1\leq k\leq n$，都有 $\operatorname{rank}((A-\mu I)^k)=\operatorname{rank}((B-\mu I)^k)$。
\end{theorem}
\begin{proof}
非 $\mu$ 特征值的块在减去 $\mu I$ 后可逆，对所有 $R_k$ 贡献相同的秩。属于 $\mu$ 的一个 $r$ 阶块贡献 $\max(r-k,0)$，其前后两级之差在 $r\geq k$ 时为 $1$，否则为 $0$，所以第一次差分统计“至少 $k$ 阶”，再差分就统计“恰为 $k$ 阶”。幂次到 $n$ 时核已稳定，故有限的秩数据足以确定全部块。相似必保留这些秩；反过来，相同的块数给出相同的 Jordan 型。
\end{proof}

\begin{example}
设 $N$ 为九阶幂零矩阵，且 $\operatorname{rank}N=5$、$\operatorname{rank}N^2=2$、$N^3=0$，求其 Jordan 型。
\end{example}
\begin{solve}
秩序列为 $9,5,2,0,0$，相邻差为 $4,3,2,0$，所以总块数为 $4$，至少二阶的有 $3$ 个，至少三阶的有 $2$ 个。再取差得到一阶、二阶、三阶块数分别为 $1,1,2$，故标准型为 $\operatorname{diag}(J_1(0),J_2(0),J_3(0),J_3(0))$。检验总阶数 $1+2+3+3=9$，极小多项式为 $\lambda^3$。
\end{solve}

\begin{example}
\label{ex:xie-ex-sim-rank-one}
设 $n>1$ 且 $\operatorname{rank}A=1$，按 $\operatorname{tr}A$ 分类求 Jordan 型。
\end{example}
\begin{solve}
写 $A=uv^{\mathsf T}$，令 $c=v^{\mathsf T}u=\operatorname{tr}A$，则 $A^2=cA$。若 $c\ne0$，多项式 $\lambda(\lambda-c)$ 无重根，故 $A$ 可对角化，且其秩为一，标准型为 $\operatorname{diag}(c,0,\ldots,0)$。若 $c=0$，则 $A^2=0$ 而 $A\ne0$，所以各块至多二阶；每个二阶零特征值块贡献秩一，故恰有一个 $J_2(0)$，其余为一阶零块。这也说明秩一矩阵幂等当且仅当迹为 $1$，幂零当且仅当迹为 $0$。
\end{solve}

\begin{example}
设
\[
 A=\begin{pmatrix}
 1&a+2&5&7\\0&1&3&6\\0&0&1&b+4\\0&0&0&1
 \end{pmatrix},
\]
按参数确定 Jordan 型。
\end{example}
\begin{solve}
令 $N=A-I_4$、$\alpha=a+2$、$\beta=b+4$。若 $\alpha\beta\ne0$，则 $\operatorname{rank}N=3$，唯一特征值 $1$ 的几何重数为一，故标准型为 $J_4(1)$。若 $\alpha\beta=0$，则 $\operatorname{rank}N=2$：第三、四列在前两行构成的子式为 $5\cdot6-7\cdot3=9$，而秩三的子式消失。又有
\[
 N^2_{13}=3\alpha,\quad N^2_{14}=6\alpha+5\beta,
 \quad N^2_{24}=3\beta,\quad N^3_{14}=3\alpha\beta.
\]
若 $\alpha,\beta$ 恰有一个为零，则 $N^2\ne0$、$N^3=0$，两块的最大阶为三，故标准型为 $\operatorname{diag}(J_3(1),1)$；若二者均为零，则 $N^2=0$，两块均为二阶，标准型为 $\operatorname{diag}(J_2(1),J_2(1))$。
\end{solve}

\begin{proposition}[Jordan 块的整数幂]
\label{prop:xie-ex-sim-jordan-power}
若 $a\ne0$、$m$ 为非零整数，则 $J_n(a)^m$ 的 Jordan 型为 $J_n(a^m)$。若 $m\geq1$ 且 $N=J_n(0)$，则 $m\geq n$ 时 $N^m=0$；$m<n$ 时写 $n=mq+r$，$0\leq r<m$，其 Jordan 型由 $r$ 个 $J_{q+1}(0)$ 与 $m-r$ 个 $J_q(0)$ 组成。
\end{proposition}
\begin{proof}
对非零 $a$，在 $aI+N$ 的有限 Taylor 展开中，$J_n(a)^m-a^mI=NH(N)$ 且 $H(0)=ma^{m-1}\ne0$，所以其各次幂的秩与 $N$ 相同。负整数幂也成立，因为 $a\ne0$ 保证 $z^m$ 在 $a$ 附近解析。对零特征值，$N^me_j=e_{j-m}$，约定下标非正时为零；按下标模 $m$ 分组，就把原来一条长度为 $n$ 的链拆为 $m$ 条链，其中 $r$ 条长 $q+1$，其余长 $q$。各条链互不重叠且覆盖全部基向量，得到所述标准型。
\end{proof}

特别地，若 $A$ 的特征值全为 $1$，则对每个正整数 $m$ 都有 $A^m\sim A$；若其特征值全属于 $\{1,-1\}$，则 $A^{-1}\sim A$。这两个结论都保留了每个块的阶数，但一般的矩阵幂可能合并不同的特征值，也可能使零特征值块分裂。

\section{过渡矩阵的求法}
\label{sec:xie-ex-sim-transition}

已知 Jordan 型 $J$ 后，还需构造可逆矩阵 $P$ 满足 $AP=PJ$，才能把标准型用于具体计算。若 $J$ 中有块 $J_r(\mu)$，$P$ 中对应的列必须构成一条链
\begin{equation}
 (A-\mu I)v_1=0,\qquad
 (A-\mu I)v_j=v_{j-1}\quad(2\leq j\leq r).
 \label{eq:xie-ex-sim-chain}
\end{equation}
这里 $v_j$ 是广义特征向量；本文以 $v_1$ 为零级，$v_j$ 为 $j-1$ 级。每条链应按 $v_1,\ldots,v_r$ 排列为 $P$ 的列，才能对应上次对角线为 $1$ 的约定。

\subsection{从多项式相抵变换提取常数矩阵}

\begin{proposition}
\label{prop:xie-ex-sim-polynomial-transition}
设幺模矩阵 $U(\lambda),W(\lambda)$ 满足
\[
 U(\lambda)(\lambda I-A)W(\lambda)=\lambda I-J.
\]
将 $W$ 作右侧除法，写成 $W(\lambda)=S(\lambda)(\lambda I-J)+P$，其中 $P$ 为常数矩阵，则 $P$ 可逆，且 $AP=PJ$。
\end{proposition}
\begin{proof}
右侧带余除法成立，是因为 $\lambda I-J$ 的最高次系数为单位矩阵，可以逐次消去最高项。代入相抵等式并整理，得
\[
 (\lambda I-A)P=
 \bigl(U^{-1}(\lambda)-(\lambda I-A)S(\lambda)\bigr)(\lambda I-J).
\]
左边次数至多为一，右边最后一因子首项为 $\lambda I$，故括号内只能是常数矩阵 $R$；比较一次项与常数项得 $P=R$、$AP=PJ$。若 $P$ 奇异，$\ker P$ 是非零的 $J$-不变子空间，在复数域中可取 $0\ne v\in\ker P$ 满足 $Jv=\mu v$，于是 $W(\mu)v=S(\mu)(\mu I-J)v+Pv=0$，与 $W(\mu)$ 可逆矛盾。原矩阵在较小数域时，也可在复数域中作这一反证，所得常数矩阵 $P$ 的元素仍属于原数域。
\end{proof}

\subsection{解链方程与选择链顶向量}

直接求解式~\eqref{eq:xie-ex-sim-chain} 时，必须保证下一步方程有解。任取一个特征向量未必能继续延长成所需链：例如 $A=\operatorname{diag}(J_2(0),0)$，特征向量 $e_3$ 不属于 $\operatorname{Im}A$，所以 $Av=e_3$ 无解。\emph[tip]{可先选择链顶 $w$，满足 $(A-\mu I)^rw=0$、$(A-\mu I)^{r-1}w\ne0$，再按倒序列出其像；这样链方程自动成立。}不同链之间还要检验线性无关，不能只逐条保证链内无关。

\begin{example}
\label{ex:xie-ex-sim-transition}
对
\[
 A=\begin{pmatrix}2&6&-15\\1&1&-5\\1&2&-6\end{pmatrix},
 \qquad J=\operatorname{diag}(-1,J_2(-1)),
\]
构造 $P$ 使 $P^{-1}AP=J$。
\end{example}
\begin{solve}
令 $N=A+I_3$，则 $N^2=0$、$\operatorname{rank}N=1$。可取短链向量 $u=(-2,1,0)^{\mathsf T}$；又取长链顶 $w=e_1$，则 $v=Nw=(3,1,1)^{\mathsf T}$。因此
\[
 P=(u,v,w)=\begin{pmatrix}-2&3&1\\1&1&0\\0&1&0\end{pmatrix},
 \qquad \det P=1.
\]
由 $Nu=Nv=0$、$Nw=v$ 可直接核对 $AP=PJ$。在 $\ker N$ 的常见基 $(-2,1,0)^{\mathsf T}$、$(5,0,1)^{\mathsf T}$ 中，任意单取一个都不能作为长链的底端，而二者之和恰为 $v$；这一例子说明保留参数、检查可解性是必要的。
\end{solve}

对一般根子空间，记 $N=A-\mu I$ 在该空间上的限制，并令
\[
 F_k=\ker N\cap\operatorname{Im}N^{k-1},\qquad
 F_1\supseteq F_2\supseteq\cdots.
\]
从最大的非零 $F_k$ 开始，逐层将 $F_{k+1}$ 的基扩张为 $F_k$ 的基；每个本层新增的向量 $b$ 都可以提升为 $N^{k-1}w=b$，并满足 $N^kw=0$。由这个 $w$ 构造长度为 $k$ 的链，最终得到全部链。其独立性与维数论证见第~\ref{sec:xie-ex-sim-geometry} 节。这种方法把“选择合适的特征向量”改为选择与核、像滤过相容的基，尤其适合多个不同长度块同时出现的情形。

\begin{example}
设 $A_{ij}\in M_n(\mathbb K)$（$1\leq i,j\leq n$）均非零，并满足 $A_{ij}A_{kl}=\delta_{jk}A_{il}$，其中 $\delta_{jk}$ 为 Kronecker 符号。证明存在同一个可逆矩阵 $P$，使所有 $P^{-1}A_{ij}P=E_{ij}$ 同时成立。
\end{example}
\begin{solve}
取 $v$ 使 $u_1=A_{11}v\ne0$，则 $A_{11}u_1=u_1$。令 $u_i=A_{i1}u_1$，有 $A_{1i}u_i=u_1$，故每个 $u_i\ne0$，且 $A_{ij}u_k=\delta_{jk}u_i$。若 $\sum_kc_ku_k=0$，作用 $A_{1j}$ 得 $c_ju_1=0$，所以全部系数为零。这 $n$ 个向量构成一组基，取 $P=(u_1,\ldots,u_n)$，便有 $A_{ij}P=PE_{ij}$，得到所求。乘法关系在此直接规定了新基下每个算子的作用。
\end{solve}

\section{Jordan 标准型的应用}
\label{sec:xie-ex-sim-applications}

只要问题在相似或同时相似变换下保持不变，就可以依次处理单个 Jordan 块、分块对角矩阵和一般矩阵。对涉及两个矩阵的问题，通常先化简其中一个，并把另一个按相同方式分块；不能独立改变两矩阵的基而仍声称保留它们之间的乘法关系。

\subsection{秩、核像分解与矩阵方程}

\begin{proposition}
\label{prop:xie-ex-sim-fitt-zero}
复矩阵 $A$ 相似于 $\operatorname{diag}(N,C)$，其中 $N$ 幂零、$C$ 可逆，允许其中一块为空。若 $\mu$ 的代数重数为 $a$，其最大 Jordan 块阶数为 $s$，则
\[
 \operatorname{rank}((A-\mu I)^k)=n-a
 \quad\Longleftrightarrow\quad k\geq s.
\]
特别地，$\operatorname{rank}A=\operatorname{rank}A^2$ 当且仅当零特征值的块全为一阶；若零不是特征值，这一秩等式也成立。
\end{proposition}
\begin{proof}
将零特征值与非零特征值的块分别归组，便得到第一结论。在 $A-\mu I$ 中，非 $\mu$ 块始终贡献秩 $n-a$，其余块各贡献 $\max(r-k,0)$；这些贡献全部消失恰好要求 $k$ 不小于最大块阶数。
\end{proof}

\begin{example}
设 $AB=BA=0$ 且 $\operatorname{rank}A=\operatorname{rank}A^2$，证明 $\operatorname{rank}(A+B)=\operatorname{rank}A+\operatorname{rank}B$。
\end{example}
\begin{solve}
在相同的基下，$A=\operatorname{diag}(0,C)$，其中 $C$ 可逆。将 $B$ 对应分块，两个乘积为零迫使其非左上块全部为零，故 $B=\operatorname{diag}(B_0,0)$。此时 $A+B=\operatorname{diag}(B_0,C)$，秩的加法公式立即成立。这一分块论证也适用于一般数域，所需核像分解见第~\ref{sec:xie-ex-sim-general-field} 节。
\end{solve}

\begin{theorem}[Sylvester 方程]
\label{thm:xie-ex-sim-sylvester}
设 $A\in M_a(\mathbb C)$、$B\in M_b(\mathbb C)$。方程 $AX-XB=0$ 只有零解，当且仅当 $A,B$ 无公共特征值；这也等价于对每个 $C\in M_{a\times b}(\mathbb C)$，方程 $AX-XB=C$ 有唯一解。
\end{theorem}
\begin{proof}
分别换基把 $B$ 化为 Jordan 型后，可按其各块独立解方程。对一个块 $J_r(\mu)$，将未知块写为 $X=(x_1,\ldots,x_r)$，右端块写为 $(c_1,\ldots,c_r)$，得到
\[
 (A-\mu I)x_1=c_1,\qquad
 (A-\mu I)x_j=x_{j-1}+c_j\quad(2\leq j\leq r).
\]
若无公共特征值，$A-\mu I$ 可逆，故依次唯一确定各列。若存在公共特征值 $\mu$，取 $A$ 的右特征向量 $u\ne0$ 与 $B$ 的左特征行向量 $w\ne0$，则 $X=uw\ne0$ 满足 $AX=XB$，所以相应线性映射有非零核，不可能有唯一解。由于定义域与值域的维数均为 $ab$，单射、满射和可逆等价，从而得到全部结论。
\end{proof}

若 $A,B$ 无公共特征值，解 $AX-XB=C$ 后取 $S=\left(\begin{smallmatrix}I&-X\\0&I\end{smallmatrix}\right)$，可直接验证
\[
 S^{-1}\begin{pmatrix}A&C\\0&B\end{pmatrix}S
 =\begin{pmatrix}A&0\\0&B\end{pmatrix}.
\]
因此两个谱不相交的对角块之间的非对角块可以通过相似变换消去。若两块有公共特征值，则这种消去一般不成立，非对角块可能把两个短链连接成长链。

\begin{example}
设 $C=AB-BA$ 且 $AC=CA$、$BC=CB$，证明 $C$ 幂零。
\end{example}
\begin{solve}
把 $C$ 按不同特征值的根子空间分块。因为 $A,B$ 都与 $C$ 交换，它们保持每个根子空间；在其中一个维数为 $d$、对应特征值为 $\mu$ 的分量上有 $C_0=A_0B_0-B_0A_0$。取迹得 $d\mu=\operatorname{tr}C_0=0$，而特征为零，故 $\mu=0$。因此 $C$ 的全部特征值为零，由 Cayley--Hamilton 定理可知 $C$ 幂零。
\end{solve}

\subsection{矩阵幂、方根与对称矩阵分解}

当 $\mu\ne0$ 时，对非负整数 $k$，Jordan 块的二项式展开为
\begin{equation}
 J_r(\mu)^k=
 \sum_{j=0}^{\min(k,r-1)}\binom{k}{j}\mu^{k-j}N_r^j.
 \label{eq:xie-ex-sim-matrix-power}
\end{equation}
当 $\mu=0$ 时可以直接使用 $N_r^k$，避免写出负指数或含混的 $0^0$。将各块计算结果拼接，再用 $A^k=PJ^kP^{-1}$ 回到原坐标。

\begin{proposition}
\label{prop:xie-ex-sim-complex-root}
每个可逆复矩阵都存在任意给定正整数次的复矩阵方根。
\end{proposition}
\begin{proof}
只需对 $J_r(\mu)=\mu I+N_r$ 构造，因可逆性知 $\mu\ne0$。取 $\nu^m=\mu$，令
\[
 B_r=\nu\sum_{j=0}^{r-1}\binom{1/m}{j}\mu^{-j}N_r^j,
 \qquad
 \binom{\alpha}{j}=\frac{\alpha(\alpha-1)\cdots(\alpha-j+1)}{j!},
 \quad \binom{\alpha}{0}=1.
\]
形式幂级数恒等式 $((1+z)^{1/m})^m=1+z$ 在模 $z^r$ 意义下成立，代入幂零矩阵 $\mu^{-1}N_r$ 后高次项全为零，故 $B_r^m=J_r(\mu)$。各块直和后再作相似变换，即得所求。奇异情形没有普遍保证，例如 $J_2(0)$ 无二阶矩阵平方根：若 $B^2=J_2(0)$，则 $B$ 的特征值全为零，二阶 Cayley--Hamilton 定理给出 $B^2=0$，矛盾。
\end{proof}

\begin{proposition}
每个复矩阵都能分解为两个复对称矩阵的乘积，并且可以指定其中任意一个因子可逆；还存在可逆复对称矩阵 $Q$ 满足 $Q^{-1}AQ=A^{\mathsf T}$。
\end{proposition}
\begin{proof}
设 $A=PJP^{-1}$，对每个 Jordan 块取反序单位矩阵 $R_i$，其反对角线元素为 $1$，其余为零。令 $R=\operatorname{diag}(R_i)$，则 $R^{\mathsf T}=R$、$R^2=I$，且 $JR=RJ^{\mathsf T}$。于是 $Q=PRP^{\mathsf T}$ 可逆且对称，并有 $AQ=QA^{\mathsf T}$，即 $Q^{-1}AQ=A^{\mathsf T}$。由此 $AQ$ 和 $Q^{-1}A$ 都对称，故
\[
 A=(AQ)Q^{-1}=Q(Q^{-1}A)
\]
给出两种分解，分别让第二个或第一个因子可逆。这里的对称指普通转置，不能替换为共轭转置。
\end{proof}

\section{矩阵函数}
\label{sec:xie-ex-sim-functions}

\subsection{Jordan 块上的有限 Taylor 公式}

\begin{theorem}
\label{thm:xie-ex-sim-functional-calculus}
若 $f$ 在 $A$ 的每个特征值附近解析，且 $A=P\operatorname{diag}(J_{r_i}(\mu_i))P^{-1}$，则可通过
\begin{equation}
 f(J_r(\mu))=\sum_{j=0}^{r-1}\frac{f^{(j)}(\mu)}{j!}N_r^j,
 \qquad
 f(A)=P\operatorname{diag}\bigl(f(J_{r_i}(\mu_i))\bigr)P^{-1}
 \label{eq:xie-ex-sim-taylor}
\end{equation}
定义矩阵函数。若 $f$ 是多项式或在包含全部特征值的圆盘内收敛的幂级数，此定义与直接代入一致。$f(A)$ 的特征值为各 $f(\mu_i)$，按相应代数重数计。
\end{theorem}
\begin{proof}
多项式在 $\mu$ 处的 Taylor 展开代入 $\mu I+N_r$ 后，$N_r^r=0$ 使高次项全部消失，得到第一式。对收敛幂级数，在收敛圆盘内部，其各阶导数级数都收敛，故部分和矩阵的每条上对角线分别趋于 $f^{(j)}(\mu)/j!$，得到相同公式。对局部解析函数，设属于互异特征值 $\mu$ 的最大块阶数为 $s_\mu$，多项式插值给出唯一的 $h$，满足 $\deg h<\deg m_A$ 及
\[
 h^{(j)}(\mu)=f^{(j)}(\mu),\qquad 0\leq j<s_\mu.
\]
这可由模 $(z-\mu)^{s_\mu}$ 的多项式中国剩余定理得到。公式定义的矩阵就是 $h(A)$，因此不依赖 Jordan 基或块的排列。各块是对角元全为 $f(\mu)$ 的上三角矩阵，所以得到特征值结论。
\end{proof}

这一公式也说明谱半径与收敛半径的关系：若 $\rho(A)<R$，全部所需导数级数都在收敛圆盘内部收敛，故矩阵级数收敛；若 $\rho(A)>R$，取一个模大于 $R$ 的特征值及其特征向量，矩阵部分和在该向量上的作用就是不收敛的标量部分和，故矩阵级数不收敛。在边界上，不仅 $f(\mu)$ 的级数，还可能需要相应的导数级数收敛，不能仅凭 $\rho(A)=R$ 判定。

\begin{example}
设 $N$ 幂零，证明 $e^N$ 与 $I+N$ 相似。
\end{example}
\begin{solve}
若 $N^s=0$，则 $e^N=\sum_{j=0}^{s-1}N^j/j!$，且
\[
 e^N-I=NH(N),\qquad H(N)=\sum_{j=0}^{s-2}\frac{N^j}{(j+1)!}
\]
在 $s\geq2$ 时成立；$N=0$ 的结论直接成立。$H(N)$ 的常数项为 $I$，故可逆，并与 $N$ 交换。于是 $\operatorname{rank}((e^N-I)^k)=\operatorname{rank}N^k$，而两矩阵 $e^N$、$I+N$ 的特征值全为 $1$，用定理~\ref{thm:xie-ex-sim-rank-sequence} 即得相似。这是“函数在特征值处的导数非零便保留块阶数”的具体应用。
\end{solve}

\subsection{矩阵指数与可交换条件}

定义 Frobenius 范数 $\|A\|_F=(\sum_{i,j}|a_{ij}|^2)^{1/2}$。Cauchy--Schwarz 不等式给出 $\|AB\|_F\leq\|A\|_F\|B\|_F$，而三角不等式来自矩阵元组成的向量的欧氏范数。因此矩阵指数 $e^A=\sum_{k=0}^{\infty}A^k/k!$ 对每个有限阶矩阵绝对收敛，矩阵级数的乘法可以用对应范数级数控制。

\begin{proposition}
\label{prop:xie-ex-sim-exp}
若 $AB=BA$，则 $e^Ae^B=e^{A+B}=e^Be^A$。任意复矩阵 $A$ 都满足
\[
 (e^A)^{-1}=e^{-A},\qquad
 \det(e^A)=e^{\operatorname{tr}A}.
\]
由幂级数定义的矩阵正弦、余弦还满足 $\sin^2A+\cos^2A=I$、$\sin(2A)=2\sin A\cos A$。
\end{proposition}
\begin{proof}
两个指数级数绝对收敛，且双重级数的范数和受 $\bigl(\sum_i\|A^i\|_F/i!\bigr)\bigl(\sum_j\|B^j\|_F/j!\bigr)$ 控制，故可按总次数重排；交换性使同次项之和为 $\sum_{i+j=k}A^iB^j/(i!j!)=(A+B)^k/k!$，于是得到指数乘法公式。取 $B=-A$ 即得逆矩阵公式。若 $\mu_1,\ldots,\mu_n$ 为按重数计的特征值，式~\eqref{eq:xie-ex-sim-taylor} 给出 $\det(e^A)=\prod_i e^{\mu_i}=e^{\sum_i\mu_i}=e^{\operatorname{tr}A}$。最后由 $\cos A=(e^{iA}+e^{-iA})/2$、$\sin A=(e^{iA}-e^{-iA})/(2i)$，及 $iA,-iA$ 的交换性，展开相乘即得两个三角恒等式。
\end{proof}

\begin{example}
说明矩阵指数的乘法公式不能省去交换性。
\end{example}
\begin{solve}
取 $A=E_{12}$、$B=E_{21}$ 为二阶基础矩阵，则 $A^2=B^2=0$，故
\[
 e^Ae^B=\begin{pmatrix}2&1\\1&1\end{pmatrix},
 \qquad e^Be^A=\begin{pmatrix}1&1\\1&2\end{pmatrix}.
\]
它们不相等，因而不可能都等于 $e^{A+B}$。对于不同矩阵参数的标量恒等式，必须先检查其推导是否依赖乘法交换性。
\end{solve}

\subsection{矩阵幂的极限}

\begin{theorem}
\label{thm:xie-ex-sim-power-limit}
$\lim_{k\to\infty}A^k$ 存在，当且仅当每个特征值满足 $|\mu|<1$，或满足 $\mu=1$ 且属于 $1$ 的 Jordan 块全为一阶。此时极限是在 $1$ 的特征子空间上的投影，其核为其他特征值的根子空间的直和。特别地，$A^k\to0$ 当且仅当 $\rho(A)<1$。
\end{theorem}
\begin{proof}
若 $|\mu|<1$，式~\eqref{eq:xie-ex-sim-matrix-power} 中每一项都是至多固定次数的 $k$ 的多项式乘以几何衰减项，因而趋于零；$\mu=0$ 时各块的高次幂直接为零。若 $|\mu|>1$，块的对角元不收敛；若 $|\mu|=1$ 且 $\mu\ne1$，序列 $\mu^k$ 也不收敛，因为相邻两项之差的模恒为 $|\mu-1|>0$。最后，$J_r(1)^k$ 在 $r\geq2$ 时的上次对角元为 $k$，不收敛；$r=1$ 时恒为 $1$。逐块判断再作相似变换，便得到结论。若 $P$ 中先排 $1$ 的基向量，则极限为 $P\operatorname{diag}(I,0)P^{-1}$。
\end{proof}

\section{Jordan 标准型的几何}
\label{sec:xie-ex-sim-geometry}

\subsection{主分解与根子空间}

\begin{theorem}[主分解]
\label{thm:xie-ex-sim-primary}
设 $m_T=\prod_{i=1}^{t}p_i^{s_i}$，其中 $p_i$ 为 $\mathbb K[\lambda]$ 中互异的首一不可约多项式，则
\begin{equation}
 V=\bigoplus_{i=1}^{t}V_i,
 \qquad V_i=\ker\bigl(p_i(T)^{s_i}\bigr).
 \label{eq:xie-ex-sim-primary}
\end{equation}
若 $\chi_T=\prod_i p_i^{a_i}$，则 $T|_{V_i}$ 的特征多项式为 $p_i^{a_i}$，极小多项式为 $p_i^{s_i}$，且 $\dim V_i=a_i\deg p_i$。特别地，$\ker(p_i(T)^{a_i})=\ker(p_i(T)^{s_i})$。
\end{theorem}
\begin{proof}
令 $h_i=m_T/p_i^{s_i}$。中国剩余定理给出多项式 $e_i=u_ih_i$，使 $e_i\equiv1\pmod{p_i^{s_i}}$ 且 $e_i\equiv0\pmod{p_j^{s_j}}$（$j\ne i$）。由于 $\sum_i e_i\equiv1\pmod{m_T}$，有 $\sum_i e_i(T)=I$；又 $p_i^{s_i}e_i$ 是 $m_T$ 的倍式，故 $e_i(T)V\subseteq V_i$。在 $V_i$ 上 $e_i(T)=I$，在其他 $V_j$ 上它为零，因此这些算子是相互正交的代数投影，给出直和分解。

各限制变换由 $p_i^{s_i}$ 零化，所以其特征多项式只含不可约因子 $p_i$；整体特征多项式是各限制特征多项式的乘积，故指数必为 $a_i$。整体极小多项式是各限制极小多项式的最小公倍式，而各分量只含不同的 $p_i$，所以其指数必为 $s_i$。维数由特征多项式次数得出。在 $V_i$ 上 $p_i(T)^{s_i}=0$，而在其余分量上 $p_i(T)$ 可逆，故用任何不小于 $s_i$ 的指数取核都得到同一个 $V_i$，特别可取 $a_i$。
\end{proof}

在复数域上 $p_i(\lambda)=\lambda-\mu_i$，空间 $V_i$ 称为对应的\emph{根子空间（generalized eigenspace）}，也等于 $\ker(T-\mu_iI)^n$。限制变换 $N_i=T|_{V_i}-\mu_iI$ 幂零，因此 Jordan 型的几何构造分成两步：先分离不同特征值的根子空间，再在每个根子空间上构造幂零链。特征子空间 $\ker(T-\mu_iI)$ 只是根子空间中的第一级核，二者相等恰好表示该特征值的块全为一阶。

\subsection{幂零链基的构造}

\begin{proposition}
\label{prop:xie-ex-sim-nilpotent-chains}
有限维空间上的幂零变换 $N$ 存在由若干条链 $w,Nw,\ldots,N^{r-1}w$ 组成的基，且 $N^rw=0$、$N^{r-1}w\ne0$。按每条链的逆序排列后，其表示矩阵是零特征值 Jordan 块的直和。
\end{proposition}
\begin{proof}
采用第~\ref{sec:xie-ex-sim-transition} 节的 $F_k=\ker N\cap\operatorname{Im}N^{k-1}$。选取与下降链 $F_1\supseteq F_2\supseteq\cdots$ 相容的基，并把每个 $F_k$ 中相对 $F_{k+1}$ 新增的基向量 $b$ 提升为 $N^{k-1}w=b$，产生一条长度为 $k$ 的链。所有链底 $b$ 合在一起是 $\ker N$ 的一组基。

证明所得全部链向量无关：若有线性关系，按向量距链底的层数计，设其中系数非零的最高层为 $h$，作用 $N^{h-1}$ 后，仅这一层留下各链底的线性组合。链底无关使该层系数全为零，逐层下降即得所有系数为零。另一方面，映射 $N:\operatorname{Im}N^{k-1}\to\operatorname{Im}N^k$ 的核为 $F_k$，所以 $\dim F_k=\operatorname{rank}N^{k-1}-\operatorname{rank}N^k$。由幂零性及有限求和，全部链向量的数目为
\[
 \sum_{k\geq1}k(\dim F_k-\dim F_{k+1})
 =\sum_{k\geq1}\dim F_k=\dim V.
\]
因此它们确为一组基。在每条链的逆序 $N^{r-1}w,\ldots,Nw,w$ 下，$N$ 的矩阵就是 $J_r(0)$。
\end{proof}

结合主分解与幂零链基，就得到 Jordan 标准型的几何存在性证明。每个块 $J_r(\mu)$ 对应一个循环子空间，链顶为最后一个基向量。对 $N=T-\mu I$，有 $\mathcal C(N,w)=\mathcal C(T,w)$，因为 $N$ 与 $T$ 互为对方的一次多项式；Jordan 链进一步要求 $N^rw=0$。因此两种循环空间在此相同，差别在于选用使幂零作用显露出来的链基。一般有理标准型的一个循环块则可能包含多个不同特征值，不能简单地把它与一个 Jordan 块等同。

\subsection{不变子空间的结构与个数}

\begin{proposition}
\label{prop:xie-ex-sim-invariant-count-complex}
单个 Jordan 块 $J_n(\mu)$ 恰有 $n+1$ 个不变子空间，即 $\operatorname{span}(e_1,\ldots,e_k)$，$0\leq k\leq n$。若复线性变换 $T$ 是循环的，且 $\chi_T=\prod_{i=1}^{t}(\lambda-\mu_i)^{a_i}$，则其不变子空间恰有 $\prod_i(a_i+1)$ 个；若 $T$ 非循环，则有无穷多个不变子空间。
\end{proposition}
\begin{proof}
对单块令 $N=J_n(\mu)-\mu I$。任取非零不变子空间 $U$，设 $k$ 为其中向量可能出现的最高非零坐标位置，则 $U\subseteq\operatorname{span}(e_1,\ldots,e_k)$。取第 $k$ 坐标非零的 $u\in U$，向量 $u,Nu,\ldots,N^{k-1}u$ 无关且全部属于右边这个 $k$ 维空间，故 $U$ 等于它。

循环时，每个特征值只有一个 Jordan 块。不变子空间 $U$ 被主分解中的多项式投影保持，故 $U=\bigoplus_i(U\cap V_i)$；每个分量有 $a_i+1$ 种选择，乘法原理给出总数。非循环时，某个特征值至少有两个 Jordan 块，所以其特征子空间至少二维；取其中独立向量 $u,v$，直线 $\mathbb C(u+cv)$（$c\in\mathbb C$）两两不同，且均为不变子空间。
\end{proof}

\section{一般数域上的相似标准型}
\label{sec:xie-ex-sim-general-field}

\subsection{核像分解与基域问题}

\begin{proposition}[Fitting 分解]
\label{prop:xie-ex-sim-fitting}
对任意数域上的有限维线性变换 $T$，当整数 $s$ 足够大（例如 $s=n$）时，有
\[
 V=\ker T^s\oplus\operatorname{Im}T^s.
\]
两个分量均为不变子空间，$T$ 在前者上幂零，在后者上可逆。并且 $V=\ker T\oplus\operatorname{Im}T$ 当且仅当 $\operatorname{rank}T=\operatorname{rank}T^2$；若 $T$ 奇异，这也等价于 $0$ 是 $m_T$ 的单根。
\end{proposition}
\begin{proof}
有限维性使核、像链在不超过 $n$ 步后稳定。若 $v=T^su\in\ker T^s$，则 $T^{2s}u=0$，稳定性给出 $u\in\ker T^s$，故 $v=0$。两空间维数之和为 $n$，所以得到直和。$T$ 在核空间上幂零，而在稳定像空间上是满射，有限维性保证可逆。对于一级核像分解，映射 $T:\operatorname{Im}T\to\operatorname{Im}T^2$ 的核为 $\ker T\cap\operatorname{Im}T$，故交为零与两像秩相同等价。若奇异且一级分解成立，$m_T$ 是 $\lambda$ 与可逆限制变换的极小多项式的最小公倍式，因后者常数项非零，$0$ 是单根。反之若 $m_T=\lambda g$、$g(0)\ne0$，则主分解给出 $V=\ker T\oplus\ker g(T)$，且 $T$ 在后者上可逆，所以 $\ker g(T)=\operatorname{Im}T$。
\end{proof}

此证明完全留在 $\mathbb K$ 内，无需构造复 Jordan 基。另一方面，如果问题只涉及秩、特征多项式、极小多项式或两个原域矩阵是否相似，也可先在复数域中证明，再依据这些量在扩域下的不变性返回原域。\emph[caution]{某个方程在复数域中存在解，并不能据此断言它在原数域中存在解。}例如一阶实矩阵 $(-1)$ 有复平方根，却无实平方根；存在性问题必须检查构造是否仍在给定数域内。

\subsection{不可约多项式对应的广义 Jordan 块}

设 $p$ 是 $\mathbb K[\lambda]$ 中 $d$ 次首一不可约多项式，记 $F=F(p)=C(p)^{\mathsf T}$，并令 $E=e_de_1^{\mathsf T}$，即第 $(d,1)$ 元为 $1$、其余为零的 $d$ 阶矩阵。对正整数 $e$，定义两种 $de$ 阶块
\[
 G_e(p)=\begin{pmatrix}
 F&I_d&&0\\&F&\ddots&\\&&\ddots&I_d\\0&&&F
 \end{pmatrix},
 \qquad
 H_e(p)=\begin{pmatrix}
 F&E&&0\\&F&\ddots&\\&&\ddots&E\\0&&&F
 \end{pmatrix}.
\]
当 $e=1$ 时都按 $F$ 理解。当 $d=1$、$p(\lambda)=\lambda-\mu$ 时，两种块都退化为通常的 $J_e(\mu)$。

\begin{proposition}
\label{prop:xie-ex-sim-generalized-block}
$G_e(p)$ 与 $H_e(p)$ 的特征多项式、极小多项式都为 $p^e$，故二者都是循环矩阵，不变因子组均为 $1,\ldots,1,p^e$。
\end{proposition}
\begin{proof}
分块上三角结构给出特征多项式 $p^e$。先看 $G=G_e(p)$：$p(G)$ 的对角块为 $p(F)=0$，第一条上分块对角线为 $p'(F)$。数域特征为零，故不可约多项式 $p$ 与 $p'$ 互素，由 Bézout 等式知 $p'(F)$ 可逆。于是 $p(G)^e=0$，而 $e\geq2$ 时 $p(G)^{e-1}$ 的第 $(1,e)$ 块为 $p'(F)^{e-1}\ne0$，所以 $m_G=p^e$；$e=1$ 时 $m_F=p$。再看 $H=H_e(p)$：它的普通上次对角元全为 $1$，其余更高的对角线全为零，故 $\lambda I-H$ 的前 $de-1$ 行、后 $de-1$ 列构成一个主对角元全为 $-1$ 的三角子式，其值非零。当 $de>1$ 时 $D_{de-1}=1$，所以唯一非常数不变因子为 $\chi_H=p^e$；$de=1$ 时结论直接成立。
\end{proof}

若 $A$ 在 $\mathbb K$ 上的初等因子组为 $p_1^{e_1},\ldots,p_t^{e_t}$，则它在 $\mathbb K$ 上同时相似于
\[
 \operatorname{diag}(G_{e_1}(p_1),\ldots,G_{e_t}(p_t)),
 \qquad
 \operatorname{diag}(H_{e_1}(p_1),\ldots,H_{e_t}(p_t)).
\]
理由是各块恰贡献对应的一个初等因子，整体初等因子多重集合与 $A$ 相同。这就是一般数域上的\emph{广义 Jordan 标准型（generalized Jordan form）}。同一初等因子还可用 $C(p_i^{e_i})$ 表示；与这个友阵相比，广义 Jordan 块显露了同一不可约多项式对应的层次结构。

在 $\mathbb R$ 上，不可约多项式只有一次多项式和形如 $(\lambda-a)^2+b^2$（$b>0$）的二次多项式。后者对应的常见实块为
\[
 R_e(a,b)=\begin{pmatrix}
 R(a,b)&I_2&&0\\&R(a,b)&\ddots&\\&&\ddots&I_2\\0&&&R(a,b)
 \end{pmatrix},
 \qquad R(a,b)=\begin{pmatrix}a&-b\\b&a\end{pmatrix}.
\]
$F((\lambda-a)^2+b^2)$ 与 $R(a,b)$ 有相同的二次不可约特征多项式，故实相似，将同一相似变换用于各对角块，且 $I_2$ 保持不变，便得到 $G_e((\lambda-a)^2+b^2)\sim R_e(a,b)$。在复数域中，这个实块分解为 $J_e(a+ib)$ 与 $J_e(a-ib)$ 的直和；这解释了实矩阵的非实特征值及对应块为何成共轭对出现。

\subsection{分解、不变子空间与实矩阵方根}

\begin{proposition}
\label{prop:xie-ex-sim-jc-general}
每个 $A\in M_n(\mathbb K)$ 唯一分解为 $A=D+N$，其中 $D,N\in M_n(\mathbb K)$，$D$ 在复数域上可对角化，$N$ 幂零且 $DN=ND$。
\end{proposition}
\begin{proof}
在每个 $G_e(p)$ 内，将对角块组成的矩阵记为 $D_0$，将上分块次对角线上的单位矩阵组成的矩阵记为 $N_0$，则 $G_e(p)=D_0+N_0$，两者交换且 $N_0$ 幂零。不可约多项式 $p$ 在复数域内无重根，所以 $F(p)$ 在复数域上可对角化，$D_0$ 也可对角化。逐块构造后用原数域上的过渡矩阵变换回去，即得到 $D,N\in M_n(\mathbb K)$；唯一性由复数域上的定理~\ref{thm:xie-ex-sim-jc} 保证。这里的 $D$ 未必在 $\mathbb K$ 上可对角化，例如实数域上的非实共轭特征值块。
\end{proof}

\begin{proposition}
\label{prop:xie-ex-sim-invariant-count-general}
若 $T$ 是循环变换，且 $m_T=\chi_T=\prod_i p_i^{a_i}$，则每个不变子空间都唯一写成 $\ker d(T)$，其中 $d$ 遍历 $\chi_T$ 的首一因子，因而恰有 $\prod_i(a_i+1)$ 个不变子空间。若 $T$ 非循环，则在任意数域上都有无穷多个不变子空间。
\end{proposition}
\begin{proof}
取循环向量 $v$，映射 $f\mapsto f(T)v$ 把 $\mathbb K[\lambda]/(\chi_T)$ 与 $V$ 对应，乘以 $\lambda$ 对应作用 $T$。一个 $T$-不变子空间在此对应下的原像是包含 $(\chi_T)$ 的理想，故为 $(g)$，其中 $g$ 是 $\chi_T$ 的唯一首一因子。令 $d=\chi_T/g$，则 $d(T)f(T)v=0$ 当且仅当 $\chi_T\mid df$，当且仅当 $g\mid f$，所以此子空间恰为 $\ker d(T)$，并有维数 $\deg d$。数出首一因子即得有限个数。

若非循环，整除链中至少有两个非常数不变因子，它们共享某个不可约因子 $p$。在两个相应的循环分量内，分别存在一个与 $\mathbb K[\lambda]/(p)$ 同构的非零不变子空间 $U_1,U_2$：在 $\mathbb K[\lambda]/(q)$ 中可取由 $q/p$ 生成的子空间。取与 $T$ 交换的模同构 $\varphi:U_1\to U_2$，则 $\{u+c\varphi(u):u\in U_1\}$ 随 $c\in\mathbb K$ 给出两两不同的不变子空间。数域为无限域，故有无穷多个。
\end{proof}

\begin{proposition}
\label{prop:xie-ex-sim-real-root}
设 $A$ 为可逆实矩阵，$m$ 为正整数。若 $m$ 为奇数，则存在实矩阵 $B$ 满足 $B^m=A$；若 $m$ 为偶数，则存在这样的实矩阵 $B$，当且仅当对每个负实特征值，属于它的每种阶数的 Jordan 块均出现偶数次。
\end{proposition}
\begin{proof}
正实特征值可以选取实数 $m$ 次方根，非实共轭特征值可以选取互为共轭的方根；用命题~\ref{prop:xie-ex-sim-complex-root} 的有限展开，可分别得到实块或可组成实矩阵的共轭块。负实特征值在 $m$ 为奇数时也有实数方根。若 $m$ 为偶数而同阶负特征值块成对出现，就给这两个块选取共轭的非实方根，仍可组成实矩阵，证明充分性。

反之，若 $B$ 为实矩阵且 $B^m=A$，当 $m$ 为偶数时，$B$ 的实特征值的 $m$ 次幂不可能是负数；所有映为某负特征值的 $B$ 的特征值都非实，并成共轭对。由于 $B$ 可逆，$z^m$ 在这些特征值处的导数非零，故取幂保留 Jordan 块阶数，每一对共轭块贡献两个同阶的负特征值块。这给出必要性。
\end{proof}

\section{基础训练与方法整合}
\label{sec:xie-ex-sim-training}

基础训练应检验“由哪些数据能恢复全部块”以及“每种操作在哪个环、哪个数域中合法”。下列题解归纳原章训练中的主要方法，不逐项重抄选择题选项。对于具体矩阵，先检查总阶数、特征多项式、最大块与几何重数是否一致；对含参数题，还要检查参数为零或特征值碰合时是否产生新的分支。

\begin{example}
设某八阶矩阵在复数域上的初等因子为 $(\lambda-1)^2,(\lambda-1)^2,(\lambda+1)^3,\lambda$，求不变因子及 Jordan 型。
\end{example}
\begin{solve}
按不可约因子分别排列指数，$\lambda-1$ 对应 $(2,2)$，$\lambda+1$ 对应 $(0,3)$，$\lambda$ 对应 $(0,1)$；逐列相乘，得到非常数不变因子
\[
 q_1=(\lambda-1)^2,\qquad
 q_2=\lambda(\lambda-1)^2(\lambda+1)^3.
\]
加上六个 $1$，即为完整的八项不变因子。Jordan 型为 $\operatorname{diag}(J_2(1),J_2(1),J_3(-1),0)$，共四块，而有理标准型只有两个块。次数检验为 $2+2+3+1=8$。
\end{solve}

\begin{example}
设 $A$ 为 $n$ 阶实矩阵，$n\geq2$，且 $\deg D_{n-1}=n-2$，描述其不变因子组。
\end{example}
\begin{solve}
由 $m_A=D_n/D_{n-1}$ 得 $\deg m_A=2$。若 $m_A$ 在 $\mathbb R$ 上不可约，则每个非常数不变因子都是 $m_A$，故 $n$ 为偶数，有 $n/2$ 个这样的因子，其余为 $1$。若 $m_A=(\lambda-a)(\lambda-b)$，允许 $a=b$，则整除链中的非常数因子只能先是若干个同一个一次因子 $\lambda-c$（$c=a$ 或 $b$），再是若干个 $m_A$；不能在链中同时出现两个不同且互不整除的一次因子。设一次因子出现 $u$ 次、二次因子出现 $v$ 次，则 $u\geq0$、$v\geq1$、$u+2v=n$，再补 $n-u-v$ 个 $1$。这些形式也都使 $\deg D_{n-1}=n-2$，故分类完整。
\end{solve}

\begin{example}
证明任意 $A\in M_n(\mathbb K)$ 都可写成一个纯量矩阵与两个幂零矩阵之和。
\end{example}
\begin{solve}
令 $c=\operatorname{tr}A/n$，则 $A-cI_n$ 迹为零。命题~\ref{prop:xie-ex-sim-commutator} 的证明已说明，可将它相似化为主对角元全为零的矩阵 $M$。将 $M$ 分成严格上三角部分 $M_1$ 与严格下三角部分 $M_2$，二者均幂零；再变换回原基，就得到 $A=cI_n+N_1+N_2$。两个幂零部分一般不交换，故不能由此推断它们的和也幂零。
\end{solve}

\begin{example}
设 $A,B$ 都与同一循环矩阵 $T$ 交换，$v$ 是循环向量，证明 $A=B$ 当且仅当 $Av=Bv$。并说明相似矩阵是否一定能用同一个过渡矩阵同时化为 Jordan 型。
\end{example}
\begin{solve}
由 $Av=Bv$ 及交换性，得 $AT^jv=T^jAv=T^jBv=BT^jv$，$0\leq j<n$；这些向量构成一组基，故 $A=B$。后一问题的答案是否定的：取 $n\geq2$、$A=J_n(0)$、$B=2J_n(0)$，二者都相似于唯一的标准块 $J_n(0)$。若同一个 $P$ 同时把两者化为该块，就会有 $P^{-1}BP=2P^{-1}AP=2J_n(0)=J_n(0)$，矛盾。相同的标准型保证分别存在换基矩阵，不保证换基矩阵相同。
\end{solve}

\begin{example}
说明多项式矩阵乘积的次数为何可能降低，并证明对固定 $A$，表示 $M(\lambda)=(\lambda I-A)Q(\lambda)+R$ 中的常数余项与商都唯一。
\end{example}
\begin{solve}
若最高次系数的乘积为零，乘积中的最高次项就消失。例如 $M(\lambda)=\lambda E_{11}$、$N(\lambda)=\lambda E_{22}$ 是二阶一次矩阵多项式，但乘积为零矩阵。对第二问，若有两个表示，相减得 $(\lambda I-A)(Q_1-Q_2)=R_2-R_1$。若 $Q_1-Q_2\ne0$，左边的次数比它高一，因为 $\lambda I-A$ 的最高次系数为单位矩阵，不会消去首项；但右边是常数矩阵，矛盾。因此两商相等，余项也相等。这种保留乘法次序的带余除法正是命题~\ref{prop:xie-ex-sim-polynomial-transition} 的基础。
\end{solve}

\begin{example}
设 $g\in\mathbb C[z]$ 是非常数多项式。若对 $A$ 的每个特征值 $\mu$，方程 $g(z)=\mu$ 都有一个根 $\nu$ 满足 $g'(\nu)\ne0$，证明存在复矩阵 $B$ 使 $g(B)=A$。
\end{example}
\begin{solve}
对 $A$ 的每个块 $J_r(\mu)$ 选取这样的 $\nu$。有限 Taylor 公式给出 $g(J_r(\nu))-\mu I=N_rH(N_r)$，且 $H(0)=g'(\nu)\ne0$，所以其幂的秩与 $N_r$ 一样，$g(J_r(\nu))$ 相似于 $J_r(\mu)$。于是适当相似变换 $J_r(\nu)$ 就得到一个 $B_r$ 使 $g(B_r)=J_r(\mu)$；把这些块合并并返回原基即可。原训练中条件 $\gcd(\chi_A(g(z)),g'(z))=1$ 保证每个原像点都是这样的单根，所以是一个充分条件；实际构造只需每个特征值至少有一个导数非零的原像。
\end{solve}

\begin{example}
证明每个可逆复矩阵 $A$ 都有复矩阵对数，即存在 $L$ 使 $e^L=A$。
\end{example}
\begin{solve}
对每个 $J_r(\mu)$，取复数 $\ell$ 使 $e^{\ell}=\mu$，并令
\[
 L_r=\ell I_r+\sum_{j=1}^{r-1}\frac{(-1)^{j+1}}{j}\mu^{-j}N_r^j.
\]
空和按零理解。形式幂级数 $\exp(\log(1+z))=1+z$ 在模 $z^r$ 意义下成立，故代入 $\mu^{-1}N_r$，结合可交换的指数乘法公式得 $e^{L_r}=\mu I_r+N_r$。各块直和并作相似变换便得到 $L$。不同的标量对数分支可能给出不同的矩阵对数，这里只证明存在性。
\end{solve}

复习时可把各类信息归于同一条链：Smith 形给出不变因子，不可约分解给出初等因子，复数域上的一次幂因子给出 Jordan 块，块的阶数决定核、像和秩序列，而链基给出实际过渡矩阵。矩阵函数只需在每个特征值处保留有限阶导数；数域上的存在性问题则还必须检查所用特征值、方根及过渡矩阵是否位于指定数域。这样既能从代数不变量读取结构，也能用循环向量与空间分解核查计算结果。

% BEGIN textbook-xie-exercises-ch8
% 教材来源：《高等代数习题（谢启鸿）》第8章，印刷页435--485；高代习题-8.pdf。
% 按原书九节主题整合学习笔记；例题为选讲，非训练题逐题誊录。
\chapter{二次型}
\label{chap:xie-ex-quadratic-forms}

二次型将二次齐次多项式与对称矩阵联系起来：对变量作可逆线性替换，对应于对系数矩阵作合同变换。研究的中心是消去交叉项，再利用标准型说明哪些量不随变量替换而改变。本章依照教材第八章的主题，先建立合同、惯性和正定性的基本概念，再以对称分块消元为线索，讨论归纳证明、多变元计算以及正定、半正定矩阵的性质。计算中既要判断表达式的正负，也要检查变量替换是否可逆；证明中则要区分合同不变量与仅在相似变换下保持的量。

\section{基本概念}
\label{sec:xie-ex-qf-basics}

\subsection{二次型、合同与对角化}

本章的数域 $\mathbb F$ 为 $\mathbb C$ 的子域，例如 $\mathbb Q$、$\mathbb R$ 或 $\mathbb C$ 本身；$x$、$y$ 均表示列向量，$A^{\mathsf T}$ 表示转置。实矩阵的正定性与惯性指数均在实数域上讨论，复矩阵的相应概念另用共轭转置定义。记 $I_k$ 为 $k$ 阶单位矩阵，零矩阵的尺寸由上下文确定，$\|x\|^2=x^{\mathsf T}x$ 为实向量的欧几里得范数平方。

\begin{definition}
数域 $\mathbb F$ 上的 $n$ 元\emph{二次型（quadratic form）}是二次齐次多项式
\[
 f(x)=\sum_{i=1}^n a_{ii}x_i^2+2\sum_{1\le i<j\le n}a_{ij}x_ix_j
     =x^{\mathsf T}Ax,
 \qquad A=A^{\mathsf T}\in\mathbb F^{n\times n}.
\]
矩阵 $A$ 称为它的系数矩阵或相伴矩阵。交叉项 $x_ix_j$ 的系数是 $2a_{ij}$，故由多项式写矩阵时应将交叉项系数除以 $2$。
\end{definition}

对称性保证系数矩阵唯一。事实上，若对称矩阵 $H$ 满足 $x^{\mathsf T}Hx=0$ 对所有 $x$ 成立，依次取标准单位向量 $e_i$ 和 $e_i+e_j$，得到 $h_{ii}=0$ 和 $2h_{ij}=0$，于是 $H=0$。若允许系数矩阵不对称，则 $x^{\mathsf T}Mx=x^{\mathsf T}(M+M^{\mathsf T})x/2$，其反对称部分不影响二次型，矩阵表示便不再唯一。

\begin{definition}
设 $A,B\in\mathbb F^{n\times n}$。若存在可逆矩阵 $C$ 使
\begin{equation}
 B=C^{\mathsf T}AC,
 \label{eq:xie-ex-qf-congruence}
\end{equation}
则称 $A$ 与 $B$ 在 $\mathbb F$ 上\emph{合同（congruent）}，记为 $A\sim_{\mathrm c}B$。
\end{definition}

将 $x=Cy$ 代入二次型，恰得 $x^{\mathsf T}Ax=y^{\mathsf T}(C^{\mathsf T}AC)y$。因此合同描述同一二次型在不同坐标下的系数矩阵。合同是等价关系，并保持秩；由
\[
 \det B=(\det C)^2\det A
\]
可知，实合同还保持行列式的符号，但一般不保持行列式的数值、迹或特征值。相似变换为 $C^{-1}AC$，与合同的表达式不同；当 $C$ 正交时，两者才在这一变换上重合。

\begin{theorem}
\label{thm:xie-ex-qf-diagonalization}
数域 $\mathbb F$ 上的每个对称矩阵均合同于对角矩阵，因而每个二次型均可经可逆线性替换化为只含平方项的\emph{标准型（diagonal form）}。
\end{theorem}
\begin{proof}
对矩阵阶数归纳。零矩阵已经对角化。若某个对角元非零，先同时置换行、列，将其移至首位，写成
\[
 A=\begin{pmatrix}a&b^{\mathsf T}\\b&D\end{pmatrix},\qquad a\ne0.
\]
配方给出
\[
 x^{\mathsf T}Ax
 =a\left(x_1+a^{-1}b^{\mathsf T}x'\right)^2
  +(x')^{\mathsf T}(D-a^{-1}bb^{\mathsf T})x',
 \qquad x'=(x_2,\ldots,x_n)^{\mathsf T}.
\]
所用替换为可逆的三角替换，余下系数矩阵仍对称，可用归纳假设。若所有对角元为零但 $A\ne0$，取 $a_{ij}\ne0$，则 $(e_i+e_j)^{\mathsf T}A(e_i+e_j)=2a_{ij}\ne0$；将 $e_i+e_j$ 扩充为一组基后，就得到非零对角元，归入前一种情况。
\end{proof}

\subsection{惯性定理与正负定性}

\begin{theorem}[Sylvester 惯性定理]
\label{thm:xie-ex-qf-inertia}
任意 $n$ 阶实对称矩阵均合同于
\begin{equation}
 \operatorname{diag}(I_p,-I_q,0_{n-p-q}),\qquad p,q\ge0,\quad p+q\le n,
 \label{eq:xie-ex-qf-inertia-form}
\end{equation}
其中 $p,q$ 由原矩阵唯一确定，分别称为正、负惯性指数，记为 $p(A),q(A)$。相应二次型称为\emph{规范标准型}；其秩为 $r=p+q$，符号差为 $s=p-q$。
\end{theorem}
\begin{proof}
先由定理~\ref{thm:xie-ex-qf-diagonalization} 对角化，再将非零对角系数的绝对值开平方并吸收到变量中，得到式~\eqref{eq:xie-ex-qf-inertia-form}。为证唯一性，说明 $p$ 等于二次型在其上为正定的子空间的最大维数。标准型的前 $p$ 个坐标张成一个正定子空间；若某子空间维数超过 $p$，它与由后 $n-p$ 个坐标张成的非正子空间必有非零交向量，故不可能正定。这一最大维数在可逆坐标替换下不变，因而 $p$ 唯一；对二次型取负，同理得到 $q$ 唯一。
\end{proof}

实对称矩阵合同的充要条件是阶数及正、负惯性指数相同。在复数域上，每个非零数都有平方根，复对称矩阵的合同标准型为 $\operatorname{diag}(I_r,0)$，只由秩决定；这一结论不能直接搬到有理数域，也不能与下文 Hermite 矩阵的复相合混同。

\begin{definition}
实对称矩阵 $A$ 及二次型 $x^{\mathsf T}Ax$ 称为\emph{正定（positive definite）}，若对一切 $x\ne0$ 有 $x^{\mathsf T}Ax>0$；称为\emph{半正定（positive semidefinite）}，若对一切 $x$ 有 $x^{\mathsf T}Ax\ge0$，分别记为 $A\succ0$、$A\succeq0$。负定、半负定通过 $-A\succ0$、$-A\succeq0$ 定义；若二次型既取正值又取负值，则称为不定型。
\end{definition}

由惯性定理，正定等价于 $p=n$，半正定等价于 $q=0$；负定等价于 $q=n$，半负定等价于 $p=0$；不定等价于 $p,q$ 均大于零。半正定允许零矩阵与正定矩阵在内，不能把“半正定”理解成“恰好有零方向”。若 $x^{\mathsf T}Ax$ 对每个非零 $x$ 均不为零，则它必正定或负定：退化标准型在零坐标方向取零，不定标准型在一正一负坐标均取 $1$ 时取零，两者都被假设排除。

\subsection{Hermite 型与复相合}

沿用前文记号，对复矩阵用 $A^{\mathsf H}=\overline A^{\mathsf T}$ 表示共轭转置，用 $\operatorname{adj}(A)$ 表示伴随矩阵；二者是不同的运算。对复向量定义 $\|z\|^2=z^{\mathsf H}z$。
\begin{definition}
若 $A=A^{\mathsf H}\in\mathbb C^{n\times n}$，则称 $A$ 为\emph{厄米矩阵（Hermitian matrix，亦称 Hermite 矩阵）}，实值函数
\[
 h(z)=z^{\mathsf H}Az=\sum_{i,j=1}^n\overline z_i a_{ij}z_j
\]
称为 Hermite 型。若 $B=C^{\mathsf H}AC$ 且 $C$ 为可逆复矩阵，则称 $A,B$ \emph{复相合（$*$-congruent）}。
\end{definition}

Hermite 型也有唯一的规范标准型 $\operatorname{diag}(I_p,-I_q,0)$。其对角化仍可通过配方归纳完成：非零 Hermite 矩阵总有 $v$ 使 $v^{\mathsf H}Av\ne0$，因为若所有对角元为零，依次用 $e_i+e_j$ 和 $e_i+\mathrm i e_j$ 检验，就可检测非零非对角元的实部与虚部；对角化后各对角系数为实数，再按绝对值开平方缩放。惯性唯一性的子空间维数证明同样成立，但子空间维数取复维数。正定、半正定的定义改为对所有非零复列向量检验 $z^{\mathsf H}Az$；实对称正定矩阵在这一意义下仍然正定，因为 $z=u+\mathrm i v$ 时
\[
 z^{\mathsf H}Az=u^{\mathsf T}Au+v^{\mathsf T}Av.
\]

\section{对称初等变换与矩阵合同}
\label{sec:xie-ex-qf-block-congruence}

对矩阵作初等行变换 $A\mapsto EA$ 后，再作对应的列变换 $EA\mapsto EAE^{\mathsf T}$，所得矩阵与 $A$ 合同，并保持对称性。对应操作包括同时对换两行、两列，同时用同一非零数乘某行、某列，以及将某行的倍数加到另一行后再作对应列加法。分块时，左乘块矩阵必须配合右乘其转置，不能把矩阵乘法次序当作标量运算处理。若要同时求出变量替换矩阵，应记录总乘子，使最后结果写成 $C^{\mathsf T}AC$。

\begin{theorem}[对称分块消元]
\label{thm:xie-ex-qf-schur}
设
\[
 M=\begin{pmatrix}A&B\\B^{\mathsf T}&D\end{pmatrix},\qquad
 A=A^{\mathsf T}\in\mathbb F^{r\times r},\quad
 D=D^{\mathsf T}\in\mathbb F^{s\times s},\quad B\in\mathbb F^{r\times s}.
\]
若 $A$ 可逆，则
\begin{equation}
 C^{\mathsf T}MC=
 \begin{pmatrix}A&0\\0&D-B^{\mathsf T}A^{-1}B\end{pmatrix},
 \qquad C=\begin{pmatrix}I_r&-A^{-1}B\\0&I_s\end{pmatrix}.
 \label{eq:xie-ex-qf-schur}
\end{equation}
矩阵 $D-B^{\mathsf T}A^{-1}B$ 称为关于 $A$ 的\emph{Schur 补（Schur complement）}。在实对称情形下，
\begin{align}
 p(M)&=p(A)+p(D-B^{\mathsf T}A^{-1}B),\\
 q(M)&=q(A)+q(D-B^{\mathsf T}A^{-1}B),
 \label{eq:xie-ex-qf-schur-inertia}
\end{align}
并且 $\det M=\det A\det(D-B^{\mathsf T}A^{-1}B)$。
\end{theorem}
\begin{proof}
直接分块相乘即得式~\eqref{eq:xie-ex-qf-schur}，其中 $C$ 可逆且行列式为 $1$。分块对角矩阵可分别对各块进行合同对角化，所以其正、负惯性指数分别是各块相应指数之和；同时取行列式便得到最后一个等式。
\end{proof}

若 $D$ 也可逆，换用 $D$ 消元，可得
\[
 p(A)+p(D-B^{\mathsf T}A^{-1}B)
 =p(D)+p(A-BD^{-1}B^{\mathsf T}),
\]
负惯性指数亦有相同形式。这就是教材中正、负惯性指数降阶公式的实质。特别地，$M\succ0$ 当且仅当 $A\succ0$ 且 $D-B^{\mathsf T}A^{-1}B\succ0$；若已知 $A\succ0$，则 $M\succeq0$ 当且仅当该 Schur 补半正定。\emph[caution]{当主对角块奇异时不能写出其逆矩阵；半正定情形下的消元需要额外检验非对角块的列空间条件。}

\begin{example}
\label{ex:xie-ex-qf-reflection}
设 $u\in\mathbb R^n$ 且 $u^{\mathsf T}u=1$，求 $I_n-2uu^{\mathsf T}$ 的正、负惯性指数。
\end{example}
\begin{solve}
考虑 $M=\begin{pmatrix}I_n&\sqrt2u\\\sqrt2u^{\mathsf T}&1\end{pmatrix}$。用 $I_n$ 消元，得 $M\sim_{\mathrm c}\operatorname{diag}(I_n,-1)$；用右下角的 $1$ 消元，得 $M\sim_{\mathrm c}\operatorname{diag}(I_n-2uu^{\mathsf T},1)$。比较惯性指数，得到 $p=n-1$、$q=1$。也可将 $u$ 扩充成正交基：矩阵在 $\mathbb Ru$ 上作用为 $-1$，在 $u^\perp$ 上作用为 $1$。
\end{solve}

\begin{example}
设 $B\in\mathbb R^{n\times n}$ 可逆，求 $H=\begin{pmatrix}0&B\\B^{\mathsf T}&0\end{pmatrix}$ 的惯性；另设 $A\succ0$，判断 $K=\begin{pmatrix}A&-I_n\\-I_n&A^{-1}\end{pmatrix}$ 的正负定性。
\end{example}
\begin{solve}
对 $H$ 用 $\operatorname{diag}(B^{-\mathsf T},I_n)$ 作合同变换，化为 $\begin{pmatrix}0&I_n\\I_n&0\end{pmatrix}$；再令 $u=(y+z)/\sqrt2$、$v=(y-z)/\sqrt2$，二次型 $2u^{\mathsf T}v$ 化为 $y^{\mathsf T}y-z^{\mathsf T}z$，故 $p(H)=q(H)=n$。对 $K$ 用 $A$ 消元，其 Schur 补为零，故 $K\sim_{\mathrm c}\operatorname{diag}(A,0)$，从而 $K\succeq0$ 且秩为 $n$。
\end{solve}

同时置换行、列还可将任意主子阵移到左上角。若 $A\succ0$ 或 $A\succeq0$，把主子阵上的向量用零补足成完整向量，就知每个主子阵也相应正定或半正定。由二阶主子阵可得
\begin{equation}
 |a_{ij}|^2\le a_{ii}a_{jj}\quad(A\succeq0),\qquad
 |a_{ij}|^2<a_{ii}a_{jj}\quad(A\succ0,\ i\ne j).
 \label{eq:xie-ex-qf-entry-bound}
\end{equation}
因此正定矩阵的元素绝对值的最大值只能在主对角线上取得。更一般地，若对称或反对称矩阵秩为 $r$，它总有非零的 $r$ 阶主子式：同时置换行、列，使前 $r$ 行构成行空间的一组基，则后面的行都是前面行的线性组合；前 $r$ 列也因对称或反对称而线性无关，所以它们在前 $r$ 行上的限制仍线性无关，否则该限制的一个非零线性关系会扩展为全列的关系。

\section{归纳法的应用}
\label{sec:xie-ex-qf-induction}

\subsection{顺序主子式与三角分解}

记 $A_k=A[1\!:\!k,1\!:\!k]$ 为前 $k$ 行、列构成的顺序主子阵，$\Delta_k=\det A_k$ 为\emph{顺序主子式（leading principal minor）}，并约定 $\Delta_0=1$。Schur 补将一个问题分成较低阶的问题，因此适合同时证明判据和构造分解。

\begin{theorem}[Jacobi 对角化与 Sylvester 判据]
\label{thm:xie-ex-qf-jacobi}
设 $A$ 为 $n$ 阶实对称矩阵，若 $\Delta_1,\ldots,\Delta_{n-1}$ 均非零，则存在主对角元均为 $1$ 的上三角矩阵 $U$，使
\begin{equation}
 A=U^{\mathsf T}\operatorname{diag}(d_1,\ldots,d_n)U,
 \qquad d_k=\frac{\Delta_k}{\Delta_{k-1}}.
 \label{eq:xie-ex-qf-jacobi}
\end{equation}
特别地，$A\succ0$ 当且仅当 $\Delta_k>0$ 对每个 $1\le k\le n$ 成立。
\end{theorem}
\begin{proof}
当 $n=1$ 时直接成立。对 $n>1$，写 $A=\begin{pmatrix}A_{n-1}&b\\b^{\mathsf T}&a\end{pmatrix}$，用式~\eqref{eq:xie-ex-qf-schur} 消元得到 $\operatorname{diag}(A_{n-1},a-b^{\mathsf T}A_{n-1}^{-1}b)$。归纳假设处理第一块，而行列式等式给出末项为 $\Delta_n/\Delta_{n-1}$；消元矩阵及归纳中的变换都是单位上三角矩阵，故合并后得到所述分解。若所有 $\Delta_k>0$，每个 $d_k>0$，合同标准型表明 $A\succ0$。反之，正定矩阵的各顺序主子阵均正定，其行列式由惯性定理为正。
\end{proof}

若某个中间的 $\Delta_k=0$，比值公式便不适用，不能据此判断矩阵退化；应改用对称置换或一个可逆的二阶主块消元。例如 $\begin{pmatrix}0&1\\1&0\end{pmatrix}$ 的一阶顺序主子式为零，但矩阵可逆。

\begin{theorem}[正定矩阵的 Cholesky 分解]
\label{thm:xie-ex-qf-cholesky}
对实对称矩阵 $A$，下列条件等价：$A\succ0$；$A=U^{\mathsf T}DU$，其中 $U$ 为单位上三角矩阵、$D$ 为正对角矩阵；$A=R^{\mathsf T}R$，其中 $R$ 为主对角元均为正的上三角矩阵。后两种分解分别唯一，后一分解称为\emph{Cholesky 分解}。
\end{theorem}
\begin{proof}
存在性由定理~\ref{thm:xie-ex-qf-jacobi} 及 $R=D^{1/2}U$ 得到；反之，$R$ 可逆，$x^{\mathsf T}Ax=\|Rx\|^2>0$。若 $R^{\mathsf T}R=S^{\mathsf T}S$，则 $Q=RS^{-1}$ 既上三角又正交。其逆矩阵既上三角又等于 $Q^{\mathsf T}$，所以 $Q$ 为对角矩阵；正的对角元再给出 $Q=I_n$。于是 $R=S$，而 $D$ 及 $U$ 分别由 $R$ 的对角元平方和行缩放唯一确定。
\end{proof}

\begin{example}
设 $A\succ0$、$A^{-1}=(b_{ij})$。证明 $a_{ii}b_{ii}\ge1$，并说明何时等号成立；若 $A$ 的所有非对角元都严格为负，证明 $A^{-1}$ 的所有元素都严格为正。
\end{example}
\begin{solve}
同时置换行、列可将所考察的 $a_{ii}$ 移至末位，写 $A=\begin{pmatrix}H&u\\u^{\mathsf T}&a\end{pmatrix}$。令 $s=a-u^{\mathsf T}H^{-1}u>0$，分块求逆得
\begin{equation}
 A^{-1}=\begin{pmatrix}
 H^{-1}+s^{-1}H^{-1}uu^{\mathsf T}H^{-1}&-s^{-1}H^{-1}u\\
 -s^{-1}u^{\mathsf T}H^{-1}&s^{-1}
 \end{pmatrix}.
 \label{eq:xie-ex-qf-block-inverse}
\end{equation}
所以 $ab_{nn}=a/s\ge1$，等号当且仅当 $u=0$，即原来第 $i$ 行、列除对角元外全为零。对于第二问，对阶数归纳；$n=1$ 时成立。若 $n>1$，归纳假设给出 $H^{-1}$ 的每个元素为正，$u$ 的每个分量为负，故 $H^{-1}u$ 的每个分量为负。式~\eqref{eq:xie-ex-qf-block-inverse} 的各块因而都逐项为正。
\end{solve}

\subsection{反对称矩阵的合同标准型}

\begin{theorem}
\label{thm:xie-ex-qf-skew-form}
数域 $\mathbb F$ 上每个反对称矩阵 $S^{\mathsf T}=-S$ 均合同于
\[
 \operatorname{diag}(J,\ldots,J,0),\qquad
 J=\begin{pmatrix}0&1\\-1&0\end{pmatrix}.
\]
因此反对称矩阵的秩为偶数；实反对称矩阵的行列式非负。
\end{theorem}
\begin{proof}
若 $S\ne0$，将一个非零非对角元通过对称置换移到 $(1,2)$，再同时缩放相应行、列，将首个二阶主块化为 $J$。此时
\[
 S=\begin{pmatrix}J&B\\-B^{\mathsf T}&D\end{pmatrix}
 \sim_{\mathrm c}\operatorname{diag}(J,D+B^{\mathsf T}J^{-1}B).
\]
这一等式由与式~\eqref{eq:xie-ex-qf-schur} 相同的分块乘法验证，余块仍反对称，故每次将阶数降低 $2$，直至余块为零。每个 $J$ 的秩为 $2$、行列式为 $1$，于是秩为偶数。在实数域上，奇异时行列式为零，可逆时由合同等式得 $(\det C)^2\det S=1$，故 $\det S>0$。
\end{proof}

\begin{corollary}
\label{cor:xie-ex-qf-skew-det}
若 $S$ 为 $n$ 阶实反对称矩阵，则
\[
 \det(I_n+S)\ge1+\det S\ge1.
\]
第一处等号当且仅当 $n\le2$ 或 $S=0$；$\det(I_n+S)=1$ 当且仅当 $S=0$。
\end{corollary}
\begin{proof}
按主子式展开，$\det(I_n+S)$ 等于常数 $1$ 加上所有非空主子式之和。各主子阵仍反对称，其行列式均非负。当 $n\ge3$，第一处等号要求所有二阶主子式 $s_{ij}^2$ 为零，故 $S=0$；$n\le2$ 可直接验证。第二个等号判定也由所有二阶主子式必须为零得到，$n=1$ 时反对称矩阵本来就是零矩阵。
\end{proof}

\section{合同标准型的应用}
\label{sec:xie-ex-qf-standard-applications}

\subsection{分解、逆矩阵与伴随矩阵}

标准型的应用应分为两步：先在对角矩阵上验证，再说明条件和结论怎样随合同变换转回。不能仅因原矩阵可对角化，就把不在合同下保持的乘积、迹或特征值一并替换。

\begin{proposition}
\label{prop:xie-ex-qf-gram}
秩为 $r$ 的对称矩阵可以写成 $r$ 个秩一对称矩阵之和。复对称矩阵总可写成 $A=T^{\mathsf T}T$，且 $\operatorname{rank}T=\operatorname{rank}A$。对于实对称矩阵，
\[
 A\succeq0\ \Longleftrightarrow\ A=T^{\mathsf T}T\text{，其中 $T$ 为某实矩阵};
 \qquad
 A\succ0\ \Longleftrightarrow\ A=T^{\mathsf T}T\text{，其中 $T$ 为可逆实方阵}.
\]
\end{proposition}
\begin{proof}
由合同对角化可写 $A=C^{\mathsf T}\operatorname{diag}(d_1,\ldots,d_r,0)C$，其中 $d_i\ne0$。若 $c_i^{\mathsf T}$ 为 $C$ 的第 $i$ 行，则 $A=\sum_{i=1}^r d_ic_ic_i^{\mathsf T}$，每项秩为一。在复数域中选取 $t_i^2=d_i$，令 $T=\operatorname{diag}(t_1,\ldots,t_r,0)C$ 即得分解；在实半正定情形可令前 $r$ 个 $d_i$ 均为 $1$。反之，实 Gram 形式满足 $x^{\mathsf T}T^{\mathsf T}Tx=\|Tx\|^2\ge0$，若 $T$ 可逆则对非零 $x$ 严格为正。
\end{proof}

在实数域上 $\ker(T^{\mathsf T}T)=\ker T$，故二者秩相同；在复数域中若只有转置而没有共轭，这一秩等式对任意 $T$ 未必成立，上述复对称分解的秩结论来自具体构造。若 $A\succeq0$ 且秩为 $r>0$，还可删去构造中的零行，使 $T$ 为满行秩的 $r\times n$ 矩阵；$r=0$ 时就是 $A=0$。

若 $A$ 为可逆实对称矩阵，则 $A^{-1}= (A^{-1})^{\mathsf T}A A^{-1}$，所以 $A^{-1}\sim_{\mathrm c}A$，具有相同惯性。因而 $A\succ0$ 时 $A^{-1}$ 和 $\operatorname{adj}(A)=\det(A)A^{-1}$ 都正定；若 $A\prec0$，则 $A^{-1}\prec0$，且伴随矩阵在 $n$ 为奇数时正定、在 $n$ 为偶数时负定。对 $A\succeq0$，先对 $A+\varepsilon I_n\succ0$ 使用上述结论，再令 $\varepsilon\downarrow0$；伴随矩阵的各元素是原矩阵元素的多项式，故极限给出 $\operatorname{adj}(A)\succeq0$。

复对称矩阵的分解还说明，每个复方阵都相似于某个复对称矩阵。为说明这一点，可先取该矩阵的 Jordan 标准形 $J$；对每个 Jordan 块选取反序置换矩阵，组成分块对角矩阵 $E$，则 $E=E^{\mathsf T}=E^{-1}$ 且 $JE$ 对称。写 $E=T^{\mathsf T}T$，其中 $T$ 可逆，于是
\[
 TJT^{-1}=T(JE)T^{\mathsf T}
\]
为复对称矩阵。这里应用的是复对称分解，不能据此宣称任意实方阵都相似于实对称矩阵。

\subsection{二次型不等式与迹判据}

\begin{proposition}
\label{prop:xie-ex-qf-weighted-inequalities}
若 $A\succ0$，$u,v\in\mathbb R^n$，则
\begin{align}
 u^{\mathsf T}Au+v^{\mathsf T}A^{-1}v&\ge2u^{\mathsf T}v,
 \label{eq:xie-ex-qf-young}\\
 (u^{\mathsf T}v)^2&\le(u^{\mathsf T}Au)(v^{\mathsf T}A^{-1}v).
 \label{eq:xie-ex-qf-cauchy}
\end{align}
前式等号当且仅当 $Au=v$，后式等号当且仅当 $Au$ 与 $v$ 线性相关。
\end{proposition}
\begin{proof}
写 $A=R^{\mathsf T}R$，则式~\eqref{eq:xie-ex-qf-young} 两边之差为 $\|Ru-R^{-\mathsf T}v\|^2$，等号条件为 $Ru=R^{-\mathsf T}v$。对向量 $Ru$ 与 $R^{-\mathsf T}v$ 使用 Cauchy--Schwarz 不等式，便得到式~\eqref{eq:xie-ex-qf-cauchy} 及其等号条件。
\end{proof}

\begin{proposition}
\label{prop:xie-ex-qf-trace}
设 $A$ 为实对称矩阵，则 $A\succeq0$ 当且仅当对每个实对称 $B\succeq0$ 有 $\operatorname{tr}(AB)\ge0$。若另外已知 $A$ 可逆，则 $A\succ0$ 当且仅当对每个 $B\succ0$ 有 $\operatorname{tr}(AB)>0$。当 $A,B\succeq0$ 时，还有
\begin{equation}
 \operatorname{tr}(AB)=0\quad\Longleftrightarrow\quad AB=0.
 \label{eq:xie-ex-qf-trace-zero}
\end{equation}
\end{proposition}
\begin{proof}
若 $A=C^{\mathsf T}C$、$B=D^{\mathsf T}D$，由迹的循环性质，
\[
 \operatorname{tr}(AB)=\operatorname{tr}(CD^{\mathsf T}DC^{\mathsf T})
 =\|DC^{\mathsf T}\|_{\mathrm F}^2\ge0,
\]
其中 $\|X\|_{\mathrm F}^2=\sum_{i,j}x_{ij}^2$。等号时 $DC^{\mathsf T}=0$，其转置也为零，所以 $AB=C^{\mathsf T}(CD^{\mathsf T})D=0$；反向直接成立。若对所有 $B\succeq0$ 满足迹条件，取 $B=xx^{\mathsf T}$ 即得 $x^{\mathsf T}Ax\ge0$。对于正定判据，必要性由 $CBC^{\mathsf T}\succ0$ 的迹为正得到，其中 $C$ 可逆。反之，对每个 $x$ 取 $B=xx^{\mathsf T}+tI_n\succ0$，令 $t\downarrow0$ 得 $A\succeq0$，再由可逆性排除零惯性方向。
\end{proof}

正定迹判据中的可逆性不能省略。例如 $A=\operatorname{diag}(1,0)$ 满足对每个二阶 $B\succ0$ 都有 $\operatorname{tr}(AB)=b_{11}>0$，但 $A$ 只有半正定性。另一个常见误区是把行列式为正当成正定的充分条件；例如 $-I_2$ 的行列式为正，却为负定矩阵。行列式为负的实对称矩阵一定有负惯性方向，但未必为负定。

\section{多变元二次型的计算}
\label{sec:xie-ex-qf-computation}

多变元题目不一定需要逐项配方。链状耦合可用二阶分块递推，平方和可先求核空间，具有求和或均值结构的表达式可用秩一矩阵描述。若题目只问规范标准型，确定正、负惯性指数和秩就已足够；若要求具体变量替换，则还应给出变换及其可逆性验证。

\begin{example}
设实二次型 $f$ 的秩大于 $1$。证明它能分解为两个实齐次一次式之积，当且仅当它的惯性指数为 $(1,1)$。
\end{example}
\begin{solve}
若 $f=(a^{\mathsf T}x)(b^{\mathsf T}x)$，则 $a,b$ 必线性无关，否则 $f$ 的秩至多为一。将这两个线性形式扩充为对偶空间的一组基，得到可逆替换 $y_1=a^{\mathsf T}x$、$y_2=b^{\mathsf T}x$ 及其他坐标；再用 $y_1=u+v$、$y_2=u-v$，便有 $f=u^2-v^2$。反之，惯性指数为 $(1,1)$ 时，规范标准型为 $y_1^2-y_2^2=(y_1+y_2)(y_1-y_2)$，代回原变量就得到分解。等价地，条件为秩等于 $2$ 且符号差等于零。
\end{solve}

\begin{example}
\label{ex:xie-ex-qf-chain}
求链式二次型 $f_n(x)=\sum_{i=1}^{n-1}x_ix_{i+1}$ 的规范标准型，其中 $n\ge1$。
\end{example}
\begin{solve}
令 $T_n$ 为主对角元为零、相邻两条副对角线均为 $1$ 的三对角矩阵，则 $2f_n=x^{\mathsf T}T_nx$。$n\ge3$ 时将前两个坐标分出，写
\[
 T_n=\begin{pmatrix}H&B\\B^{\mathsf T}&T_{n-2}\end{pmatrix},
 \qquad H=\begin{pmatrix}0&1\\1&0\end{pmatrix}.
\]
$B$ 只有第 $(2,1)$ 个元素为 $1$，其余为零，故 $B^{\mathsf T}H^{-1}B=0$，于是 $T_n\sim_{\mathrm c}H\oplus T_{n-2}$。由 $T_1=(0)$、$T_2=H$，以及 $H\sim_{\mathrm c}\operatorname{diag}(1,-1)$，得到
\[
 p(f_n)=q(f_n)=\left\lfloor\frac n2\right\rfloor,
 \qquad \operatorname{rank}f_n=2\left\lfloor\frac n2\right\rfloor.
\]
故规范标准型有 $\lfloor n/2\rfloor$ 个正平方与同样多个负平方，$n$ 为奇数时另有一个零方向。乘以正数 $2$ 不改变惯性，故这一结论也适用于 $f_n$。
\end{solve}

\begin{example}
\label{ex:xie-ex-qf-all-cross}
求 $f(x)=\sum_{i=1}^n x_i^2+\sum_{i<j}x_ix_j$ 的正负定性与标准型。
\end{example}
\begin{solve}
利用求和结构，
\[
 f(x)=\frac12\sum_{i=1}^n x_i^2+\frac12\left(\sum_{i=1}^n x_i\right)^2>0\qquad(x\ne0).
\]
所以规范标准型为 $\sum_{i=1}^n y_i^2$。若要得到一个具体的对角型，系数矩阵为 $A=(I_n+\mathbf1\mathbf1^{\mathsf T})/2$，其中 $\mathbf1=(1,\ldots,1)^{\mathsf T}$。矩阵 $I_k+\mathbf1\mathbf1^{\mathsf T}$ 在常数向量方向的特征值为 $k+1$，在坐标和为零的子空间上作用为恒等，故 $\Delta_k=(k+1)/2^k$。由定理~\ref{thm:xie-ex-qf-jacobi}，存在可逆替换将 $f$ 化为 $\sum_{k=1}^n\frac{k+1}{2k}y_k^2$。
\end{solve}

\begin{example}
\label{ex:xie-ex-qf-centering}
设 $\bar x=(x_1+\cdots+x_n)/n$，求 $f(x)=\sum_{i=1}^n(x_i-\bar x)^2$ 的规范标准型。
\end{example}
\begin{solve}
每项均非负，而 $f(x)=0$ 当且仅当 $x_1=\cdots=x_n$，故零集是维数为 $1$ 的线性空间。令 $P=I_n-\mathbf1\mathbf1^{\mathsf T}/n$，则 $f(x)=\|Px\|^2$，且 $\ker P=\mathbb R\mathbf1$。由 $\ker(P^{\mathsf T}P)=\ker P$，可得秩为 $n-1$，正、负惯性指数分别为 $n-1$、$0$，规范标准型为 $y_1^2+\cdots+y_{n-1}^2$。当 $n=1$ 时这表示零型；当 $n\ge2$ 时也有
\[
 \sum_{i<j}(x_i-x_j)^2=n\sum_{i=1}^n(x_i-\bar x)^2,
\]
所以这一差分平方和有相同惯性。
\end{solve}

\begin{example}
\label{ex:xie-ex-qf-cycle}
设 $n\ge2$、$a_1,\ldots,a_n\in\mathbb R$，求
\[
 f(x)=\sum_{i=1}^{n-1}(x_i-a_ix_{i+1})^2+(x_n-a_nx_1)^2
\]
的规范标准型。
\end{example}
\begin{solve}
写 $f=\|Lx\|^2$，其中 $Lx=(x_1-a_1x_2,\ldots,x_{n-1}-a_{n-1}x_n,x_n-a_nx_1)^{\mathsf T}$。展开可得 $\det L=1-a_1\cdots a_n$，而 $L$ 的左上角 $n-1$ 阶子式为 $1$。因此乘积不等于 $1$ 时 $L$ 可逆，$f\succ0$，规范标准型为 $\sum_{i=1}^n y_i^2$；乘积等于 $1$ 时 $\operatorname{rank}L=n-1$，故规范标准型为 $\sum_{i=1}^{n-1}y_i^2$。这里 $f$ 的系数矩阵为 $L^{\mathsf T}L$，而不是 $L$ 本身。
\end{solve}

\begin{example}
以实矩阵 $X=(x_{ij})\in\mathbb R^{n\times n}$ 的 $n^2$ 个元素为变量，求二次型 $f(X)=\operatorname{tr}(X^2)$ 的惯性指数。
\end{example}
\begin{solve}
展开得到 $f=\sum_i x_{ii}^2+2\sum_{i<j}x_{ij}x_{ji}$。对每一对 $i<j$，令 $u_{ij}=(x_{ij}+x_{ji})/\sqrt2$、$v_{ij}=(x_{ij}-x_{ji})/\sqrt2$，这是可逆替换，且 $2x_{ij}x_{ji}=u_{ij}^2-v_{ij}^2$。所以
\[
 p(f)=\frac{n(n+1)}2,\qquad q(f)=\frac{n(n-1)}2.
\]
它在对称矩阵子空间上为正定，在反对称矩阵子空间上为负定；二者维数正好给出上述指数。这里 $\operatorname{tr}(X^2)$ 一般不同于恒非负的 $\operatorname{tr}(X^{\mathsf T}X)$。
\end{solve}

\begin{proposition}
\label{prop:xie-ex-qf-inertia-compression}
若 $A$ 为 $m$ 阶实对称矩阵、$C\in\mathbb R^{m\times n}$，则
\[
 p(C^{\mathsf T}AC)\le p(A),\qquad q(C^{\mathsf T}AC)\le q(A).
\]
从而对同阶实对称矩阵 $A,B$，有 $p(A+B)\le p(A)+p(B)$ 及 $q(A+B)\le q(A)+q(B)$。
\end{proposition}
\begin{proof}
若 $W$ 是 $C^{\mathsf T}AC$ 的正定子空间，则 $C$ 在 $W$ 上单射，否则某非零 $x\in W$ 满足 $Cx=0$，与正定性矛盾；其像 $CW$ 是 $A$ 的正定子空间，且维数等于 $\dim W$。使用定理~\ref{thm:xie-ex-qf-inertia} 中的最大维数刻画即得第一式，对 $-A$ 同理。最后写成
\[
 A+B=\begin{pmatrix}I_n\\I_n\end{pmatrix}^{\mathsf T}
       \operatorname{diag}(A,B)
       \begin{pmatrix}I_n\\I_n\end{pmatrix}
\]
并用分块惯性的可加性。
\end{proof}

\section{矩阵与二次型}
\label{sec:xie-ex-qf-matrix-view}

\subsection{由矩阵求二次型的极值与标准型}

\begin{example}
\label{ex:xie-ex-qf-minimum}
设 $A\succ0$、$b\in\mathbb R^n$、$c\in\mathbb R$，求 $g(x)=x^{\mathsf T}Ax+2b^{\mathsf T}x+c$ 的最小值。
\end{example}
\begin{solve}
矩阵配方给出
\[
 g(x)=(x+A^{-1}b)^{\mathsf T}A(x+A^{-1}b)+c-b^{\mathsf T}A^{-1}b.
\]
故唯一最小点为 $x=-A^{-1}b$，最小值为 $c-b^{\mathsf T}A^{-1}b$。这也就是将增广对称矩阵 $\begin{pmatrix}A&b\\b^{\mathsf T}&c\end{pmatrix}$ 关于 $A$ 消元后的末项；增广矩阵本身未必正定。
\end{solve}

如果只有 $A\succeq0$，则不能直接套用逆矩阵公式。此时最小值存在当且仅当 $b\in\operatorname{im}A$：若 $Az=b$，则 $g(x)=(x+z)^{\mathsf T}A(x+z)+c-z^{\mathsf T}Az$，所有最小点为 $-z+\ker A$；若 $b\notin\operatorname{im}A=(\ker A)^\perp$，可取 $u\in\ker A$ 使 $b^{\mathsf T}u\ne0$，沿 $x=tu$ 便得函数无下界。

\begin{example}
设 $A=(\max\{i,j\})_{1\le i,j\le n}$，求 $x^{\mathsf T}Ax$ 的惯性。
\end{example}
\begin{solve}
对第 $k$ 阶顺序主子阵，依次在 $i=1,\ldots,k-1$ 时作“第 $i$ 行减去第 $i+1$ 行”。此时被减的行尚未变动，新第 $i$ 行的前 $i$ 个元素为 $-1$，其余为零，末行全为 $k$，所以得到下三角矩阵，行列式为 $\Delta_k=(-1)^{k-1}k$。于是 $\Delta_1=1$，而 $\Delta_k/\Delta_{k-1}=-k/(k-1)<0$ 对 $k\ge2$ 成立。由式~\eqref{eq:xie-ex-qf-jacobi}，规范标准型为 $y_1^2-y_2^2-\cdots-y_n^2$。
\end{solve}

\subsection{用二次型证明矩阵性质}

\begin{example}
设 $n\ge2$，求距离矩阵 $D=(|i-j|)_{i,j=1}^n$ 的惯性。
\end{example}
\begin{solve}
令 $S=\sum_{i=1}^n x_i$、$S_k=\sum_{i=1}^k x_i$，将每个距离 $|i-j|$ 分解为它跨过的相邻区间数，得到
\[
 x^{\mathsf T}Dx=2\sum_{k=1}^{n-1}S_k(S-S_k).
\]
在 $S=0$ 的 $n-1$ 维子空间上，它等于 $-2\sum_{k=1}^{n-1}S_k^2$，且只有 $x=0$ 才为零，故 $q(D)\ge n-1$。另一方面，取 $x=e_1+e_n$ 得值 $2(n-1)>0$，故 $p(D)\ge1$。两者之和不超过 $n$，所以 $p(D)=1$、$q(D)=n-1$，规范标准型也为一正、$n-1$ 负。
\end{solve}

\begin{theorem}[Hadamard 乘积的保正性]
\label{thm:xie-ex-qf-hadamard-product}
若 $A=(a_{ij})$、$B=(b_{ij})$ 为同阶实半正定矩阵，则逐元素乘积 $A\circ B=(a_{ij}b_{ij})$ 也半正定；若二者均正定，则 $A\circ B$ 正定。
\end{theorem}
\begin{proof}
写 $B=C^{\mathsf T}C$，令 $D_k=\operatorname{diag}(c_{k1},\ldots,c_{kn})$。逐项计算得到
\[
 A\circ B=\sum_{k=1}^nD_kAD_k,
 \qquad x^{\mathsf T}(A\circ B)x=\sum_{k=1}^n(D_kx)^{\mathsf T}A(D_kx)\ge0.
\]
若 $A,B\succ0$，可取 $C$ 可逆。对任一 $x\ne0$，至少存在一个 $D_kx\ne0$；否则，取 $x_j\ne0$，将推出 $C$ 的第 $j$ 列全为零。故上式至少有一项严格为正。
\end{proof}

\begin{example}
若 $A$ 为可逆实对称矩阵，$S$ 为实反对称矩阵，并且 $AS=SA$，证明 $A+S$ 可逆。
\end{example}
\begin{solve}
利用交换性消去交叉项，
\[
 (A+S)^{\mathsf T}(A+S)=A^{\mathsf T}A+S^{\mathsf T}S.
\]
右边第一项正定、第二项半正定，故整个矩阵正定，因而 $A+S$ 的核为零。这里 $A$ 不必正定，但交换性参与了交叉项相消；若已知 $A\succ0$，则无需交换性，见第~\ref{sec:xie-ex-qf-positive} 节的行列式结论。
\end{solve}

\section{正定型与正定矩阵}
\label{sec:xie-ex-qf-positive}

\subsection{判定、构造与矩阵次序}

综合前面的结论，实对称矩阵正定等价于合同于 $I_n$、存在可逆实 Gram 分解、所有顺序主子式为正，也等价于所有主子式为正。最后一个条件的必要性来自主子阵的正定性，充分性只需选取其中的顺序主子式。由实对称矩阵的正交对角化，它还等价于所有特征值严格为正；这一谱判据使用的是正交相似标准型，不能由一般合同变换保持特征值来解释。Hermite 矩阵具有对应的正定判据，其中转置换成共轭转置，主子式仍为实数。

\begin{proposition}
\label{prop:xie-ex-qf-diagonal-dominance}
若实对称矩阵满足 $a_{ii}>\sum_{j\ne i}|a_{ij}|$ 对所有 $i$ 成立，则 $A\succ0$。此外，对任意实对称矩阵 $A$，存在 $K>0$ 使 $|x^{\mathsf T}Ax|\le K\|x\|^2$ 对所有 $x$ 成立。
\end{proposition}
\begin{proof}
由 $2|x_ix_j|\le x_i^2+x_j^2$，
\[
 x^{\mathsf T}Ax\ge\sum_i\left(a_{ii}-\sum_{j\ne i}|a_{ij}|\right)x_i^2>0\qquad(x\ne0).
\]
同一估计取绝对值得 $|x^{\mathsf T}Ax|\le\sum_i(\sum_j|a_{ij}|)x_i^2$，故取 $K=1+\max_i\sum_j|a_{ij}|$ 即可。
\end{proof}

\begin{example}
证明 Cauchy 型矩阵 $H=(1/(i+j))_{i,j=1}^n$ 正定。
\end{example}
\begin{solve}
利用 $1/(i+j)=\int_0^1t^{i+j-1}\,\mathrm dt$，有限求和可与积分交换，得到
\[
 x^{\mathsf T}Hx=\int_0^1t\left(\sum_{i=1}^n x_it^{i-1}\right)^2\,\mathrm dt.
\]
$x\ne0$ 时括号内是非零多项式，在 $(0,1)$ 上不可能恒为零，并由连续性在某个区间上平方严格为正，故积分为正。这给出了教材中 Cauchy 行列式判据之外的一个直接验证。
\end{solve}

对同阶实对称矩阵，约定 $A\succ B$ 表示 $A-B\succ0$，$A\succeq B$ 表示 $A-B\succeq0$，称为由半正定性诱导的矩阵次序。它不是逐元素大小关系。正定矩阵的和、正数倍、可逆合同像及逆矩阵仍正定；若 $A\succ0$、$B\succeq0$，则 $A+B\succ0$。两个正定矩阵的乘积却未必对称，例如
\[
 \begin{pmatrix}2&1\\1&2\end{pmatrix}
 \begin{pmatrix}1&0\\0&2\end{pmatrix}
 =\begin{pmatrix}2&2\\1&4\end{pmatrix},
\]
所以不能一般地将其称为正定矩阵。

\begin{proposition}
\label{prop:xie-ex-qf-inverse-order}
若 $A\succ B\succ0$，则 $B^{-1}\succ A^{-1}$。
\end{proposition}
\begin{proof}
令 $H=A-B\succ0$。分块求逆或直接相乘可验证
\[
 B^{-1}-A^{-1}=(B+BH^{-1}B)^{-1}.
\]
右边括号内正定，因此其逆也正定。
\end{proof}

\begin{example}
\label{ex:xie-ex-qf-positive-construction}
给定非零实列向量 $u,v$，证明存在 $A\succ0$ 使 $Av=u$ 当且仅当 $u^{\mathsf T}v>0$。
\end{example}
\begin{solve}
必要性来自 $u^{\mathsf T}v=v^{\mathsf T}Av>0$。若 $u^{\mathsf T}v>0$，构造
\[
 A=I_n-\frac{vv^{\mathsf T}}{v^{\mathsf T}v}
     +\frac{uu^{\mathsf T}}{u^{\mathsf T}v}.
\]
直接计算得 $Av=u$。第一项与第二项之和是到 $v^\perp$ 的正交投影，第三项半正定；若 $x^{\mathsf T}Ax=0$，前一部分要求 $x=tv$，后一部分要求 $u^{\mathsf T}x=0$，由 $u^{\mathsf T}v>0$ 得 $t=0$。故 $A\succ0$。更一般地，教材中的修正
\[
 H=G-\frac{Gvv^{\mathsf T}G}{v^{\mathsf T}Gv}
       +\frac{uu^{\mathsf T}}{u^{\mathsf T}v}\qquad(G\succ0)
\]
同样正定：加权 Cauchy--Schwarz 不等式表明前两项的二次型非负，等号恰在 $x\in\mathbb Rv$ 时成立，再用第三项排除非零零方向。
\end{solve}

\subsection{行列式不等式}

\begin{proposition}[Hadamard 与分块行列式不等式]
\label{prop:xie-ex-qf-hadamard-det}
若 $A=(a_{ij})\succ0$，则
\begin{equation}
 0<\det A\le\prod_{i=1}^n a_{ii},
 \label{eq:xie-ex-qf-hadamard-det}
\end{equation}
等号当且仅当 $A$ 为对角矩阵。若 $M=\begin{pmatrix}A&B\\B^{\mathsf T}&D\end{pmatrix}\succ0$，则 $\det M\le\det A\det D$，等号当且仅当 $B=0$。
\end{proposition}
\begin{proof}
写 $A=\begin{pmatrix}A_{n-1}&u\\u^{\mathsf T}&a_{nn}\end{pmatrix}$，由 Schur 补有
\[
 \det A=\det A_{n-1}(a_{nn}-u^{\mathsf T}A_{n-1}^{-1}u)
 \le a_{nn}\det A_{n-1}.
\]
等号当且仅当 $u=0$。反复降阶便得到式~\eqref{eq:xie-ex-qf-hadamard-det} 及等号条件。对于分块不等式，分别用可逆合同将 $A,D$ 化为单位矩阵，得到 $\widetilde M=\begin{pmatrix}I&K\\K^{\mathsf T}&I\end{pmatrix}\succ0$。它的对角元全为 $1$，故 $\det\widetilde M\le1$，等号当且仅当 $K=0$，即 $B=0$；按合同的行列式缩放关系还原即可。
\end{proof}

若实方阵 $X=(Y,Z)$ 按列分块，则
\[
 (\det X)^2\le\det(Y^{\mathsf T}Y)\det(Z^{\mathsf T}Z).
\]
当 $X$ 可逆时，对 $X^{\mathsf T}X$ 使用分块不等式即得；奇异时左边为零，右边由两个 Gram 矩阵的半正定性非负。若 $X$ 可逆，等号条件为 $Y^{\mathsf T}Z=0$，即两组列向量张成的子空间正交。

\begin{proposition}
\label{prop:xie-ex-qf-positive-skew-det}
若 $A\succ0$、$S^{\mathsf T}=-S$ 均为 $n$ 阶实矩阵，则
\[
 \det(A+S)\ge\det A+\det S\ge\det A>0.
\]
第一处等号当且仅当 $n\le2$ 或 $S=0$，且 $\det(A+S)=\det A$ 当且仅当 $S=0$。
\end{proposition}
\begin{proof}
取可逆实矩阵 $C$ 使 $C^{\mathsf T}AC=I_n$，令 $T=C^{\mathsf T}SC$，则 $T$ 仍反对称。对 $I_n+T$ 使用推论~\ref{cor:xie-ex-qf-skew-det}，再将各行列式同时除以 $(\det C)^2$，便得到不等式及等号条件。
\end{proof}

\subsection{非对称矩阵的对称部分与复向量方法}

负定矩阵的判定可统一转化为对 $-A$ 使用正定判据。特别地，$A\prec0$ 当且仅当 $(-1)^k\Delta_k>0$ 对所有 $k$ 成立。若 $A\succ0$，则关于 $x$ 的加边行列式
\[
 h(x)=\det\begin{pmatrix}A&x\\x^{\mathsf T}&0\end{pmatrix}
     =-\det(A)x^{\mathsf T}A^{-1}x
     =-x^{\mathsf T}\operatorname{adj}(A)x
\]
是负定二次型。这一计算同时说明，加边行列式问题往往可转化为逆矩阵或伴随矩阵的正定性问题。

\begin{definition}
教材将满足 $x^{\mathsf T}Mx>0$ 对所有非零实向量 $x$ 成立的实矩阵 $M$ 称为\emph{亚正定矩阵}，将处处非负的情形称为\emph{亚半正定矩阵}。这里沿用教材术语，其本质是矩阵的对称部分正定或半正定。
\end{definition}

唯一分解
\[
 M=H+S,\qquad H=\frac{M+M^{\mathsf T}}2,\quad
 S=\frac{M-M^{\mathsf T}}2
\]
满足 $x^{\mathsf T}Mx=x^{\mathsf T}Hx$，所以亚正定等价于 $H\succ0$。由命题~\ref{prop:xie-ex-qf-positive-skew-det} 可得亚正定矩阵的行列式为正，因而可逆；其转置与逆矩阵也亚正定，因为对 $y=M^{-1}x\ne0$，
\[
 x^{\mathsf T}M^{-1}x=(My)^{\mathsf T}y=y^{\mathsf T}M^{\mathsf T}y>0.
\]
伴随矩阵为正数 $\det M$ 乘以 $M^{-1}$，故也亚正定；正数倍、和及可逆实合同像亦保持这一性质。若 $Mz=\lambda z$、$z\in\mathbb C^n\setminus\{0\}$，则 $z^{\mathsf H}Sz$ 为纯虚数，从而
\[
 \operatorname{Re}\lambda=\frac{z^{\mathsf H}Hz}{z^{\mathsf H}z}>0.
\]
这里使用复向量时必须改用共轭转置。反过来，仅有特征值实部为正不足以保证亚正定，例如 $\begin{pmatrix}1&4\\0&1\end{pmatrix}$ 的对称部分不定。

\begin{example}
设 $A,B,C\succ0$，证明 $g(t)=\det(t^2A+tB+C)$ 的所有复根实部都小于零。
\end{example}
\begin{solve}
若 $t_0$ 为根，存在非零复向量 $z$ 使 $(t_0^2A+t_0B+C)z=0$。左乘 $z^{\mathsf H}$，得到 $at_0^2+bt_0+c=0$，其中 $a=z^{\mathsf H}Az$、$b=z^{\mathsf H}Bz$、$c=z^{\mathsf H}Cz$ 均为正实数。若 $t_0$ 为实根，则非负数不可能满足此方程，故 $t_0<0$；若为非实根，其共轭也是该实系数二次方程的根，由根和得 $\operatorname{Re}t_0=-b/(2a)<0$。三个系数依赖于所选的根与向量，但对该根的这一结论已经足够。
\end{solve}

\section{半正定型与半正定矩阵}
\label{sec:xie-ex-qf-semidefinite}

\subsection{极限性质与主子式判据}

\begin{theorem}
\label{thm:xie-ex-qf-psd-minors}
设 $A$ 为 $n$ 阶实对称矩阵，则下列条件等价：$A\succeq0$；对每个 $\varepsilon>0$ 都有 $A+\varepsilon I_n\succ0$；$A$ 的所有主子式非负。它们也等价于存在实 Gram 分解，以及所有特征值非负。
\end{theorem}
\begin{proof}
若 $A\succeq0$，则对 $x\ne0$ 有 $x^{\mathsf T}(A+\varepsilon I_n)x\ge\varepsilon\|x\|^2>0$；反之，让此不等式中的 $\varepsilon\downarrow0$ 即得半正定性。半正定矩阵的每个主子阵都半正定，其行列式由 Gram 分解非负。若所有主子式非负，则对每个顺序主子阵 $A_k$，展开得
\[
 \det(A_k+\varepsilon I_k)
 =\varepsilon^k+\sum_{j=1}^k e_j(A_k)\varepsilon^{k-j}>0,
\]
其中 $e_j(A_k)$ 为 $A_k$ 的所有 $j$ 阶主子式之和。每个系数非负，首项严格为正，所以 $A+\varepsilon I_n$ 的所有顺序主子式为正，由 Sylvester 判据得其正定，再取极限。Gram 条件已在命题~\ref{prop:xie-ex-qf-gram} 中证明，谱条件则由实对称矩阵的正交对角化得到。
\end{proof}

\begin{notealert}[type=caution,title={顺序主子式的非负性不够}]
矩阵 $\operatorname{diag}(1,0,-1)$ 的顺序主子式为 $1,0,0$，但它有负方向。因此，将正定判据中的每个严格不等号直接改为非严格不等号，会得到错误结论；半正定的主子式判据要求\emph[caution]{所有主子式}非负。
\end{notealert}

半正定矩阵是正定矩阵的极限，且半正定性对矩阵元素的极限封闭：若 $A_j\succeq0$ 且逐元素趋于 $A$，则固定任意 $x$ 后由 $x^{\mathsf T}A_jx\ge0$ 取极限即可。这种方法能延拓由连续表达式构成的不等式，但不能未经检查就延拓严格不等号、等号条件或含逆矩阵的表达式。

例如对 $A\succeq0$，式~\eqref{eq:xie-ex-qf-hadamard-det} 延拓为 $0\le\det A\le\prod_i a_{ii}$。等号当且仅当 $A$ 为对角矩阵或至少有一个 $a_{ii}=0$：若 $A$ 非奇异，就用正定情形的等号条件；若 $A$ 奇异，则等号等价于对角元之积为零。类似地，若 $S$ 实反对称，对 $A+\varepsilon I_n$ 使用命题~\ref{prop:xie-ex-qf-positive-skew-det} 并取极限，得到
\begin{equation}
 \det(A+S)\ge\det A+\det S\ge0\qquad(A\succeq0).
 \label{eq:xie-ex-qf-psd-skew-det}
\end{equation}

\begin{example}
设 $A,B\succeq0$，证明 $\operatorname{tr}(AB)/n\ge\sqrt[n]{\det A\det B}$，并求等号条件。
\end{example}
\begin{solve}
写 $B=C^{\mathsf T}C$，令 $T=CAC^{\mathsf T}\succeq0$。则 $\operatorname{tr}T=\operatorname{tr}(AB)$、$\det T=\det A\det B$，由半正定的 Hadamard 不等式和非负数的算术几何平均不等式，
\[
 \sqrt[n]{\det A\det B}
 =\sqrt[n]{\det T}
 \le\sqrt[n]{\prod_i t_{ii}}
 \le\frac1n\sum_i t_{ii}.
\]
若行列式之积为零，等号要求迹为零，由式~\eqref{eq:xie-ex-qf-trace-zero} 等价于 $AB=0$。若行列式之积为正，则 $C$ 可逆，两个等号要求 $T=cI_n$、$c>0$，这等价于 $AB=cI_n$。故统一的等号条件是 $AB=cI_n$，其中 $c\ge0$。
\end{solve}

\subsection{零对角元与零二次型值}

\begin{proposition}
\label{prop:xie-ex-qf-psd-kernel}
若 $A\succeq0$，则 $a_{ii}=0$ 蕴含第 $i$ 行、列全为零；对任意实向量 $x$，有
\begin{equation}
 x^{\mathsf T}Ax=0\quad\Longleftrightarrow\quad Ax=0.
 \label{eq:xie-ex-qf-psd-kernel}
\end{equation}
因此二次型的零集为 $\ker A$，其维数为 $n-\operatorname{rank}A$。
\end{proposition}
\begin{proof}
第一条由式~\eqref{eq:xie-ex-qf-entry-bound} 的 $|a_{ij}|^2\le a_{ii}a_{jj}$ 得到。第二条写 $A=C^{\mathsf T}C$，则 $x^{\mathsf T}Ax=\|Cx\|^2=0$ 等价于 $Cx=0$，从而蕴含 $Ax=0$；逆向直接成立。最后使用秩与零空间维数公式。
\end{proof}

这两条性质分别是教材的第二、第三个核心性质。第一条将一个零对角元扩展为零行、零列；第二条将一个标量等式提升为向量等式。例如例~\ref{ex:xie-ex-qf-centering} 与例~\ref{ex:xie-ex-qf-cycle} 的求解，只需找出平方和的零集并计算维数，即可得到秩和规范标准型。\emph[caution]{一般不定二次型的零集未必是线性空间，也不能从 $x^{\mathsf T}Ax=0$ 推出 $Ax=0$。}例如 $A=\operatorname{diag}(1,-1)$、$x=(1,1)^{\mathsf T}$ 就是反例。

更精确地，对实对称矩阵 $A$，性质“对每个 $x$，由 $x^{\mathsf T}Ax=0$ 均能推出 $Ax=0$”等价于 $A$ 半正定或半负定。若 $A$ 半正定或半负定，只需对 $A$ 或 $-A$ 使用命题~\ref{prop:xie-ex-qf-psd-kernel}；反之，若 $A$ 不定，规范标准型的一正、一负坐标均取 $1$ 就给出零二次型值却非零矩阵像，经可逆替换转回便与所述性质矛盾。

\begin{proposition}
\label{prop:xie-ex-qf-kernel-sums}
若 $A,B\succeq0$、$S^{\mathsf T}=-S$ 均为 $n$ 阶实矩阵，则
\begin{align*}
 \ker(A+B)&=\ker A\cap\ker B,\\
 \ker(A+S)&=\ker A\cap\ker S.
\end{align*}
故
\[
 \operatorname{rank}(A+B)=\operatorname{rank}\begin{pmatrix}A\\B\end{pmatrix},
 \qquad
 \operatorname{rank}(A+S)=\operatorname{rank}\begin{pmatrix}A\\S\end{pmatrix}.
\]
特别地，$A+B\succ0$ 当且仅当 $\ker A\cap\ker B=\{0\}$；$\det(A+S)>0$ 当且仅当 $\ker A\cap\ker S=\{0\}$。
\end{proposition}
\begin{proof}
若 $(A+B)x=0$，则 $0=x^{\mathsf T}Ax+x^{\mathsf T}Bx$，两项非负，故分别为零，再用式~\eqref{eq:xie-ex-qf-psd-kernel} 得 $Ax=Bx=0$。若 $(A+S)x=0$，左乘 $x^{\mathsf T}$，利用 $x^{\mathsf T}Sx=0$ 得 $x^{\mathsf T}Ax=0$，从而 $Ax=0$，继而 $Sx=0$。两个逆向都直接成立。秩等式来自相同的核；最后，$A+B$ 已经半正定，而 $\det(A+S)\ge0$ 由式~\eqref{eq:xie-ex-qf-psd-skew-det} 得到，所以非奇异分别等价于正定与行列式严格为正。
\end{proof}

由对称性与反对称性，堆叠矩阵的转置具有相同秩，上述秩也可写成横向拼接矩阵 $(A,B)$、$(A,S)$ 的秩。如果 $0\preceq B\preceq A$，则对 $Ax=0$ 的 $x$，等式 $A=(A-B)+B$ 表明 $Bx=0$；因此 $\ker A\subseteq\ker B$，并有 $\operatorname{rank}(A,B)=\operatorname{rank}A$。

\begin{example}
设 $A,B$ 为实对称矩阵，$B\succeq0$。若 $\det(A+\mathrm iB)=0$，证明存在非零实向量 $u$ 使 $Au=Bu=0$。
\end{example}
\begin{solve}
取非零复向量 $z$ 满足 $(A+\mathrm iB)z=0$，左乘 $z^{\mathsf H}$。由于 $z^{\mathsf H}Az,z^{\mathsf H}Bz$ 都为实数，虚部给出 $z^{\mathsf H}Bz=0$。实半正定矩阵对复向量仍半正定，Gram 分解于是给出 $Bz=0$，继而 $Az=0$。将 $z$ 分成实部与虚部，二者都被 $A,B$ 消去，且至少一个非零，取作 $u$ 即可。
\end{solve}

\subsection{主对角块的列空间控制与奇异消元}

\begin{theorem}
\label{thm:xie-ex-qf-singular-schur}
设 $M=\begin{pmatrix}A&B\\B^{\mathsf T}&D\end{pmatrix}$ 为实对称矩阵，其中 $A$ 为 $r$ 阶。则 $M\succeq0$ 当且仅当 $A\succeq0$、$\operatorname{im}B\subseteq\operatorname{im}A$，并且对满足 $AX=B$ 的一个矩阵 $X$，有 $D-X^{\mathsf T}AX\succeq0$。此时
\begin{equation}
 \begin{pmatrix}I&0\\-X^{\mathsf T}&I\end{pmatrix}
 M\begin{pmatrix}I&-X\\0&I\end{pmatrix}
 =\operatorname{diag}(A,D-X^{\mathsf T}AX).
 \label{eq:xie-ex-qf-singular-schur}
\end{equation}
特别地，半正定性蕴含 $\operatorname{rank}(A,B)=\operatorname{rank}A$。
\end{theorem}
\begin{proof}
若 $M\succeq0$，则 $A\succeq0$。对 $u\in\ker A$，向量 $(u^{\mathsf T},0)^{\mathsf T}$ 在 $M$ 的二次型下取零，由命题~\ref{prop:xie-ex-qf-psd-kernel} 可得 $B^{\mathsf T}u=0$。所以 $\ker A\subseteq\ker B^{\mathsf T}$，取正交补得到 $\operatorname{im}B\subseteq\operatorname{im}A$，从而方程 $AX=B$ 可解。直接相乘验证式~\eqref{eq:xie-ex-qf-singular-schur}，右边的第二块作为半正定合同像的主子阵必半正定。反之，所述三个条件使右边分块对角矩阵半正定，故 $M$ 半正定。若 $AX_1=AX_2=B$，则 $A(X_1-X_2)=0$，所以 $X_1^{\mathsf T}AX_1=X_2^{\mathsf T}AX_2$，条件不依赖于解的选择。
\end{proof}

这是教材第四个核心性质“主对角块可消去同行、同列非对角块”的精确形式。它不要求主对角块可逆；真正需要的是非对角块的每一列都属于该主对角块的像空间。零对角元迫使零行、零列，是这一性质在一阶主块上的特殊情形。

\begin{proposition}
\label{prop:xie-ex-qf-semicholesky}
每个实半正定矩阵均可写成 $A=U^{\mathsf T}DU$，其中 $U$ 为单位上三角矩阵、$D$ 为非负对角矩阵；也可写成 $A=R^{\mathsf T}R$，其中 $R$ 为对角元非负的上三角矩阵。反之，两种分解都蕴含半正定性，但一般不唯一。
\end{proposition}
\begin{proof}
对阶数归纳。若首元 $a_{11}=0$，则第一行、列全为零，直接对余下主子阵使用归纳假设。若 $a_{11}>0$，用它作首个主元消去首行、列的非对角元，得到非负数 $a_{11}$ 与一个较低阶的半正定 Schur 补，再使用归纳假设。消元乘子均为单位上三角矩阵，故得到 $U^{\mathsf T}DU$；令 $R=D^{1/2}U$ 即得后一分解，逆向由平方范数非负得到。非唯一性可由
\[
 A=\begin{pmatrix}0&0\\0&1\end{pmatrix},\quad
 R_1=\begin{pmatrix}0&1\\0&0\end{pmatrix},\quad
 R_2=\begin{pmatrix}0&0\\0&1\end{pmatrix}
\]
验证：两者都是对角元非负的上三角矩阵，且 $R_1^{\mathsf T}R_1=R_2^{\mathsf T}R_2=A$。
\end{proof}

对于亚半正定矩阵 $M=H+S$，其中 $H\succeq0$、$S^{\mathsf T}=-S$，式~\eqref{eq:xie-ex-qf-psd-skew-det} 给出 $\det M\ge0$；同第~\ref{sec:xie-ex-qf-positive} 节的复向量计算，可知所有特征值实部非负。其和、非负数倍、转置以及任意实合同像仍亚半正定；伴随矩阵的亚半正定性可由 $M+\varepsilon I_n$ 的伴随矩阵取极限得到。但可逆的亚半正定矩阵未必亚正定：例如 $J=\begin{pmatrix}0&1\\-1&0\end{pmatrix}$ 可逆，却对所有实向量满足 $x^{\mathsf T}Jx=0$。

\section{基础训练与综合应用}
\label{sec:xie-ex-qf-training}

\subsection{概念辨析}

系数矩阵、坐标替换与惯性之间的对应关系是基础题的核心。若一个 $n$ 阶实对称矩阵的秩为 $r$、符号差为 $s$，则
\[
 p=\frac{r+s}{2},\qquad q=\frac{r-s}{2},\qquad
 |s|\le r,\quad r\equiv s\pmod2.
\]
固定阶数 $n$ 时，合同类由非负整数对 $(p,q)$、$p+q\le n$ 决定，共有 $\sum_{r=0}^n(r+1)=(n+1)(n+2)/2$ 类。任意两个同阶正定实对称矩阵都合同；若非退化二次型与其相反数合同，则 $p=q=n/2$，所以 $n$ 必为偶数。可逆实对称矩阵与其逆矩阵合同，但与伴随矩阵未必合同，例如 $\operatorname{diag}(1,1,-1)$ 的伴随为 $\operatorname{diag}(-1,-1,1)$，惯性不同。

正定性要求检验所有非零向量，不能只检验各分量均非零的向量。$\operatorname{diag}(1,0)$ 对各分量均非零的向量给出正值，却非正定；负惯性指数为零也只保证半正定。正定矩阵不要求所有元素为正，$\begin{pmatrix}2&-1\\-1&2\end{pmatrix}$ 就是反例；仅有所有元素为正也不足够，$\begin{pmatrix}1&2\\2&1\end{pmatrix}$ 就不定。对于实矩阵 $B\in\mathbb R^{m\times n}$，$B^{\mathsf T}B$ 总半正定，而其正定等价于 $B$ 满列秩；更一般地，若 $A\succ0$ 为 $m$ 阶矩阵，则 $B^{\mathsf T}AB\succ0$ 当且仅当 $\operatorname{rank}B=n$，因为它的二次型取零恰等价于 $Bx=0$。

\begin{example}
\label{ex:xie-ex-qf-dependent-squares}
解释为什么写成若干个正平方的和，仍不能直接断言原二次型正定。
\end{example}
\begin{solve}
例如 $f=(x_1-x_2)^2+(x_2-x_3)^2+(x_3-x_1)^2$。三个平方内部的线性形式之和为零，所定义的替换矩阵奇异；原变量 $x=(1,1,1)^{\mathsf T}$ 非零却使 $f=0$。其零空间为 $\mathbb R(1,1,1)^{\mathsf T}$，所以 $p=2$、$q=0$。平方和证明的是半正定；要得到正定，必须补证全部线性形式同时为零时只有零解。
\end{solve}

\subsection{综合训练选讲}

\begin{example}
% 对应原书基础训练解答题1。
证明整数反对称矩阵的行列式是某个整数的平方。
\end{example}
\begin{solve}
若矩阵奇异，行列式为 $0^2$。若非奇异，在有理数域上按定理~\ref{thm:xie-ex-qf-skew-form} 消元，得到可逆有理矩阵 $C$，使 $C^{\mathsf T}SC$ 为若干 $J$ 的直和，于是 $\det S=(1/\det C)^2$ 为有理数的平方。又 $\det S$ 为整数；若既约有理数 $a/b$ 的平方为整数，则 $b^2\mid a^2$ 与 $\gcd(a,b)=1$ 迫使 $b=1$，所以 $\det S$ 必为整数平方。
\end{solve}

\begin{example}
% 对应原书基础训练填空题6。
确定实参数 $\lambda$ 的范围，使
\[
 f(x)=\lambda(x_1^2+x_2^2+x_3^2)+x_4^2
       +2x_1x_2+2x_1x_3-2x_2x_3
\]
正定。
\end{example}
\begin{solve}
系数矩阵为 $A=\begin{pmatrix}\lambda&1&1&0\\1&\lambda&-1&0\\1&-1&\lambda&0\\0&0&0&1\end{pmatrix}$，顺序主子式依次为
\[
 \Delta_1=\lambda,\quad \Delta_2=\lambda^2-1,\quad
 \Delta_3=(\lambda+1)^2(\lambda-2),\quad \Delta_4=\Delta_3.
\]
四者均为正当且仅当 $\lambda>2$。边界 $\lambda=2$ 时行列式为零，故不能纳入正定范围。
\end{solve}

\begin{example}
% 对应原书基础训练解答题7的秩一更新主题。
若 $A\succ0$，证明对任意 $x\in\mathbb R^n$，有 $0\le x^{\mathsf T}(A+xx^{\mathsf T})^{-1}x<1$；若 $A\succeq0$，证明存在 $x$ 使 $A+xx^{\mathsf T}\succ0$ 当且仅当 $\operatorname{rank}A\ge n-1$。
\end{example}
\begin{solve}
记 $K=A+xx^{\mathsf T}$。考虑 $M=\begin{pmatrix}K&x\\x^{\mathsf T}&1\end{pmatrix}$，用右下角消元得到 $M\sim_{\mathrm c}\operatorname{diag}(A,1)$。若 $A\succ0$，则 $M\succ0$，再用 $K\succ0$ 消元，得到 $1-x^{\mathsf T}K^{-1}x>0$；下界由 $K^{-1}\succ0$ 得到。若 $A\succeq0$，则
\[
 \ker(A+xx^{\mathsf T})=\ker A\cap x^\perp.
\]
一个线性条件至多将核空间的维数降低 $1$，所以必要条件为 $\dim\ker A\le1$。若 $A$ 已正定，任取 $x$ 即可；若核为一维，选取非零 $x\in\ker A$，则 $\ker A\cap x^\perp=\{0\}$，从而 $K\succ0$。特别地，在 $A$ 奇异且存在这样的 $x$ 时，$\operatorname{rank}A=n-1$，上面的分块矩阵秩为 $n$；再按 $K$ 消元得末项为零，即 $x^{\mathsf T}K^{-1}x=1$。
\end{solve}

\begin{example}
% 对应原书基础训练解答题10；用整数矩阵构造反例。
若 $A\succ B\succ0$，是否必有 $A^2\succ B^2$？
\end{example}
\begin{solve}
取
\[
 B=\begin{pmatrix}1&0\\0&10\end{pmatrix},\quad
 A=\begin{pmatrix}2&1\\1&12\end{pmatrix},\quad
 A-B=\begin{pmatrix}1&1\\1&2\end{pmatrix}.
\]
三个矩阵均由顺序主子式判据正定，但
\[
 A^2-B^2=\begin{pmatrix}4&14\\14&45\end{pmatrix},\qquad
 \det(A^2-B^2)=-16<0,
\]
所以它不定。矩阵乘法一般不交换，不能将 $A^2-B^2$ 当作两个可交换正定因子的乘积；逆矩阵的反序结论也不能迁移到平方函数。
\end{solve}

\begin{example}
% 对应原书基础训练解答题13。
设 $A\succ0$、$B\in\mathbb R^{n\times n}$，且 $M=\begin{pmatrix}A&B^{\mathsf T}\\B&A^{-1}\end{pmatrix}\succeq0$。证明 $B$ 的每个复特征值 $\lambda$ 均满足 $|\lambda|\le1$。
\end{example}
\begin{solve}
用右下角 $A^{-1}$ 消元得 $A-B^{\mathsf T}AB\succeq0$。取非零复向量 $z$ 满足 $Bz=\lambda z$，利用实半正定性对复向量的延拓，得到
\[
 0\le z^{\mathsf H}(A-B^{\mathsf T}AB)z
 =(1-|\lambda|^2)z^{\mathsf H}Az.
\]
$z^{\mathsf H}Az>0$，故 $|\lambda|\le1$。证明只使用特征向量，不要求 $B$ 可对角化。
\end{solve}

\begin{example}
% 对应原书基础训练解答题15，并补充最大维数。
设 $f$ 为非退化的 $n$ 元实二次型，符号差为 $s$。求满足 $f$ 在其上恒为零的线性子空间的最大维数。
\end{example}
\begin{solve}
取可逆替换 $x=Cy$ 使 $f=\sum_{i=1}^p y_i^2-\sum_{j=1}^q y_{p+j}^2$，其中 $p+q=n$、$p-q=s$。记 $m=\min(p,q)=(n-|s|)/2$，则向量 $e_i+e_{p+i}$（$1\le i\le m$）线性无关，其张成空间中的任一向量都使正负平方相消，故其经 $C$ 映出的子空间具有维数 $m$ 且 $f$ 恒为零。反之，若 $W$ 上 $f$ 恒零，则从 $W$ 到前 $p$ 个坐标的投影单射：一个向量若前 $p$ 个坐标为零，由负平方和为零也只能是零向量。同理到后 $q$ 个坐标的投影也单射，故 $\dim W\le p,q$，最大维数正为 $m$。这里强调的是整个子空间上恒零，单独列出若干零点并不能保证其线性组合仍为零。
\end{solve}
% END textbook-xie-exercises-ch8


% BEGIN textbook-xie-exercises-ch9
% 教材：高代习题-9.pdf，PDF 第 1--118 页，原书第 486--603 页。
% 覆盖第 9 章各节；基础训练按知识点整合，不逐题抄录题目与答案。
\chapter{内积空间}
\label{chap:xie-ip}

线性空间的基、维数和线性映射只涉及加法与数乘；要讨论长度、夹角、垂直和最佳逼近，还需要内积这一度量结构。本章沿着“内积与 Gram 矩阵、正交分解与伴随、特殊算子的标准型、矩阵分解及其应用”的路线整理教材第 9 章。矩阵问题中首先应辨明允许的变换：换一般基对应相似，改变二次型的坐标对应合同，换标准正交基则对应正交相似或酉相似。它们保留的信息不同，不能仅因都得到对角矩阵就相互替代。

\section{基本概念}
\label{sec:xie-ip-basics}
% 对应教材 9.1；各类算子与分解在后续各节展开。

\begin{definition}
设 $\mathbb F=\mathbb R$ 或 $\mathbb C$，$V$ 为 $\mathbb F$ 上的线性空间。若二元运算 $\langle\cdot,\cdot\rangle:V\times V\to\mathbb F$ 满足共轭对称性、第一变量的线性及正定性，即
\[
\langle y,x\rangle=\overline{\langle x,y\rangle},\qquad
\langle ax+bz,y\rangle=a\langle x,y\rangle+b\langle z,y\rangle,
\qquad \langle x,x\rangle\geq0,
\]
且 $\langle x,x\rangle=0$ 当且仅当 $x=0$，则称其为\emph{内积（inner product）}。带有内积的有限维实线性空间称为欧氏空间（Euclidean space），有限维复线性空间称为酉空间（unitary space）。定义范数 $\|x\|=\sqrt{\langle x,x\rangle}$ 和距离 $d(x,y)=\|x-y\|$；若 $\langle x,y\rangle=0$，称 $x,y$ 正交，记为 $x\perp y$。
\end{definition}

本章沿用教材\emph[caution]{第一变量线性、第二变量共轭线性}的约定，因而 $\langle x,ay+bz\rangle=\bar a\langle x,y\rangle+\bar b\langle x,z\rangle$。对列向量，标准内积为 $\langle x,y\rangle=y^{\mathsf H}x$，其中 $A^{\mathsf T}$ 表示转置，$A^{\mathsf H}=\bar A^{\mathsf T}$ 表示共轭转置；对实矩阵，$A^{\mathsf H}=A^{\mathsf T}$。沿用前文的矩阵记号，上标 $\mathsf H$ 不与伴随矩阵混用；抽象算子 $T$ 的伴随则记为 $T^*$。实空间中非零向量的夹角 $\theta\in[0,\pi]$ 由 $\cos\theta=\langle x,y\rangle/(\|x\|\|y\|)$ 定义。若采用教材的复空间夹角约定，则以 $|\langle x,y\rangle|/(\|x\|\|y\|)$ 定义余弦，所得角在 $[0,\pi/2]$ 内；它不记录复相位。

\begin{theorem}
\label{thm:xie-ip-cauchy}
对内积空间中的任意 $x,y$，有 Cauchy--Schwarz 不等式与三角不等式
\[
|\langle x,y\rangle|\leq\|x\|\|y\|,
\qquad \|x+y\|\leq\|x\|+\|y\|.
\]
前一个不等式取等号当且仅当 $x,y$ 线性相关；当 $x,y$ 都非零时，后一个取等号当且仅当 $x=ty$，其中 $t>0$ 为实数。
\end{theorem}
\begin{proof}
若 $y\ne0$，令 $c=\langle x,y\rangle/\|y\|^2$，则
\[
0\leq\|x-cy\|^2
=\|x\|^2-\frac{|\langle x,y\rangle|^2}{\|y\|^2}.
\]
这给出第一个不等式，取等号等价于 $x=cy$；$y=0$ 时结论也成立。再由
\[
\|x+y\|^2=\|x\|^2+\|y\|^2+2\operatorname{Re}\langle x,y\rangle
\leq(\|x\|+\|y\|)^2
\]
得到三角不等式。当两向量非零时，取等号要求同时有线性相关及 $\langle x,y\rangle=\|x\|\|y\|>0$，故比例系数为正实数。
\end{proof}

内积的正定性还给出常用的\emph[tip]{非退化性检验}：若 $\langle x,y\rangle=0$ 对所有 $y$ 成立，取 $y=x$ 即得 $x=0$；若 $x-y$ 与一组基的每个向量都正交，则 $x=y$。范数所包含的信息也足以恢复内积：实空间中
\[
\langle x,y\rangle=\frac14\bigl(\|x+y\|^2-\|x-y\|^2\bigr),
\]
复空间中则为
\begin{equation}
\label{eq:xie-ip-polarization}
\langle x,y\rangle=\frac14\bigl(\|x+y\|^2-\|x-y\|^2
+i\|x+iy\|^2-i\|x-iy\|^2\bigr).
\end{equation}
直接展开即可验证；同时得到平行四边形恒等式 $\|x+y\|^2+\|x-y\|^2=2\|x\|^2+2\|y\|^2$。式~\eqref{eq:xie-ip-polarization} 中虚部的符号依赖于内积约定。

\section{内积空间与 Gram 矩阵}
\label{sec:xie-ip-gram}
% 对应教材 9.2，例 9.1--9.10，并整合例 9.14、9.33 的基础结论。

\begin{definition}
复方阵 $H$ 满足 $H^{\mathsf H}=H$ 时称为 Hermite 矩阵（Hermitian matrix）；实情形就是对称矩阵。若对所有非零列向量 $z$，均有 $z^{\mathsf H}Hz>0$，称其正定；若对所有 $z$，均有 $z^{\mathsf H}Hz\geq0$，称其半正定。Hermite 条件保证这些二次型取实数值。
\end{definition}

\begin{example}
教材中的几类典型内积可以统一理解为“把坐标或函数的乘积相加”。在实多项式空间中，$\langle\sum a_it^i,\sum b_it^i\rangle=\sum a_ib_i$ 是系数内积；在 $C[a,b]$（$a<b$）及其多项式子空间中，$\langle f,g\rangle=\int_a^bf(t)g(t)\,dt$ 是积分内积；在复矩阵空间中，
\[
\langle X,Y\rangle_F=\operatorname{tr}(Y^{\mathsf H}X)
=\sum_{i,j}x_{ij}\overline{y_{ij}}
\]
是 Frobenius 内积（Frobenius inner product），对应范数为 $\|X\|_F=(\sum_{i,j}|x_{ij}|^2)^{1/2}$。
\end{example}
\begin{solve}
对称性或共轭对称性及线性都可由运算性质验证。正定性需要分别落实：系数或矩阵内积为各系数模平方之和，只有全部为零时才为零；积分内积中，若连续函数 $f\ne0$，则在某个非退化子区间上 $|f|$ 有正下界，故 $\int_a^bf(t)^2\,dt>0$。这里连续性不可随意删去，否则仅在单点非零的函数也可能具有零积分。
\end{solve}

\begin{definition}
向量组 $v_1,\ldots,v_m$ 的 \emph{Gram 矩阵（Gram matrix）}定义为
\[
G(v_1,\ldots,v_m)=(\langle v_i,v_j\rangle)_{i,j=1}^m.
\]
当向量组为一组基时，也称度量矩阵。若 $e=(e_1,\ldots,e_n)$ 为基，$x=e\xi,y=e\eta$，则
\begin{equation}
\label{eq:xie-ip-gram-coordinates}
\langle x,y\rangle=\xi^{\mathsf T}G\bar\eta.
\end{equation}
若另一组基为 $f=eC$，其 Gram 矩阵为 $H$，则
\begin{equation}
\label{eq:xie-ip-gram-change}
H=C^{\mathsf T}G\bar C;
\qquad \mathbb F=\mathbb R\text{ 时 }H=C^{\mathsf T}GC.
\end{equation}
\end{definition}

上述公式均由逐项展开内积得到。\emph[caution]{在第一变量线性且 $G_{ij}=\langle e_i,e_j\rangle$ 的约定下，复情形不能把式~\eqref{eq:xie-ip-gram-coordinates} 直接写成 $\xi^{\mathsf H}G\eta$。}若改用 $M=\bar G$，同一公式可写成 $\eta^{\mathsf H}M\xi$，换基公式相应变为 $\bar H=C^{\mathsf H}\bar G C$。固定一组基后，内积与正定实对称矩阵或正定 Hermite 矩阵一一对应；例如给定正定 Hermite 矩阵 $G$，以 $\xi^{\mathsf T}G\bar\eta$ 定义内积，其正定性可由 $\xi^{\mathsf T}G\bar\xi=(\bar\xi)^{\mathsf H}G\bar\xi>0$ 验证。

\begin{proposition}
\label{prop:xie-ip-gram-rank}
Gram 矩阵 $G=G(v_1,\ldots,v_m)$ 是半正定 Hermite 矩阵，且
\[
\operatorname{rank}G=\dim\operatorname{span}\{v_1,\ldots,v_m\}.
\]
因此向量组线性无关当且仅当 $G$ 正定，也当且仅当 $\det G\ne0$。反之，任意 $m$ 阶半正定 Hermite 矩阵，只要其秩不超过 $\dim V$，都能成为 $V$ 中某组向量的 Gram 矩阵。
\end{proposition}
\begin{proof}
令 $L(c)=\sum_jc_jv_j$。由定义，$\|L(c)\|^2=c^{\mathsf T}G\bar c\geq0$；共轭对称性给出 $G^{\mathsf H}=G$，故 $G$ 半正定。又
\[
L(c)=0\quad\Longleftrightarrow\quad G\bar c=0.
\]
向右的蕴含由逐项内积得到；向左则由 $c^{\mathsf T}G\bar c=0$ 及正定性得到。因此 $\ker L$ 的维数为 $m-\operatorname{rank}G$，秩的结论随之成立。对于反向构造，$\bar G$ 也是半正定 Hermite 矩阵，故存在 $r\times m$ 矩阵 $X$ 使 $\bar G=X^{\mathsf H}X$，其中 $r=\operatorname{rank}G$。取 $V$ 中 $r$ 个标准正交向量，以 $X$ 的各列为坐标构造 $v_j$，其 Gram 矩阵就是 $X^{\mathsf T}\bar X=G$。这里的半正定矩阵分解也可由后文的谱定理构造。
\end{proof}

\begin{example}
在 $\mathbb R[t]$ 上取 $\langle f,g\rangle=\int_0^1f(t)g(t)\,dt$。向量组 $1,t,\ldots,t^{n-1}$ 的 Gram 矩阵为 Hilbert 矩阵
\[
H_n=\left(\frac1{i+j-1}\right)_{i,j=1}^n.
\]
证明 $H_n$ 正定。
\end{example}
\begin{solve}
对任意非零 $a=(a_1,\ldots,a_n)^{\mathsf T}$，非零多项式 $p(t)=\sum_{j=1}^na_jt^{j-1}$ 不可能在整个 $[0,1]$ 上恒为零，故
\[
a^{\mathsf T}H_na=\int_0^1p(t)^2\,dt>0.
\]
这种证明把分母相互交错的矩阵转化为向量组的度量问题，无须展开行列式。
\end{solve}

\begin{example}
在 $n$ 维欧氏空间中，两两夹角大于直角的非零向量至多有 $n+1$ 个，且该上界可达到。
\end{example}
\begin{solve}
若存在 $n+2$ 个这样的向量 $v_1,\ldots,v_{n+2}$，则前 $n+1$ 个满足非平凡关系 $\sum_{i=1}^{n+1}c_iv_i=0$。将正系数项移到一边，得 $w=\sum_{c_i>0}c_iv_i=\sum_{c_j<0}(-c_j)v_j$，并可令正系数项非空。若 $w=0$，与 $v_{n+2}$ 作内积即得 $0<0$；故 $w\ne0$，负系数项也非空。再将两种表达式作内积，得到
\[
0<\|w\|^2=\sum_{c_i>0,\,c_j<0}c_i(-c_j)\langle v_i,v_j\rangle<0,
\]
矛盾。取 $\mathbb R^n$ 中 $v_i=(n+1)e_i-\boldsymbol1$（$1\leq i\leq n$）及 $v_{n+1}=-\boldsymbol1$，其中 $\boldsymbol1=(1,\ldots,1)^{\mathsf T}$，任意两个不同向量的内积分别为 $-n-2$ 或 $-1$，均小于零，故上界达到。
\end{solve}

\section{Gram--Schmidt 正交化与正交补}
\label{sec:xie-ip-orthogonalization}
% 对应教材 9.3，例 9.11--9.22。

\begin{definition}
由两两正交向量组成的一组基称为正交基（orthogonal basis）；若每个基向量还具有单位长度，则称为标准正交基（orthonormal basis）。任意两两正交的非零向量组都线性无关：若 $\sum_jc_jv_j=0$，依次与 $v_i$ 作内积，便有 $c_i\|v_i\|^2=0$，从而所有 $c_i=0$。
\end{definition}

\begin{theorem}
\label{thm:xie-ip-gs}
对线性无关向量组 $u_1,\ldots,u_m$，递推定义
\begin{equation}
\label{eq:xie-ip-gs}
v_k=u_k-\sum_{j<k}\frac{\langle u_k,v_j\rangle}{\|v_j\|^2}v_j,
\qquad e_k=\frac{v_k}{\|v_k\|},\qquad 1\leq k\leq m.
\end{equation}
则 $v_1,\ldots,v_m$ 两两正交且非零，并且对每个 $k$，$\operatorname{span}(u_1,\ldots,u_k)=\operatorname{span}(v_1,\ldots,v_k)$。因此每个有限维内积空间都有标准正交基，任一子空间的标准正交基都可以扩张为全空间的标准正交基。
\end{theorem}
\begin{proof}
设前 $k-1$ 个向量已满足结论。对 $i<k$，式~\eqref{eq:xie-ip-gs} 与 $v_i$ 作内积，求和中只有 $j=i$ 的项保留，故 $\langle v_k,v_i\rangle=0$。该式同时给出两组前 $k$ 个向量张成空间相同。若 $v_k=0$，则 $u_k$ 属于前 $k-1$ 个向量的张成空间，与线性无关矛盾。归纳完成证明。扩张结论可先将子空间的基扩张为一般基，再从已有正交向量继续正交化。
\end{proof}

\begin{theorem}
\label{thm:xie-ip-projection}
设 $U$ 是有限维内积空间 $V$ 的子空间，定义 $U^\perp=\{x:\langle x,u\rangle=0,\ \forall u\in U\}$。则
\[
V=U\mathbin{\oplus^\perp}U^\perp,\qquad \dim U^\perp=\dim V-\dim U.
\]
若 $e_1,\ldots,e_r$ 是 $U$ 的标准正交基，则到 $U$ 上的正交投影（orthogonal projection）为
\begin{equation}
\label{eq:xie-ip-projection}
P_Ux=\sum_{j=1}^r\langle x,e_j\rangle e_j,
\qquad \|x-u\|^2=\|x-P_Ux\|^2+\|P_Ux-u\|^2\quad(u\in U).
\end{equation}
因而 $P_Ux$ 是 $U$ 中唯一使 $\|x-u\|$ 最小的向量，并且
\[
\sum_{j=1}^r|\langle x,e_j\rangle|^2\leq\|x\|^2.
\]
最后的不等式称为 Bessel 不等式，等号成立当且仅当 $x\in U$；$U=V$ 时给出 Parseval 等式。
\end{theorem}
\begin{proof}
将 $U$ 的标准正交基扩张为 $V$ 的标准正交基，新增向量恰好张成 $U^\perp$，由此得到正交直和与维数公式。逐项验证 $x-P_Ux\perp U$，再对 $x-u=(x-P_Ux)+(P_Ux-u)$ 使用勾股恒等式，即得式~\eqref{eq:xie-ip-projection} 及唯一性。取 $u=0$ 得到 Bessel 不等式。
\end{proof}

在有限维空间中，正交补满足
\[
(U^\perp)^\perp=U,\qquad
(U+W)^\perp=U^\perp\cap W^\perp,\qquad
(U\cap W)^\perp=U^\perp+W^\perp.
\]
第二式由定义直接验证，第三式对第二式取正交补得到。对任意子集 $S$，有 $S^\perp=(\operatorname{span}S)^\perp$ 及 $(S^\perp)^\perp=\operatorname{span}S$。在 $\mathbb F^n$ 中，$\ker A=(\operatorname{Im}A^{\mathsf H})^\perp$，故
\begin{equation}
\label{eq:xie-ip-fredholm}
Ax=b\text{ 有解}\quad\Longleftrightarrow\quad b\perp\ker A^{\mathsf H}.
\end{equation}
这些等式依赖有限维条件；一般无限维空间中的双重正交补不能在不加条件时直接换成原子空间。

\begin{example}
在 $\mathbb R[t]_{\leq n}$ 上取积分内积 $\langle f,g\rangle=\int_{-1}^1f(t)g(t)\,dt$。由 $1,t,\ldots,t^n$ 正交化得到的标准正交多项式为
\[
p_k(t)=\sqrt{\frac{2k+1}{2}}P_k(t),\qquad
P_k(t)=\frac1{2^kk!}\frac{d^k}{dt^k}(t^2-1)^k.
\]
其中 $P_k$ 是 Legendre 多项式（Legendre polynomial）。
\end{example}
\begin{solve}
多项式 $(t^2-1)^k$ 在端点的零点重数为 $k$，故对任意 $\deg q<k$，分部积分 $k$ 次且所有边界项均为零，得到 $\int_{-1}^1P_kq=0$。$P_k$ 的首项系数为 $L_k=(2k)!/(2^k(k!)^2)>0$；再次分部积分得
\[
\int_{-1}^1P_k(t)^2\,dt
=\frac{L_k}{2^k}\int_{-1}^1(1-t^2)^k\,dt
=\frac2{2k+1}.
\]
这里 $I_k=\int_{-1}^1(1-t^2)^k\,dt$ 满足 $I_0=2$、$I_k=\frac{2k}{2k+1}I_{k-1}$，故 $I_k=2^{2k+1}(k!)^2/(2k+1)!$。每个 $p_k$ 均为首项系数为正的单位向量，且属于 $\mathbb R[t]_{\leq k}$ 对 $\mathbb R[t]_{\leq k-1}$ 的一维正交补，因此正是按原顺序正交化所得的向量。特别地，
\[
p_0=\frac1{\sqrt2},\quad p_1=\sqrt{\frac32}t,\quad
p_2=\sqrt{\frac58}(3t^2-1),\quad p_3=\sqrt{\frac78}(5t^3-3t).
\]
\end{solve}

\begin{example}
求次数不超过 $3$ 的实多项式 $f$，使 $\int_{-1}^1(e^t-f(t))^2\,dt$ 最小。
\end{example}
\begin{solve}
由定理~\ref{thm:xie-ip-projection}，唯一最优多项式为 $f_*=\sum_{k=0}^3c_kp_k$，其中
\[
(c_0,c_1,c_2,c_3)=\left(
\frac{e-e^{-1}}{\sqrt2},\ \sqrt6\,e^{-1},\
\sqrt{\frac52}(e-7e^{-1}),\
\sqrt{\frac72}(37e^{-1}-5e)\right).
\]
各系数由分部积分计算 $c_k=\int_{-1}^1e^tp_k(t)\,dt$ 得到。最小误差平方为
\[
\frac{e^2-e^{-2}}2-\sum_{k=0}^3c_k^2
\approx 2.228887\times10^{-5}.
\]
最佳逼近的条件是残差与整个多项式子空间正交，不能仅令几个采样点上的误差为零。
\end{solve}

\begin{proposition}
\label{prop:xie-ip-qr}
若 $A\in\mathbb F^{m\times n}$ 列满秩，则 $m\geq n$，并有唯一分解 $A=QR$，其中 $Q\in\mathbb F^{m\times n}$ 满足 $Q^{\mathsf H}Q=I_n$，$R\in\mathbb F^{n\times n}$ 是对角元全为正实数的上三角矩阵。这称为薄 QR 分解（thin QR decomposition）。实情形中，$A^{\mathsf T}A=R^{\mathsf T}R$ 正是其 Gram 矩阵的 Cholesky 分解。
\end{proposition}
\begin{proof}
依次对 $A$ 的列向量作标准正交化，将每列表示为所得前几个标准正交向量的线性组合，即得 $A=QR$，且 $r_{kk}=\|v_k\|>0$。若还有 $A=\widetilde Q\widetilde R$，两组 $Q$ 的列空间均为 $\operatorname{Im}A$，故 $Q^{\mathsf H}\widetilde Q=R\widetilde R^{-1}$ 为酉上三角矩阵。酉上三角矩阵必为对角矩阵，而其对角元又是正实数，故只能为 $I_n$，唯一性得证。两边计算 $A^{\mathsf H}A$ 得 $R^{\mathsf H}R$；这也说明正定矩阵的 Cholesky 分解与正交化之间的联系。
\end{proof}

对奇异方阵，补选正交向量仍可得到 $A=QR$，其中 $Q$ 为正交或酉矩阵，$R$ 为对角元非负的上三角矩阵，但不再保证唯一。无论实复，正交化中的单位上三角过渡矩阵均使 Gram 行列式保持不变，因而
\begin{equation}
\label{eq:xie-ip-gram-hadamard}
0\leq\det G(u_1,\ldots,u_m)\leq\prod_{j=1}^m\|u_j\|^2.
\end{equation}
当向量组线性无关时，左边等于正交化残差范数平方的乘积，而每个残差范数不超过原向量范数；线性相关时左边为零。右侧等号成立当且仅当向量两两正交或至少有一个零向量。将此式用于方阵的列向量，得到 Hadamard 不等式 $|\det A|\leq\prod_j\|a_j\|$。实向量组的 $\sqrt{\det G}$ 还等于其张成的平行多面体体积。

\section{伴随}
\label{sec:xie-ip-adjoint}
% 对应教材 9.4，例 9.23--9.31。

\begin{theorem}
设 $V,W$ 为有限维内积空间，$T:V\to W$ 为线性映射，则存在唯一线性映射 $T^*:W\to V$，使
\begin{equation}
\label{eq:xie-ip-adjoint}
\langle Tx,y\rangle_W=\langle x,T^*y\rangle_V
\qquad(x\in V,\ y\in W).
\end{equation}
$T^*$ 称为\emph{伴随映射（adjoint map）}。在定义域和值域的标准正交基下，若 $T$ 的矩阵为 $A$，则 $T^*$ 的矩阵为 $A^{\mathsf H}$。
\end{theorem}
\begin{proof}
取两空间的标准正交基，以 $A^{\mathsf H}$ 定义 $T^*$，由 $\langle A\xi,\eta\rangle=\eta^{\mathsf H}A\xi=\langle\xi,A^{\mathsf H}\eta\rangle$ 验证存在性。如果 $S$ 也满足式~\eqref{eq:xie-ip-adjoint}，则 $\langle x,(T^*-S)y\rangle=0$ 对任意 $x$ 成立，由非退化性得 $(T^*-S)y=0$，唯一性成立。
\end{proof}

\begin{definition}
内积空间上的线性算子 $T$ 称为正规算子（normal operator），若 $T^*T=TT^*$；称为自伴随算子（self-adjoint operator），若 $T=T^*$；称为斜自伴随算子（skew-adjoint operator），若 $T^*=-T$。复标准坐标中对应正规矩阵、Hermite 矩阵与斜 Hermite 矩阵；实自伴随算子也称对称变换。自伴随算子 $T$ 若对所有非零 $x$ 满足 $\langle Tx,x\rangle>0$，称为正定；若对所有 $x$ 满足 $\langle Tx,x\rangle\geq0$，称为半正定。
\end{definition}

在允许复合的情况下，伴随满足
\[
(S+T)^*=S^*+T^*,\quad(cT)^*=\bar cT^*,\quad
(ST)^*=T^*S^*,\quad(T^*)^*=T,
\]
以及可逆时 $(T^{-1})^*=(T^*)^{-1}$。这些式子既可由定义逐一验证，也可在标准正交基中转化为共轭转置的恒等式。对于一般基，若定义域和值域的 Gram 矩阵分别为 $G_V,G_W$，由 $A^{\mathsf T}G_W=G_V\bar B$ 得到伴随矩阵
\[
B=\bar G_V^{-1}A^{\mathsf H}\bar G_W.
\]
实空间自同态的公式简化为 $B=G^{-1}A^{\mathsf T}G$。因此\emph[caution]{伴随的矩阵等于转置或共轭转置，必须以前述标准正交基为条件}。

\begin{proposition}
\label{prop:xie-ip-adjoint-kernel}
有限维线性映射 $T:V\to W$ 满足
\[
\ker T=(\operatorname{Im}T^*)^\perp,\qquad
\operatorname{Im}T=(\ker T^*)^\perp.
\]
若 $T:V\to V$ 且 $U$ 是 $T$ 的不变子空间，则 $U^\perp$ 是 $T^*$ 的不变子空间。
\end{proposition}
\begin{proof}
$x\in\ker T$ 等价于 $\langle Tx,y\rangle=0$ 对所有 $y\in W$ 成立，再由伴随定义即等价于 $x\perp\operatorname{Im}T^*$。对 $T^*$ 使用这一等式并取正交补得到第二式。若 $u\in U,z\in U^\perp$，则 $\langle u,T^*z\rangle=\langle Tu,z\rangle=0$，故 $T^*z\in U^\perp$。
\end{proof}

伴随还使特征多项式与极小多项式的系数取共轭：若 $p_T(t)$、$m_T(t)$ 分别为二者，则 $p_{T^*}(t)=\overline{p_T}(t)$、$m_{T^*}(t)=\overline{m_T}(t)$，这里横线只对系数取共轭。特征多项式结论由行列式与共轭转置得到；极小多项式结论则由 $q(T)=0\Longleftrightarrow\bar q(T^*)=0$ 及首一最小次数得到。特别地，$T^*$ 的特征值是 $T$ 的特征值的共轭，重数相同。

\begin{example}
在 $M_n(\mathbb F)$ 上取 Frobenius 内积，设 $T(X)=PXQ$，求 $T^*$。
\end{example}
\begin{solve}
利用迹的循环性质，
\[
\langle PXQ,Y\rangle_F=\operatorname{tr}(Y^{\mathsf H}PXQ)
=\operatorname{tr}(QY^{\mathsf H}PX)=\langle X,P^{\mathsf H}YQ^{\mathsf H}\rangle_F,
\]
故 $T^*(Y)=P^{\mathsf H}YQ^{\mathsf H}$。例如左乘算子 $X\mapsto PX$ 自伴随当且仅当 $P=P^{\mathsf H}$，保持内积当且仅当 $P^{\mathsf H}P=I$，两者均可取 $X=I$ 检验必要性。在实对称矩阵空间上，$X\mapsto R^{\mathsf T}XR$ 的伴随是 $Y\mapsto RYR^{\mathsf T}$，此时还需确认所得矩阵仍属于该子空间。
\end{solve}

\begin{example}
说明有限维条件在伴随存在性中不可直接删去。在 $\mathbb R[t]$ 上取系数内积，定义 $D(f)=f'$ 及 $S(t^k)=1+t+\cdots+t^k$，线性延拓到所有多项式。
\end{example}
\begin{solve}
对单项式核对可得 $D^*(t^j)=(j+1)t^{j+1}$，故 $D^*g=t^2g'+tg$，它确实把多项式映为多项式。然而若 $S^*$ 存在，则
\[
\langle t^k,S^*1\rangle=\langle S(t^k),1\rangle=1\qquad(k\geq0),
\]
要求多项式 $S^*1$ 的每个系数均为 $1$，矛盾。因此无限维内积空间中的线性算子可能有伴随，也可能没有，不能只照搬有限维矩阵的构造。
\end{solve}

\section{保积同构、正交变换与镜像变换}
\label{sec:xie-ip-isometry}
% 对应教材 9.5，例 9.32--9.47。

\begin{definition}
内积空间之间保持内积的线性同构称为保积同构（isometric isomorphism）。欧氏空间上的保积自同构称为正交变换（orthogonal transformation），酉空间上的保积自同构称为酉变换（unitary transformation）。对应的标准正交基表示矩阵分别满足 $Q^{\mathsf T}Q=I$ 与 $U^{\mathsf H}U=I$。
\end{definition}

对维数相同的有限维内积空间之间的线性映射，保持内积、保持范数、把一组标准正交基映为标准正交基三者等价：前两者由极化恒等式联系，后一条件经坐标展开即可验证。保持内积的线性映射是单射，维数相同才进一步保证满射。因此同一基域上的两个有限维内积空间保积同构当且仅当维数相同。对自同态 $T$，上述条件还等价于 $T^*T=I$，也等价于 $T^*=T^{-1}$。

\begin{proposition}
\label{prop:xie-ip-equal-gram}
设 $V,W$ 为维数相同的有限维内积空间。存在保积同构 $T:V\to W$ 使 $T(v_j)=w_j$（$1\leq j\leq m$）的充要条件是两组向量的 Gram 矩阵相同。特别地，对 $A,B\in\mathbb F^{m\times n}$，
\[
A^{\mathsf H}A=B^{\mathsf H}B\quad\Longleftrightarrow\quad A=UB
\]
其中 $U$ 为某个 $m$ 阶酉矩阵；实情形中可取正交矩阵。
\end{proposition}
\begin{proof}
必要性由保积性得到。若两组 Gram 矩阵相等，则对应线性组合的范数相等，故 $\sum c_jv_j=0$ 当且仅当 $\sum c_jw_j=0$。于是
\[
T_0\left(\sum_jc_jv_j\right)=\sum_jc_jw_j
\]
是两张成空间之间良定义的保积同构。分别在其正交补中选标准正交基并一一对应，可将 $T_0$ 扩张到全空间。矩阵结论应用于列向量组；按本章约定，列向量的 Gram 矩阵是 $\overline{A^{\mathsf H}A}$，所以条件仍等价于两组 Gram 矩阵相等。
\end{proof}

\begin{proposition}
\label{prop:xie-ip-householder}
设 $u$ 是欧氏空间中的单位向量，则
\[
H_u(x)=x-2\langle x,u\rangle u
\]
是关于超平面 $u^\perp$ 的镜像变换（reflection）。它在 $\operatorname{span}\{u\}$ 上为 $-I$，在 $u^\perp$ 上为 $I$，故 $H_u^*=H_u$、$H_u^2=I$、$\det H_u=-1$。在标准坐标中，其矩阵为 Householder 矩阵 $I-2uu^{\mathsf T}$。
\end{proposition}
\begin{proof}
由正交分解 $x=\langle x,u\rangle u+z$（$z\perp u$），有 $H_u(x)=-\langle x,u\rangle u+z$，其全部性质由这一分解得到。若 $a\ne b$ 且 $\|a\|=\|b\|$，取 $u=(a-b)/\|a-b\|$，利用
\[
2\langle a,a-b\rangle=\|a-b\|^2
\]
可核对 $H_u(a)=b$，这给出将等长向量互相映送的具体构造。
\end{proof}

\begin{theorem}
$n$ 维欧氏空间中的任意正交变换都能表示为至多 $n$ 个镜像变换之积。
\end{theorem}
\begin{proof}
对 $n$ 归纳。恒等变换是零个镜像变换之积。若正交变换 $T\ne I$，取单位向量 $e$ 使 $Te\ne e$，由命题~\ref{prop:xie-ip-householder} 构造 $H$ 满足 $HTe=e$。正交变换 $HT$ 保持 $e^\perp$，在这一 $n-1$ 维空间上由归纳假设是至多 $n-1$ 个镜像之积。将这些镜像按固定 $e$ 的方式延拓到全空间，再在左边乘 $H=H^{-1}$，即得结论。这就是实欧氏空间中的 Cartan--Dieudonn\'e 定理。
\end{proof}

正交或酉矩阵的列向量与行向量都是标准正交向量组。由 $Ux=\lambda x$ 及范数不变性，任一特征值满足 $|\lambda|=1$；实正交矩阵的行列式为 $\pm1$，酉矩阵的行列式模长为 $1$。此外，任一方形子阵 $B$ 都满足 $\|Bx\|\leq\|x\|$：将 $x$ 嵌入原矩阵相应列的位置，再选取相应行，只会减小范数。因此 $B$ 的特征值模长不超过 $1$。这些性质都不意味着正交矩阵的特征值必须是实数。

\begin{proposition}
\label{prop:xie-ip-conformal}
设 $T$ 是有限维实或复内积空间上的非零线性变换。$T$ 把每对正交向量映为正交向量，当且仅当存在 $c>0$ 使 $T^*T=cI$。因此这类映射恰为 $T=\sqrt c\,U$，其中 $U$ 为正交或酉变换。
\end{proposition}
\begin{proof}
取标准正交基 $e_1,\ldots,e_n$。像向量 $Te_i$ 两两正交；对 $i\ne j$，再由 $e_i+e_j\perp e_i-e_j$ 得到 $\|Te_i\|^2=\|Te_j\|^2$。若 $n=1$，同一结论只涉及唯一的基向量。由于 $T\ne0$，这一公共值 $c>0$，故 $T/\sqrt c$ 将标准正交基映为标准正交基。逆向结论由 $\langle Tx,Ty\rangle=\langle x,T^*Ty\rangle=c\langle x,y\rangle$ 得到。
\end{proof}

实空间中，命题~\ref{prop:xie-ip-conformal} 还表明可逆线性变换保持正交性当且仅当保持所有非零向量的夹角。但保持夹角只限定统一缩放，并不强制 $c=1$；保持长度才给出正交变换。

\section{用正交变换化简二次型}
\label{sec:xie-ip-quadratic}
% 对应教材 9.6，例 9.48--9.51 及前章例题的新解法。

实二次型 $q(x)=x^{\mathsf T}Ax$ 的系数矩阵应取为实对称矩阵，混合项 $x_ix_j$ 的系数是 $2a_{ij}$。实对称矩阵的谱定理说明存在正交矩阵 $Q$，使 $Q^{\mathsf T}AQ=\operatorname{diag}(\lambda_1,\ldots,\lambda_n)$，于是正交代换 $x=Qy$ 将二次型化为 $\sum_i\lambda_iy_i^2$。该定理将在定理~\ref{thm:xie-ip-spectral-normal} 中证明。实际计算时，先求各特征子空间，在同一特征子空间内部正交化并单位化，再按特征值顺序把单位特征向量排成 $Q$ 的列；不同特征值的特征向量已经正交。

正交相似同时是相似和合同：它既保留所有特征值，也保留正负惯性指数。一般合同只保留惯性，允许把非零对角元进一步缩放为 $\pm1$。因此含参题中常先用 $\operatorname{tr}A=\sum\lambda_i$、$\det A=\prod\lambda_i$ 以及 $\operatorname{rank}(A-\lambda I)=n-\dim\ker(A-\lambda I)$ 求参数，再求正交矩阵；不能只凭正负惯性指数确定正交标准型。

\begin{example}
已知 $q(x_1,x_2,x_3)=x_1^2+ax_2^2+x_3^2+2bx_1x_2+2x_1x_3+2x_2x_3$ 可经正交变换化为 $y_2^2+4y_3^2$，求 $a,b$ 及一组正交变换。
\end{example}
\begin{solve}
系数矩阵为 $A=\begin{pmatrix}1&b&1\\b&a&1\\1&1&1\end{pmatrix}$，目标特征值为 $0,1,4$。比较迹得 $a=3$，再由 $\det A=-(b-1)^2=0$ 得 $b=1$。此时相应特征向量可取 $(1,0,-1)^{\mathsf T}$、$(1,-1,1)^{\mathsf T}$、$(1,2,1)^{\mathsf T}$，它们已经两两正交。令
\[
Q=\begin{pmatrix}
1/\sqrt2&1/\sqrt3&1/\sqrt6\\
0&-1/\sqrt3&2/\sqrt6\\
-1/\sqrt2&1/\sqrt3&1/\sqrt6
\end{pmatrix},
\]
则 $Q^{\mathsf T}Q=I$、$Q^{\mathsf T}AQ=\operatorname{diag}(0,1,4)$，故 $x=Qy$ 满足要求。
\end{solve}

\begin{example}
四阶实对称矩阵 $A$ 的特征值为 $0,0,0,4$，且零特征子空间由 $(-1,1,0,0)^{\mathsf T}$、$(-1,0,1,0)^{\mathsf T}$、$(-1,0,0,1)^{\mathsf T}$ 张成，求 $A$。
\end{example}
\begin{solve}
特征值 $4$ 的特征子空间是上述空间的正交补，故单位特征向量为 $u=\tfrac12(1,1,1,1)^{\mathsf T}$。于是对任意 $x$，有 $Ax=4\langle x,u\rangle u$，从而 $A=4uu^{\mathsf T}$ 是所有元素均为 $1$ 的矩阵。只需构造非零特征值对应的正交投影，无须把三维零特征子空间的基逐一正交化。
\end{solve}

\section{实对称矩阵的标准型及其应用}
\label{sec:xie-ip-symmetric}
% 对应教材 9.7，例 9.52--9.74。

\begin{theorem}
\label{thm:xie-ip-rayleigh}
设 $A=A^{\mathsf H}\in\mathbb F^{n\times n}$，特征值为 $\lambda_1\leq\cdots\leq\lambda_n$。则对所有 $x$，
\[
\lambda_1\|x\|^2\leq x^{\mathsf H}Ax\leq\lambda_n\|x\|^2.
\]
左右等号分别等价于 $x\in\ker(A-\lambda_1I)$、$x\in\ker(A-\lambda_nI)$。因此 Rayleigh 商 $R_A(x)=x^{\mathsf H}Ax/(x^{\mathsf H}x)$（$x\ne0$）的最小值与最大值恰为这两个极端特征值。
\end{theorem}
\begin{proof}
以标准正交特征基表示 $x=\sum_i c_ie_i$，则 $x^{\mathsf H}Ax=\sum_i\lambda_i|c_i|^2$、$\|x\|^2=\sum_i|c_i|^2$。上界等号成立等价于所有满足 $\lambda_i<\lambda_n$ 的 $c_i$ 均为零；下界同理。
\end{proof}

由定理~\ref{thm:xie-ip-rayleigh}，若 $A,B$ 都是 Hermite 矩阵，则 $A+B$ 的特征值落在 $[\lambda_{\min}(A)+\lambda_{\min}(B),\lambda_{\max}(A)+\lambda_{\max}(B)]$ 内。对一般矩阵 $C$ 及其特征值 $\lambda$，在特征向量上估计 $C^{\mathsf H}C$ 的二次型还得到
\begin{equation}
\label{eq:xie-ip-eigen-singular-bound}
\sqrt{\lambda_{\min}(C^{\mathsf H}C)}\leq|\lambda|
\leq\sqrt{\lambda_{\max}(C^{\mathsf H}C)}.
\end{equation}
同样，由 $H=(C+C^{\mathsf H})/2$、$K=(C-C^{\mathsf H})/(2i)$ 的 Rayleigh 商，可以分别估计 $\operatorname{Re}\lambda$ 与 $\operatorname{Im}\lambda$。把实矩阵当作复矩阵考察时，必须使用 $x^{\mathsf H}$ 与模平方，不能继续把复特征向量当作实向量。

对 Hermite 矩阵，以 $A\succeq0$、$A\succ0$ 分别表示半正定与正定，并以 $A\succeq B$ 表示 $A-B\succeq0$，这称为 Loewner 序（Loewner order）。它是对二次型的比较，不是逐元素比较。谱定理给出
\[
A\succeq0\Longleftrightarrow\lambda_i(A)\geq0\ (\forall i),
\qquad A\succ0\Longleftrightarrow\lambda_i(A)>0\ (\forall i).
\]
例如 $A\succ0$ 时，$A+A^{-1}\succeq2I$，等号当且仅当 $A=I$；其依据是对每个正特征值使用 $t+t^{-1}\geq2$。

\begin{proposition}
\label{prop:xie-ip-root}
半正定 Hermite 矩阵 $A$ 对每个正整数 $k\geq2$ 都有唯一半正定 Hermite $k$ 次方根
\[
A^{1/k}=U\operatorname{diag}(\lambda_1^{1/k},\ldots,\lambda_n^{1/k})U^{\mathsf H},
\]
其中 $A=U\operatorname{diag}(\lambda_1,\ldots,\lambda_n)U^{\mathsf H}$。实对称矩阵的根仍可取为实对称矩阵。若 $C$ 与 $A$ 交换，则 $C$ 与 $A^{1/k}$ 也交换。
\end{proposition}
\begin{proof}
公式直接给出存在性。若 $B\succeq0$ 且 $B^k=A$，则 $AB=BA$，故每个 $A$ 的特征子空间都在 $B$ 下不变。在特征值 $\lambda$ 的子空间中，$B$ 的所有特征值非负且满足 $\mu^k=\lambda$，所以其限制只能是 $\lambda^{1/k}I$，唯一性成立。对 $A$ 的不同特征值作 Lagrange 插值，可得多项式 $p$ 满足 $p(\lambda)=\lambda^{1/k}$，故 $A^{1/k}=p(A)$，交换性随之成立。
\end{proof}

\begin{proposition}
若 $A\succeq0$，则对任意列向量 $x,y$，有
\[
|x^{\mathsf H}Ay|^2\leq(x^{\mathsf H}Ax)(y^{\mathsf H}Ay),
\]
等号当且仅当 $A^{1/2}x,A^{1/2}y$ 线性相关。若另有 $B\succeq0$，则 $\operatorname{tr}(AB)\geq0$，且等号当且仅当 $AB=0$。
\end{proposition}
\begin{proof}
第一个不等式对 $A^{1/2}x,A^{1/2}y$ 使用 Cauchy--Schwarz 不等式即得；即使 $A$ 奇异，这一方法也仍适用，取等条件同样由这两个像向量决定。第二个结论利用
\[
\operatorname{tr}(AB)=\operatorname{tr}(B^{1/2}AB^{1/2})
=\|A^{1/2}B^{1/2}\|_F^2\geq0.
\]
若迹为零，则 $A^{1/2}B^{1/2}=0$，从而 $AB=A^{1/2}(A^{1/2}B^{1/2})B^{1/2}=0$；逆向由取迹直接得到。
\end{proof}

\begin{proposition}
\label{prop:xie-ip-positive-products}
若 $A\succ0$、$B=B^{\mathsf H}$，则 $AB$ 相似于 Hermite 矩阵 $A^{1/2}BA^{1/2}$，因而可对角化且特征值全实。若 $A,B\succeq0$，则 $AB$ 的特征值非负，且仍可对角化。特别地，对半正定 Hermite 矩阵 $A,B$，$AB$ 是半正定 Hermite 矩阵当且仅当 $AB=BA$。
\end{proposition}
\begin{proof}
第一种情形中，$A^{-1/2}(AB)A^{1/2}=A^{1/2}BA^{1/2}$，结论由谱定理得到。对实半正定 $A,B$，取可逆 $C$ 使 $C^{\mathsf T}AC=\operatorname{diag}(I_r,0)$，同时将 $B$ 变为 $C^{-1}BC^{-\mathsf T}$，它仍半正定；此时两矩阵的乘积为 $C^{\mathsf T}(AB)C^{-\mathsf T}$，只发生相似变化。以下沿用 $A,B$ 表示变换后的矩阵。写成分块后，乘积为 $\begin{pmatrix}B_{11}&B_{12}\\0&0\end{pmatrix}$。若 $z\in\ker B_{11}$，则半正定性使 $B\binom{z}{0}=0$，从而 $B_{21}z=0$。因此 $\operatorname{Im}B_{12}\subseteq\operatorname{Im}B_{11}$，可写 $B_{12}=B_{11}M$，进而
\[
\begin{pmatrix}I&M\\0&I\end{pmatrix}
\begin{pmatrix}B_{11}&B_{11}M\\0&0\end{pmatrix}
\begin{pmatrix}I&-M\\0&I\end{pmatrix}
=\begin{pmatrix}B_{11}&0\\0&0\end{pmatrix}.
\]
右端半正定且可对角化。复情形以共轭转置代替转置，证明相同。最后，$AB$ 为 Hermite 矩阵等价于 $AB=(AB)^{\mathsf H}=BA$，再结合已证的非负特征值结论即可。
\end{proof}

这里使用的事实 $H\succeq0$ 且 $x^{\mathsf H}Hx=0\Rightarrow Hx=0$，可由 $x^{\mathsf H}Hx=\|H^{1/2}x\|^2$ 证明。上述正定或半正定条件各有用途：若只有 $A\succeq0$ 而 $B=B^{\mathsf H}$，乘积不一定可对角化。例如
\[
A=\begin{pmatrix}1&0\\0&0\end{pmatrix},\quad
B=\begin{pmatrix}0&1\\1&0\end{pmatrix},\quad
AB=\begin{pmatrix}0&1\\0&0\end{pmatrix}.
\]
$AB$ 的特征值都是零，却不是零矩阵，因而不可对角化。

\begin{example}
设 $A,C$ 为正定实对称矩阵，证明矩阵方程 $AX+XA=C$ 有唯一解，且该解正定实对称。
\end{example}
\begin{solve}
取正交矩阵 $Q$ 使 $Q^{\mathsf T}AQ=\operatorname{diag}(\lambda_1,\ldots,\lambda_n)$，其中 $\lambda_i>0$。令 $Y=Q^{\mathsf T}XQ$、$D=Q^{\mathsf T}CQ$，方程等价于
\[
(\lambda_i+\lambda_j)y_{ij}=d_{ij},\qquad
 y_{ij}=\frac{d_{ij}}{\lambda_i+\lambda_j}.
\]
这给出唯一解且 $Y=Y^{\mathsf T}$。若该实对称解 $X$ 有特征向量 $u\ne0$、特征值 $\mu$，则
\[
0<u^{\mathsf T}Cu=u^{\mathsf T}(AX+XA)u
=2\mu\,u^{\mathsf T}Au,
\]
故 $\mu>0$，于是 $X\succ0$。由 $A,X\succ0$ 不能反过来直接断定 $AX+XA\succ0$，上述论证是从已给定的正定右端推出解的正定性。
\end{solve}

\begin{proposition}
若实矩阵 $T$ 满足 $T^{\mathsf T}G=GT$，其中 $G\succ0$ 为实对称矩阵，则 $T$ 在实数域上可对角化且特征值全实。
\end{proposition}
\begin{proof}
$G^{1/2}TG^{-1/2}$ 的转置等于自身，故它是实对称矩阵；而它与 $T$ 相似。几何上，$T$ 对由 $G$ 定义的内积自伴随，这说明某个矩阵在普通坐标中不对称，并不排除它是另一种内积下的自伴随算子。
\end{proof}

\section{同时合同对角化}
\label{sec:xie-ip-congruence}
% 对应教材 9.8，例 9.75--9.84。

\begin{theorem}
\label{thm:xie-ip-simultaneous-congruence}
设 $A\succ0$ 为实对称矩阵，$B$ 为同阶实对称矩阵，则存在可逆矩阵 $C$，使
\begin{equation}
\label{eq:xie-ip-congruence}
C^{\mathsf T}AC=I,\qquad C^{\mathsf T}BC=\operatorname{diag}(\lambda_1,\ldots,\lambda_n),
\end{equation}
其中 $\lambda_i$ 是 $A^{-1}B$ 的特征值，也是广义特征值问题 $Bx=\lambda Ax$ 的全部特征值。
\end{theorem}
\begin{proof}
矩阵 $A^{-1/2}BA^{-1/2}$ 实对称，取正交矩阵 $Q$ 将它对角化，令 $C=A^{-1/2}Q$ 即得式~\eqref{eq:xie-ip-congruence}。又
\[
A^{1/2}(A^{-1}B)A^{-1/2}=A^{-1/2}BA^{-1/2},
\]
故二者特征值相同。这里 $C$ 通常不是正交矩阵，也不要求 $AB=BA$。
\end{proof}

同时合同对角化的关键是先用 $A$ 把度量化为单位矩阵，再在这一度量中正交对角化 $B$。行列式不等式尤其适合此法，但须同时保留每个行列式引入的 $\det(C)^2$。例如当 $A\succ0,B\succeq0$ 时，上式中的 $\lambda_i\geq0$，于是
\[
\det(A+B)=\det A\prod_i(1+\lambda_i)
\geq\det A+\det B.
\]
$n\geq2$ 时，等号要求所有中间阶的非负乘积项为零，特别要求 $\sum_i\lambda_i=0$，故等价于 $B=0$；$n=1$ 时恒取等号。若 $A,B\succ0$，则逐项使用 $1+\lambda_i\geq2\sqrt{\lambda_i}$ 得
\begin{equation}
\label{eq:xie-ip-determinant-amgm}
\det(A+B)\geq2^n\sqrt{\det A\det B},
\end{equation}
等号当且仅当 $A=B$。

\begin{proposition}
若 $A\succeq B\succ0$，则 $B^{-1}\succeq A^{-1}$ 且 $\det A\geq\det B$。若 $A,B\succ0$ 且 $A^2\succeq B^2$，则 $A\succeq B$。
\end{proposition}
\begin{proof}
对前一结论，在式~\eqref{eq:xie-ip-congruence} 中 $0<\lambda_i\leq1$，于是 $\lambda_i^{-1}\geq1$。对逆矩阵使用 $C^{-1}(\cdot)(C^{\mathsf T})^{-1}$ 的合同变换，得到 $B^{-1}-A^{-1}\succeq0$；行列式结论来自 $\det B=\det A\prod_i\lambda_i$。

对后一结论，仍写 $A=C^{-\mathsf T}C^{-1}$、$B=C^{-\mathsf T}\Lambda C^{-1}$，其中 $\Lambda=\operatorname{diag}(\lambda_i)\succ0$，并令 $D=C^{-1}C^{-\mathsf T}\succ0$。假设化为 $D-\Lambda D\Lambda\succeq0$，所以其每个对角元 $d_{ii}(1-\lambda_i^2)\geq0$。由于 $d_{ii}>0$，有 $\lambda_i\leq1$，故 $I-\Lambda\succeq0$，再作合同变换得 $A-B\succeq0$。
\end{proof}

\begin{notealert}[type=caution,title={矩阵序不能任意平方}]
一般不能由 $A\succeq B\succ0$ 推出 $A^2\succeq B^2$。例如 $B=\operatorname{diag}(2,1)$、$A=\begin{pmatrix}3&1\\1&2\end{pmatrix}$ 满足 $A-B=\begin{pmatrix}1&1\\1&1\end{pmatrix}\succeq0$，但 $A^2-B^2=\begin{pmatrix}6&5\\5&4\end{pmatrix}$ 的行列式为 $-1$，不是半正定矩阵。若另有 $AB=BA$，则可在公共特征基中逐项比较正特征值，平方等价才成立。
\end{notealert}

\begin{example}
设 $M=\begin{pmatrix}A&B\\B^{\mathsf T}&D\end{pmatrix}\succ0$，证明 $\det M\leq\det A\det D$，并确定等号条件。
\end{example}
\begin{solve}
对 $M$ 作对称分块消元，得到合同矩阵
\[
\begin{pmatrix}I&0\\-B^{\mathsf T}A^{-1}&I\end{pmatrix}
M
\begin{pmatrix}I&-A^{-1}B\\0&I\end{pmatrix}
=\begin{pmatrix}A&0\\0&D-B^{\mathsf T}A^{-1}B\end{pmatrix}.
\]
故 Schur 补 $S=D-B^{\mathsf T}A^{-1}B\succ0$，且 $D\succeq S$，于是 $\det M=\det A\det S\leq\det A\det D$。正定矩阵行列式单调性的等号要求 $D=S$，即 $B^{\mathsf T}A^{-1}B=0$，等价于 $A^{-1/2}B=0$，故等价于 $B=0$。
\end{solve}

两个半正定实对称矩阵仍可同时合同对角化。一个简洁的证明是在 $(\ker A\cap\ker B)^\perp$ 上考察 $A+B$：它在该空间正定，先将 $A+B$ 化为 $I$，再将 $A$ 正交对角化，$B=I-A$ 也随之对角化；在公共核上二者均为零。由 $\ker(A+B)=\ker A\cap\ker B$ 还得到
\[
\operatorname{rank}(A+B)=\operatorname{rank}(A\ B).
\]
但“一个半正定、另一个仅对称”不足以保证同时合同对角化；命题~\ref{prop:xie-ip-positive-products} 后的两个二阶矩阵即构成反例。若有公共对角化基，则 $A=\operatorname{diag}(1,0)$ 在此基中的对角元恰有一个非零，故两个基向量中必有一个为 $v=(0,v_2)^{\mathsf T}$，$v_2\ne0$，另一个 $u=(u_1,u_2)^{\mathsf T}$ 满足 $u_1\ne0$。然而对 $B=\begin{pmatrix}0&1\1&0\end{pmatrix}$，有 $u^{\mathsf T}Bv=u_1v_2\ne0$，与公共对角化要求的非对角元为零矛盾。

\section{Schur 定理}
\label{sec:xie-ip-schur}
% 对应教材 9.9，例 9.85--9.90。

\begin{theorem}
\label{thm:xie-ip-schur}
任一复方阵 $A$ 都酉相似于上三角矩阵：存在酉矩阵 $U$，使 $U^{\mathsf H}AU=R$ 为上三角矩阵，且其对角元为 $A$ 的全部特征值。
\end{theorem}
\begin{proof}
对阶数归纳。复矩阵存在特征值 $\lambda$ 及单位特征向量 $u_1$。把 $u_1$ 扩张为标准正交基，所得酉矩阵 $U_1$ 满足
\[
U_1^{\mathsf H}AU_1=\begin{pmatrix}\lambda&*\\0&A_1\end{pmatrix}.
\]
由归纳假设，存在酉矩阵 $U_2$ 把 $A_1$ 化为上三角矩阵；用 $U_1\operatorname{diag}(1,U_2)$ 即完成归纳。上三角矩阵的特征多项式为对角元对应一次因子的乘积，故对角元就是全部特征值。
\end{proof}

教材还给出 QR 分解的证明：先取可逆 $P$ 使 $P^{-1}AP=M$ 上三角，写 $P=UR$，则 $U^{\mathsf H}AU=RMR^{-1}$ 仍上三角。这个方法说明，正交化不会破坏由前几个基向量张成的不变子空间链。Schur 定理适用于任意复方阵，并不要求可对角化。

实矩阵未必存在实特征值，故其一般形式只能是\emph{实 Schur 分解（real Schur decomposition）}：经正交相似化为分块上三角矩阵，对角块为实的一阶块或有共轭非实特征值的二阶块。其构造依赖于实不变子空间：若 $A(u+iv)=(a+ib)(u+iv)$ 且 $b\ne0$，比较实虚部得 $Au=au-bv$、$Av=bu+av$，且 $u,v$ 线性无关。若二者相关，则实特征向量会有非实特征值，矛盾。因此可先在一维或二维不变子空间内选正交基，再对其余的压缩矩阵归纳。只有当全部特征值为实时，才保证可以化为实上三角矩阵；一般的二阶 Schur 块也尚未具有正规矩阵的旋转伸缩形式。

\begin{lemma}
\label{lem:xie-ip-normal-block}
若分块上三角矩阵 $N=\begin{pmatrix}A&C\\0&B\end{pmatrix}$ 为正规矩阵，即 $N^{\mathsf H}N=NN^{\mathsf H}$，则 $C=0$，并且 $A,B$ 均正规。
\end{lemma}
\begin{proof}
比较左上角分块得 $A^{\mathsf H}A=AA^{\mathsf H}+CC^{\mathsf H}$。两边取迹，因 $\operatorname{tr}(A^{\mathsf H}A)=\operatorname{tr}(AA^{\mathsf H})$，有 $\operatorname{tr}(CC^{\mathsf H})=\sum_{i,j}|c_{ij}|^2=0$，从而 $C=0$。再比较对角分块即得 $A,B$ 正规。
\end{proof}

\section{复正规算子与复正规矩阵}
\label{sec:xie-ip-complex-normal}
% 对应教材 9.10，例 9.91--9.102。



\begin{theorem}
\label{thm:xie-ip-spectral-normal}
复矩阵 $A$ 正规当且仅当它可酉对角化。自伴随算子的特征值全为实数；实自伴随算子有由实特征向量组成的标准正交基，即任一实对称矩阵可正交对角化。
\end{theorem}
\begin{proof}
先对复矩阵应用 Schur 定理，再反复应用引理~\ref{lem:xie-ip-normal-block}，可知正规上三角矩阵必为对角矩阵。反之，对角矩阵正规，而正规性在酉相似下保持。若 $Ax=\lambda x$ 且 $A=A^{\mathsf H}$，则
\[
\lambda\|x\|^2=\langle Ax,x\rangle
=\langle x,Ax\rangle=\bar\lambda\|x\|^2,
\]
故 $\lambda\in\mathbb R$。实对称矩阵视为 Hermite 矩阵后，所有特征值均实；每个实特征值的复特征向量的实部或虚部中至少一个是非零实特征向量。取一个实单位特征向量，其正交补仍为不变子空间，限制后仍自伴随，对维数归纳即得实标准正交特征基。
\end{proof}

不同特征值的正规算子特征子空间互相正交；这既可由标准正交特征基看出，也可利用下述伴随性质验证。自伴随算子的特征值全实，酉算子的特征值模长全为 $1$，斜自伴随算子的特征值为纯虚数或零。\emph[caution]{“可对角化且特征值全实”不自动意味着自伴随；必须要求标准正交特征基。}例如 $\begin{pmatrix}1&1\\0&2\end{pmatrix}$ 可对角化，但不是 Hermite 矩阵，也不是正规矩阵。

\begin{proposition}
\label{prop:xie-ip-normal-tests}
有限维实或复内积空间上的算子 $T$ 正规当且仅当 $\|Tx\|=\|T^*x\|$ 对所有 $x$ 成立。正规时，
\begin{equation}
\label{eq:xie-ip-eigen-adjoint}
\ker(T-\lambda I)=\ker(T^*-\bar\lambda I)
\end{equation}
对基域中的所有 $\lambda$ 成立，故 $\ker T=\ker T^*$。在复空间中，正规性还等价于对每个特征值 $\lambda$ 有
\[
V=\ker(T-\lambda I)\mathbin{\oplus^\perp}\operatorname{Im}(T-\lambda I).
\]
\end{proposition}
\begin{proof}
正规时 $\|Tx\|^2=\langle x,T^*Tx\rangle=\langle x,TT^*x\rangle=\|T^*x\|^2$。反之，$H=T^*T-TT^*$ 自伴随，范数相等说明 $\langle x,Hx\rangle=0$ 对所有 $x$ 成立；在 $H$ 的标准正交特征基上检验可知每个特征值为零，故 $H=0$。由于 $T-\lambda I$ 仍正规，对其使用范数等式即得式~\eqref{eq:xie-ip-eigen-adjoint}。

复空间中，必要性由特征子空间的正交直和得到。反之，题设正交分解意味着 $\operatorname{Im}(T-\lambda I)=(\ker(T-\lambda I))^\perp$，再由命题~\ref{prop:xie-ip-adjoint-kernel} 得到 $\ker(T-\lambda I)=\ker(T^*-\bar\lambda I)$。取单位特征向量 $e$，它同时为 $T^*$ 的特征向量，故 $e^\perp$ 同时在 $T,T^*$ 下不变。限制后每个特征向量仍满足这一伴随关系，归纳即可构造标准正交特征基，从而 $T$ 正规。
\end{proof}

\begin{proposition}
\label{prop:xie-ip-normal-polynomial}
复正规算子 $T$ 满足 $T^*=p(T)$，其中 $p$ 是某个复系数多项式；反之，若 $T^*=p(T)$，则 $T$ 正规。因而任何与 $T$ 交换的算子也与 $T^*$ 交换。此外，$T=H+iK$ 的分解中，
\[
H=\frac{T+T^*}{2},\qquad K=\frac{T-T^*}{2i}
\]
均自伴随，且 $T$ 正规当且仅当 $HK=KH$。
\end{proposition}
\begin{proof}
设 $T$ 的不同特征值为 $\lambda_1,\ldots,\lambda_s$，取 Lagrange 插值多项式
\[
p(t)=\sum_{j=1}^s\bar\lambda_j
\prod_{\ell\ne j}\frac{t-\lambda_\ell}{\lambda_j-\lambda_\ell}.
\]
在标准正交特征基下有 $p(T)=T^*$。逆向结论来自 $T$ 与自身多项式交换。关于 $H,K$ 的结论可将 $T^*T$ 与 $TT^*$ 展开，相减后得到 $T^*T-TT^*=2i(HK-KH)$。
\end{proof}

\begin{proposition}
设复方阵 $A=(a_{ij})$ 的特征值按代数重数记为 $\lambda_1,\ldots,\lambda_n$，则
\begin{equation}
\label{eq:xie-ip-frobenius-spectrum}
\sum_i|\lambda_i|^2\leq\|A\|_F^2=\sum_{i,j}|a_{ij}|^2,
\end{equation}
等号当且仅当 $A$ 正规。
\end{proposition}
\begin{proof}
Frobenius 范数在左右乘酉矩阵时不变。将 $A$ 酉相似为 Schur 上三角矩阵 $R$ 后，
\[
\|A\|_F^2=\|R\|_F^2
=\sum_i|\lambda_i|^2+\sum_{i<j}|r_{ij}|^2.
\]
等号成立当且仅当 $R$ 对角，由定理~\ref{thm:xie-ip-spectral-normal} 即得结论。若 $A$ 是实矩阵且所有特征值均实，则等号进一步等价于 $A$ 实对称。
\end{proof}

\begin{example}
斜 Hermite 矩阵 $K$ 与不含特征值 $-1$ 的酉矩阵 $U$ 通过 Cayley 变换（Cayley transform）一一对应：
\[
U=(I-K)(I+K)^{-1},\qquad K=(I-U)(I+U)^{-1}.
\]
\end{example}
\begin{solve}
$K^{\mathsf H}=-K$ 使其特征值为纯虚数或零，故 $I+K$ 可逆；利用 $I-K$ 与 $I+K$ 交换，直接计算得 $U^{\mathsf H}U=I$，且 $I+U=2(I+K)^{-1}$ 可逆。反过来，若 $U$ 酉且 $I+U$ 可逆，以 $U^{\mathsf H}=U^{-1}$ 代入可得所定义的 $K^{\mathsf H}=-K$。再把两式相互代入，得到原矩阵，故为双射。实矩阵的情形即为实反对称矩阵与不含特征值 $-1$ 的正交矩阵的对应。
\end{solve}

\section{实正规算子与实正规矩阵}
\label{sec:xie-ip-real-normal}
% 对应教材 9.11，例 9.103--9.113。

\begin{theorem}
\label{thm:xie-ip-real-normal}
任一实正规矩阵 $A$ 都正交相似于分块对角矩阵
\begin{equation}
\label{eq:xie-ip-real-normal-form}
\operatorname{diag}\bigl(C(a_1,b_1),\ldots,C(a_r,b_r),c_1,\ldots,c_s\bigr),
\qquad C(a,b)=\begin{pmatrix}a&-b\\b&a\end{pmatrix},
\end{equation}
其中 $b_j\ne0$、$2r+s=n$；各二阶块对应共轭特征值 $a_j\pm ib_j$。因此实正规矩阵在复数域可对角化，在实数域却可能只能分成一维与二维的正交不变子空间。
\end{theorem}
\begin{proof}
实 Schur 分解与引理~\ref{lem:xie-ip-normal-block} 已使所有非对角分块消失。若有非实特征值 $a+ib$ 及特征向量 $z=u+iv$，则 $\bar z$ 对应 $a-ib$。把 $A$ 看作复正规矩阵，由不同特征子空间正交得到 $\langle z,\bar z\rangle=0$，亦即
\[
\|u\|^2-\|v\|^2+2i\langle u,v\rangle=0.
\]
故 $u\perp v$ 且 $\|u\|=\|v\|>0$。在实标准正交基 $u/\|u\|,-v/\|v\|$ 下，限制矩阵恰为 $C(a,b)$。其二维空间同时在 $A,A^{\mathsf T}$ 下不变，正交补亦然；对维数归纳即可得到所需分解。
\end{proof}

复正规算子的若干判定有实版本，但须照顾非实特征值。实算子总能唯一写成 $T=H+K$，其中 $H=(T+T^*)/2$ 自伴随、$K=(T-T^*)/2$ 斜自伴随；$T$ 正规当且仅当 $HK=KH$。对实正规矩阵，把命题~\ref{prop:xie-ip-normal-polynomial} 中插值多项式写成 $p=g+ih$，其中 $g,h$ 为实系数多项式，则 $A^{\mathsf T}=g(A)+ih(A)$ 的两边比较实虚部可知 $A^{\mathsf T}=g(A)$。故与实正规矩阵交换也自动与其转置交换。若只检查“实特征向量是否为伴随的特征向量”，当矩阵没有实特征值时条件可能完全为空，不能据此判定正规性。

\begin{proposition}
\label{prop:xie-ip-reducing}
有限维内积空间中，正规算子 $T$ 的每个不变子空间 $U$ 都同时在 $T^*$ 下不变，且 $U^\perp$ 也同时在 $T,T^*$ 下不变；$T|_U$ 仍为正规算子。
\end{proposition}
\begin{proof}
取 $U$ 及 $U^\perp$ 的标准正交基，$T$ 的表示矩阵为 $\begin{pmatrix}A&C\\0&B\end{pmatrix}$。由引理~\ref{lem:xie-ip-normal-block} 得 $C=0$，矩阵成为分块对角形；取伴随后仍保持此分块，且两限制矩阵正规。这个证明同时适用于实空间与复空间。
\end{proof}

\begin{example}
比较实、复空间中条件 $Tx\perp x$ 对算子的约束。
\end{example}
\begin{solve}
实空间中，若 $\langle Tx,x\rangle=0$ 对所有 $x$ 成立，把 $x+y$ 代入并减去两个对角项，得 $\langle Tx,y\rangle+\langle Ty,x\rangle=0$，故 $T^*=-T$；反向也由对称性直接成立。因此非零实反对称矩阵给出非零实例，例如 $\begin{pmatrix}0&-1\\1&0\end{pmatrix}$。

复空间中，同样代入 $x+y$ 得 $a+b=0$，其中 $a=\langle Tx,y\rangle$、$b=\langle Ty,x\rangle$；再代入 $x+iy$，扣除对角项，得 $-ia+ib=0$，故 $a=b=0$。任意固定 $x$ 后取 $y=Tx$，即得 $Tx=0$，所以只能 $T=0$。实空间的斜自伴随判据不能直接搬到复空间；复斜自伴随算子的 $\langle Tx,x\rangle$ 一般是纯虚数，并非必为零。
\end{solve}

\section{实正规矩阵标准型的应用}
\label{sec:xie-ip-real-normal-applications}
% 对应教材 9.12，例 9.114--9.122。

实正交矩阵在式~\eqref{eq:xie-ip-real-normal-form} 中的二阶块满足 $a^2+b^2=1$，故为平面旋转块
\[
R(\theta)=\begin{pmatrix}\cos\theta&-\sin\theta\\\sin\theta&\cos\theta\end{pmatrix},
\]
一阶块为 $1$ 或 $-1$。实反对称矩阵则由 $\begin{pmatrix}0&-b_j\\b_j&0\end{pmatrix}$ 与零块构成，故秩为偶数，非零特征值成对为 $\pm ib_j$。若正交矩阵的标准型中有 $t$ 个 $-1$ 块，则 $\det Q=(-1)^t$；奇数维且 $\det Q=1$ 时必有特征值 $1$，而 $\det Q=-1$ 时必有特征值 $-1$。在三维，行列式为 $1$ 的情形有固定轴，行列式为 $-1$ 的一般情形是轴向反向与垂直平面旋转的复合，不一定是单独一个超平面镜像。

\begin{proposition}
若 $A\succ0$ 实对称、$S$ 实反对称，则存在可逆 $C$，使 $C^{\mathsf T}AC=I$，且 $C^{\mathsf T}SC$ 为上述反对称分块标准型。于是
\[
\det(A+S)\geq\det A>0,
\]
等号当且仅当 $S=0$。
\end{proposition}
\begin{proof}
先取 $A^{-1/2}$ 将 $A$ 化为 $I$，再正交标准化 $A^{-1/2}SA^{-1/2}$，即得 $C$。比较合同变换后的行列式，有
\[
\frac{\det(A+S)}{\det A}=\prod_j(1+b_j^2)\geq1,
\]
等号要求所有非零二阶块消失，即 $S=0$。
\end{proof}

\begin{example}
若实矩阵 $M$ 的对称部分 $H=(M+M^{\mathsf T})/2$ 正定，证明 $\det(M+M^{\mathsf T})\leq2^n\det M$，并求等号条件。
\end{example}
\begin{solve}
写 $M=H+S$，其中 $S=(M-M^{\mathsf T})/2$ 反对称。由上一命题，$\det M\geq\det H>0$，因此 $\det(M+M^{\mathsf T})=2^n\det H\leq2^n\det M$。等号当且仅当 $S=0$，即 $M$ 本身对称。
\end{solve}

\begin{proposition}
任一实正交矩阵都能表示为至多两个实对称矩阵之积；任一实方阵都能表示为至多三个实对称矩阵之积。
\end{proposition}
\begin{proof}
在每个旋转块上有
\[
R(\theta)=
\begin{pmatrix}\cos\theta&\sin\theta\\\sin\theta&-\cos\theta\end{pmatrix}
\begin{pmatrix}1&0\\0&-1\end{pmatrix},
\]
两因子均实对称；在 $\pm1$ 的一阶块上也可如此配成两个对称因子。将各块组合，再用原正交矩阵作相同的正交相似变换，得到整个矩阵的分解。对任意实方阵，由定理~\ref{thm:xie-ip-polar} 的极分解 $A=QS$，$Q$ 正交且 $S$ 半正定实对称，便得到三个对称因子。
\end{proof}

还有两类常见结论可直接从标准型读取：对正交矩阵 $Q$，有 $\operatorname{rank}(I-Q)=\operatorname{rank}((I-Q)^2)$；两个实正规矩阵若相似，则必正交相似，因为含重数的复特征值完全确定一阶块与共轭二阶块，块内 $b$ 的符号又可由正交换基改变。另一方面，对称正交矩阵只有 $\pm1$ 特征值，故按负特征值个数分成 $n+1$ 个正交相似类。

\section{同时正交对角化与同时正交标准化}
\label{sec:xie-ip-simultaneous-orthogonal}
% 对应教材 9.13，例 9.123--9.130。

\begin{theorem}
\label{thm:xie-ip-simultaneous-normal}
有限个实对称矩阵可以同时正交对角化，当且仅当它们两两交换；有限个复正规矩阵可以同时酉对角化，当且仅当它们两两交换。
\end{theorem}
\begin{proof}
必要性由对角矩阵交换得到。对充分性，若所有矩阵均为纯量矩阵则结论成立。否则任选一个非纯量矩阵，以其特征子空间作正交直和分解。与它交换的各矩阵均保持每个特征子空间，且由命题~\ref{prop:xie-ip-reducing}，各限制仍自伴随或正规，并保持两两交换。每个子空间维数均严格小于原空间，故对维数归纳，在各子空间选公共标准正交特征基，合并即得全空间所需的基。
\end{proof}

对于两两交换的实正规矩阵，同样存在公共的实标准正交基，使每个矩阵均由一阶块和 $C(a,b)$ 二阶块组成，并且各矩阵采用同一组一维、二维子空间。其理由是先在复数域同时酉对角化：若公共特征值全实，可在相应公共特征子空间内选实基；若某个公共特征值非实，就把公共特征向量的实部和虚部配成二维实空间，由定理~\ref{thm:xie-ip-real-normal} 的论证得到该空间的标准正交基。正交补对所有算子及其伴随均不变，再作归纳。这里称为同时正交标准化，不能一般要求所有二阶块也化成实对角块。

同时合同与同时正交对角化的条件不同。对多个实对称矩阵 $A_1\succ0,A_2,\ldots,A_m$，存在可逆 $C$ 使 $C^{\mathsf T}A_1C=I$ 且其余全部对角化，当且仅当
\begin{equation}
\label{eq:xie-ip-multiple-congruence}
A_iA_1^{-1}A_j=A_jA_1^{-1}A_i\qquad(2\leq i,j\leq m).
\end{equation}
事实上，将 $A_1$ 化为单位矩阵后，待对角化的是实对称矩阵 $A_1^{-1/2}A_iA_1^{-1/2}$；它们两两交换正好等价于式~\eqref{eq:xie-ip-multiple-congruence}。若其余矩阵为实反对称矩阵，同一交换条件则给出共同的一阶、二阶合同标准化。

\begin{example}
设 $A\succ0$、$B,C\succeq0$ 均为实对称矩阵，并满足 $BA^{-1}C=CA^{-1}B$。证明
\[
\det A\det(A+B+C)\leq\det(A+B)\det(A+C),
\]
并求等号条件。
\end{example}
\begin{solve}
由式~\eqref{eq:xie-ip-multiple-congruence} 同时合同对角化，使三者变为 $I,\operatorname{diag}(b_i),\operatorname{diag}(c_i)$，其中 $b_i,c_i\geq0$。两边有相同的行列式变换因子，故问题化为
\[
\prod_i(1+b_i+c_i)\leq\prod_i(1+b_i)(1+c_i),
\]
由 $1+b_i+c_i\leq(1+b_i)(1+c_i)$ 即得。所有因子均正，等号当且仅当每个 $b_ic_i=0$，等价于 $BA^{-1}C=0$。
\end{solve}

\begin{example}
设 $A,B$ 实对称，$A\succeq0$，且 $AB+BA=0$，证明 $AB=BA=0$，并说明二者可以同时正交对角化。
\end{example}
\begin{solve}
先把 $A$ 正交化为 $\operatorname{diag}(\lambda_i)$，其中 $\lambda_i\geq0$。矩阵方程逐项为 $(\lambda_i+\lambda_j)b_{ij}=0$，故只要 $\lambda_i$ 或 $\lambda_j$ 非零，相应的 $b_{ij}$ 必为零。因此 $B$ 仅在 $\ker A$ 上可能非零，从而 $AB=BA=0$；再在 $\ker A$ 中正交对角化 $B$ 即得到公共标准型。若把 $B$ 换为反对称矩阵 $S$，同一逐项论证也给出 $AS=SA=0$，随后应使用同时正交标准化。
\end{solve}

\section{谱分解、极分解、奇异值分解及其应用}
\label{sec:xie-ip-decompositions}
% 对应教材 9.14，例 9.131--9.140。

\subsection{谱分解与函数演算}

\begin{theorem}
\label{thm:xie-ip-spectral-resolution}
设 $T$ 为复内积空间上的正规算子，或实内积空间上的自伴随算子，$\lambda_1,\ldots,\lambda_s$ 为其不同特征值，$E_j$ 为到 $\ker(T-\lambda_jI)$ 的正交投影，则
\begin{equation}
\label{eq:xie-ip-spectral-resolution}
T=\sum_{j=1}^s\lambda_jE_j,\qquad
E_j^*=E_j,\quad E_j^2=E_j,\quad
E_iE_j=0\ (i\ne j),\quad \sum_jE_j=I.
\end{equation}
这些投影由 $T$ 及特征值唯一确定，并且
\begin{equation}
\label{eq:xie-ip-spectral-projector}
E_j=\prod_{\ell\ne j}\frac{T-\lambda_\ell I}{\lambda_j-\lambda_\ell}.
\end{equation}
\end{theorem}
\begin{proof}
标准正交特征基把全空间分解为各特征子空间的正交直和。在第 $j$ 个子空间上，$T$ 为 $\lambda_jI$，相应 $E_j$ 为恒等，其余投影为零，故式~\eqref{eq:xie-ip-spectral-resolution} 逐子空间成立。同样检验式~\eqref{eq:xie-ip-spectral-projector}，它在第 $j$ 个子空间上为 $I$，在其他子空间上为零，因此它等于 $E_j$，也给出唯一性。
\end{proof}

谱分解（spectral decomposition）把算子计算化为有限个标量计算：对定义在谱上的任意函数 $f$，定义 $f(T)=\sum_jf(\lambda_j)E_j$；多项式情形与通常的代入一致。自伴随算子正定、半正定与否可逐特征值判断；半正定算子的算术平方根是 $\sum_j\sqrt{\lambda_j}E_j$；正规算子为酉算子当且仅当全部 $|\lambda_j|=1$。对实正规但非自伴随的算子，一般不能把实空间分解为实特征子空间的正交直和，应使用一维与二维的实标准型，或明确在复数域中讨论。

反过来，若一组算子 $E_j$ 满足式~\eqref{eq:xie-ip-spectral-resolution} 中的自伴随、幂等、两两乘积为零及和为 $I$ 的条件，则对任意复数 $\lambda_j$，$T=\sum_j\lambda_jE_j$ 均为正规算子，因为 $T^*T=TT^*=\sum_j|\lambda_j|^2E_j$。若所有 $\lambda_j$ 为实数，$T$ 还自伴随。这说明谱分解的条件本身也是一个正规性判据。

\subsection{极分解与唯一性}

\begin{theorem}
\label{thm:xie-ip-polar}
任一 $n$ 阶实或复矩阵 $A$ 均有\emph{极分解（polar decomposition）}
\begin{equation}
\label{eq:xie-ip-polar}
A=US,\qquad S=(A^{\mathsf H}A)^{1/2}\succeq0,
\end{equation}
其中 $U$ 为正交或酉矩阵。$S$ 总是唯一；若 $A$ 可逆，则 $U=AS^{-1}$ 也唯一。另有左极分解 $A=HU$，其中 $H=(AA^{\mathsf H})^{1/2}$。若 $A$ 奇异，全空间上的酉因子一般不唯一，但其在 $(\ker A)^\perp$ 上的作用已由 $A$ 唯一确定。
\end{theorem}
\begin{proof}
$S^2=A^{\mathsf H}A$ 使 $\|Sx\|=\|Ax\|$，故 $\ker S=\ker A$，且 $S$ 在 $E=(\ker A)^\perp$ 上可逆。定义 $U_0(Sx)=Ax$，其中 $x\in E$，则由范数等式及极化恒等式可知 $U_0:E\to\operatorname{Im}A$ 为保积同构。两边正交补维数相等，把 $U_0$ 扩张为全空间的酉或正交矩阵 $U$，即可得到 $A=US$。若还有 $A=\widetilde U\widetilde S$，则 $A^{\mathsf H}A=\widetilde S^2$，由半正定平方根唯一性得 $\widetilde S=S$。当 $S$ 可逆时再得 $U=AS^{-1}$；当其有核时，酉因子在核上的延拓仍有选择。由 $AA^{\mathsf H}=US^2U^{\mathsf H}$ 得 $H=USU^{\mathsf H}$，所以 $A=HU$。
\end{proof}

正规矩阵还可以选择彼此交换的极因子。复情形中，若 $A=\sum_j\lambda_jE_j$，取
\[
S=\sum_j|\lambda_j|E_j,\qquad
U=\sum_{\lambda_j\ne0}\frac{\lambda_j}{|\lambda_j|}E_j
+\sum_{\lambda_j=0}E_j,
\]
则 $A=US$ 且 $US=SU$。实正规矩阵按每个旋转伸缩块取模与旋转、按每个实特征值取绝对值与符号，也得到实的交换极分解。反过来，若 $A=US$ 且 $U$ 酉、$S\succeq0$、$US=SU$，则 $A^{\mathsf H}A=S^2=AA^{\mathsf H}$，所以 $A$ 正规。一般矩阵的极分解并不要求两因子交换。

\subsection{奇异值分解}

\begin{theorem}
\label{thm:xie-ip-svd}
设 $A\in\mathbb F^{m\times n}$，$\operatorname{rank}A=r$，则存在 $m$ 阶、$n$ 阶酉矩阵 $U,V$，实情形中可取正交矩阵，使
\begin{equation}
\label{eq:xie-ip-svd}
A=U\Sigma V^{\mathsf H},\qquad
\Sigma=\begin{pmatrix}D&0\\0&0\end{pmatrix}\in\mathbb R^{m\times n},
\qquad D=\operatorname{diag}(\sigma_1,\ldots,\sigma_r),
\end{equation}
其中 $\sigma_1\geq\cdots\geq\sigma_r>0$。这称为\emph{奇异值分解（singular value decomposition，SVD）}；这些正数是 $A^{\mathsf H}A$ 的正特征值的算术平方根，称为正奇异值。零块的大小由 $m,n,r$ 决定，$r=0$ 时 $A=0$。
\end{theorem}
\begin{proof}
对半正定 Hermite 矩阵 $A^{\mathsf H}A$ 取标准正交特征基 $v_1,\ldots,v_n$，令其正特征值为 $\sigma_j^2$。对 $1\leq j\leq r$ 定义 $u_j=Av_j/\sigma_j$，则
\[
u_i^{\mathsf H}u_j=\frac{v_i^{\mathsf H}A^{\mathsf H}Av_j}{\sigma_i\sigma_j}=\delta_{ij}.
\]
其中 $\delta_{ij}$ 是 Kronecker 符号。把 $u_1,\ldots,u_r$ 扩张为 $\mathbb F^m$ 的标准正交基。对 $j>r$，$\|Av_j\|^2=v_j^{\mathsf H}A^{\mathsf H}Av_j=0$，故 $Av_j=0$。令 $U,V$ 分别由这两组基组成，便有 $AV=U\Sigma$，即式~\eqref{eq:xie-ip-svd}。同时，$A^{\mathsf H}u_j=\sigma_jv_j$，说明 $u_j,v_j$ 分别是相应的左、右奇异向量。
\end{proof}

SVD 的定义域与值域通常使用\emph[caution]{两组独立选择的标准正交基}；它是左右酉等价分解，一般不是相似分解。即使是方阵，奇异值也不等于特征值的绝对值，只有正规矩阵才有这一对应。例如 $\begin{pmatrix}0&1\\0&0\end{pmatrix}$ 的特征值全零，正奇异值却为 $1$。按式~\eqref{eq:xie-ip-svd} 还可直接读取
\[
\operatorname{Im}A=\operatorname{span}(u_1,\ldots,u_r),\quad
\ker A=\operatorname{span}(v_{r+1},\ldots,v_n),\quad
\|A\|_F^2=\sum_{j=1}^r\sigma_j^2.
\]
定义算子 $2$-范数为 $\|A\|_2=\max_{\|x\|=1}\|Ax\|$，则非零矩阵满足 $\|A\|_2=\sigma_1$。记 $\operatorname{spec}(A)$ 为方阵 $A$ 的特征值集合，则其谱半径 $\rho(A)=\max_{\lambda\in\operatorname{spec}(A)}|\lambda|$ 满足 $\rho(A)\leq\sigma_1$，也可由式~\eqref{eq:xie-ip-eigen-singular-bound} 得到。

方阵的 SVD 与极分解可以互相联系：由 $A=U\Sigma V^{\mathsf H}$ 得 $A=(UV^{\mathsf H})(V\Sigma V^{\mathsf H})$，后一个因子正是 $(A^{\mathsf H}A)^{1/2}$。但 SVD 的左右奇异向量在重奇异值空间内可共同旋转，零奇异值空间内也可选不同基；奇异值有序后唯一，分解矩阵一般并不唯一。

\begin{example}
设 $A$ 为秩 $r$ 的实方阵，证明 $(\operatorname{tr}A)^2\leq r\operatorname{tr}(A^{\mathsf T}A)$，并求等号条件。
\end{example}
\begin{solve}
令 $P$ 为到 $\operatorname{Im}A$ 的正交投影，则 $PA=A$、$\|P\|_F^2=\operatorname{tr}P=r$，并有
\[
\langle A,P\rangle_F=\operatorname{tr}(PA)=\operatorname{tr}A.
\]
由 Frobenius 内积的 Cauchy--Schwarz 不等式得所需结论。若 $r>0$，等号当且仅当 $A=cP$，其中 $c\ne0$ 为实数，亦即 $A$ 实对称且所有非零特征值都相同；$A=0$ 时也取等号。该证明把教材的奇异值分析改写为正交投影上的一次内积估计。
\end{solve}

\subsection{Moore--Penrose 广义逆与最小二乘}

\begin{theorem}
\label{thm:xie-ip-pseudoinverse}
对任意 $A\in\mathbb F^{m\times n}$，存在唯一 $A^\dagger\in\mathbb F^{n\times m}$，满足
\begin{equation}
\label{eq:xie-ip-penrose}
AA^\dagger A=A,\qquad A^\dagger AA^\dagger=A^\dagger,
\qquad(AA^\dagger)^{\mathsf H}=AA^\dagger,\qquad
(A^\dagger A)^{\mathsf H}=A^\dagger A.
\end{equation}
它称为 Moore--Penrose 广义逆（Moore--Penrose pseudoinverse）。若 $A=U\Sigma V^{\mathsf H}$ 为 SVD，则
\[
A^\dagger=V\Sigma^\dagger U^{\mathsf H},\qquad
\Sigma^\dagger=\begin{pmatrix}D^{-1}&0\\0&0\end{pmatrix}\in\mathbb R^{n\times m}.
\]
特别地，$AA^\dagger$ 是到 $\operatorname{Im}A$ 的正交投影，$A^\dagger A$ 是到 $(\ker A)^\perp$ 的正交投影。
\end{theorem}
\begin{proof}
用 SVD 代入可逐块验证四个等式，同时得到两个投影结论。对唯一性，设 $B$ 也满足四个条件。$AB$ 自伴随且幂等，$\operatorname{Im}(AB)=\operatorname{Im}A$：一侧由乘积形式得到，另一侧由 $ABA=A$ 得到。因此 $AB$ 必为到 $\operatorname{Im}A$ 的正交投影。同样，$BA$ 自伴随幂等，且由 $ABA=A$ 有 $\ker(BA)=\ker A$，所以它是到 $(\ker A)^\perp$ 的正交投影。若 $b=Ax$，则 $Bb=BAx$ 是 $x$ 在 $(\ker A)^\perp$ 上的唯一投影；若 $b\perp\operatorname{Im}A$，则 $ABb=0$，再由 $BAB=B$ 得 $Bb=BABb=0$。由于这两个子空间正交直和为整个值域空间，$B$ 的作用处处唯一，故 $B=A^\dagger$。
\end{proof}

\begin{theorem}
\label{thm:xie-ip-least-squares}
对任意 $A\in\mathbb F^{m\times n}$ 与 $b\in\mathbb F^m$，向量 $x_*=A^\dagger b$ 是所有使 $\|Ax-b\|$ 最小的向量中唯一范数最小者。全体最小二乘解为
\begin{equation}
\label{eq:xie-ip-least-squares}
x=A^\dagger b+(I-A^\dagger A)z,\qquad z\in\mathbb F^n,
\end{equation}
等价地，它们是正规方程 $A^{\mathsf H}Ax=A^{\mathsf H}b$ 的全部解。若 $Ax=b$ 有解，则上述向量恰为其全部解，$A^\dagger b$ 为唯一最小范数解。
\end{theorem}
\begin{proof}
令 $p=AA^\dagger b$，它是 $b$ 在 $\operatorname{Im}A$ 上的正交投影，因此
\[
\|Ax-b\|^2=\|Ax-p\|^2+\|p-b\|^2.
\]
故最小二乘解恰满足 $Ax=p$，等价于 $x-x_*\in\ker A$，从而得到式~\eqref{eq:xie-ip-least-squares}。又 $x_*\in(\ker A)^\perp$，故对 $v\in\ker A$，$\|x_*+v\|^2=\|x_*\|^2+\|v\|^2$，唯一最小范数性质成立。另一方面，$Ax-b\perp\operatorname{Im}A$ 等价于 $A^{\mathsf H}(Ax-b)=0$，所以正规方程总有解，即使原方程无解也是如此。
\end{proof}

当 $A$ 列满秩时，$A^{\mathsf H}A$ 正定，$A^\dagger=(A^{\mathsf H}A)^{-1}A^{\mathsf H}$，最小二乘解唯一；当 $A$ 行满秩时，$A^\dagger=A^{\mathsf H}(AA^{\mathsf H})^{-1}$，所有右端都有解，而该公式选出最小范数解。若方阵 $A$ 可逆，则 $A^\dagger=A^{-1}$；零矩阵的广义逆仍为零。这里的“正规方程”是最小二乘的驻点方程，与“正规矩阵”是不同概念。

\begin{example}
设 $A=\begin{pmatrix}1&1\\1&1\\0&0\end{pmatrix}$、$b=(1,0,1)^{\mathsf T}$，求全部最小二乘解及其中范数最小者。
\end{example}
\begin{solve}
$A$ 的秩为 $1$，唯一正奇异值为 $2$，且
\[
A^\dagger=\frac14\begin{pmatrix}1&1&0\\1&1&0\end{pmatrix},
\qquad x_*=\begin{pmatrix}1/4\\1/4\end{pmatrix},
\qquad Ax_*=\begin{pmatrix}1/2\\1/2\\0\end{pmatrix}.
\]
残差 $b-Ax_*=(1/2,-1/2,1)^{\mathsf T}$ 与列空间 $\operatorname{span}\{(1,1,0)^{\mathsf T}\}$ 正交。因此全部最小二乘解为 $(1/4,1/4)^{\mathsf T}+t(1,-1)^{\mathsf T}$（$t\in\mathbb R$），最小残差平方为 $3/2$；其中 $t=0$ 时的解范数最小。这个例子同时区分了“残差最小”与“解向量范数最小”两个目标。
\end{solve}

\section{基础训练：综合方法与易错辨析}
\label{sec:xie-ip-training}
% 对应教材 9.15；将选择、填空与解答训练按概念整合。

基础训练首先检验\emph[tip]{条件是否足以保证所用的标准型}。实对称矩阵有标准正交特征基，但不一定有互异特征值，同一重特征值的任意特征向量也未必两两正交；自伴随算子在一般基中的矩阵未必对称，实情形中正确判据是 $A^{\mathsf T}G=GA$。正交变换在标准正交基中的表示矩阵正交，在一般基中则只保证 $A^{\mathsf T}GA=G$。任意可逆矩阵作相似变换并不保持矩阵的正交性，例如 $Q=\begin{pmatrix}0&-1\\1&0\end{pmatrix}$、$P=\operatorname{diag}(2,1)$ 时，$P^{-1}QP=\begin{pmatrix}0&-1/2\\2&0\end{pmatrix}$ 已不正交。

分类问题还要区分特征多项式与极小多项式：实对称矩阵由含重数的特征值决定其正交相似类，因此特征多项式相同就正交相似；但 $\operatorname{diag}(1,1,2)$ 与 $\operatorname{diag}(1,2,2)$ 的极小多项式同为 $(t-1)(t-2)$，却不相似。类似地，正规且幂零的算子必为零，复正规算子特征值全实时必自伴随；如果只知道在某组一般基中可对角化，这两类正规性结论都不能无条件使用。

\begin{example}
设 $u_1,\ldots,u_r$ 线性无关，$U=\operatorname{span}(u_1,\ldots,u_r)$。用 Gram 行列式表示 $x$ 到 $U$ 的距离。
\end{example}
\begin{solve}
写 $x=P_Ux+z$，其中 $z\perp U$。从向量组 $(u_1,\ldots,u_r,x)$ 的最后一个向量减去 $P_Ux$，对应行列式为 $1$ 的基变换，故 Gram 行列式不变。新 Gram 矩阵是 $G(u_1,\ldots,u_r)$ 与 $\|z\|^2$ 的分块对角矩阵，因此
\[
d(x,U)^2=\|z\|^2
=\frac{\det G(u_1,\ldots,u_r,x)}{\det G(u_1,\ldots,u_r)}.
\]
分母正定非零；当 $x\in U$ 时分子为零，公式仍成立。实情形中，它就是“体积等于底面积乘以高”的平方形式。
\end{solve}

\begin{example}
设 $n$ 维欧氏空间中有 $n+1$ 个点 $a_0,\ldots,a_n$，任意两点的距离均为 $d>0$。证明 $a_i-a_0$（$1\leq i\leq n$）构成一组基。
\end{example}
\begin{solve}
令 $v_i=a_i-a_0$，则 $\|v_i\|=d$，而 $\|v_i-v_j\|^2=d^2$ 给出 $\langle v_i,v_j\rangle=d^2/2$（$i\ne j$）。所以
\[
G(v_1,\ldots,v_n)=\frac{d^2}{2}(I_n+\boldsymbol1\boldsymbol1^{\mathsf T}),
\qquad \det G=(n+1)\left(\frac{d^2}{2}\right)^n>0.
\]
由命题~\ref{prop:xie-ip-gram-rank} 即得线性无关。这说明这 $n+1$ 个等距点确实张成一个非退化的 $n$ 维单纯形。
\end{solve}

\begin{example}
设 $a,b$ 是内积空间中的非零向量，定义秩一算子 $T(x)=\langle x,a\rangle b$。求其伴随，并判定它何时正规。
\end{example}
\begin{solve}
由
\[
\langle Tx,y\rangle=\langle x,a\rangle\langle b,y\rangle
=\langle x,\langle y,b\rangle a\rangle
\]
得 $T^*(y)=\langle y,b\rangle a$。进而
\[
T^*Tx=\|b\|^2\langle x,a\rangle a,\qquad
TT^*x=\|a\|^2\langle x,b\rangle b.
\]
二者相等必使像空间 $\operatorname{span}\{a\}$ 与 $\operatorname{span}\{b\}$ 相同，故 $a,b$ 线性相关；反之若 $b=ca$，则 $T$ 是秩一正交投影的某个标量倍数，确实正规。特别地，$a,b$ 为正交单位向量时，在以它们开头的标准正交基下，$T,T^*$ 的矩阵分别是 $E_{21},E_{12}$，该算子不正规。
\end{solve}

\begin{example}
设 $P^2=P$。证明 $\|Px\|\leq\|x\|$ 对所有 $x$ 成立，当且仅当 $P$ 是正交投影，也当且仅当 $P=P^*$。
\end{example}
\begin{solve}
幂等性给出 $V=\operatorname{Im}P\oplus\ker P$。若 $u\in\operatorname{Im}P$、$v\in\ker P$，则 $P(u+tv)=u$，收缩性使 $\|u\|^2\leq\|u+tv\|^2$ 对所有实数 $t$ 成立，因而展开后一次项必须为零，即 $\operatorname{Re}\langle u,v\rangle=0$；在复空间中再以 $iv$ 代替 $v$，得到虚部也为零。因此两子空间正交，$P$ 即正交投影，其矩阵在相应标准正交基下为 $\operatorname{diag}(I,0)$，所以自伴随。反之，自伴随幂等矩阵具有这一标准型，因而是正交投影并满足范数收缩。
\end{solve}

\begin{example}
证明八阶实正交矩阵中不可能出现所有元素均为 $1/\sqrt8$ 的三阶子矩阵。
\end{example}
\begin{solve}
若该子矩阵为 $B$，将所在的三行在其他五列上的元素记为 $C$。由这三行标准正交，
\[
CC^{\mathsf T}=I_3-BB^{\mathsf T}
=I_3-\frac38\boldsymbol1\boldsymbol1^{\mathsf T}.
\]
右边在 $\boldsymbol1=(1,1,1)^{\mathsf T}$ 上的特征值为 $1-9/8=-1/8$，与 $CC^{\mathsf T}\succeq0$ 矛盾。利用子向量组的 Gram 矩阵，可以将“某些元素能否同时出现”转化为半正定性检验。
\end{solve}

\begin{example}
设 $A\in\mathbb F^{m\times n}$ 列满秩，证明 $P=A(A^{\mathsf H}A)^{-1}A^{\mathsf H}$ 可写为 $P=QQ^{\mathsf H}$，其中 $Q^{\mathsf H}Q=I_n$。
\end{example}
\begin{solve}
取 $Q=A(A^{\mathsf H}A)^{-1/2}$，则 $Q^{\mathsf H}Q=I_n$ 且 $QQ^{\mathsf H}=P$。因而 $P$ 是到 $\operatorname{Im}A$ 上的正交投影；也可由薄 QR 分解 $A=QR$ 得到同一个投影。这里 $A$ 的列向量未必标准正交，而乘以 $(A^{\mathsf H}A)^{-1/2}$ 正是把其度量归一化。
\end{solve}

\begin{example}
不预先假设线性。设映射 $F:V\to V$ 满足 $\langle F(x),F(y)\rangle=\langle x,y\rangle$ 对所有 $x,y$ 成立，其中 $V$ 为有限维实或复内积空间。证明 $F$ 自动是线性的保积自同构。
\end{example}
\begin{solve}
取标准正交基 $e_1,\ldots,e_n$，则 $F(e_1),\ldots,F(e_n)$ 也是标准正交基。任意 $x$ 的像在此基下展开为
\[
F(x)=\sum_{j=1}^n\langle F(x),F(e_j)\rangle F(e_j)
=\sum_{j=1}^n\langle x,e_j\rangle F(e_j).
\]
右端关于 $x$ 线性，故 $F$ 线性；它将一组基映为一组基，因而为同构。这个证明先利用有限维性把像基变成全空间的基，再由坐标恢复线性，避免把题目未给出的线性当作已知条件。
\end{solve}
% END textbook-xie-exercises-ch9


\chapter{双线性型}
\label{chap:laex-bilinear}

% 教材来源：高等代数/textbooks/高等代数习题（谢启鸿）/高代习题-10.pdf。
% 覆盖原书第 10 章 §§10.1--10.6，PDF 第 1--23 页，印刷页 604--626。
% 按主题整合基本概念、例 10.1--10.30 及基础训练，不逐项抄录选择题答案。
双线性型用一个标量记录两个向量之间的关系，对偶空间则把其中一个变量固定后得到的线性函数组织起来。本章先建立对偶基与零化子，再用矩阵表示双线性型，最后研究交错型和对称型所确定的几何结构。解题时要区分三种操作：线性映射换基产生相似变换，两个空间分别换基产生相抵变换，同一空间中的双线性型同时换基产生合同变换。它们所保留的信息不同，不能相互替代。

以下 $\mathbb F$ 为数域，各空间均在 $\mathbb F$ 上；除明确说明外，空间有限维，坐标写成列向量。$V^*=\mathcal L(V,\mathbb F)$ 中的星号表示代数对偶，$A^{\mathsf T}$ 表示矩阵转置，$A^{\mathsf H}$ 表示共轭转置。本章讨论的复数域上的双线性型对两个变量都线性，不含共轭，因此也不同于复内积。

\section{基本概念}
\label{sec:laex-bilinear-basics}

\subsection{对偶空间、对偶基与自然配对}

\begin{definition}
\label{def:laex-dual}
线性函数（linear functional）是线性映射 $f:V\to\mathbb F$。所有线性函数按逐点加法和数乘构成的空间称为 $V$ 的\emph{代数对偶空间（algebraic dual space）}，记为 $V^*$，自然配对记为 $\langle f,v\rangle=f(v)$。若 $e_1,\ldots,e_n$ 是 $V$ 的基，定义坐标函数 $e^i$ 为
\[
 e^i\left(\sum_{j=1}^n x_je_j\right)=x_i,
 \qquad e^i(e_j)=\delta_{ij}.
\]
则 $e^1,\ldots,e^n$ 称为这组基的\emph{对偶基（dual basis）}。
\end{definition}

任意 $f\in V^*$ 满足 $f=\sum_i f(e_i)e^i$，因为两边在每个基向量上的值相同；若 $\sum_i c_ie^i=0$，在 $e_j$ 处取值得 $c_j=0$。因此对偶基确为基，$\dim V^*=\dim V=n$，且有相互对应的展开式
\begin{equation}
 v=\sum_{i=1}^n e^i(v)e_i,\qquad
 f=\sum_{i=1}^n f(e_i)e^i.
 \label{eq:laex-dual-expansion}
\end{equation}

\begin{proposition}
\label{prop:laex-double-dual}
映射 $J_V:V\to V^{**}$，$J_V(v)(f)=f(v)$，是无需选基即可定义的线性同构。
\end{proposition}
\begin{proof}
评价对 $v$ 和 $f$ 分别线性，故 $J_V$ 有定义且线性。若 $J_V(v)=0$，则所有坐标函数 $e^i$ 在 $v$ 上均为零，由式~\eqref{eq:laex-dual-expansion} 得 $v=0$。所以 $J_V$ 单射，再由 $\dim V^{**}=\dim V$ 得满射。
\end{proof}

$V$ 与 $V^*$ 虽然维数相同，但仅凭向量空间结构没有指定的同构；$V$ 与 $V^{**}$ 的上述同构则由评价自然确定。无限维时，坐标函数通常不能张成整个代数对偶，不能照搬有限维的维数等式或双重对偶满射结论。

\subsection{对偶映射及其矩阵}

\begin{definition}
若 $T:V\to U$ 线性，定义其\emph{对偶映射（dual map）} $T^*:U^*\to V^*$ 为
\begin{equation}
 T^*f=f\circ T,\qquad
 \langle T^*f,v\rangle=\langle f,Tv\rangle.
 \label{eq:laex-dual-map}
\end{equation}
\end{definition}

评价恒等式唯一确定 $T^*$，并直接给出
\[
 (aT+bS)^*=aT^*+bS^*,\qquad
 (RT)^*=T^*R^*,\qquad I_V^*=I_{V^*}.
\]
第一式要求 $T,S$ 有相同的定义域和陪域，第二式要求 $T:V\to U$、$R:U\to W$。若 $T$ 可逆，则 $(T^{-1})^*=(T^*)^{-1}$。复数域上这里仍是线性的 $aT^*$，不会出现 $\overline a$。

\begin{theorem}
\label{thm:laex-dual-transpose}
若 $T:V\to U$ 在基 $e_1,\ldots,e_n$ 与 $u_1,\ldots,u_m$ 下的矩阵为 $A$，则 $T^*$ 在相应对偶基 $u^1,\ldots,u^m$ 与 $e^1,\ldots,e^n$ 下的矩阵为 $A^{\mathsf T}$。因而
\[
 \operatorname{rank}T^*=\operatorname{rank}T,
 \qquad \dim\ker T^*=m-\operatorname{rank}T,
\]
且 $T$ 单射当且仅当 $T^*$ 满射，$T$ 满射当且仅当 $T^*$ 单射。
\end{theorem}
\begin{proof}
由 $Te_j=\sum_i a_{ij}u_i$ 得 $(T^*u^i)(e_j)=a_{ij}$，故 $T^*u^i=\sum_j a_{ij}e^j$，其坐标列正是 $A$ 第 $i$ 行的转置。其余结论由转置不改变秩及维数公式推出。
\end{proof}

\subsection{双线性型、根子空间与非退化性}

\begin{definition}
\label{def:laex-bilinear}
映射 $g:U\times V\to\mathbb F$ 若对每个变量分别线性，则称为\emph{双线性型（bilinear form）}。给定 $U$ 的基 $u_1,\ldots,u_m$ 与 $V$ 的基 $v_1,\ldots,v_n$，矩阵 $G=(g(u_i,v_j))$ 称为其表示矩阵；若 $u,v$ 的坐标列分别为 $x,y$，则
\begin{equation}
 g(u,v)=x^{\mathsf T}Gy.
 \label{eq:laex-bilinear-matrix}
\end{equation}
定义左、右\emph{根子空间（left and right radicals）}为
\[
 L_g=\{u\in U:g(u,V)=0\},\qquad
 R_g=\{v\in V:g(U,v)=0\}.
\]
若 $L_g=0$ 且 $R_g=0$，称 $g$ \emph{非退化（nondegenerate）}。
\end{definition}

若两组基分别按可逆矩阵 $P,Q$ 变换，即新基向量的旧坐标列组成 $P,Q$，则 $x=Px'$、$y=Qy'$，所以
\begin{equation}
 G'=P^{\mathsf T}GQ.
 \label{eq:laex-bilinear-change}
\end{equation}
因此 $\operatorname{rank}G$ 与基无关，称为型的秩；由相抵标准型可分别选基使 $G'=\begin{pmatrix}I_r&0\\0&0\end{pmatrix}$。又 $L_g$ 的坐标满足 $G^{\mathsf T}x=0$，$R_g$ 的坐标满足 $Gy=0$，故
\begin{equation}
 \dim L_g=m-r,\qquad \dim R_g=n-r.
 \label{eq:laex-radical-dimensions}
\end{equation}
于是 $g$ 非退化当且仅当 $m=n=r$，即表示矩阵为可逆方阵。只要求一侧根子空间为零时，两个空间的维数可以不同。

\subsection{纯量积、对称型与交错型}

\begin{definition}
\label{def:laex-scalar-form}
同一空间上 $g:V\times V\to\mathbb F$ 的双线性型，教材称为\emph{纯量积（scalar product）}，这里不附加正定性。若 $g(x,y)=g(y,x)$，称为\emph{对称型（symmetric bilinear form）}；若 $g(x,x)=0$ 对每个 $x$ 成立，称为\emph{交错型（alternating bilinear form）}。由于数域的特征为零，后一条件等价于 $g(x,y)=-g(y,x)$。
\end{definition}

等价性来自 $0=g(x+y,x+y)=g(x,y)+g(y,x)$；反向令 $y=x$ 并除以 $2$。相应矩阵分别满足 $G^{\mathsf T}=G$ 与 $G^{\mathsf T}=-G$。任意纯量积唯一分解为
\begin{equation}
 g=g_{\mathrm s}+g_{\mathrm a},\qquad
 g_{\mathrm s}(x,y)=\frac{g(x,y)+g(y,x)}2,\qquad
 g_{\mathrm a}(x,y)=\frac{g(x,y)-g(y,x)}2.
 \label{eq:laex-form-split}
\end{equation}
两部分分别对称与交错，唯一性来自一个型若同时具有两种性质就只能为零。纯量积换基必须对两个变量使用同一个 $P$，故矩阵变为 $P^{\mathsf T}GP$，这称为\emph{合同变换（congruence transformation）}。

一般纯量积的正交关系有方向：$g(x,y)=0$ 不必推出 $g(y,x)=0$。对称型与交错型的正交关系都对称，因而可统一记
\[
 W^{\perp_g}=\{v\in V:g(v,w)=0\text{ 对所有 }w\in W\}.
\]
下标可避免将它与对偶空间中的零化子混淆。

\begin{proposition}
\label{prop:laex-reciprocal-orthogonality}
若纯量积 $g$ 满足 $g(x,y)=0$ 当且仅当 $g(y,x)=0$，则 $g$ 为对称型或交错型。
\end{proposition}
\begin{proof}
若所有 $g(x,x)$ 均为零，则按定义为交错型。否则选 $a$ 使 $d=g(a,a)\ne0$。任意 $x$ 都能唯一写成 $x=\alpha a+u$，其中 $g(a,u)=0$，由假设也有 $g(u,a)=0$。对任意这样的 $u,v$，置 $b=d^{-1}g(u,v)a-v$，则 $g(a+u,b)=0$；交换次序后得到 $0=g(b,a+u)=g(u,v)-g(v,u)$。故 $g$ 在 $a$ 的正交补上对称，在 $\mathbb Fa$ 上也对称，混合项为零，所以在整个 $V$ 上对称。
\end{proof}

\section{线性函数与对偶空间}
\label{sec:laex-dual-methods}

\subsection{用核空间刻画线性函数}

\begin{proposition}
\label{prop:laex-functional-kernels}
对任意维数的 $V$，两个非零线性函数 $f,h$ 线性相关，当且仅当 $\ker f=\ker h$。有限维时，每个超平面都是某个非零线性函数的核，每个子空间都能表示为有限个这类核的交。
\end{proposition}
\begin{proof}
若 $h=cf$，非零性使 $c\ne0$，两核相等。反过来，取 $a$ 使 $f(a)=1$；对任意 $v$，有 $v-f(v)a\in\ker f=\ker h$，故 $h(v)=h(a)f(v)$。若 $h(a)=0$ 就有 $h=0$，故比例系数非零。

有限维时，取子空间 $W$ 的基 $e_1,\ldots,e_k$ 并扩张成 $V$ 的基，则
\[
 W=\bigcap_{i=k+1}^n\ker e^i.
\]
特别地 $k=n-1$ 时只需一个核；$W=V$ 时把空交理解为 $V$。
\end{proof}

\begin{example}
\label{ex:laex-functional-decomposition}
设 $f,h\in V^*$ 线性无关。证明任意 $v\in V$ 可写为 $v=u+w$，满足 $f(w)=f(v)$、$h(u)=h(v)$，并说明分解是否唯一。
\end{example}
\begin{solve}
条件等价于 $u\in\ker f$、$w\in\ker h$。映射 $v\mapsto(f(v),h(v))$ 的秩为 $2$，否则两个分量线性相关；所以 $\dim(\ker f\cap\ker h)=n-2$，由维数公式得 $\ker f+\ker h=V$，所求分解存在。两个分解之差满足 $u'-u=-(w'-w)\in\ker f\cap\ker h$，故 $n=2$ 时唯一，$n>2$ 时不唯一。
\end{solve}

\subsection{零化子把子空间转成线性约束}

\begin{definition}
\label{def:laex-annihilator}
对 $W\le V$，其\emph{零化子（annihilator）}定义为
\[
 W^\circ=\{f\in V^*:f(w)=0\text{ 对所有 }w\in W\}.
\]
对 $F\le V^*$，相应定义 $F^\circ=\{v\in V:f(v)=0\text{ 对所有 }f\in F\}$，即通过 $J_V$ 把双重零化子视为 $V$ 的子空间。
\end{definition}

\begin{theorem}
\label{thm:laex-annihilator}
若 $\dim V=n$，则对 $W,U\le V$ 有
\begin{equation}
 \dim W^\circ=n-\dim W,\qquad (W^\circ)^\circ=W,
 \label{eq:laex-annihilator-dim}
\end{equation}
以及
\begin{equation}
 (U+W)^\circ=U^\circ\cap W^\circ,\qquad
 (U\cap W)^\circ=U^\circ+W^\circ.
 \label{eq:laex-annihilator-sum-intersection}
\end{equation}
包含关系在取零化子后反向：$U\le W$ 蕴含 $W^\circ\le U^\circ$。
\end{theorem}
\begin{proof}
取 $W$ 的基并扩张成 $e_1,\ldots,e_n$，若前 $k$ 个生成 $W$，则 $W^\circ=\operatorname{span}\{e^{k+1},\ldots,e^n\}$，得到维数式。$W\subseteq(W^\circ)^\circ$ 由定义成立，两边同维即相等。一个函数在 $U+W$ 上为零，当且仅当它在 $U,W$ 上分别为零，得到第一个和交公式；将它应用于 $U^\circ,W^\circ$ 并取双重零化子，得到第二个公式。反向包含性质直接来自定义。
\end{proof}

这些关系把几何问题转成约束问题：生成空间增加时，能同时消去它的线性函数减少。计算 $W^\circ$ 时，若 $W$ 由坐标列 $w_1,\ldots,w_k$ 生成，令 $B=(w_1,\ldots,w_k)$，则 $f$ 的对偶基坐标列 $c$ 恰满足 $B^{\mathsf T}c=0$。

\subsection{对偶映射的像、核与不变子空间}

\begin{theorem}
\label{thm:laex-dual-image-kernel}
若 $T:V\to U$ 线性，则
\begin{equation}
 \ker T^*=(\operatorname{im}T)^\circ\subseteq U^*,\qquad
 \operatorname{im}T^*=(\ker T)^\circ\subseteq V^*.
 \label{eq:laex-dual-image-kernel}
\end{equation}
若 $T:V\to V$ 且 $W\le V$，则 $W$ 为 $T$ 的不变子空间，当且仅当 $W^\circ$ 为 $T^*$ 的不变子空间。
\end{theorem}
\begin{proof}
$T^*f=0$ 等价于 $f(Tv)=0$ 对全部 $v$ 成立，即第一个等式。若 $f=T^*h$，则 $f$ 在 $\ker T$ 上为零，得到第二式的一个包含。反向证明不必只数维数：若 $f\in(\ker T)^\circ$，在 $\operatorname{im}T$ 上定义 $h_0(Tv)=f(v)$。若 $Tv=Tv'$，则 $v-v'\in\ker T$，故定义无歧义。把 $\operatorname{im}T$ 的基扩张成 $U$ 的基，即可将 $h_0$ 延拓为 $h\in U^*$，于是 $f=T^*h$。

若 $TW\subseteq W$，则对 $f\in W^\circ,w\in W$，有 $(T^*f)(w)=f(Tw)=0$，所以 $T^*W^\circ\subseteq W^\circ$。反过来，该包含给出所有 $f\in W^\circ$ 在 $Tw$ 上都为零，于是 $Tw\in(W^\circ)^\circ=W$。
\end{proof}

\begin{notealert}[type=caution,title={代数对偶与无限维条件}]
上述像、核公式使用的是全体线性函数组成的代数对偶。若允许无限维空间中的基扩张，延拓论证仍成立，因而 $\operatorname{im}T^*=(\ker T)^\circ$ 也成立；不能把它说成一般无法推广到无限维。连续对偶还要求函数连续，须另查拓扑条件。本章其余涉及维数的推导仍按有限维理解。
\end{notealert}

对任意维数的幂等变换 $P^2=P$，第二个等式还有无需基扩张的直接证明。若 $f$ 在 $\ker P$ 上为零，则 $v-Pv\in\ker P$，故 $f(v)=f(Pv)$，即 $f=P^*f$。这说明 $\operatorname{im}P^*$ 中恰为在 $\ker P$ 上消失的函数。

\subsection{直和、双重对偶与欧氏空间中的识别}

外直和的函数可分别限制在两分量上，从而有自然同构
\[
 (U\oplus V)^*\longrightarrow U^*\oplus V^*,\qquad
 F\longmapsto\bigl(u\mapsto F(u,0),\ v\mapsto F(0,v)\bigr).
\]
其逆把 $(f,h)$ 送到 $(u,v)\mapsto f(u)+h(v)$，两者复合为恒等。这一论证不需要维数有限，但这里仅讨论两个分量的直和。

对 $T:V\to U$，自然评价给出 $T^{**}J_V=J_UT$：两边在任意 $f\in U^*$ 上的值都为 $f(Tv)$。有限维时通过 $J_V,J_U$ 识别双重对偶，才可简记为 $T^{**}=T$。

若 $V$ 是实欧氏空间，内积额外给出同构 $\iota:V\to V^*$，$\iota(u)(v)=(u,v)$。其单射性由正定性得出，有限维下再得满射。标准正交基经 $\iota$ 对应自身的对偶基；若 $T^\dagger$ 是满足 $(T^\dagger u,v)=(u,Tv)$ 的内积伴随，则
\[
 \iota^{-1}T^*\iota=T^\dagger.
\]
因此对偶变换与伴随变换的对应依赖内积提供的识别。一般基下不能直接把原基当作对偶基；复内积还涉及共轭线性，不能直接套用这里的实线性公式。

\section{双线性型与纯量积}
\label{sec:laex-bilinear-methods}

\subsection{非退化配对与相对对偶变换}

非退化配对 $g:U\times V\to\mathbb F$ 把 $U$ 通过 $u\mapsto g(u,-)$ 同构到 $V^*$。给定 $V$ 的任意基 $v_1,\ldots,v_n$，存在唯一一组 $U$ 的基 $u_1,\ldots,u_n$ 使 $g(u_i,v_j)=\delta_{ij}$，称为关于 $g$ 的一对对偶基。

\begin{proposition}
\label{prop:laex-relative-dual}
对 $T\in\mathcal L(V)$，存在唯一的 $T^\vee\in\mathcal L(U)$ 满足
\begin{equation}
 g(T^\vee u,v)=g(u,Tv).
 \label{eq:laex-relative-dual}
\end{equation}
在关于 $g$ 的对偶基下，若 $T$ 的矩阵为 $A$，则 $T^\vee$ 的矩阵为 $A^{\mathsf T}$。并有
\[
 (aT+bS)^\vee=aT^\vee+bS^\vee,\qquad
 (TS)^\vee=S^\vee T^\vee,
\]
且 $T$ 可逆时 $(T^{-1})^\vee=(T^\vee)^{-1}$。
\end{proposition}
\begin{proof}
令 $\iota_g:U\to V^*$，$\iota_g(u)=g(u,-)$，则取 $T^\vee=\iota_g^{-1}T^*\iota_g$ 即得存在性，非退化性保证唯一性。若 $T^\vee$ 的矩阵为 $B$，则 $b_{ji}=g(T^\vee u_i,v_j)=g(u_i,Tv_j)=a_{ij}$，所以 $B=A^{\mathsf T}$。运算法则由式~\eqref{eq:laex-relative-dual} 连续使用两次并利用唯一性得到。
\end{proof}

再次取相对对偶时，要交换 $U,V$ 并使用反向配对 $g^{\mathrm{op}}(v,u)=g(u,v)$，这时 $(T^\vee)^\vee=T$。若 $U=V$ 却始终使用同一个不对称配对，则不能无条件写出这个等式。

\begin{example}
\label{ex:laex-form-operator-matrix}
设非退化纯量积 $g$ 与纯量积 $h$ 在同一组基下的矩阵分别为 $G,H$。求满足 $h(x,y)=g(Tx,y)$ 的变换 $T$，并与作用在第二个变量上的情形比较。
\end{example}
\begin{solve}
若 $C$ 为 $T$ 的矩阵，则 $x^{\mathsf T}Hy=(Cx)^{\mathsf T}Gy=x^{\mathsf T}C^{\mathsf T}Gy$，故
\begin{equation}
 C^{\mathsf T}G=H,\qquad C=G^{-\mathsf T}H^{\mathsf T},
 \label{eq:laex-left-form-operator}
\end{equation}
其中 $G^{-\mathsf T}=(G^{-1})^{\mathsf T}$。若改为 $h(x,y)=g(x,Ty)$，则 $GC=H$，所以 $C=G^{-1}H$。两者一般不同。若 $h$ 也非退化，则两种构造中的 $T$ 均可逆；这也证明两个非退化配对可通过其中一个变量的唯一可逆变换相互联系。
\end{solve}

\subsection{抽取非退化部分与分别变换两变量}

\begin{proposition}
\label{prop:laex-bilinear-core}
设 $g:U\times V\to\mathbb F$ 的秩为 $r$。存在分解
\[
 U=U_0\oplus L_g,\qquad V=V_0\oplus R_g,
 \qquad \dim U_0=\dim V_0=r,
\]
使 $g$ 在 $U_0\times V_0$ 上非退化。两个相同定义空间上的双线性型 $g,h$ 秩相同，当且仅当存在 $S\in\operatorname{GL}(U)$、$T\in\operatorname{GL}(V)$ 使
\[
 h(u,v)=g(Su,Tv).
\]
\end{proposition}
\begin{proof}
将表示矩阵按式~\eqref{eq:laex-bilinear-change} 化为相抵标准型，取前 $r$ 个基向量分别生成 $U_0,V_0$，剩余向量分别生成左右根子空间，限制矩阵为 $I_r$。对第二个结论，分别变换两变量得到 $H=C^{\mathsf T}GD$，因左右乘子可逆，必要性成立；若两矩阵同秩，由相抵标准型可取这样的可逆 $C,D$，从而得到 $S,T$。
\end{proof}

即使 $U=V$，这里仍允许 $S,T$ 不同，不能据此推出同秩纯量积合同。例如实对称型 $x_1y_1+x_2y_2$ 与 $x_1y_1-x_2y_2$ 同秩，却有不同的正负惯性指数。对一般不对称型，也不能把上述两个补空间强行选成同一个子空间。

若 $W=U\oplus V$，在两分量上分别给定纯量积 $g,h$，则
\[
 q(u+v,u'+v')=g(u,u')+h(v,v')
\]
定义直和型，合并基下的矩阵为 $\operatorname{diag}(G,H)$。所以它非退化当且仅当两分量均非退化；两分量都对称或都交错时，直和型保留相应性质。

\subsection{非退化限制与正交直和}

\begin{theorem}
\label{thm:laex-orthogonal-splitting}
设 $g$ 为 $V$ 上的对称型或交错型，$W\le V$。其限制 $g|_{W\times W}$ 非退化，当且仅当 $W\cap W^{\perp_g}=0$；这时
\begin{equation}
 V=W\oplus W^{\perp_g},\qquad g(W,W^{\perp_g})=g(W^{\perp_g},W)=0.
 \label{eq:laex-orthogonal-splitting}
\end{equation}
若 $g$ 在整个 $V$ 上也非退化，则对任意子空间 $W$ 有
\[
 \dim W^{\perp_g}=\dim V-\dim W,\qquad
 (W^{\perp_g})^{\perp_g}=W.
\]
\end{theorem}
\begin{proof}
限制的根子空间就是 $W\cap W^{\perp_g}$，得到第一个等价关系。限制非退化时，映射 $\Phi:V\to W^*$，$\Phi(v)=g(v,-)|_W$，在 $W$ 上的限制已经是同构，故 $\Phi$ 满射，其核为 $W^{\perp_g}$。于是 $\dim V=\dim W+\dim W^{\perp_g}$，结合零交即得直和。

若全空间上的型非退化，则 $V\to V^*$ 的配对同构与限制映射 $V^*\to W^*$ 的复合仍满射，后者满射由基扩张延拓函数得到。因此维数公式对任意 $W$ 成立；双重正交补包含 $W$ 且与它同维，所以相等。
\end{proof}

\subsection{矩阵空间中的两种典型配对}

矩阵空间 $M_n(\mathbb F)$ 上的 $b(A,B)=\operatorname{tr}(AB)$ 是对称非退化双线性型。若对所有 $B$ 都有 $\operatorname{tr}(AB)=0$，取 $B=E_{ji}$ 得 $a_{ij}=0$，故 $A=0$。关于这一配对，基础矩阵 $E_{ij}$ 的对偶基为 $E_{ji}$。实数域上它不正定，例如 $A=\begin{pmatrix}0&1\\-1&0\end{pmatrix}$ 满足 $\operatorname{tr}(A^2)=-2$；Frobenius 内积则是 $(A,B)_{\mathrm F}=\operatorname{tr}(A^{\mathsf T}B)$，满足 $(A,A)_{\mathrm F}=\sum_{i,j}a_{ij}^2$。

\begin{example}
\label{ex:laex-frobenius-dual-identity}
设 $A_1,\ldots,A_{n^2}$ 与 $B_1,\ldots,B_{n^2}$ 是 $M_n(\mathbb R)$ 关于 Frobenius 内积的对偶基，即 $(A_i,B_j)_{\mathrm F}=\delta_{ij}$。证明
\[
 \sum_{i=1}^{n^2}A_iB_i=I_n.
\]
\end{example}
\begin{solve}
把基础矩阵依次记为 $E_1,\ldots,E_{n^2}$，写 $A_i=\sum_k p_{ki}E_k$、$B_i=\sum_l q_{li}E_l$。因 $E_k$ 为标准正交基，对偶条件为 $P^{\mathsf T}Q=I_{n^2}$，故 $Q=P^{-\mathsf T}$，也有 $PQ^{\mathsf T}=I_{n^2}$。于是
\[
 \sum_i A_iB_i
 =\sum_{k,l}\left(\sum_i p_{ki}q_{li}\right)E_kE_l
 =\sum_k E_k^2
 =\sum_{j=1}^n E_{jj}=I_n.
\]
最后一步用 $E_{ij}^2=0$（$i\ne j$）及 $E_{ii}^2=E_{ii}$。这里乘积没有转置；把配对公式中的转置机械移到结论里会得到另一恒等式。
\end{solve}

\section{交错型与辛几何}
\label{sec:laex-symplectic}

\subsection{交错标准型及其构造}

\begin{theorem}
\label{thm:laex-alternating-normal}
任意有限维空间上的交错型 $g$，都存在一组基使表示矩阵为
\begin{equation}
 \operatorname{diag}(\underbrace{S,\ldots,S}_{r\text{ 个}},0_s),
 \qquad S=\begin{pmatrix}0&1\\-1&0\end{pmatrix}.
 \label{eq:laex-alternating-normal}
\end{equation}
因此交错型的秩为偶数 $2r$；同阶交错矩阵合同当且仅当秩相同。
\end{theorem}
\begin{proof}
若 $g=0$，取任意基即可。否则存在 $e,f$ 使 $g(e,f)\ne0$，缩放 $f$ 后使 $g(e,f)=1$。在 $H=\operatorname{span}\{e,f\}$ 上的矩阵为 $S$，所以 $H$ 非退化。对任意 $v$，向量
\begin{equation}
 v'=v+g(f,v)e-g(e,v)f
 \label{eq:laex-symplectic-elimination}
\end{equation}
同时满足 $g(e,v')=g(f,v')=0$，从而 $V=H\oplus H^{\perp_g}$。在正交补上递归；每一步减少二维，最终剩下零型，得到标准型。合同保持秩，反之同秩对应相同标准型。
\end{proof}

\begin{definition}
配有非退化交错型的空间 $(V,g)$ 称为\emph{辛空间（symplectic vector space）}。其维数必为偶数 $2r$；使矩阵为 $J=\operatorname{diag}(S,\ldots,S)$ 的有序基 $e_1,f_1,\ldots,e_r,f_r$ 称为\emph{辛基（symplectic basis）}。即 $g(e_i,f_j)=\delta_{ij}$，而 $g(e_i,e_j)=g(f_i,f_j)=0$。
\end{definition}

\begin{example}
\label{ex:laex-three-dimensional-alternating}
设 $\dim V=3$，$h$ 是非零交错型。证明存在 $f,k\in V^*$ 使 $h(x,y)=f(x)k(y)-f(y)k(x)$。
\end{example}
\begin{solve}
非零性及偶数秩说明 $h$ 的秩为 $2$。取标准型对应基 $e_1,f_1,z$，使 $h(e_1,f_1)=1$ 且 $z$ 在根子空间中。定义 $f(x)=h(x,f_1)$、$k(x)=h(e_1,x)$，它们分别抽取 $e_1,f_1$ 坐标，代入标准型即得结论。这个表示依赖三维中的秩至多为 $2$；高维交错型一般需要多个这样的项相加。
\end{solve}

\begin{example}
\label{ex:laex-symplectic-basis}
四维空间中的交错型在基 $e_1,e_2,e_3,e_4$ 下的矩阵为
\[
 A=\begin{pmatrix}
 0&2&-1&3\\-2&0&4&-2\\1&-4&0&1\\-3&2&-1&0
 \end{pmatrix}.
\]
求一组辛基。
\end{example}
\begin{solve}
先取 $u_1=e_1$、$v_1=\tfrac12e_2$，使 $g(u_1,v_1)=1$。用式~\eqref{eq:laex-symplectic-elimination} 消去另外两向量与这一对的配对，得到
\[
 w_3=e_3+2e_1+\tfrac12e_2,\qquad
 w_4=e_4-e_1-\tfrac32e_2.
\]
两向量均在 $\operatorname{span}\{u_1,v_1\}^{\perp_g}$ 中，且直接计算得 $g(w_3,w_4)=6$。于是取 $u_2=w_3$、$v_2=\tfrac16w_4$，所求过渡矩阵为
\[
 C=\begin{pmatrix}
 1&0&2&-\tfrac16\\
 0&\tfrac12&\tfrac12&-\tfrac14\\
 0&0&1&0\\
 0&0&0&\tfrac16
 \end{pmatrix},\qquad
 C^{\mathsf T}AC=\operatorname{diag}(S,S).
\]
$C$ 可逆同时说明 $A$ 非退化。这一过程每次先归一化一对向量，再消去混合配对，适用于任意交错型。
\end{solve}

\subsection{辛变换及其矩阵判据}

\begin{definition}
辛空间上的可逆线性变换 $T$ 若满足 $g(Tx,Ty)=g(x,y)$，称为\emph{辛变换（symplectic transformation）}。在辛基下，其矩阵 $M$ 称为辛矩阵，并满足
\begin{equation}
 M^{\mathsf T}JM=J.
 \label{eq:laex-symplectic-matrix}
\end{equation}
\end{definition}

这一矩阵等式由 $g(Tx,Ty)=x^{\mathsf T}M^{\mathsf T}JMy$ 对所有坐标列成立得到；反过来也直接给出保型性，并由 $J$ 可逆推出 $M$ 可逆。辛变换恰把一组辛基变成另一组辛基：基向量的全部配对保持即保证任意线性组合的配对保持。恒等、复合及逆变换仍保型，因而组成辛群。由式~\eqref{eq:laex-symplectic-matrix} 得
\[
 M^{-1}=J^{-1}M^{\mathsf T}J,\qquad (\det M)^2=1.
\]
后一个等式只是取行列式得到的必要条件，不能取代完整的辛矩阵判据。二维时直接计算 $M^{\mathsf T}SM=(\det M)S$，故二维辛矩阵恰为行列式等于 $1$ 的矩阵。

\section{对称型与正交几何}
\label{sec:laex-orthogonal-geometry}

\subsection{正交基、惯性指数与迷向性}

\begin{theorem}
\label{thm:laex-symmetric-normal}
任意数域上的对称型都有一组正交基，使表示矩阵为 $\operatorname{diag}(d_1,\ldots,d_r,0,\ldots,0)$，其中 $d_i\ne0$。实数域上可归一化为 $\operatorname{diag}(I_p,-I_q,0_s)$，其中 $p,q,s$ 由型唯一确定；复数域上可归一化为 $\operatorname{diag}(I_r,0_s)$。
\end{theorem}
\begin{proof}
若 $g\ne0$，必有 $v$ 使 $g(v,v)\ne0$，否则极化恒等式
\[
 2g(x,y)=g(x+y,x+y)-g(x,x)-g(y,y)
\]
将使 $g=0$。一维空间 $\mathbb Fv$ 非退化，由定理~\ref{thm:laex-orthogonal-splitting} 分裂出其正交补，再归纳得到对角形。实数域中分别缩放正、负对角元使其成为 $1,-1$；复数域中每个非零数有平方根，故可全部化为 $1$。

实数域中，$p$ 是正定限制子空间的最大维数：前 $p$ 个坐标方向实现它，而任何维数大于 $p$ 的子空间都与其余非正方向的空间有非零交，不可能正定。负定情形同理确定 $q$，根子空间维数确定 $s$，所以三者不依赖选基。
\end{proof}

\begin{definition}
配有非退化对称型的 $(V,g)$ 称为\emph{正交空间（orthogonal space）}。非零向量 $v$ 若满足 $g(v,v)=0$，称为\emph{迷向向量（isotropic vector）}，亦称各向同性向量。含有迷向向量的子空间称为迷向子空间；若子空间 $W$ 上 $g(w,w)=0$ 对所有 $w$ 成立，则称为\emph{全迷向子空间（totally isotropic subspace）}。由极化恒等式，全迷向等价于 $g(W,W)=0$，也等价于 $W\subseteq W^{\perp_g}$。
\end{definition}

非退化意味着没有与所有向量正交的非零向量，并不排除迷向向量。例如双曲平面的矩阵 $\begin{pmatrix}0&1\\1&0\end{pmatrix}$ 可逆，但两个基向量都迷向。实正定空间没有迷向向量；实不定型则存在迷向向量，所以正交空间的几何不能一律按欧氏空间处理。

\begin{theorem}
\label{thm:laex-witt-real}
设实非退化对称型的正、负惯性指数为 $p,q$。全迷向子空间的最大维数为 $\min\{p,q\}$；每个按包含关系极大的全迷向子空间也有这一维数。
\end{theorem}
\begin{proof}
按正交标准型写 $V=V_+\oplus V_-$。若 $W$ 全迷向，向 $V_+$ 的投影在 $W$ 上单射：投影为零的非零向量会属于负定空间 $V_-$，不可能迷向。同理向 $V_-$ 的投影也单射，故 $\dim W\le\min\{p,q\}$。取正标准向量 $e_i$ 与负标准向量 $f_i$，$e_i+f_i$（$1\le i\le\min\{p,q\}$）张成全迷向子空间，达到上界。

还需排除维数较小但无法扩大的情形。设 $W$ 的基为 $w_1,\ldots,w_k$。非退化性使可选 $z_1,\ldots,z_k$ 满足 $g(w_i,z_j)=\delta_{ij}$；这是因为这些线性函数独立，所确定的 $V\to\mathbb R^k$ 映射满射。将
\[
 z_i'=z_i-\frac12\sum_{j=1}^k g(z_i,z_j)w_j
\]
代入可得 $g(z_i',z_j')=0$，而 $g(w_i,z_j')=\delta_{ij}$ 不变。因此 $H=\operatorname{span}\{w_i,z_i':1\le i\le k\}$ 非退化，其矩阵为 $\begin{pmatrix}0&I_k\\I_k&0\end{pmatrix}$，惯性指数为 $(k,k)$。在分解 $V=H\oplus H^{\perp_g}$ 中，正交补的惯性指数为 $(p-k,q-k)$。若 $k<\min\{p,q\}$，其中仍能取到非零迷向向量 $a$，于是 $W+\mathbb Ra$ 是更大的全迷向子空间，矛盾。
\end{proof}

\subsection{正交变换、伴随与镜像}

\begin{definition}
正交空间上的可逆线性变换 $T$ 若满足 $g(Tx,Ty)=g(x,y)$，称为\emph{正交变换（orthogonal transformation）}。任意基下，若型的矩阵为可逆对称矩阵 $G$，变换的矩阵为 $M$，则等价条件是
\begin{equation}
 M^{\mathsf T}GM=G.
 \label{eq:laex-orthogonal-matrix}
\end{equation}
\end{definition}

关于 $g$ 的伴随 $T^\dagger$ 由 $g(T^\dagger x,y)=g(x,Ty)$ 唯一确定，其矩阵为 $G^{-1}M^{\mathsf T}G$。因此 $T$ 正交当且仅当 $T^\dagger T=I$，并有 $\det T=\pm1$。恒等、复合与逆仍为正交变换，它们组成正交群；只有在 $G=I$ 的标准正交基下，判据才简化为 $M^{\mathsf T}M=I$。

\begin{proposition}
\label{prop:laex-reflection}
若 $a$ 满足 $g(a,a)\ne0$，则
\begin{equation}
 s_a(x)=x-\frac{2g(x,a)}{g(a,a)}a
 \label{eq:laex-reflection}
\end{equation}
是正交变换，称为关于 $a^{\perp_g}$ 的\emph{镜像变换（reflection）}。它在 $\mathbb Fa$ 上作用为 $-I$，在 $a^{\perp_g}$ 上作用为 $I$，故 $s_a^2=I$、$\det s_a=-1$。
\end{proposition}
\begin{proof}
将式~\eqref{eq:laex-reflection} 代入两个变量，两个交叉项之和为 $-4g(x,a)g(y,a)/g(a,a)$，二次修正项为其相反数，因此 $g(s_ax,s_ay)=g(x,y)$。其在两分量上的作用由公式得到，再用 $V=\mathbb Fa\oplus a^{\perp_g}$ 即得其余结论。
\end{proof}

\begin{example}
\label{ex:laex-hyperbolic-reflection}
设双曲平面有基 $u,v$ 满足 $g(u,u)=g(v,v)=0$、$g(u,v)=1$。证明其中行列式为 $-1$ 的正交变换必为镜像变换。
\end{example}
\begin{solve}
矩阵 $\begin{pmatrix}a&b\\c&d\end{pmatrix}$ 的保型条件给出 $ac=bd=0$、$ad+bc=1$。结合可逆性，它或为 $\operatorname{diag}(t,t^{-1})$，或为 $\begin{pmatrix}0&t\\t^{-1}&0\end{pmatrix}$，其中 $t\ne0$。前者行列式为 $1$，故题设对应后者。令 $a_0=tu-v$，则 $g(a_0,a_0)=-2t\ne0$，由式~\eqref{eq:laex-reflection} 计算得 $s_{a_0}(u)=t^{-1}v$、$s_{a_0}(v)=tu$，正是所给变换。这个结论依赖二维条件。
\end{solve}

\subsection{分块保距同构与不动子空间}

若 $V=U_1\mathbin{\perp_g}\cdots\mathbin{\perp_g}U_r=W_1\mathbin{\perp_g}\cdots\mathbin{\perp_g}W_r$ 是两组正交直和分解，且各有保距线性同构 $T_i:U_i\to W_i$，则
\[
 T(u_1+\cdots+u_r)=T_1u_1+\cdots+T_ru_r
\]
是 $V$ 上的正交变换。直和保证定义唯一及可逆，各分量间正交保证 $g(Tx,Ty)=\sum_i g(T_ix_i,T_iy_i)=g(x,y)$。

任意线性变换 $T$ 的不动子空间为 $\operatorname{Fix}(T)=\ker(I-T)$，维数公式总给出
\[
 \dim V=\dim\operatorname{Fix}(T)+\dim\operatorname{im}(I-T).
\]
若 $T$ 为正交变换，则还有 $\operatorname{im}(I-T)=\operatorname{Fix}(T)^{\perp_g}$：对 $w=Tw$，$g((I-T)v,w)=g(v,w)-g(Tv,Tw)=0$，再比较维数即可。但不定型下二者可能有非零交，\emph[caution]{维数相加不等于已经得到直和分解}；正定情形下限制非退化，才可直接得到正交直和。

\subsection{迹为零与迷向标准正交基}

\begin{theorem}
\label{thm:laex-isotropic-orthonormal}
设 $V$ 是 $n$ 维实欧氏空间，$g$ 是另一个对称双线性型。存在一组关于欧氏内积的标准正交基 $u_1,\ldots,u_n$，使每个 $u_i$ 都是关于 $g$ 的迷向向量，当且仅当 $g$ 在一组欧氏标准正交基下的矩阵 $A$ 满足 $\operatorname{tr}A=0$。
\end{theorem}
\begin{proof}
欧氏标准正交基之间的过渡矩阵 $Q$ 满足 $Q^{\mathsf T}Q=I$，故换基矩阵 $Q^{\mathsf T}AQ=Q^{-1}AQ$ 保持迹。若全部 $g(u_i,u_i)=0$，新矩阵对角元全零，必要性成立。

对充分性按 $n$ 归纳。$n=1$ 时 $A=0$，结论成立。若 $A=0$，任意标准正交基均可。否则实对称矩阵的谱定理给出标准正交特征向量；其特征值之和为零且不全为零，故有 $\lambda>0>\mu$ 及相应单位向量 $e,f$。令
\[
 u=\sqrt{\frac{-\mu}{\lambda-\mu}}\,e
   +\sqrt{\frac{\lambda}{\lambda-\mu}}\,f.
\]
则 $(u,u)=1$、$g(u,u)=0$。把 $u$ 扩张成欧氏标准正交基，此时 $A$ 的首个对角元为零，余下主块仍实对称且迹为零；在欧氏正交补上应用归纳假设即可。
\end{proof}

这里的两种关系由不同形式给出：标准正交性来自既定的正定欧氏内积，迷向性来自 $g$。不能将结论误读成非退化对称型存在一组关于自身既正交又全为迷向的基。

\subsection{Minkowski 空间中的正交限制}

\begin{example}
\label{ex:laex-minkowski}
在 $\mathbb R^4$ 上取
\[
 g(x,y)=x_1y_1+x_2y_2+x_3y_3-x_4y_4.
\]
按教材约定，$g(x,x)>0$ 称为空间向量，$g(x,x)<0$ 称为时间向量，非零且 $g(x,x)=0$ 称为光向量。证明时间向量不能与时间向量或光向量正交，两个光向量正交当且仅当它们线性相关。
\end{example}
\begin{solve}
写时间向量为 $x=(a,t)$，其中 $a\in\mathbb R^3$ 且 $\lVert a\rVert^2<t^2$，故 $t\ne0$。若 $y=(b,s)\in x^{\perp_g}$，则 $s=(a\cdot b)/t$，由 Cauchy--Schwarz 不等式得
\[
 g(y,y)=\lVert b\rVert^2-\frac{(a\cdot b)^2}{t^2}
 \ge\left(1-\frac{\lVert a\rVert^2}{t^2}\right)\lVert b\rVert^2.
\]
当 $y\ne0$ 时必有 $b\ne0$，故右边严格为正，得到前两个结论。

若 $x=(a,t)$、$y=(b,s)$ 为光向量，则 $\lVert a\rVert=|t|>0$、$\lVert b\rVert=|s|>0$。正交时 $a\cdot b=ts$，Cauchy--Schwarz 不等式取等，故 $b=ca$；代入 $a\cdot b=ts$ 及 $\lVert a\rVert^2=t^2$，得 $s=ct$，于是 $y=cx$。反向由 $g(x,x)=0$ 立即得到。
\end{solve}

\section{基础训练与方法整合}
\label{sec:laex-bilinear-training}

\subsection{对偶基的逆转置与双线性型的维数}

\begin{example}
\label{ex:laex-dual-basis-change}
若新旧基满足 $(e'_1,\ldots,e'_n)=(e_1,\ldots,e_n)P$，求对偶基的过渡矩阵，并计算三维基 $e'_1=e_1+e_2+e_3$、$e'_2=e_2+e_3$、$e'_3=e_3$ 的对偶基。
\end{example}
\begin{solve}
设 $(e'^1,\ldots,e'^n)=(e^1,\ldots,e^n)Q$。由 $e'^i(e'_j)=\delta_{ij}$ 得 $Q^{\mathsf T}P=I$，故
\begin{equation}
 Q=P^{-\mathsf T}.
 \label{eq:laex-dual-basis-change}
\end{equation}
在所给三维例子中，$P=\begin{pmatrix}1&0&0\\1&1&0\\1&1&1\end{pmatrix}$，所以 $e'^1=e^1$、$e'^2=-e^1+e^2$、$e'^3=-e^2+e^3$。可逐一在新基上取值检查。对偶映射的矩阵为转置，对偶基的过渡矩阵为逆转置，两者来自不同的问题。
\end{solve}

若 $\dim U=m$、$\dim V=n$，给定基后 $g\mapsto G$ 是双线性型空间到 $M_{m\times n}(\mathbb F)$ 的线性同构，故其维数为 $mn$。在 $n$ 维空间上，对称型的自由参数为主对角元与上三角元，维数为 $n(n+1)/2$；交错型对角元为零、下三角元由上三角元确定，维数为 $n(n-1)/2$。两者维数之和为 $n^2$，对应式~\eqref{eq:laex-form-split} 的直和分解。

\subsection{由表达式识别矩阵、退化性与变量次序}

判定一个函数是否双线性，应分别固定一个变量检验齐次性与可加性。形如 $\sum_{i,j}a_{ij}x_iy_j$ 的表达式是双线性的；常数项、仅依赖一个变量的非零一次项、$x_i^2y_j$ 等项一般破坏条件。对复数变量，$\overline{x_i}y_j$ 是半双线性项，不是这里的双线性项。矩阵 $G=(a_{ij})$ 的行指标对应第一个变量，列指标对应第二个变量。

若 $T$ 的矩阵为 $B$，则 $h(x,y)=g(Tx,y)$ 的矩阵为 $B^{\mathsf T}G$，$h(x,y)=g(x,Ty)$ 的矩阵为 $GB$。例如 $g(x,y)=\langle f,x\rangle\langle h,y\rangle$ 的矩阵为 $ab^{\mathsf T}$，其中 $a,b$ 为两个函数的对偶基坐标列，非零时秩为 $1$。又若固定非退化配对 $b$ 并定义 $g(x,y)=b(x,Ty)$，则 $R_g=\ker T$；交换变量中 $T$ 的位置后，对应的则是左根子空间。

\begin{example}
\label{ex:laex-idempotent-form}
设 $P\in M_n(\mathbb F)$ 满足 $P^2=P$，且 $P\ne I_n$。说明 $g(x,y)=x^{\mathsf T}Py$ 必退化。
\end{example}
\begin{solve}
若 $P$ 可逆，$P^2=P$ 左乘 $P^{-1}$ 得 $P=I_n$，矛盾。因此 $P$ 奇异，$\ker P\ne0$，这正是 $g$ 的右根子空间；由式~\eqref{eq:laex-radical-dimensions}，方阵情形的左根子空间也非零。幂等性本身并不保证对称性或正定性。
\end{solve}

\subsection{积分配对与 Gram 矩阵}

\begin{example}
\label{ex:laex-polynomial-integral-form}
设 $V=\{f\in\mathbb R[t]:\deg f<3\}$，定义 $g(f,h)=\int_0^1 f(t)h(t)\,dt$。求它在基 $1,t,t^2$ 下的矩阵并判定非退化性。
\end{example}
\begin{solve}
积分对两个变量分别线性，且 $g(t^i,t^j)=1/(i+j+1)$（$0\le i,j\le2$），所以表示矩阵为
\[
 H=\begin{pmatrix}
 1&\tfrac12&\tfrac13\\
 \tfrac12&\tfrac13&\tfrac14\\
 \tfrac13&\tfrac14&\tfrac15
 \end{pmatrix}.
\]
对非零实多项式 $f$，它不可能在整个 $[0,1]$ 上为零；由连续性，在某个正长度子区间上 $f^2$ 严格为正，故 $g(f,f)>0$。因此这是正定对称型，$H$ 可逆；直接计算也有 $\det H=1/2160>0$。非退化性证明来自积分的严格正性，比只计算一个行列式更能解释其原因。
\end{solve}

本章的计算可以沿同一条路线检查：先写出两个变量各自所在的空间及其坐标，再确认换基是分别进行还是同时进行；涉及对偶时用评价恒等式确定映射方向，涉及非退化性时计算根子空间，涉及几何分解时检查限制是否非退化。交错型通过成对选基消去混合项，对称型通过非迷向方向分裂出正交补；这两种选基过程都依赖非退化限制，而不只依赖维数相加。

\end{document}
